% !TeX program = xelatex
% Standalone English volume. Use XeLaTeX until the references stabilize.
\documentclass[a4paper,11pt]{article}
\usepackage[T1]{fontenc}
\usepackage{lmodern}
\usepackage[english]{babel}
\usepackage[left=26mm,right=26mm,top=25mm,bottom=25mm,headheight=15pt,headsep=7mm,footskip=11mm]{geometry}
\usepackage{amsmath,amssymb,mathtools,amsthm}
\newtheorem{proposition}{Proposition}[section]
\newtheorem{lemma}[proposition]{Lemma}
\newtheorem{theorem}[proposition]{Theorem}
\newtheorem{corollary}[proposition]{Corollary}
\newtheorem{maintheorem}[proposition]{Main Theorem}
\theoremstyle{definition}
\newtheorem{definition}[proposition]{Definition}
\theoremstyle{remark}
\newtheorem{remark}[proposition]{Remark}
\makeatletter
\renewcommand*\l@section{\@dottedtocline{1}{0em}{2.4em}}
\renewcommand*\l@subsection{\@dottedtocline{2}{1.45em}{3.45em}}
\renewcommand*\l@subsubsection{\@dottedtocline{3}{3.8em}{4.55em}}
\makeatother
\usepackage{xcolor,fancyhdr,titlesec,needspace,xurl}
\usepackage{longtable,multicol}
\usepackage{hyperref}
\definecolor{ink}{HTML}{1A3542}
\definecolor{muted}{HTML}{4F606A}
\hypersetup{unicode=true,colorlinks=true,linkcolor=ink,filecolor=ink,urlcolor=ink,citecolor=ink,bookmarksnumbered=true,
 pdfdisplaydoctitle=true,pdflang={en},
 pdftitle={The generalized Riemann hypothesis for Dirichlet L-functions. Edition 930, Volume XX},
 pdfsubject={A natural-language reconstruction of the proposed proof},
 pdfauthor={A mathematical reconstruction}}
\titleformat{\section}{\Large\bfseries\color{ink}}{\thesection}{.65em}{}
\titlespacing*{\section}{0pt}{.8em}{.65em}
\setlength{\parindent}{1.2em}
\linespread{1.045}
\setlength{\parskip}{.22em plus .05em minus .03em}
\setlength{\emergencystretch}{2em}
\setcounter{topnumber}{2}
\setcounter{tocdepth}{3}\setcounter{secnumdepth}{3}
\widowpenalty=4500
\clubpenalty=4500
\pagestyle{fancy}\fancyhf{}
\fancyhead[L]{\small\color{muted}Dirichlet GRH \textperiodcentered\ Volume I}
\fancyhead[R]{\small\color{muted}Dirichlet zeros and power intervals}
\fancyfoot[C]{\small\thepage}
\renewcommand{\headrulewidth}{.3pt}
\newcommand{\lead}[1]{\par\smallskip\noindent{\bfseries\color{ink}#1.}\enspace}
\newcommand{\sour}[2]{\unskip\penalty150\hspace{.35em}{\small\hyperlink{ref#1}{[#2]}}}
\titleformat{\subsection}{\large\bfseries}{\thesubsection}{.6em}{}
\titlespacing*{\subsection}{0pt}{1.15em plus .15em}{.4em}
\setlength{\abovedisplayskip}{9pt plus 2pt minus 3pt}
\setlength{\belowdisplayskip}{9pt plus 2pt minus 3pt}
\newcommand{\RR}{\mathbb R}\newcommand{\CC}{\mathbb C}\newcommand{\QQ}{\mathbb Q}
\newcommand{\ZZ}{\mathbb Z}\newcommand{\ord}{\operatorname{ord}}
\newcommand{\Frob}{\operatorname{Frob}}\newcommand{\Gal}{\operatorname{Gal}}
\newcommand{\len}{\operatorname{length}}\newcommand{\Res}{\operatorname*{Res}}
\newcommand{\wmean}{\operatorname{av}}\newcommand{\Ker}{\operatorname{ker}}
\newcommand{\Lh}{\mathcal L}\newcommand{\Gg}{\mathcal G}
\newcommand{\real}{\operatorname{Re}}\newcommand{\imag}{\operatorname{Im}}
\newcommand{\abs}[1]{\left|#1\right|}\newcommand{\eps}{\varepsilon}
\newcommand{\anchoritem}[2]{\par\Needspace{2.8\baselineskip}\hypertarget{ref#1}{}\noindent{\bfseries\color{ink}#2.}\enspace}
\newcommand{\srcloc}[2]{{\small\nolinkurl{#1.lean}}, line~#2}
\numberwithin{equation}{section}
\raggedbottom

\newcommand{\Sel}{\mathcal S}
\newcommand{\Cl}{\operatorname{Cl}}
\newcommand{\Aut}{\operatorname{Aut}}
\newcommand{\GL}{\operatorname{GL}}
\newcommand{\OO}{\mathcal O}
\hypersetup{bookmarksdepth=3}


\addto\captionsenglish{\renewcommand{\partname}{Volume}}


\makeatletter
\expandafter\def\csname r@cor:normaldegree347\endcsname{{127.2}{27}{}{proposition.127.2}{Dirichlet-GRH-v930-Volume-XXIV.pdf}}
\expandafter\def\csname r@cor:original200\endcsname{{83.7}{57}{}{proposition.83.7}{Dirichlet-GRH-v930-Volume-XVI.pdf}}
\expandafter\def\csname r@cor:vectors169\endcsname{{72.7}{51}{}{proposition.72.7}{Dirichlet-GRH-v930-Volume-XIV.pdf}}
\expandafter\def\csname r@eq:letters111\endcsname{{9.1}{117}{}{equation.9.1}{Dirichlet-GRH-v930-Volume-I.pdf}}
\expandafter\def\csname r@eq:seed127\endcsname{{92.16}{20}{}{equation.92.16}{Dirichlet-GRH-v930-Volume-XVIII.pdf}}
\expandafter\def\csname r@eq:src109\endcsname{{87.4}{37}{}{equation.87.4}{Dirichlet-GRH-v930-Volume-XVII.pdf}}
\expandafter\def\csname r@eq:tail114\endcsname{{19.15}{17}{}{equation.19.15}{Dirichlet-GRH-v930-Volume-IV.pdf}}
\expandafter\def\csname r@eq:window\endcsname{{{W}}{2}{}{AMS.2}{Dirichlet-GRH-v930-Volume-I.pdf}}
\expandafter\def\csname r@lem:blockphysical295\endcsname{{100.12}{73}{}{proposition.100.12}{Dirichlet-GRH-v930-Volume-XIX.pdf}}
\expandafter\def\csname r@lem:cas96\endcsname{{121.25}{55}{}{proposition.121.25}{Dirichlet-GRH-v930-Volume-XXIII.pdf}}
\expandafter\def\csname r@lem:forms203\endcsname{{110.10}{46}{}{proposition.110.10}{Dirichlet-GRH-v930-Volume-XXI.pdf}}
\expandafter\def\csname r@lem:good175\endcsname{{72.18}{60}{}{proposition.72.18}{Dirichlet-GRH-v930-Volume-XIV.pdf}}
\expandafter\def\csname r@lem:hook86\endcsname{{92.4}{5}{}{proposition.92.4}{Dirichlet-GRH-v930-Volume-XVIII.pdf}}
\expandafter\def\csname r@lem:monomials177\endcsname{{72.20}{63}{}{proposition.72.20}{Dirichlet-GRH-v930-Volume-XIV.pdf}}
\expandafter\def\csname r@lem:project271\endcsname{{71.5}{19}{}{proposition.71.5}{Dirichlet-GRH-v930-Volume-XIV.pdf}}
\expandafter\def\csname r@lem:returns261\endcsname{{111.3}{59}{}{proposition.111.3}{Dirichlet-GRH-v930-Volume-XXI.pdf}}
\expandafter\def\csname r@lem:root94\endcsname{{121.20}{51}{}{proposition.121.20}{Dirichlet-GRH-v930-Volume-XXIII.pdf}}
\expandafter\def\csname r@lem:signs286\endcsname{{95.5}{69}{}{proposition.95.5}{Dirichlet-GRH-v930-Volume-XVIII.pdf}}
\expandafter\def\csname r@lem:size97\endcsname{{87.1}{33}{}{proposition.87.1}{Dirichlet-GRH-v930-Volume-XVII.pdf}}
\expandafter\def\csname r@lem:threeminors349\endcsname{{120.3}{16}{}{proposition.120.3}{Dirichlet-GRH-v930-Volume-XXIII.pdf}}
\expandafter\def\csname r@lem:topchar266\endcsname{{88.10}{65}{}{proposition.88.10}{Dirichlet-GRH-v930-Volume-XVII.pdf}}
\expandafter\def\csname r@lem:translate163\endcsname{{86.19}{16}{}{proposition.86.19}{Dirichlet-GRH-v930-Volume-XVII.pdf}}
\expandafter\def\csname r@nav:I-28\endcsname{{6.6}{32}{}{subsection.6.6}{Dirichlet-GRH-v930-Volume-I.pdf}}
\expandafter\def\csname r@nav:VIII-02\endcsname{{72.3}{50}{}{subsection.72.3}{Dirichlet-GRH-v930-Volume-XIV.pdf}}
\expandafter\def\csname r@nav:VIII-06\endcsname{{72.7}{60}{}{subsection.72.7}{Dirichlet-GRH-v930-Volume-XIV.pdf}}
\expandafter\def\csname r@nav:VIII-07\endcsname{{72.8}{63}{}{subsection.72.8}{Dirichlet-GRH-v930-Volume-XIV.pdf}}
\expandafter\def\csname r@prop:besorbit265\endcsname{{88.4}{56}{}{proposition.88.4}{Dirichlet-GRH-v930-Volume-XVII.pdf}}
\expandafter\def\csname r@prop:block295\endcsname{{100.7}{69}{}{proposition.100.7}{Dirichlet-GRH-v930-Volume-XIX.pdf}}
\expandafter\def\csname r@prop:blockfinite295\endcsname{{100.2}{64}{}{proposition.100.2}{Dirichlet-GRH-v930-Volume-XIX.pdf}}
\expandafter\def\csname r@prop:blockfinite297\endcsname{{101.7}{86}{}{proposition.101.7}{Dirichlet-GRH-v930-Volume-XIX.pdf}}
\expandafter\def\csname r@prop:blockmarks296\endcsname{{100.3}{65}{}{proposition.100.3}{Dirichlet-GRH-v930-Volume-XIX.pdf}}
\expandafter\def\csname r@prop:composition298\endcsname{{100.9}{71}{}{proposition.100.9}{Dirichlet-GRH-v930-Volume-XIX.pdf}}
\expandafter\def\csname r@prop:corecusp251\endcsname{{71.8}{22}{}{proposition.71.8}{Dirichlet-GRH-v930-Volume-XIV.pdf}}
\expandafter\def\csname r@prop:cost177\endcsname{{72.23}{65}{}{proposition.72.23}{Dirichlet-GRH-v930-Volume-XIV.pdf}}
\expandafter\def\csname r@prop:dominantexchange344\endcsname{{86.41}{32}{}{proposition.86.41}{Dirichlet-GRH-v930-Volume-XVII.pdf}}
\expandafter\def\csname r@prop:extend300\endcsname{{102.3}{90}{}{proposition.102.3}{Dirichlet-GRH-v930-Volume-XIX.pdf}}
\expandafter\def\csname r@prop:field12\endcsname{{6.4}{27}{}{proposition.6.4}{Dirichlet-GRH-v930-Volume-I.pdf}}
\expandafter\def\csname r@prop:flagunmarked264\endcsname{{102.1}{87}{}{proposition.102.1}{Dirichlet-GRH-v930-Volume-XIX.pdf}}
\expandafter\def\csname r@prop:jacorbit266\endcsname{{75.6}{118}{}{proposition.75.6}{Dirichlet-GRH-v930-Volume-XIV.pdf}}
\expandafter\def\csname r@prop:lattice282\endcsname{{99.11}{51}{}{proposition.99.11}{Dirichlet-GRH-v930-Volume-XIX.pdf}}
\expandafter\def\csname r@prop:marked291\endcsname{{99.3}{42}{}{proposition.99.3}{Dirichlet-GRH-v930-Volume-XIX.pdf}}
\expandafter\def\csname r@prop:plane281\endcsname{{111.15}{71}{}{proposition.111.15}{Dirichlet-GRH-v930-Volume-XXI.pdf}}
\expandafter\def\csname r@prop:satake287\endcsname{{101.2}{81}{}{proposition.101.2}{Dirichlet-GRH-v930-Volume-XIX.pdf}}
\expandafter\def\csname r@prop:shift203\endcsname{{110.11}{47}{}{proposition.110.11}{Dirichlet-GRH-v930-Volume-XXI.pdf}}
\expandafter\def\csname r@prop:sphere289\endcsname{{101.3}{82}{}{proposition.101.3}{Dirichlet-GRH-v930-Volume-XIX.pdf}}
\expandafter\def\csname r@prop:split287\endcsname{{100.10}{72}{}{proposition.100.10}{Dirichlet-GRH-v930-Volume-XIX.pdf}}
\expandafter\def\csname r@prop:tracelat253\endcsname{{122.7}{63}{}{proposition.122.7}{Dirichlet-GRH-v930-Volume-XXIII.pdf}}
\expandafter\def\csname r@prop:transfer286\endcsname{{89.6}{80}{}{proposition.89.6}{Dirichlet-GRH-v930-Volume-XVII.pdf}}
\expandafter\def\csname r@prop:weight163\endcsname{{86.20}{17}{}{proposition.86.20}{Dirichlet-GRH-v930-Volume-XVII.pdf}}
\expandafter\def\csname r@sec:absolute163\endcsname{{86.7}{15}{}{subsection.86.7}{Dirichlet-GRH-v930-Volume-XVII.pdf}}
\expandafter\def\csname r@sec:char94\endcsname{{121.16}{51}{}{subsection.121.16}{Dirichlet-GRH-v930-Volume-XXIII.pdf}}
\expandafter\def\csname r@sec:lowerlattice347\endcsname{{83.13}{73}{}{subsection.83.13}{Dirichlet-GRH-v930-Volume-XVI.pdf}}
\expandafter\def\csname r@sec:missing\endcsname{{6.6}{32}{}{subsection.6.6}{Dirichlet-GRH-v930-Volume-I.pdf}}
\expandafter\def\csname r@sec:realtower169\endcsname{{72.3}{50}{}{subsection.72.3}{Dirichlet-GRH-v930-Volume-XIV.pdf}}
\expandafter\def\csname r@sec:success135\endcsname{{86.3}{8}{}{subsection.86.3}{Dirichlet-GRH-v930-Volume-XVII.pdf}}
\expandafter\def\csname r@seriesguide\endcsname{{1.1}{4}{}{subsection.1.1}{Dirichlet-GRH-v930-Volume-I.pdf}}
\expandafter\def\csname r@thm:allflags298\endcsname{{111.14}{69}{}{proposition.111.14}{Dirichlet-GRH-v930-Volume-XXI.pdf}}
\expandafter\def\csname r@thm:bestransfer265\endcsname{{88.16}{72}{}{proposition.88.16}{Dirichlet-GRH-v930-Volume-XVII.pdf}}
\expandafter\def\csname r@thm:both184\endcsname{{20.12}{32}{}{proposition.20.12}{Dirichlet-GRH-v930-Volume-IV.pdf}}
\expandafter\def\csname r@thm:bounds282\endcsname{{101.6}{85}{}{proposition.101.6}{Dirichlet-GRH-v930-Volume-XIX.pdf}}
\expandafter\def\csname r@thm:capelli250\endcsname{{71.7}{21}{}{proposition.71.7}{Dirichlet-GRH-v930-Volume-XIV.pdf}}
\expandafter\def\csname r@thm:codimension294\endcsname{{111.13}{68}{}{proposition.111.13}{Dirichlet-GRH-v930-Volume-XXI.pdf}}
\expandafter\def\csname r@thm:common178\endcsname{{72.30}{71}{}{proposition.72.30}{Dirichlet-GRH-v930-Volume-XIV.pdf}}
\expandafter\def\csname r@thm:commonfischer348\endcsname{{120.2}{15}{}{proposition.120.2}{Dirichlet-GRH-v930-Volume-XXIII.pdf}}
\expandafter\def\csname r@thm:compact248\endcsname{{71.15}{31}{}{proposition.71.15}{Dirichlet-GRH-v930-Volume-XIV.pdf}}
\expandafter\def\csname r@thm:completeplane293\endcsname{{111.12}{68}{}{proposition.111.12}{Dirichlet-GRH-v930-Volume-XXI.pdf}}
\expandafter\def\csname r@thm:conductor232\endcsname{{19.15}{12}{}{proposition.19.15}{Dirichlet-GRH-v930-Volume-IV.pdf}}
\expandafter\def\csname r@thm:corecopy251\endcsname{{71.1}{13}{}{proposition.71.1}{Dirichlet-GRH-v930-Volume-XIV.pdf}}
\expandafter\def\csname r@thm:covectorpaths352\endcsname{{73.1}{73}{}{proposition.73.1}{Dirichlet-GRH-v930-Volume-XIV.pdf}}
\expandafter\def\csname r@thm:decay96\endcsname{{121.26}{55}{}{proposition.121.26}{Dirichlet-GRH-v930-Volume-XXIII.pdf}}
\expandafter\def\csname r@thm:direction249\endcsname{{71.6}{19}{}{proposition.71.6}{Dirichlet-GRH-v930-Volume-XIV.pdf}}
\expandafter\def\csname r@thm:dominantflag344\endcsname{{9.21}{125}{}{proposition.9.21}{Dirichlet-GRH-v930-Volume-I.pdf}}
\expandafter\def\csname r@thm:finite281\endcsname{{99.10}{50}{}{proposition.99.10}{Dirichlet-GRH-v930-Volume-XIX.pdf}}
\expandafter\def\csname r@thm:first250\endcsname{{71.10}{25}{}{proposition.71.10}{Dirichlet-GRH-v930-Volume-XIV.pdf}}
\expandafter\def\csname r@thm:first284\endcsname{{111.7}{63}{}{proposition.111.7}{Dirichlet-GRH-v930-Volume-XXI.pdf}}
\expandafter\def\csname r@thm:fullperiod266\endcsname{{89.5}{79}{}{proposition.89.5}{Dirichlet-GRH-v930-Volume-XVII.pdf}}
\expandafter\def\csname r@thm:global175\endcsname{{72.19}{61}{}{proposition.72.19}{Dirichlet-GRH-v930-Volume-XIV.pdf}}
\expandafter\def\csname r@thm:global176\endcsname{{72.29}{70}{}{proposition.72.29}{Dirichlet-GRH-v930-Volume-XIV.pdf}}
\expandafter\def\csname r@thm:grid354\endcsname{{99.22}{62}{}{proposition.99.22}{Dirichlet-GRH-v930-Volume-XIX.pdf}}
\expandafter\def\csname r@thm:invar148\endcsname{{20.11}{31}{}{proposition.20.11}{Dirichlet-GRH-v930-Volume-IV.pdf}}
\expandafter\def\csname r@thm:jactransfer266\endcsname{{88.14}{69}{}{proposition.88.14}{Dirichlet-GRH-v930-Volume-XVII.pdf}}
\expandafter\def\csname r@thm:localmass156\endcsname{{87.24}{48}{}{proposition.87.24}{Dirichlet-GRH-v930-Volume-XVII.pdf}}
\expandafter\def\csname r@thm:lowerlattice347\endcsname{{83.28}{74}{}{proposition.83.28}{Dirichlet-GRH-v930-Volume-XVI.pdf}}
\expandafter\def\csname r@thm:lowesttwo261\endcsname{{75.4}{116}{}{proposition.75.4}{Dirichlet-GRH-v930-Volume-XIV.pdf}}
\expandafter\def\csname r@thm:markedstrict337\endcsname{{22.11}{70}{}{proposition.22.11}{Dirichlet-GRH-v930-Volume-IV.pdf}}
\expandafter\def\csname r@thm:norm177\endcsname{{72.21}{64}{}{proposition.72.21}{Dirichlet-GRH-v930-Volume-XIV.pdf}}
\expandafter\def\csname r@thm:norm91\endcsname{{94.9}{44}{}{proposition.94.9}{Dirichlet-GRH-v930-Volume-XVIII.pdf}}
\expandafter\def\csname r@thm:ordinary163\endcsname{{86.22}{18}{}{proposition.86.22}{Dirichlet-GRH-v930-Volume-XVII.pdf}}
\expandafter\def\csname r@thm:packing246\endcsname{{111.17}{72}{}{proposition.111.17}{Dirichlet-GRH-v930-Volume-XXI.pdf}}
\expandafter\def\csname r@thm:phase282\endcsname{{111.5}{61}{}{proposition.111.5}{Dirichlet-GRH-v930-Volume-XXI.pdf}}
\expandafter\def\csname r@thm:pointlattice353\endcsname{{99.20}{60}{}{proposition.99.20}{Dirichlet-GRH-v930-Volume-XIX.pdf}}
\expandafter\def\csname r@thm:prime210\endcsname{{93.17}{34}{}{proposition.93.17}{Dirichlet-GRH-v930-Volume-XVIII.pdf}}
\expandafter\def\csname r@thm:product281\endcsname{{111.4}{60}{}{proposition.111.4}{Dirichlet-GRH-v930-Volume-XXI.pdf}}
\expandafter\def\csname r@thm:projection86\endcsname{{92.5}{6}{}{proposition.92.5}{Dirichlet-GRH-v930-Volume-XVIII.pdf}}
\expandafter\def\csname r@thm:range239\endcsname{{88.1}{52}{}{proposition.88.1}{Dirichlet-GRH-v930-Volume-XVII.pdf}}
\expandafter\def\csname r@thm:rangenew50\endcsname{{6.29}{62}{}{proposition.6.29}{Dirichlet-GRH-v930-Volume-I.pdf}}
\expandafter\def\csname r@thm:ratio248\endcsname{{89.3}{76}{}{proposition.89.3}{Dirichlet-GRH-v930-Volume-XVII.pdf}}
\expandafter\def\csname r@thm:raw97\endcsname{{87.2}{33}{}{proposition.87.2}{Dirichlet-GRH-v930-Volume-XVII.pdf}}
\expandafter\def\csname r@thm:realrow257\endcsname{{74.1}{100}{}{proposition.74.1}{Dirichlet-GRH-v930-Volume-XIV.pdf}}
\expandafter\def\csname r@thm:rest96\endcsname{{121.24}{54}{}{proposition.121.24}{Dirichlet-GRH-v930-Volume-XXIII.pdf}}
\expandafter\def\csname r@thm:sectors149\endcsname{{86.11}{10}{}{proposition.86.11}{Dirichlet-GRH-v930-Volume-XVII.pdf}}
\expandafter\def\csname r@thm:short246\endcsname{{111.1}{57}{}{proposition.111.1}{Dirichlet-GRH-v930-Volume-XXI.pdf}}
\expandafter\def\csname r@thm:split281\endcsname{{101.1}{80}{}{proposition.101.1}{Dirichlet-GRH-v930-Volume-XIX.pdf}}
\expandafter\def\csname r@thm:src109\endcsname{{87.7}{37}{}{proposition.87.7}{Dirichlet-GRH-v930-Volume-XVII.pdf}}
\expandafter\def\csname r@thm:sum287\endcsname{{111.8}{64}{}{proposition.111.8}{Dirichlet-GRH-v930-Volume-XXI.pdf}}
\expandafter\def\csname r@thm:tail283\endcsname{{111.6}{62}{}{proposition.111.6}{Dirichlet-GRH-v930-Volume-XXI.pdf}}
\expandafter\def\csname r@thm:terminaldual325\endcsname{{122.6}{62}{}{proposition.122.6}{Dirichlet-GRH-v930-Volume-XXIII.pdf}}
\expandafter\def\csname r@thm:terminalformula347\endcsname{{127.1}{26}{}{proposition.127.1}{Dirichlet-GRH-v930-Volume-XXIV.pdf}}
\expandafter\def\csname r@thm:tower169\endcsname{{72.6}{50}{}{proposition.72.6}{Dirichlet-GRH-v930-Volume-XIV.pdf}}
\expandafter\def\csname r@thm:tower170\endcsname{{72.26}{68}{}{proposition.72.26}{Dirichlet-GRH-v930-Volume-XIV.pdf}}
\expandafter\def\csname r@thm:vectorchain353\endcsname{{73.2}{74}{}{proposition.73.2}{Dirichlet-GRH-v930-Volume-XIV.pdf}}
\expandafter\def\csname r@thm:wordlower348\endcsname{{83.29}{75}{}{proposition.83.29}{Dirichlet-GRH-v930-Volume-XVI.pdf}}
\makeatother
\begin{document}
\setcounter{part}{19}
\fancyhead[L]{\small\color{muted}Dirichlet GRH \textperiodcentered\ Volume XX}
\fancyhead[R]{\small\color{muted}Real orbital integrals}
\thispagestyle{plain}\begin{center}
{\LARGE\bfseries The generalized Riemann hypothesis\\[2pt]for Dirichlet \(L\)-functions}\\[6pt]
{\small 27 November 2026 \quad\textperiodcentered\quad Edition 930}\end{center}
\part{Real orbital integrals and matrix coefficients}\label{part:XX}
\begin{abstract}\noindent We estimate the archimedean factors in the orbital expansion. We compute matrix coefficients and height moments for the first-plane and higher-rank terms. For a rational flag with integral unmarked root ratios we bound the actual Bessel integral at its displaced archimedean frequency.\end{abstract}
\setcounter{section}{103}
\noindent The \hyperref[seriesguide]{guide in Volume~I} lists all twenty-seven volumes.
\begingroup\small\setlength{\parskip}{0pt}\tableofcontents\endgroup
\clearpage

\section{Real theta norms and projected coefficients}\label{nav:g-X-0}\label{nav:volume-XX}

\subsection{Positive generators and the real scale of the theta tests}\label{nav:VIII-08}
\label{sec:units179}

The choices of the additive character and of the skew form in
Lemma~\ref{lem:good175} involve totally positive generators of
two ideals of \(F=E^+\). Their norms are prescribed by those
ideals. The real theta kernels use every conjugate of each
generator. We bound these conjugates separately, and then
compute how a change of scale enters the real norm of the
monomials in Theorem~\ref{thm:norm177}.

\begin{lemma}[A bound using positive cyclotomic units]
\label{lem:unit179}
Suppose \(F\) is totally real, \(F\subset\QQ(\zeta_N)\),
and \(f=[F:\QQ]\). Put \(B_F=(f-1)\varphi(N)\log(2N)\).
For each totally positive \(u\in F^\times\) there is a
totally positive unit \(v\in\OO_F^\times\) such that
\[
   \max_{\tau:F\hookrightarrow\RR}
   \left|\log\tau(uv)-\frac1f\log N_{F/\QQ}(u)\right|\le B_F.
\]
\end{lemma}
\begin{proof}
Let \(G=\operatorname{Gal}(F/\QQ)\), labelling its real
embeddings by a fixed inclusion and the elements of \(G\).
We construct a collection of positive units whose logarithm
vectors span the hyperplane of coordinate sum zero. If
\(f=1\) this space is zero and \(v=1\) works.

Take a nontrivial character of \(G\). Its primitive Dirichlet
character \(\chi\) has conductor \(h\mid N\) and is even
because complex conjugation fixes \(F\). The subfield fixed by the kernel of this character lies in both
\(F\) and \(\QQ(\zeta_h)\). Write
\(F_h=F\cap\QQ(\zeta_h)\), choose \(a\) prime to \(h\)
with \(\chi(a)\ne1\), and put
\[
 y_a=\frac{1-\zeta_h^a}{1-\zeta_h},\qquad
 v_a=N_{\QQ(\zeta_h)/F_h}(y_a\overline{y_a})\in F_h^\times.
\]
A geometric sum shows that \(y_a\) is integral. If \(ab\equiv
1\bmod h\), the corresponding sum in powers of \(\zeta_h^a\)
is \(y_a^{-1}\), hence \(y_a\) is a unit. All conjugates of
\(v_a\) are positive. Each \(|1-\zeta_h^j|\) with
\((j,h)=1\) is between \(2\sin(\pi/h)\) and two. The lower
bound is at least \(4/h\) for \(h\ge2\). It follows that
each coordinate of the log vector of \(v_a\) has absolute
value at most \(2\varphi(h)\log(2h)\).

This vector has nonzero coefficient in the \(\chi\)-direction.
Indeed the primitivity and even parity of \(\chi\) give
\[
 \sum_{j\bmod h}^{*}\overline{\chi(j)}\log|1-\zeta_h^j|
              =-\mathfrak g(\bar\chi)L(1,\chi),
 \qquad \mathfrak g(\bar\chi)=\sum_{j\bmod h}^{*}
                           \overline{\chi(j)}\zeta_h^j.
\]
To see this, expand \(-\log|1-s\zeta_h^j|\) in the
cosine series at \(0<s<1\). Evenness identifies the
two exponential sums for each cosine. The primitive Gauss
sum identity gives \(\chi(n)\mathfrak g(\bar\chi)\) for
the coefficient at \(s^n/n\), also when \((n,h)>1\) since
both sides then vanish. Letting \(s\) tend to one gives
the formula by Abel's theorem. The Gauss sum has absolute
value \(\sqrt h\), and Dirichlet's nonvanishing theorem
gives \(L(1,\chi)\ne0\). Replacing \(j\) by \(aj\) and
subtracting multiplies this nonzero coefficient by
\(\chi(a)-1\). Taking the norm defining \(v_a\) sums log
coordinates on the fibers above \(F_h\), where the character
is constant. The square of the absolute value contributes
a factor two. Inflating from \(F_h\) to the embeddings
of \(F\) contributes \([F:F_h]\). The coefficient thus
remains nonzero.

Include all \(G\)-translates of the units just constructed.
The complex span of their log vectors is a subrepresentation
of the regular representation of the finite abelian group
\(G\). Each character occurs once there. Applying the finite
character idempotents and using the preceding computation
shows that the span contains every nontrivial character
line. Unit norms give coordinate sum zero, so the real
vectors span the real hyperplane as asserted. Choose
\(f-1\) independent vectors from this collection. Their
supremum norms are at most \(2\varphi(N)\log(2N)\), since
\(h\mid N\). Express the log vector of \(u\) minus its
coordinate mean as a real linear combination of this basis.
Subtract the nearest integer from each coefficient by
multiplying by the corresponding inverse unit power. Each
remaining coefficient has absolute value at most one half.
The triangle inequality proves the stated bound.
\end{proof}

\begin{theorem}[Real and finite sizes for the same two generators]
\label{thm:balanced179}
Keep the very primes of Proposition~\ref{prop:cost177}. There
are totally positive elements \(\beta,\gamma\) with
\((\beta)\mathfrak d_F=\mathfrak p\) and
\((\gamma)\mathfrak a_0=\mathfrak q\) such that, at every
real embedding \(\tau\),
\begin{align*}
 \left|\log\tau\beta-\frac1f\log
       \frac{N\mathfrak p}{|D_F|}\right|&\le B_F,\\
 \left|\log\tau\gamma-\frac1f\log
       \frac{N\mathfrak q}{d_0^{f/2}}\right|&\le B_F.
\end{align*}
On the original intervals,
\(1+\max_\tau(|\log\tau\beta|+|\log\tau\gamma|)
= b^{o(1)}\), where \(b=\exp(x^\epsilon)\).
The two additional primes have logarithmic norms polynomial
in \(x\), and all original inert local data remain unramified.
\end{theorem}
\begin{proof}
The narrow ideal class equalities give initial positive
generators. The two means of their logarithms have the
displayed values since \(N\mathfrak d_F=|D_F|\) and
\(\mathfrak a_0^2=(d_0)\), the latter by extending to
\(E\) and using \((\sqrt{-d_0})^2=(d_0)\).
Apply the lemma separately to the two generators. Unit
multiplication preserves their ideals and real signs.
For these fields, Section \ref{nav:I-28} gives \(f\ll_\psi x^3\)
and \(\log(2N)\ll_{\psi,\epsilon} x^{\epsilon/5}
\log^2(2x)+\log(2x)\). Thus \(\log (1+B_F)=o(x^\epsilon)\).
The numerators of the two norm quotients have logarithms
bounded in Proposition~\ref{prop:cost177}. For the denominators
use \(\log|D_F|\le f\log N\) and \(d_0\mid N\). These bounds
are uniform on the chosen intervals and prove the assertion.
\end{proof}

\begin{proposition}[The real scalar under dilation]
\label{prop:dilation179}
Use one fixed compatible complex structure in each real
oscillator, and give its Fock vacuum norm one. Keep the
polynomials in these coordinates, the unit source vectors,
and the Haar measures on the real groups fixed. Let
\(d_{j,\tau}\) be the degree of the monomial on edge
\(j\to j+1\) at the place \(\tau\) in
Theorem~\ref{thm:norm177}, and write
\(D_\tau=\sum_{j=0}^{m-1}d_{j,\tau}\),
\(\bar D=f^{-1}\sum_\tau D_\tau\).
Let \(R_\infty(\beta,\gamma)\) be the product of the absolute
values of the real diagonal Rallis factors of these tests,
with the scaled additive character and skew form. Use
\(R_\infty(1,1)\) for the same calculation at the real places
with the original character and form. For
\(\lambda=\beta\gamma\) one has
\begin{equation}
 \frac{R_\infty(\beta,\gamma)}{R_\infty(1,1)}
      =\prod_\tau(\tau\lambda)^{-D_\tau},\qquad
 N_{F/\QQ}(\lambda)=
      \frac{N\mathfrak p\,N\mathfrak q}{|D_F|\,d_0^{f/2}} .
 \label{eq:dilation179}
\end{equation}
The absolute difference between the logarithm of this ratio
and \(-\bar D\log N_{F/\QQ}(\lambda)\) is at most
\(2B_F\sum_\tau|D_\tau-\bar D|\) for the generators above.
It is \(o(b/x^A)\) for every fixed \(A\) whenever the sum
of the degree deviations is bounded by a power of \(x\).
\end{proposition}
\begin{proof}
The real symplectic pairing with its additive character
acquires the positive scalar \(s=\tau\lambda\), since the
skew form is scaled by \(\tau\gamma\) and the character
by \(\tau\beta\). A compatible complex structure for the
initial pairing remains compatible for its positive multiple.
In complex dimension \(l\) the Gaussian measure for its
Fock inner product is proportional to
\(s^l\exp(-\pi s Q(z))\,dz\) for the initial positive
quadratic form \(Q\). The constant is chosen to give the
vacuum norm one. The isometry to the initial Fock space
sends a polynomial \(P(z)\) to \(P(z/\sqrt s)\), by change
of variables. The map \(z\mapsto\sqrt s\,z\) is symplectic
from the scaled space to the initial one and commutes with
both real isometry groups. Its induced Fock isometry
therefore intertwines their oscillators. Equivalently
one may transport the Heisenberg translations and their
multipliers through this linear change. The characters defining the unitary splittings are unchanged.
In the real oscillator the compact vacuum character depends on
the signatures, the sign of the real additive character and
these angular exponents. All agree after the positive dilation.
Any difference between the two group splittings would be a
character of the real unitary groups trivial on their maximal
compacts. Such a character is trivial since the abelianization of
a connected real unitary group is its determinant circle,
already covered by the maximal compact.
The metaplectic phase of an isometry has no effect on a
diagonal matrix coefficient.

For homogeneous degree \(d\), the image of the polynomial
is exactly \(s^{-d/2}P\). Its contraction gives the same
target line. Its oscillator matrix coefficient, and hence
its diagonal integral against the fixed source coefficient,
is multiplied by \(s^{-d}\). This reasoning applies at
each edge with the successive unit lines, even though the
fields used to define the real comparison alone need have
no prescribed finite normalization at \((1,1)\). Multiplication
gives the first identity. Multiplication of ideal norms
gives the second. The error after subtracting the mean is
\[
 -\sum_\tau(D_\tau-\bar D)
          \left(\log\tau\lambda-f^{-1}\log N_{F/\QQ}(\lambda)\right).
\]
Each factor in parentheses has absolute value at most
\(2B_F\) by the theorem. The stated asymptotic follows
from \(1+B_F=b^{o(1)}\).
\end{proof}

\begin{corollary}[Conjugate choices with the same finite cost]
\label{cor:mean179}
With the real row and its polynomial coordinates fixed,
the geometric mean of the ratios in \eqref{eq:dilation179}
over the choices \((\sigma\beta,\sigma\gamma)\), \(\sigma\in
\operatorname{Gal}(F/\QQ)\), is exactly
\((N_{F/\QQ}\lambda)^{-\bar D}\). Some choice attains
at most this value and some choice attains at least this
value. All choices satisfy the same estimates on the two
auxiliary prime norms and the conjugates of the generators.
\end{corollary}
\begin{proof}
The ideals \(\mathfrak d_F\) and \(\mathfrak a_0\) are
Galois invariant. For the latter this follows also from
\(\mathfrak a_0^2=(d_0)\) and unique factorization of
fractional ideals. The conjugate primes are split in \(E\)
since \(E/\QQ\) is abelian. They still lie above the same
two distinct rational primes outside the original levels.
For each fixed \(\tau\), the arithmetic mean over \(\sigma\)
of \(\log \tau\sigma\lambda\) equals
\(f^{-1}\log N_{F/\QQ}(\lambda)\). Apply the first identity
in \eqref{eq:dilation179}, keeping \(D_\tau\) at the fixed
place, and exponentiate. The two inequalities follow from
the elementary bounds on the mean of real numbers.
\end{proof}

\begin{proposition}[Dependence on positive units at fixed prime ideals]
\label{prop:unitfreedom179}
Fix the ideals, the real monomials and the vacuum normalization
of Proposition~\ref{prop:dilation179}. The real normalizing
factor is independent of the choice of positive generators
if and only if \(D_\tau\) is independent of \(\tau\). If the
degrees are nonconstant, the possible positive real factors
have infimum zero and are unbounded above even when both
prime ideals themselves are fixed.
\end{proposition}
\begin{proof}
For parallel degrees the right side of \eqref{eq:dilation179}
is a power of a fixed ideal norm. For nonparallel degrees the
vector \((D_\tau)_\tau\) has a nonzero projection onto the
hyperplane of coordinate sum zero. The spanning assertion in
Lemma~\ref{lem:unit179} supplies a positive unit \(v\) with
\(c=\sum_\tau D_\tau\log\tau v\ne0\). Replace \(\beta\) by
\(\beta v^k\), leaving \(\gamma\) fixed, for \(k\in\ZZ\).
All ideals and signs are preserved. By the dilation identity
the real factor is multiplied by \(\exp(-kc)\). Both conclusions
follow by letting \(k\) run through positive and negative
integers. The bounded representatives in
Theorem~\ref{thm:balanced179} are a particular choice in this
family of permissible generators.
\end{proof}

This controls the dependence of the real norm on the two
generators in the vacuum normalization. The reference factors
\(R_\infty(1,1)\), the degrees for the designated arithmetic
tensors, and the pairing with the required integral projection
still enter the global period.

\subsection{An exact projected norm on the last compact pair}\label{nav:VIII-09}
\label{sec:last181}

At the last forward step the source and target at each real
place are both compact. This permits an absolute calculation
of the real factor in a fixed matrix frame, including a
projection to a specified target line. For this section put
\(n=n_{m-1}\), \(N'=n_m=n+t\), and write \(a_T=-e_m+h_{m-1}\)
and \(a_S=-e_m+h_m\) at the real place in question. Write
\(p_1>\cdots>p_n\) for the entries of \(P_{m-1}\) there.
Use probability Haar measures on the two real groups. In
orthonormal matrix coordinates \(Z\in M_{N',n}(\CC)\) take
the Gaussian \(e^{-\pi\operatorname{tr}(Z^*Z)}dZ\), with
each complex coordinate given ordinary plane Lebesgue measure.
It has mass one. These are fixed coordinates for the reference
real calculation. A further positive dilation has the exact
effect of Proposition~\ref{prop:dilation179}.

\begin{lemma}[The unique degree at the last place]
\label{lem:degree181}
Put
\[
 \lambda_i=p_i-\frac{n+1}{2}+i-a_T-\frac{N'}2
       \quad(1\le i\le n),\qquad d_*=|\lambda|=\sum_i\lambda_i.
\]
These entries form a partition. After the vacuum characters
are removed, the joint compact type is
\(S_\lambda(\CC^n)\boxtimes S_\lambda(\CC^{N'})\), where
\(S_\lambda\) denotes the Schur representation. It occurs
once and only in degree \(d_*\) of the matrix polynomials.
At a place \(\tau\) with weight offset \(B_\tau\) from
equation \ref{eq:letters111} and the weight lists following that equation, the degree is
\begin{align}
 d_{*,\tau}
   &=(k-1)B_\tau+\sum_{i=1}^{k-1} v_i-\frac{nN'}2 \notag\\
   &\quad+
       \big((k-1)(2m-1)+m^2+t m(m-1)\big)H_0 .      \label{eq:lastdegree181}
\end{align}
Here the lists are equivalently specified by
\(b_{i,\tau}=B_\tau+2(i-1)H_0\) and
\(e_{i,\tau}=B_\tau+(2i-1)H_0\).
\end{lemma}
\begin{proof}
Both smaller signatures are zero by
\(s_j=\lfloor(m-j)/2\rfloor\). The source highest weight is
\((p_i-(n+1)/2+i)_i\). In the forward convention the compact
Fock vacuum has determinant exponents \(a_T+N'/2\) on the
source and \(a_S+n/2\) on the target. These are the compact
shifts of \href{https://arxiv.org/pdf/2002.09148}{Ichino,
Lemma 7.8}, with both negative signature blocks empty.
In particular both polynomial weights in the forward
pair have the same partition. Equivalently this follows
by expanding the symmetric algebra of the two defining
spaces after removing the vacuum characters. Dualizing the
entire realization gives the conjugate pair and the same
dimensions and squared norms. Subtraction on the source
gives the displayed \(\lambda_i\). The entries are integral
by the parity of the two ranks and splittings. Successive
differences are nonnegative since \(p_i-p_{i+1}\ge1\).
Moreover
\(\lambda_n=p_n-a_T-(t+1)/2>0\), since every entry of the
centered source list exceeds \((t+1)/2\) by the spacings in
Section~\ref{sec:realtower169}.

For degree \(d\), the symmetric power of the tensor of the
two defining spaces is the space of simultaneous symmetric
group invariants in the tensor of their \(d\)-fold powers.
Schur--Weyl duality decomposes each tensor power into
Schur modules times symmetric group irreducibles. The
invariants have dimension one on each pair of equal
symmetric group labels and zero on two different labels.
Thus precisely the matching partitions of size \(d\)
occur, with multiplicity one. On the target, padding \(\lambda\)
by \(t\) zeros and adding \(a_S+n/2\) also gives the highest
weight computed from \(P_m\) in Theorem~\ref{thm:tower169}.

For the sum use \(\sum_i ((n+1)/2-i)=0\). The centered
entries on the source are \(v_i+e_m\), \(e_m-b_i\), and the
strings of length \(t\) centered at \(e_m-e_i\) for \(i<m\).
The sum of \(e_m-b_i=(2m+1-2i)H_0\) over \(i\le m\) is
\(m^2H_0\). The sums of the symmetric strings give
\(tm(m-1)H_0\). Inserting \(e_m=B_\tau+(2m-1)H_0\) and
subtracting \(nN'/2\) proves \eqref{eq:lastdegree181}.
\end{proof}

Write \(D_n=\dim S_\lambda(\CC^n)\) and
\(D_{N'}=\dim S_\lambda(\CC^{N'})\). For unit source
and target vectors \(u,v\) in the respective types, let
\(e_v\) be the compact matrix-unit projector onto the line
of \(v\) in each copy of the target type, zero on the other
types. For a polynomial \(P\) the last projected local norm is
\[
 Z_v(P;u)=\int_{U(n)}
     \langle\omega(g)e_v P,e_vP\rangle
                \overline{\langle\sigma(g)u,u\rangle}\,dg.
\]
The source and target actions in this formula include their
vacuum twists. Equivariance of theta lifting identifies
projection of the target form with applying \(e_v\) to the
test. In particular selecting the line in this definition
applies directly to the global norm formula of Section~\ref{nav:VIII-07}.

\begin{theorem}[Monomials detecting a prescribed final line]
\label{thm:line181}
Write \(M=\binom{d_*+nN'-1}{nN'-1}\). For every pair of
unit vectors \(u,v\), and every coefficient-one monomial
\(Z^\alpha\) of degree \(d_*\), set
\(E_\alpha=(\pi^{d_*}/\alpha !)^{1/2}Z^\alpha\), where
\(\alpha !=\prod_{i,j}\alpha_{ij}!\). Then
\[
  \sum_{|\alpha|=d_*}Z_v(E_\alpha;u)=\frac1{D_n},\qquad
 0\le Z_v(E_\alpha;u)\le\frac1{D_n}.
\]
For some monomial of this degree one therefore has
\begin{equation}
 \frac1{M D_n}
    \le \frac{\pi^{d_*}}{\alpha !} Z_v(Z^\alpha;u)
    \le\frac1{D_n}.                                      \label{eq:last181}
\end{equation}
Over all real places, tests can be chosen with the sum of
the absolute logarithms of these normalized local factors
bounded by a fixed power of \(x\), uniformly in the real
unit lines on the intervals of Section \ref{nav:I-28}.
\end{theorem}
\begin{proof}
For probability Haar measure, compact Schur orthogonality
identifies integration against
\(\overline{\langle\sigma(g)u,u\rangle}\) with \(1/D_n\)
times the orthogonal matrix unit \(e_u\) in the oscillator.
The two matrix units commute. On the polynomial space
their product projects onto the one line \(u\otimes v\)
in the unique joint summand from Lemma~\ref{lem:degree181}.
In particular
\[
 Z_v(P;u)=D_n^{-1}\|e_u e_v P\|^2.
\]
A monomial has squared norm \(\pi^{-|\alpha|}\alpha !\)
by angular integration and the one-variable Gaussian moment.
Distinct monomials are orthogonal. Hence the \(E_\alpha\)
at the displayed degree form an orthonormal basis of the
whole degree space. Summing the diagonals computes the
trace of a rank-one projection divided by \(D_n\).
Each diagonal is nonnegative and bounded by \(1/D_n\),
and their number is \(M\), proving the inequalities.

On the given intervals the ranks and \(f=[F:\QQ]\) are
polynomially bounded in \(x\), while the logs of two plus
the absolute weights are \(O_{\psi,\epsilon}(x^\epsilon+
\log(2x))\) by the row construction. The same bound holds
for \(\log(2+d_*+nN')\), by the sum formula above. Since
\(\log M\le nN'\log(2+d_*+nN')\), it is polynomial in \(x\).
Also \(D_n\le M\) because the degree space contains the
joint summand of dimension \(D_nD_{N'}\). Summation of
\(\log(MD_n)\) over the \(f\) places proves the last assertion.
\end{proof}

\begin{proposition}[Exact value for a product of leading minors]
\label{prop:minor181}
Take matching highest source and target lines in the
polynomial realization, write \(\lambda_{n+1}=0\), and let
\(\Delta_i(Z)\) denote the determinant of the leading
\(i\)-by-\(i\) submatrix. The polynomial with integer
coefficients
\[
       H_\lambda(Z)=\prod_{i=1}^n
                      \Delta_i(Z)^{\lambda_i-\lambda_{i+1}}
\]
has projected real norm
\begin{equation}
    Z_v(H_\lambda;u)
    =\frac{\pi^{-d_*}}{D_n^2}
       \prod_{i=1}^n\frac{(\lambda_i+n-i)!}{(n-i)!}.
                                                              \label{eq:minor181}
\end{equation}
For a positive real scale \(s\) relative to these coordinates
the right side is multiplied by \(s^{-d_*}\).
\end{proposition}
\begin{proof}
The highest vector in the tensor product of
\((\bigwedge^i\CC^n)^{\otimes(\lambda_i-\lambda_{i+1})}\)
is the product of the wedges \(e_1\wedge\cdots\wedge e_i\)
with those multiplicities. It is a unit for the induced
unitary structure and spans the highest line in its irreducible
Cartan component \(S_\lambda(\CC^n)\). On an \(n\)-by-\(n\)
matrix \(X\) its diagonal coefficient for \(S_\lambda(X)\)
is the product of the leading minors above. As a polynomial
on all \(N'\)-by-\(n\) matrices this is the joint highest
line, using the same wedges in the first \(n\) target
coordinates. This can also be seen by the unipotent
invariances of each leading minor. Consequently \(e_u e_v\)
fixes this polynomial with the matching lines. Its projected
local norm is its squared Gaussian norm divided by \(D_n\).

The Fischer squared norm of this leading-minor polynomial
is computed in Lemma~\ref{lem:hook86}. In the Gaussian used
there, Gram--Schmidt on the top square block separates the
unitary part, with squared coefficient mean \(1/D_n\),
from independent gamma variables of shapes \(n-i+1\).
Their \(\lambda_i\)-th moments give the displayed factorial
ratios. Passing from that Fischer norm to the present
vacuum measure multiplies the squared norm by \(\pi^{-d_*}\).
Taking the additional factor \(1/D_n\) from the source
matrix unit now gives the formula. The dependence on \(s\)
follows from Section~\ref{nav:VIII-08}.
\end{proof}

\begin{theorem}[Integral evaluation in the last compact norm]
\label{thm:integral181}
Keep the two highest lines of Proposition~\ref{prop:minor181}.
Let \(A\) be the integer ring of a number field \(L\) and let
\(Q\in A[Z]\) have column degrees \(\lambda_i\). Assume that it
is invariant under adding any multiple of an earlier column
to a later column, as a polynomial identity. Let \(J\) have
the identity matrix in its top \(n\) rows and zeros below.
For every embedding \(\iota:L\hookrightarrow\CC\) and
every positive scale \(s\), one has the exact comparison
\begin{equation}
   \frac{ Z_v^{(s)}(Q^\iota;u)}
        { Z_v^{(s)}(H_\lambda;u)}
      = |\iota(Q(J))|^2 .                            \label{eq:lattice181}
\end{equation}
In particular it vanishes precisely when \(Q(J)=0\). For
\(Q(J)\ne0\), taking the product over all embeddings gives
\[
 \sum_{\iota:L\hookrightarrow\CC}
 \log\frac{Z_v^{(s_\iota)}(Q^\iota;u)}
               {Z_v^{(s_\iota)}(H_\lambda;u)}
      = 2\sum_{\mathfrak l}
           \operatorname{ord}_{\mathfrak l}(Q(J))\log N\mathfrak l .
\]
The scales \(s_\iota>0\) may vary with the embedding.
The same equalities hold for any sum of these polynomials
with coefficients in \(A\).
\end{theorem}
\begin{proof}
In the Cauchy decomposition of Lemma~\ref{lem:degree181}
the column conditions select the highest source line in the
unique \(\lambda\)-summand. In this space the target weight
\((\lambda_1,\ldots,\lambda_n,0,\ldots,0)\), before adding
the vacuum character, is the one-dimensional highest
weight space. Its orthogonal projector keeps exactly those
row degrees, by integration over the diagonal compact torus.
The resulting polynomial is \(Q(J)H_\lambda\). Indeed the
row-degree projection does not affect the value at \(J\):
the only monomial of those column degrees that can be
nonzero there is \(\prod_i Z_{ii}^{\lambda_i}\). The
highest vector spans the line and takes the value one at
\(J\). This reproves in this compact pair the projection
identity of Theorem~\ref{thm:projection86}, now in the
local theta norm with its source factor. The identity
\(Z_v(P;u)=D_n^{-1}\|e_u e_vP\|^2\) gives
\eqref{eq:lattice181}. Both polynomials are homogeneous
of the same degree, so their dilation factors cancel at
each embedding separately. Finally \(Q(J)\) is an algebraic
integer. The product of its absolute values over all
embeddings, counting conjugate complex embeddings separately,
is the positive absolute norm of its principal ideal.
Factoring that ideal gives the stated sum. Linearity of
evaluation applies before taking the square of the absolute
value, including for sums with cancellations.
\end{proof}

\begin{corollary}[Logarithmic size on the common interval]
\label{cor:asym181}
Use highest-line polynomials at the real places and any
positive scales \(s_\tau\) relative to the reference coordinates.
For the projected local integrals of those coefficient-one
polynomials one has
\begin{align*}
 \sum_\tau\log Z_{v_\tau}^{(s_\tau)}(H_{\lambda_\tau};u_\tau)
     ={}&(k-1)\sum_\tau B_\tau
                     (\log B_\tau-1-\log\pi)\\
        &-\sum_\tau d_{*,\tau}\log s_\tau
            + O_{\psi,\epsilon}(x^C)
\end{align*}
for a fixed \(C=C(\psi,\epsilon)\). The error is
\(o(b/x^A)\) for every fixed \(A>0\), uniformly on the
intervals of Section \ref{nav:I-28}.
\end{corollary}
\begin{proof}
In the reference norm, take the logarithm of
\eqref{eq:minor181}. Replacing
\(\log((\lambda_i+n-i)!/(n-i)!)\) by
\(\lambda_i\log\lambda_i-\lambda_i\) costs
\(O(n\log(2+n+\lambda_1))\) per index, by the elementary
Stirling estimate and the at most \(n-i\) extra factors.
We use \(0\log0=0\). The Weyl dimension product
\(D_n=\prod_{i<j}(\lambda_i-\lambda_j+j-i)/(j-i)\)
has logarithm \(O(n^2\log(2+n+\lambda_1))\).
By the weight bounds in Theorem~\ref{thm:line181}, the sum
of all these errors over the embeddings is polynomial in \(x\).

The \(k-1\) centered entries \(v_i+e_m\) are each
\(B_\tau+O(x^{C_1})\), as also are their shifted partition
entries, for a fixed \(C_1\). All the other centered
entries are bounded by \(2mH_0+t\), and their partition
entries are bounded by a power of \(x\). These assertions
use only the lists in Lemma~\ref{lem:degree181} and the
polynomial bounds for ranks, spacings and \(v_i\).
The three permitted sizes of \(B_\tau\) all exceed every
fixed power of \(x\) on the chosen intervals, as in
the construction following equation \ref{eq:letters111}. On an interval
between \(B_\tau\) and any of the first \(k-1\) partition
entries the derivative of \(y\log y-y\) is
\(O(\log(2+B_\tau))\). Replacing each such entry by
\(B_\tau\) therefore costs only a power of \(x\).
The smaller entries, their number, and the number of
places are polynomially bounded. The same replacement
in \(-d_{*,\tau}\log\pi\) costs only a power of \(x\).
Keep the scale term from Proposition~\ref{prop:minor181}
exact, sum, and use \(b=\exp(x^\epsilon)\).
\end{proof}

In Theorem~\ref{thm:norm177} the last source line is fixed by
the previous tests. Theorem~\ref{thm:line181} therefore lets
one prescribe a unit target line separately at every real
place and retain a positive common constant after that
projection. Indeed choose the last monomial for those two
lines, apply the target matrix unit, and repeat the norm
identity in Section~\ref{nav:VIII-07} with the resulting projected local
factors. The explicit finite factors and the bound
\(\Delta_x\) there are the same. The reference real norm and
the factorial in \eqref{eq:last181} concern this last edge.
The full arithmetic projection also acts on the finite
classes, and its valuation and the normalizations on the
preceding edges remain to be estimated.

The two highest lines in Proposition~\ref{prop:minor181} may
also be imposed in a complex chain if the penultimate real
test can be chosen and projected. To see this, fix the unit
source vector on that preceding edge. By Lemma~\ref{lem:monomials177}
one real test gives a nonzero local contraction in the next
compact representation. Translates of its contraction by
the compact target span that irreducible representation.
Thus a translate has nonzero projection to the desired
highest line. Acting on the preceding test by the same
translation and the compact matrix unit realizes that
projection, by equivariance. The compact action preserves
the homogeneous polynomial spaces in the Fock realization.
The nonzero diagonal factor persists after this projection:
one may use a reference cuspidal constituent, with all other
diagonal factors nonzero, and project its nonzero pure
archimedean tensor line by this same matrix unit. Its new
norm is positive. The resulting fixed archimedean choices
work for the full row as in Section~\ref{nav:VIII-07} and allow the
coefficient-one polynomial of \eqref{eq:minor181} on the last
edge. The compact averaging in the preceding test introduces
a separate question about its arithmetic normalization.

\subsection{Integral projection from the last noncompact real group}\label{nav:VIII-10}
\label{sec:next185}

Fix one real place and take the preceding edge \(j=m-2\).
Put \(n=n_j\), \(p=n-1\), \(L=n+t=n_{j+1}\). The groups
on this edge are \(G=U(p,1)\) and \(H=U(L)\).
The source parameter \(P_j\) has \(p\) plus entries
\(p_1>\cdots>p_p\) followed, in the combined decreasing
order, by its sole minus entry \(p_-\).
Put \(a_T=-e_{m-1}+h_{m-2}\) and \(a_S=-e_{m-1}+h_{m-1}\).
Coordinates in the vacuum Fock norm are
\(Z\in M_{L,p}(\CC)\) and \(W\in\CC^L\), with independent
Gaussian entries of density \(e^{-\pi |z|^2}\) relative to
plane measure. Haar measure on \(H\) is probability.
Let \(d_\sigma\) denote the formal degree of the discrete
series \(\sigma\) on \(G\) for the Haar measure used there.

\begin{lemma}[The compact types on the penultimate edge]
\label{lem:types185}
The highest weight of the lowest \(U(p)\times U(1)\)-type
of \(\sigma\), after removal of the compact vacuum characters,
is \((\alpha_1,\ldots,\alpha_p;-\beta)\), where
\[
 \alpha_i=p_i-a_T-(p+L)/2+i,\qquad
 \beta=H_0-(t+1)/2>0 .
\]
The \(\alpha_i\) form a partition, with \(\alpha_p=\beta\).
The corresponding compact target after its vacuum twist is
removed has highest weight
\((\alpha_1,\ldots,\alpha_p,0,\ldots,0,-\beta)\).
The part of its oscillator isotypic space on the lowest
source compact type lies in bidegree \((|\alpha|,\beta)\)
in \((Z,W)\).
\end{lemma}
\begin{proof}
At this edge the centered lists in Section~\ref{sec:realtower169}
have \(p\) positive plus entries and precisely one negative
minus entry \(p_--a_T=-H_0\). The least positive entry is
\(e_{m-1}-b_{m-1}=H_0\). The old strings have centers at
least \(2H_0\) and half-width \((t-1)/2<H_0\). On this chamber
the half-sum of the positive noncompact roots is
\((1/2,\ldots,1/2;-p/2)\). The compact half-sum is
\(((p+1)/2-i)_{1\le i\le p}\) on the first block and zero
on the last. The lowest compact type of the discrete series
has weight \(P_j+\rho_{\rm nc}-\rho_{\rm c}\), and occurs
once. Remove the vacuum exponents \(a_T+L/2\) and \(a_T-L/2\)
on the respective blocks. The plus entries become the given
\(\alpha_i\), and the minus entry becomes
\(-H_0+(L-p)/2=-\beta\). The plus differences are
nonnegative integers by regularity and parity of the parameter.
The least is \(H_0-(L-p)/2=\beta\).

The compact target vacuum exponent is \(a_S+(p-1)/2\),
from the same compact splitting weights as in
Ichino~\cite[Lemma~7.8]{XX-XV-ichino328}.
For the target parameter of Section~\ref{nav:VIII-02}, subtract its compact
half-sum and this vacuum exponent. The old positive entries
give \(\alpha_i\). The inserted string of length \(t\) gives
zeros, and the last old entry gives \(-\beta\). In the
polynomial action of the compact groups, \(Z\) carries the
defining target weights and the defining \(U(p)\) weights.
The \(W\) block carries dual target weights and the negative
\(U(1)\) weight. Thus the two source compact centers fix
the two degrees of the lowest type, as claimed.
\end{proof}

Write \(H(\alpha)=\|H_\alpha\|_{\rm F}^2\) for the hook
product in Lemma~\ref{lem:hook86}. In the following test
\(H_\alpha(Z)\) uses the leading minors in the first \(p\)
target rows. Let \(e_L\) be the last coordinate vector,
\(J=\binom{1_p}0\), \(d_-=|\alpha|+\beta\), and set
\[
                     P_0(Z,W)=H_\alpha(Z)\,W_L^\beta .
\]
For a unit compact target vector \(v\) let \(e_v\) be its
compact matrix unit in the oscillator. Define \(Z_v(Q;u)\)
as in Section~\ref{nav:VIII-09} using this edge's groups and the integral
over \(G\).

\begin{theorem}[A projected local norm with noncompact source]
\label{thm:pen185}
Let \(u\) span the highest line of the lowest compact
source type, and \(v\) the compact target highest line,
with unit vectors on both. Suppose that the polynomial
\(Q(Z,W)\) has column degrees \(\alpha_i\) in \(Z\), is
invariant under upper unitriangular column operations in
\(Z\), and has degree \(\beta\) homogeneously in \(W\).
For a positive real dilation \(s\) relative to the
reference frame one has
\begin{equation}
 Z_v^{(s)}(Q;u)=
   |Q(J,e_L)|^2\,
   (\pi s)^{-d_-}\,\frac{H(\alpha)\,\beta!}{d_\sigma}.
                                                        \label{eq:pen185}
\end{equation}
The normalized ratio to the value on \(P_0\) is therefore
the squared absolute value of an integral evaluation whenever
the coefficients of \(Q\) are algebraic integers.
\end{theorem}
\begin{proof}
The compact member theorem for the oscillator realizes the
isotypic space of the indicated irreducible \(H\)-type as
the Hilbert tensor product of that type and its irreducible
unitary theta partner on \(G\). This is the compact-member
decomposition used in
\href{https://arxiv.org/pdf/2002.09148}{Ichino, Section 7.4,
in the proof of Lemma 7.7 for the compact factors}.
Its hypotheses apply to \(H=U(L)\). The nonzero pair already
identified in Theorem~\ref{thm:tower169} determines this
partner as \(\sigma\) by local Howe duality. On this tensor
product the action of \(G\) is on the first factor alone.
For a test projected by \(e_v\), write the resulting vector
as \(z\otimes v\). Square integrability of the matrix
coefficients of the discrete series and Cauchy--Schwarz
give convergence of the defining local integral. The
formal-degree orthogonality identity gives
\[
        Z_v(Q;u)=d_\sigma^{-1}|\langle z,u\rangle|^2 .
\]
This uses the invariant Hilbert norms, so any unitary
identification of the compact isotypic summand gives the
same value.

To determine the target projection, separate the target
diagonal torus weights before the vacuum character is added.
In a monomial let \(z_i\) be the total number of \(Z\)-entries
from row \(i\), and \(w_i\) the exponent of \(W_i\). The target
weight is \((z_i-w_i)_i\). For the required highest weight
its last entry is \(-\beta\), and \(\sum w_i=\beta\). This
forces \(w_L=\beta\), all the other \(w_i=0\), and \(z_L=0\).
The remaining weights force row degrees
\((\alpha_1,\ldots,\alpha_p,0,\ldots,0)\) in \(Z\).
Extracting these row weights preserves the column identities.
By Theorem~\ref{thm:projection86}, or by evaluating the
highest row component at \(J\), the extracted polynomial
is exactly \(Q(J,e_L)P_0\). The polynomial \(P_0\) is
highest on the target. Its two factors are highest vectors
in the Schur module and the symmetric power of the dual
defining representation, and their highest-weight tensor
lies in the Cartan component. Therefore the weight projection
just computed already lies in the specified irreducible
highest line on the target. The compact matrix unit \(e_v\)
kills the other weights and fixes this component.

On the source the polynomial \(P_0\) has the highest
weight of the lowest compact type from the lemma and is
annihilated by its compact raising operators, by the minor
and column identities. This type occurs with multiplicity
one in \(\sigma\). Thus the vector \(z\) for \(P_0\) lies
on the line of \(u\) and has norm \(\|P_0\|\).
Its two factors use independent Gaussian variables, so
\(\|P_0\|^2=\pi^{-d_-}H(\alpha)\beta!\) by Lemma~\ref{lem:hook86}
and the one-variable Gaussian moment. Apply the formal-degree
identity and then the dilation of Section~\ref{nav:VIII-08}. This proves
\eqref{eq:pen185}, including zero evaluations.
\end{proof}

\begin{corollary}[The last two integral evaluations]
\label{cor:two185}
Choose the compact highest vector of Theorem~\ref{thm:pen185}
as the input to the last edge. For a polynomial \(Q\) on
the preceding edge and \(Q'\) on the last edge satisfying
their respective column conditions, the normalized product
of the two projected local norms is
\[
              |Q(J,e_L)\, Q'(J')|^2 ,
\]
where \(J'\) is the top identity frame of Section~\ref{nav:VIII-09}. The
normalizing denominator uses \(P_0\) and \(H_\lambda\) on
the respective edges with the same local scales. If the
product of evaluations is nonzero and the coefficients
are in the integers of a number field \(D\), the sum of
the logarithms of these ratios over all its embeddings
equals
\[
       2\sum_{\mathfrak l}
          v_{\mathfrak l}\big(Q(J,e_L)Q'(J')\big)\log N\mathfrak l .
\]
\end{corollary}
\begin{proof}
The two intermediate lines agree because the compact
representation and its positive-root ordering are the
same on both sides. A change of its vacuum determinant
character shifts every weight equally and preserves its
highest line. Apply \eqref{eq:pen185} first, take the
intermediate vector to have unit norm, and apply
\eqref{eq:lattice181} next. The homogeneity factors cancel
edge by edge. For nonzero evaluations, factor the principal
ideal of the product exactly as in
Theorem~\ref{thm:integral181}.
\end{proof}

\begin{proposition}[A quantitative choice for arbitrary compact target line]
\label{prop:mono185}
Allow any unit \(u\) in the lowest compact source type and
any unit \(v\) in the compact target. With
\[
 M_-=\binom{|\alpha|+pL-1}{pL-1}\binom{\beta+L-1}{L-1},
\]
the sum of \(Z_v(P;u)\) over all reference unit monomials
of bidegree \((|\alpha|,\beta)\) equals \(1/d_\sigma\).
For at least one such monomial,
\[
      (M_-d_\sigma)^{-1}\le Z_v(P;u)\le d_\sigma^{-1}.
\]
Here \(\log M_-\) is bounded by a power of \(x\). For any
choice of real Haar measure, the logarithmic ratio of
\(d_\sigma\) to the formal degree for the regular parameter
\(\rho_i=(n+1)/2-i\) in the same chamber is also bounded by
a power of \(x\), uniformly after summation over real places.
In the Haar measure giving the maximal compact probability
in the integration formula and inducing
\((1-\|z\|^2)^{-(p+1)}d^{2p}z\) on the unit ball, that
reference degree is \(p!/\pi^p\).
\end{proposition}
\begin{proof}
Translate \(P_0\) by \((U(p)\times U(1))\times H\). Its
translates span the tensor product of the lowest source
compact type and the entire target representation, by
irreducibility of both finite-dimensional factors. They stay
in the stated bidegree. In the oscillator the vector
\(u\otimes v\) consequently lies in that bidegree. After
compact target projection, formal-degree orthogonality from
the preceding proof shows that \(d_\sigma Z_v(P;u)\) is
the squared projection of \(P\) to the single unit line
\(u\otimes v\). Summation in its homogeneous space is the
trace one, and the number of monomials is \(M_-\).
The size of this binomial product follows from the same
interval bounds on ranks and log weights as in
Theorem~\ref{thm:line181}.

In Harish-Chandra's formal degree formula for a discrete series
the Haar-dependent constant and the root-length factors are
the same for two parameters of this group and chamber. Their
ratio is thus \(\prod_{i<k}(P_i-P_k)/(k-i)\), where the
complete parameter is in decreasing order. The reference
parameter \(\rho\) is regular and integral after the usual
half-sum shift and defines a discrete series in this chamber.
Each ratio is at least one by the integral gaps. Taking
logarithms bounds their sum by
\(O(n^2\log(2+\max_i |P_i|+n))\). The family bounds and
the polynomial number of embeddings prove the estimate.
For the formal-degree product see also
\href{https://link.springer.com/article/10.1007/s12220-026-02573-z}{\emph{Wehrl
Inequalities for Matrix Coefficients of Holomorphic Discrete
Series}, Section 4.1, equations (12)--(13)}. A compact central
factor changes only the common Haar constant.

For the reference constant, realize the scalar Bergman
representation on the unit ball in \(\CC^p\) with ordinary
Lebesgue measure, or its complex conjugate to match the
chamber above. The group acts by change of variables with
the holomorphic Jacobian, so its action is unitary. The
constant functions have multiplicity one as a compact
type. Their translates are the Bergman kernels at all ball
points, up to nonzero scalars. These span a dense space by
the reproducing property. An invariant orthogonal projection
acts by zero or one on this compact line and then on all its
translates, proving irreducibility. In the conjugate model
the line has weight \((1,\ldots,1;-p)=2\rho_{\rm nc}\):
for \(\operatorname{diag}(A,d)\) the inverse ball map is
\(z\mapsto d A^{-1}z\), whose holomorphic Jacobian is
\(d^p/\det A\). This is the scalar lowest compact type of
the reference parameter \(\rho\).

For a unit constant vector the squared absolute diagonal
coefficient at \(g\) is \((1-\|g0\|^2)^{p+1}\).
In fact its coefficient is evaluation at zero of the
Jacobian multiplier, up to the normalization of the
constant, since the ball average of a holomorphic function
is its value at zero. For a block matrix
\(g=\left(\begin{smallmatrix}A&b\\c&d\end{smallmatrix}\right)\)
the modulus of the Jacobian there is \(|d|^{-(p+1)}\),
by the fractional linear ball action and \(|\det g|=1\).
The relation \(|d|^2-\|b\|^2=1\) gives
\(1-\|g0\|^2=|d|^{-2}\), with the same modulus for the
inverse. Integrating the squared coefficient in the stated
invariant measure leaves exactly the Euclidean volume
\(\pi^p/p!\) of the ball. In particular this representation
is square integrable, and its lowest type identifies it
with the reference discrete series. Formal-degree
orthogonality now gives \(d_\rho=p!/\pi^p\).
\end{proof}

\begin{corollary}[Uniform two-step main term]
\label{cor:twomain185}
For each place \(\tau\) let \(Z_{-,\tau}\) and \(Z_{+,\tau}\)
be the projected integrals of \(P_0\) and \(H_{\lambda_\tau}\),
respectively, with reference scales \(s_{-,\tau},s_{+,\tau}>0\).
Let \(d_{\rho,\tau}\) be the reference formal degree from
Proposition~\ref{prop:mono185}, in the same Haar measure
as the source integral on the first of these edges. Uniformly
on the intervals of Section \ref{nav:I-28},
\begin{align*}
 \sum_\tau \log (d_{\rho,\tau} Z_{-,\tau}Z_{+,\tau})
   ={}&2(k-1)\sum_\tau B_\tau(\log B_\tau-1-\log\pi)\\
      &-\sum_\tau\big(d_{-,\tau}\log s_{-,\tau}
                     +d_{*,\tau}\log s_{+,\tau}\big)
         +O_{\psi,\epsilon}(x^C).
\end{align*}
\end{corollary}
\begin{proof}
On the penultimate edge the only centered source entries
containing \(B_\tau\) are the \(k-1\) entries \(v_i+e_{m-1}\).
They are \(B_\tau+O(x^{C_1})\) for a fixed \(C_1\) and
eventually exceed all the remaining entries. The shifts
in Lemma~\ref{lem:types185}, as well as the other partition
entries and \(\beta\), are bounded by powers of \(x\).
In \(\log H(\alpha)\) the factorials give
\(\sum_i(\alpha_i\log\alpha_i-\alpha_i)\) with polynomial
error, and the Weyl differences in the denominator have
a polynomial total logarithm. This is the Stirling argument
of Corollary~\ref{cor:asym181}. The same argument replaces
the \(k-1\) large terms by those for \(B_\tau\).
The extra \(\log(\beta!)\) is polynomial. The ratio of
formal degrees is covered by Proposition~\ref{prop:mono185}.
Apply \eqref{eq:pen185}, keep its scale term exact, and
add Corollary~\ref{cor:asym181} for the last edge. All
errors remain polynomial after summation over real places.
\end{proof}

To insert these tests in a complex chain one may prescribe
the penultimate source line by compact projection at the
prior edge. For fixed nonzero prior input the images of all
polynomial tests under contraction span the target module,
as in Corollary~\ref{cor:vectors169}. On a reference cuspidal
constituent choose a prior test with nonzero norm. Its
translates in the real target span a dense invariant
subspace of that irreducible unitary representation, so
some translate has nonzero projection to the prescribed
line of the multiplicity-one lowest compact type. Realize
that projection by its compact matrix unit, as in Section~\ref{nav:VIII-09}.
Continuity and density permit approximation by polynomial
tests followed by this matrix unit. A term of some fixed
homogeneous degree still has nonzero projected norm by
linearity, and compact averaging preserves that degree.
This fixes the preceding data for the full row. The last
two choices then work with the previously determined finite
tests. To pass to the full arithmetic period one still
needs the earlier integral tensors, their finite class sums
and the normalization of the earlier noncompact pairings.

\subsection{Highest lines throughout the real theta correspondence}\label{nav:VIII-11}
\label{sec:harmonic187}

We determine a single polynomial with integer coefficients on every
real pair in Section~\ref{sec:realtower169}. In each intermediate
representation we use the highest line of its lowest compact type
for both the incoming and the outgoing map. We first determine the
least polynomial degree for the actual discrete series.

Fix \(0\le j<m\), a real embedding, and the splitting conventions of
Section~\ref{sec:realtower169}. Write \(n=n_j\), \(t=t_{\rm r}\),
\(h=(t+1)/2\), \(A=k+j(t+1)\), \(y=s_{j+1}\), \(z=s_j\).
The source signature is \((p,q)=(A+y,z)\) and the target signature
is \((P,Q)=(A+t+z,y)\). Use \(\ell_i\) for the entries of the
source parameter minus \(a_T\) in combined decreasing order.
The first \(A\) entries are positive, all of plus sign. Thereafter
the signs are alternately minus and plus. Define three lists by
\begin{align*}
 a_i&=\ell_i+i-A-h &&(1\le i\le A),\\
 b_v&=\ell_{A+2v}-h &&(1\le v\le y),\\
 c_u&=\ell_{A+2u-1}+h &&(1\le u\le z).
\end{align*}
Set \(\alpha=(a_1,\ldots,a_A)\),
\(\beta=(-b_y,\ldots,-b_1)\) and
\(\gamma=(-c_z,\ldots,-c_1)\). They are partitions, with the empty
partition permitted. In fact the integral gaps in \(\ell_i\) are
at least one, the last positive entry is at least \(H_0\), and the
negative entries are at most \(-H_0\), where \(H_0>h\).

\begin{lemma}[The two compact highest weights]
\label{lem:harmweights187}
After removing the Fock vacuum characters, the highest weights
of the lowest compact types in \(\sigma_j,\sigma_{j+1}\) are
respectively
\[
              (a,b\,;\,c),\qquad (a,0^t,c\,;\,b).
\]
The two compact types each have multiplicity one in their
respective discrete series.
\end{lemma}
\begin{proof}
For a discrete series with parameter \(\lambda\) the highest
weight of the lowest compact type is
\(\lambda+\rho_{\rm nc}-\rho_{\rm c}=
\lambda+\rho_{\rm tot}-2\rho_{\rm c}\), and that type has
multiplicity one. Here positive roots use the order of the
decreasing combined parameter. At combined position \(i\) in
rank \(n\) we have \(\rho_{{\rm tot},i}=(n+1)/2-i\). At
position \(v\) of a compact block of size \(d\),
\(2\rho_{{\rm c},i}=d+1-2v\). On the source subtract the vacuum
\(a_T\pm(P-Q)/2\), with the sign of the compact block.
The first \(A\) entries have \(v=i\). For the next entries the
positions within their blocks are \(u\) at \(A+2u-1\) and
\(A+v\) at \(A+2v\). Substitution, using
\(p-q+P-Q=2A+t\), gives \(a_i,c_u,b_v\), respectively.

On the target, Section~\ref{sec:realtower169} inserts the \(t\)
entries \(h-l\) (\(1\le l\le t\)) between the positive and
negative lists, and reverses the compact signs of the old negative
entries. Remove its vacuum \(a_S\pm(p-q)/2\). A middle entry at
combined and compact position \(A+l\) becomes
\(h-l+A+l-A-h=0\). An old negative entry is now at combined
position \(A+t+2u-1\) and plus position \(A+t+u\), or at
combined position \(A+t+2v\) and minus position \(v\). The
same half-sum formulas give \(c_u\) and \(b_v\) there. The old
positive entries still give \(a_i\). The Fock vacuum conventions
and the compact harmonic correspondence for these signs are in
Ichino~\cite[Lemma~7.8]{XX-XV-ichino328}.
\end{proof}

To describe the polynomial, a defining variable for a unitary block
has torus weight \(e_i\), while a dual variable has weight \(-e_i\).
In Fock coordinates write \(X,Y,W,V\) for the four matrices whose
types on the source and target compact factors are as follows.
Their row indices always belong to the target.
\[
\begin{array}{c|c|c|c}
 &\text{size}&\text{source factor and action}
                         &\text{target factor and action}\\ \hline
 X&P\times p& U(p)\ \text{defining}&U(P)\ \text{defining}\\
 Y&Q\times p& U(p)\ \text{dual}&U(Q)\ \text{dual}\\
 W&P\times q& U(q)\ \text{dual}&U(P)\ \text{dual}\\
 V&Q\times q& U(q)\ \text{defining}&U(Q)\ \text{defining}
\end{array}
\]
These are the four tensor products of compact defining or dual
spaces in the polynomial polarization, with the constant vacuum
weight removed. Let \(X_0\) be the first \(A\) rows and columns of
\(X\). Let \(Y_0\) use the last \(y\) rows and columns of \(Y\), and
\(W_0\) the last \(z\) rows and columns of \(W\). Order the indices
in each dual square in decreasing order on both axes.
For a partition \(\nu\) of length \(d\) write
\[
 H_\nu(M)=\prod_{l=1}^d
     \det(M_{[1,l],[1,l]})^{\nu_l-\nu_{l+1}},\qquad
 \nu_{d+1}=0.
\]
An empty square contributes one. Put
\begin{equation}
 \mathcal P_j=H_\alpha(X_0)H_\beta(Y_0)H_\gamma(W_0),\qquad
 D_j=|\alpha|+|\beta|+|\gamma|.                       \label{eq:poly187}
\end{equation}
Let
\(\mathcal J_j\) have the identity in each of the three indicated
squares and zero elsewhere, including all of \(V\).

\begin{theorem}[An explicit polynomial surviving every real quotient]
\label{thm:harmonic187}
The least Fock degree among the compact types of \(\sigma_j\)
on this pair is \(D_j\). In any nonzero local theta quotient
onto \(\sigma_j\boxtimes\sigma_{j+1}\), the image of the
integer polynomial \(\mathcal P_j\) spans the tensor product of the
highest lines of their lowest compact types. In particular
the successive contractions of the \(\mathcal P_j\)'s on these lines
are nonzero through the entire real chain.
\end{theorem}
\begin{proof}
Leading minors are highest vectors for the two standard block
actions. On a dual block its last coordinates in reverse order
play the same role, since a compact raising operator on a dual
basis sends a coordinate to a multiple of a later one.
Consequently \(\mathcal P_j\) is highest for both compact groups with the
weights in Lemma~\ref{lem:harmweights187}. This also follows by
taking the top exterior powers of the nested initial or terminal
flags. The polynomial has degree \(D_j\) and
\(\mathcal P_j(\mathcal J_j)=1\). For either compact group each variable
contributes just one signed coordinate to its torus weight.
Thus any monomial of a given vacuum-stripped weight has degree
at least the sum of the absolute values of its coordinates.
For each of our two compact types this sum is \(D_j\). In
particular both types have least degree \(D_j\) in the Fock space.

It remains to compare other compact types inside the source
discrete series. In its combined ordering let \(\delta_i=1\)
on the first \(A\) coordinates and \(\delta_i=-1\) on the rest.
This vector is dominant within each compact block. It pairs
nonnegatively with every positive noncompact root \(e_i-e_l\)
because \(i<l\). Let \(C\) denote the vector of the two constant
source vacuum weights, and \(\mu\) the lowest compact highest
weight before removing \(C\). Then
\(\langle\delta,\mu-C\rangle=D_j\).
If \(\nu\) is the highest weight of any compact type in the
source discrete series, the Blattner multiplicity formula
gives a compact Weyl permutation \(w\) and nonnegative integers
\(m_\xi\) with
\[
 w(\nu+\rho_{\rm c})
      =\mu+\rho_{\rm c}+\sum_{\xi>0,\ {\rm noncompact}}m_\xi\xi .
\]
Indeed at least one summand of the signed partition formula
must be nonzero for a type of positive multiplicity.
Since both \(\nu+\rho_{\rm c}\) and \(\delta\) are compact
dominant, sorting the coordinates within each block gives
\(\langle\delta,\nu+\rho_{\rm c}\rangle
 \ge\langle\delta,w(\nu+\rho_{\rm c})\rangle\).
It follows that
\[
                \|\nu-C\|_1\ge\langle\delta,\nu-C\rangle
                      \ge\langle\delta,\mu-C\rangle=D_j .
\]
The compact case is included, with no noncompact roots.
This proves the degree assertion for the discrete series.
For the multiplicity formula see Hecht and Schmid,
\href{https://doi.org/10.1007/BF01404112}{\emph{A proof of
Blattner's conjecture}}, and its \(U(p,q)\) convention in
\href{https://webusers.imj-prg.fr/~michele.vergne/publications/discreteseries.pdf}{Baldoni
and Vergne, Section 2.2}.

We explain why this polynomial survives the quotient. Use Howe's
harmonic generation theorem and his least-degree lemma for a
quotient of the polynomial oscillator. The space of a given
source compact type in the oscillator is generated, under the
target enveloping algebra, by its polynomials of least degree.
Since \(\mu\) occurs in the quotient, that degree-\(D_j\) subspace
has nonzero image. The least-degree lemma identifies \(D_j\)
also with the minimum degree of any target compact type in
the quotient. Thus a target compact type contributing to the
image of this subspace has least degree exactly \(D_j\). It
belongs there to the joint harmonics, since polynomials of
least degree for either compact type are harmonic for it.
In the joint harmonics a given source type has one target
partner and their compact tensor occurs once. The polynomial
\(\mathcal P_j\), of least degree for both its types, is a nonzero joint
harmonic highest vector. It therefore lies in this unique
compact tensor. The quotient map is nonzero there and injective
on it by irreducibility. Multiplicity one in the two discrete
series identifies the image of \(\mathcal P_j\) with the asserted lines.
The occurrence of the discrete-series quotient itself is
Theorem~\ref{thm:tower169}. Its successive compact highest
lines agree because the lowest compact type in an intermediate
discrete series and its positive compact roots do not depend
on the edge. These arguments use
\href{https://www.math.utah.edu/~ptrapa/math-library/howe/howe-transcending-1989.pdf}{Howe,
\emph{Transcending classical invariant theory}, (3.9)--(3.10)
and Lemma 4.1}, and the explicit harmonic correspondence of
Ichino, Lemma 7.8.
\end{proof}

\begin{theorem}[Integral evaluation on all the matched lines]
\label{thm:frame187}
Let \(F\) be homogeneous of degree \(D_j\) and highest for the
source compact group of the same weight as \(\mathcal P_j\). Project it
to the target compact highest line by its compact matrix unit.
The resulting polynomial is exactly
\[
                        F(\mathcal J_j)\mathcal P_j .
\]
In particular this matrix unit on the indicated source lattice
has integral image and coefficient one on \(\mathcal P_j\). The ratio of
the squared norms of the contracted, projected theta vectors for
\(F\) and \(\mathcal P_j\), using the same nonzero intertwiner, source
unit vector on the highest line, Haar measures and real scale,
is \(|F(\mathcal J_j)|^2\).
\end{theorem}
\begin{proof}
First take the target torus weight of \(\mathcal P_j\). On the source,
equality between its weight norm and the degree forces every
used coordinate to have the sign of its column weight. For a
zero weight the coordinate cannot occur at all. Otherwise
opposite signs in one column would increase the degree by at
least two. The identical argument on the target fixes the
allowed rows. Thus in the chosen target weight only the three
squares \(X_0,Y_0,W_0\) can occur. The row and column degrees
are \(\alpha,\beta,\gamma\) on the successive squares in their
specified orders. Source and target torus actions commute.
The extracted polynomial consequently retains the column
highest-vector identities for the independent changes within
each of the three column sets.

For a square and any partition \(\nu\), the polynomials of its
given highest column weight form one Schur module for its rows,
by the Cauchy decomposition of the polynomial ring. The row
weight equal to \(\nu\) is the one-dimensional highest weight
space in this module. Its generator \(H_\nu\) is one at the
identity, by the minor formula. Reversing the order in the two
dual squares gives exactly the same assertion there. Taking
the tensor product for the three independent squares shows
that the extracted polynomial is a scalar multiple of \(\mathcal P_j\).
Evaluation at \(\mathcal J_j\) reads this scalar from \(F\)
itself. Indeed each monomial nonzero at that frame has its
rows matched to its columns on each square and hence has
the target weight just extracted. The result is already highest
of the specified target compact type. Its matrix unit fixes
this vector and kills all the other target weights. The
projection formula follows. Applying the quotient of
Theorem~\ref{thm:harmonic187} and contracting on the source
highest line gives the norm ratio. A common dilation of the
coordinates on this space of homogeneous degree preserves
the ratio.
\end{proof}

\begin{corollary}[The product over a real chain]
\label{cor:product187}
For tests \(F_{j,\tau}\) as in Theorem~\ref{thm:frame187} on
the fields and intervals of Section \ref{nav:I-28}, the ratio
of the full products of these projected real squared norms,
relative to the polynomials \(\mathcal P_{j,\tau}\), equals
\[
          \left|\prod_{\tau}\prod_{j=0}^{m-1}
                         F_{j,\tau}(\mathcal J_j)\right|^2 .
\]
For a cuspidal row present in Theorem~\ref{thm:global175}, the
same ratio holds for the projected real factors in the global
Rallis norm at the prescribed highest lines, with a nonzero
reference factor on each edge. If all coefficients lie in
\(\mathcal O_L\) and the displayed product of evaluations is
\(u\ne0\), summation of the logarithm of the ratio over all
embeddings of \(L\), acting on these polynomial coefficients,
gives \(2\sum_{\mathfrak l} v_{\mathfrak l}(u)\log N\mathfrak l\).
\end{corollary}
\begin{proof}
Take unit vectors on the successive lines. On each line an
intertwiner contributes one scalar, with conjugation where the
Hermitian contraction requires it. Squared absolute values
multiply. This proves the first assertion by the two theorems.
For an occurring cuspidal representation use one nonzero global
lift from Section~\ref{nav:VIII-06}. Fix its other local arguments so that
its local bilinear map at the chosen real place is nonzero.
Equivariance and the local irreducible theta correspondence
identify this map on oscillator polynomials with the small
theta quotient followed by contraction, up to a nonzero
multiplicity vector in the other factors. Its application to
\(\mathcal P_j\) and the prescribed source vector is nonzero by
Theorem~\ref{thm:harmonic187}. The global norm identity with
the other arguments fixed gives a nonzero diagonal real
factor for this choice. The equality of the projected
polynomials in Theorem~\ref{thm:frame187} then gives the same
ratio in that integral. The argument applies at every step
by the successive nonzero cuspidal lifts of Section~\ref{nav:VIII-06}.
Finally the frames have entries zero or one, so \(u\) belongs
to \(\mathcal O_L\). The product of \(|\iota u|\) over all
embeddings \(\iota:L\hookrightarrow\CC\) is the absolute
norm of its principal ideal. Factoring this ideal gives the
asserted sum.
\end{proof}

\begin{proposition}[The reference Fock norms]\leavevmode
\label{prop:fock187}
In unit-vacuum Gaussian coordinates with dilation \(s>0\)
as in Section~\ref{nav:VIII-08} the squared norm of \(\mathcal P_j\) in the Fock
space is
\[
    \|\mathcal P_j\|_s^2=(\pi s)^{-D_j} H(\alpha)H(\beta)H(\gamma),
\]
where \(H(\nu)\) is the hook product of Lemma~\ref{lem:hook86}.
Summed over all edges and real places on the selected intervals,
\[
 \sum_{\tau,j}\log\|\mathcal P_{j,\tau}\|^2_{s_{j,\tau}}
   =m(k-1)\sum_\tau B_\tau(\log B_\tau-1-\log\pi)
        -\sum_{\tau,j}D_{j,\tau}\log s_{j,\tau}
        +O_{\psi,\epsilon}(x^C).
\]
Here \(B_\tau\) is the weight offset of Section~\ref{nav:VIII-09}.
\end{proposition}
\begin{proof}
The Gaussian entries in the three squares are disjoint. Their moments therefore multiply. Lemma~\ref{lem:hook86} gives their Fischer
norms. The homogeneous change of scale gives the displayed power.
Of the parts of \(\alpha\) exactly the original \(k-1\) entries
from the \(v_i\) are \(B_\tau+O(x^{C_1})\), by Section~\ref{nav:VIII-02} and
the bounds on the spacings. Every other entry of the three
partitions, and the number of entries, are bounded by a fixed
power of \(x\), uniformly in \(j,\tau\). The factorial expression
for each hook product in Lemma~\ref{lem:hook86} and Stirling's
inequality give an error polynomial in \(x\) after replacing
each large part by \(B_\tau\) in \(u\log u-u\). Indeed
\(\log(2+B_\tau)\) is polynomially bounded on the intervals,
and this also bounds each logarithm of a Weyl difference up
to a polynomial error. The number of summands \(m[E^+:\QQ]\)
is polynomially bounded. Keeping each dilation term exact gives
the formula.
\end{proof}

The factor converting this last Fock norm into the projected
theta norm on a pair with two noncompact groups is a separate
absolute normalization. Sections~\ref{nav:VIII-09} and~\ref{nav:VIII-10} calculate it for
the two terminal edges. The integer evaluations above determine
the dependence on a change of test throughout the chain. The
absolute normalized global period also involves the finite
class sum and the projection at the original finite level.

\subsection{An absolute upper bound for every real theta integral}\label{nav:VIII-12}
\label{sec:absolute189}

At every real place keep the parameters, highest lines and
unit-vacuum Fock norm of Section~\ref{sec:harmonic187}.
Write \(G=U(p,q)\), \(n=p+q=n_j\), \(K=U(p)\times U(q)\),
and use an arbitrary fixed Haar measure on \(G\). For a unit
vector \(u\) on the prescribed source line let \(d_\sigma\)
be the formal degree of its discrete series. The following
bound uses the actual integral, prior to cancellation against
a reference integral.

\begin{lemma}[Integrability of the matrix unit]
\label{lem:atom189}
Every \(K\)-finite matrix coefficient of each source
\(\sigma_j\) in the chain belongs to \(L^1(G)\). If
\[
       f_u(g)=d_\sigma\,\overline{\langle\sigma_j(g)u,u\rangle}
\]
with inner products linear in the first argument, then
\(f_u=f_u^*\) and \(f_u*f_u=f_u\) in \(L^1(G)\).
For any unitary representation \(\omega\) of \(G\) its
integrated operator \(\omega(f_u)\) is an orthogonal
projection.
\end{lemma}
\begin{proof}
The group has compact center. When \(q>0\), its derived
group is \(\mathrm{SU}(p,q)\), and the product with the
central circle covers \(G\) with finite kernel.
Absolute values of the coefficients are invariant under
the unitary central character. Thus one may apply the
semisimple integrability theorem on the derived group.
Restriction of this irreducible representation to the derived
group stays irreducible, since every element of \(G\) is the
product of a derived-group element and a central element.

For a discrete series with parameter \(\lambda\), a sufficient
criterion for integrability of its compact-finite coefficients
is that every noncompact root \(\xi\) satisfies
\[
     |(\lambda,\xi)|>\frac14\sum_{\zeta\ {\rm root}}
                                         |(\zeta,\xi)|.
\]
This is the \(L^1\) consequence of
\href{https://www.math.utah.edu/~ptrapa/math-library/milicic/milicic-asymptotics-duke-1977.pdf}{D.~Mili\v ci\'c,
\emph{Asymptotic behaviour of matrix coefficients of the
discrete series}, pp.~84--85}. Use the Euclidean pairing on
the trace-zero hyperplane. In type \(A_{n-1}\), for
\(\xi=e_i-e_l\), two roots pair with absolute value two,
and \(4(n-2)\) pair with absolute value one. The right
side is \(n-1\). Subtracting a common scalar from the
coordinates of the parameter preserves every pairing.

On our source, all the positive centered entries are plus,
and the first negative entry is at most \(-H_0\).
The absolute values on the positive side are at least
\(H_0\). The negative tail is
\(-H_0,-3H_0,-5H_0,\ldots\) and its signs alternate.
Hence any two entries of opposite compact sign are
separated by at least \(2H_0\). The noncompact roots are
exactly the differences between such entries. The chosen
spacing satisfies \(H_0\ge k+m(t_{\rm r}+3)>n_m\),
since \(n_m=k+m-1+mt_{\rm r}\). This verifies the strict
inequality for all \(n_j\). For \(q=0\) integrability
follows from compactness.

Unitarity and unimodularity give \(f_u^*(g)=
\overline{f_u(g^{-1})}=f_u(g)\).
Schur orthogonality of the square-integrable representation
with formal degree \(d_\sigma\) gives
\[
 \int_G\langle\sigma_j(h)u,u\rangle
         \langle\sigma_j(h^{-1}g)u,u\rangle\,dh
              =d_\sigma^{-1}\langle\sigma_j(g)u,u\rangle .
\]
The integral exists by Cauchy--Schwarz on its two
square-integrable coefficients. Conjugation and multiplication
by \(d_\sigma^2\) give the idempotent identity. An integrated
unitary representation is a star representation of \(L^1(G)\).
The image of a self-adjoint idempotent is an orthogonal
projection, including when its image is zero.
\end{proof}

\begin{theorem}[Absolute real Rallis bound in the polynomial norm]
\label{thm:bound189}
Let \(\phi\) be a Fock test, \(e\) a compact target matrix
unit projecting orthogonally onto the specified target
highest line in each copy of its compact type, and \(\omega\)
the oscillator. Write \(Z_e(\phi;u)\) for the local integral
of the coefficient of \(e\phi\) against the conjugate
source coefficient at \(u\). Then
\[
      d_\sigma Z_e(\phi;u)
          =\|\omega(f_u)e\phi\|^2\le\|e\phi\|^2 .
\]
In particular for every \(F\) of Theorem~\ref{thm:frame187},
at dilation \(s>0\),
\begin{equation}
 0\le Z_e^{(s)}(F;u)\le
   \frac{|F(\mathcal J_j)|^2}{d_\sigma}
       (\pi s)^{-D_j}H(\alpha)H(\beta)H(\gamma).
                                                       \label{eq:abs189}
\end{equation}
For \(F=\mathcal P_j\) the ratio to the upper bound
is the squared norm of the projection of
\(\mathcal P_j/\|\mathcal P_j\|\) by \(\omega(f_u)\).
It belongs to \((0,1]\) on a surviving cuspidal row, and
is one at the last two real edges.
\end{theorem}
\begin{proof}
The oscillator coefficient is bounded by \(\|e\phi\|^2\).
Lemma~\ref{lem:atom189} makes its integral against \(f_u\)
absolutely convergent. The definition of the integrated
operator gives the inner product
\(\langle\omega(f_u)e\phi,e\phi\rangle\), which is the
squared norm of its orthogonal projection. The identity
\(eF=F(\mathcal J_j)\mathcal P_j\) in
Theorem~\ref{thm:frame187} and the three Gaussian moments
in Proposition~\ref{prop:fock187} now give \eqref{eq:abs189}.
The homogeneous dilation multiplies both squared norms by
the same power \(s^{-D_j}\).
Nonvanishing on an occurring cuspidal row follows from
Corollary~\ref{cor:product187}. The exact evaluations of
Sections~\ref{nav:VIII-09} and~\ref{nav:VIII-10} give equality on their reference tests.
Equivalently in those two cases the compact target isotypic
summand is a tensor product with the source discrete series,
and the reference polynomial lies on the single source line.
\end{proof}

\begin{proposition}[A common Haar normalization in all signatures]
\label{prop:degree189}
Give \(K\) probability in the compact disintegration and
take on \(G/K\) the matrix ball
\(\mathcal D_{p,q}=\{Z\in M_{p,q}(\CC): I_q-Z^*Z>0\}\)
with quotient measure
\[
               \det(I_q-Z^*Z)^{-n}d^{2pq}Z .
\]
For any source parameter \(\lambda_1>\cdots>\lambda_n\)
in combined order, its formal degree in this measure is
\begin{equation}
 d_\sigma=
 \pi^{-pq}\prod_{l=1}^q\frac{\Gamma(p+l)}{\Gamma(l)}
       \prod_{i<k}\frac{\lambda_i-\lambda_k}{k-i}.
                                                      \label{eq:degrees189}
\end{equation}
In particular \(\sum_{\tau,j}|\log d_{\sigma_{j,\tau}}|\)
is bounded by a fixed power of \(x\) on the common intervals.
\end{proposition}
\begin{proof}
We compute the constant at the infinitesimal character
of \(\rho_i=(n+1)/2-i\) in the holomorphic chamber with
all plus entries first. For \(q>0\) consider the scalar
Bergman space of square-integrable holomorphic functions
on the ball with Lebesgue measure. The action by the inverse
fractional linear transformation with its holomorphic
Jacobian is unitary by change of variables. The constants
form one compact line, with multiplicity one as the central
circle rotating all the entries separates degree zero from
positive degrees. Evaluation at zero is a constant multiple
of integration, by that circle action. Translates of the unit
constant therefore represent evaluations at all points of
the ball up to nonzero scalars. Their span is dense, since
a function orthogonal to all of them vanishes at every point.
An invariant orthogonal projection acts on the constant line
by zero or one, hence has the same action on its dense span
of translates. This proves irreducibility.

Write \(g=\left(\begin{smallmatrix}A&B\\C&D\end{smallmatrix}\right)\).
At zero the fractional linear action
\(Z\mapsto(AZ+B)(CZ+D)^{-1}\) has differential
\(X\mapsto(A-BD^{-1}C)XD^{-1}\).
Its Jacobian is \((\det g)^q(\det D)^{-n}\).
Since \(|\det g|=1\) and \(D^*D-B^*B=I_q\), its modulus
squared is \(\det(I_q-(g0)^*(g0))^n\). The inverse matrix
has bottom right block \(D^*\), so its Jacobian at zero
has the same modulus. More generally
\[
 I_q-(gZ)^*(gZ)=(CZ+D)^{-*}(I_q-Z^*Z)(CZ+D)^{-1}.
\]
The Jacobian at \(Z\) has modulus \(|\det(CZ+D)|^{-n}\),
by the same differential calculation with
\(A-(gZ)C\) in the left factor. Indeed the Schur complement
of \(CZ+D\) in
\(g\left(\begin{smallmatrix}I&Z\\0&I\end{smallmatrix}\right)\)
is \(A-(gZ)C\). This proves invariance of the stated
quotient measure. The squared diagonal coefficient of
the unit constant is the squared Jacobian at zero by
evaluation. Its integral in quotient measure is therefore
the ordinary Euclidean volume of the ball.

For \(a\ge0\) integrate \(\det(I-Z^*Z)^a\) over the
ball with Lebesgue measure. Write \(Z=(X,v)\) and put
\(v=(I-XX^*)^{1/2}w\). The block determinant is
\(\det(I-X^*X)(1-\|w\|^2)\), by the Schur complement,
and the Jacobian on the last complex column is
\(\det(I-XX^*)=\det(I-X^*X)\). Thus integrating the
column multiplies the integral for \(q-1\) columns and
exponent \(a+1\) by
\(\pi^p\Gamma(a+1)/\Gamma(a+p+1)\). This last factor is
the radial beta integral on the unit ball in \(\CC^p\).
Starting with \(a=0\) gives for the volume
\(\pi^{pq}\prod_{l=1}^q\Gamma(l)/\Gamma(p+l)\).
The representation is square integrable. In its conjugate
model the constant has compact weight
\((q,\ldots,q\,;\,-p,\ldots,-p)=2\rho_{\rm nc}\),
as follows by taking the Jacobian of \(Z\mapsto A^{-1}ZD\)
for \((A,D)\in K\). It is the scalar lowest type of the
holomorphic or conjugate-holomorphic discrete series with
parameter \(\rho\) in the indicated chamber. Schur
orthogonality gives reciprocal ball volume as its degree.
For \(q=0\) the reference is the trivial representation
of a compact group of mass one.

Harish-Chandra's formal degree formula for a fixed compact
Cartan gives a Haar-dependent constant times the absolute
product of the positive-root pairings of the parameter.
The constant is independent of the discrete series chamber
for that Cartan. Its value is consequently fixed by the
reference just computed. Taking the ratio to \(\rho\) gives
the last product of \eqref{eq:degrees189}, with absolute
differences taken in combined decreasing order. This also
gives the usual Weyl dimension formula for \(q=0\).
For the general degree formula see Harish-Chandra,
\emph{Discrete series for semisimple Lie groups II},
Acta Math. \textbf{116} (1966), 1--111, and the proof by
\href{https://www.math.utah.edu/~ptrapa/math-library/schmid/atiyah-schmid.pdf}{Atiyah
and Schmid, \emph{A geometric construction of the discrete
series for semisimple Lie groups}}.

Finally the gaps in \(\lambda\) are integral and at least one,
so each Weyl ratio is at least one. Its logarithm is bounded
by \(O(\log(2+n+\max|\lambda_i|))\). The logarithms of the
gamma ratios cost \(O(n^2\log(2+n))\) in total. The ranks,
\(m\), the number of embeddings and the logarithms of the
largest weights are all polynomially bounded in \(x\)
on the intervals, giving the assertion.
\end{proof}

\begin{corollary}[An absolute bound along the full chain]
\label{cor:absall189}
Use the measures of Proposition~\ref{prop:degree189}. On a
surviving cuspidal row, put \(R_\infty^\circ=
\prod_{\tau,j}Z_e^{(s_{j,\tau})}(\mathcal P_{j,\tau};u_{j,\tau})\).
Uniformly on the fields and intervals of Section \ref{nav:I-28},
\begin{align*}
 \log R_\infty^\circ\le {}&
    m(k-1)\sum_\tau B_\tau(\log B_\tau-1-\log\pi)\\
    &-\sum_{\tau,j} D_{j,\tau}\log s_{j,\tau}
       +O_{\psi,\epsilon}(x^C).
\end{align*}
For other real Haar measures the same bound holds for the
sum of \(\log(d_{\sigma_{j,\tau}} Z_e^{(s_{j,\tau})})\).
For changed tests of the stated degree and highest source
weight one adds the exact logarithm in Corollary~\ref{cor:product187}
when the evaluation product is nonzero.
\end{corollary}
\begin{proof}
Multiply \eqref{eq:abs189} for the reference polynomials
and use Proposition~\ref{prop:fock187} to sum the Gaussian
logarithms. In ball measure the formal-degree logarithms
are absorbed by the uniform error just proved. The changed
tests use exactly the same orthogonal projections and Haar
measures. Their relative factor is therefore the one in
Theorem~\ref{thm:frame187}.
\end{proof}

The absolute upper bound permits independent local loss factors
\(\|\omega(f_u)\mathcal P_j\|^2/\|\mathcal P_j\|^2\).
A lower bound for their full product and control of the finite
class sum in the designated global period still require further
arguments.

\subsection{Totally positive generators in the comparison of real norms}\label{nav:VIII-23}
In Theorem~\ref{thm:common178}, the additive character is
scaled by \(\beta\in F^\times\) and the skew form by
\(\gamma\in F^\times\), with \(F=E^+\). Both are totally
positive. Fix the two split prime ideals of that theorem.
Their relations to the absolute different \(\mathfrak d_F\)
and the ideal \(\mathfrak a_0\) below \((\sqrt{-d_0})\) are
\[
  (\beta)\mathfrak d_F=\mathfrak p,\qquad
  (\gamma)\mathfrak a_0=\mathfrak q.
\]

\begin{theorem}[Individual real scales on the period fields]
\label{thm:units180}
Put \(f=[F:\QQ]\), take the given cyclotomic multiple \(N\)
with \(E\subset\QQ(\zeta_N)\), and write
\(B_F=(f-1)\varphi(N)\log(2N)\).
The generators can be chosen with
\[
 \max_\tau\left|\log\tau u-f^{-1}\log N_{F/\QQ}u\right|
       \le B_F\quad(u=\beta,\gamma),\qquad
 1+\max_\tau(|\log\tau\beta|+|\log\tau\gamma|)=b^{o(1)}
\]
on the same intervals as Theorem~\ref{thm:invar148}.
At the real places give the Fock vacuum norm one, and keep
the polynomial coordinates and source vectors fixed. Write
\(D_\tau\) for the sum of the chosen monomial degrees through
all edges at \(\tau\) and \(\bar D=f^{-1}\sum_\tau D_\tau\).
If \(R_\infty\) is the product of the absolute diagonal real
factors, then
\[
 \frac{R_\infty(\beta,\gamma)}{R_\infty(1,1)}
       =\prod_\tau\big(\tau(\beta\gamma)\big)^{-D_\tau},\qquad
 N_{F/\QQ}(\beta\gamma)=
       \frac{N\mathfrak p\,N\mathfrak q}{|D_F|\,d_0^{f/2}} .
\]
The logarithm of the real-factor ratio differs from \(-\bar D\log N_{F/\QQ}(\beta\gamma)\)
by at most \(2 B_F\sum_\tau|D_\tau-\bar D|\). The geometric
mean of the ratios over all simultaneous conjugates of the
two generators, with the real tests fixed, is exactly
\((N_{F/\QQ}(\beta\gamma))^{-\bar D}\).
\end{theorem}
\begin{proof}
For every nontrivial character of the totally real abelian
field \(F\), take its primitive even Dirichlet character
\(\chi\) of conductor \(h\mid N\). The squared absolute
values of \((1-\zeta_h^a)/(1-\zeta_h)\) are totally
positive units in the real cyclotomic field. Norming to
\(F\cap\QQ(\zeta_h)\) gives units whose log coordinates
are at most \(2\varphi(N)\log(2N)\) in absolute value.
Taking \(\chi(a)\ne1\) gives a nonzero coefficient in the
\(\chi\)-direction because the \(\chi\)-Fourier coefficient of
\(\log|1-\zeta_h^j|\) is a nonzero Gauss sum times
\(L(1,\chi)\). The character idempotents now show that
the log vectors of these units and their conjugates span
the sum-zero hyperplane. Rounding coordinates in an
independent subset gives the bound \(B_F\) on the distance
of the logs of a positive generator from their mean after
multiplication by a positive unit. Use this for both
generators. Their norms and the bound on \(\log N\) in
Section~\ref{sec:missing} give the uniform assertion.

At a real embedding set \(s=\tau(\beta\gamma)>0\). In the
vacuum normalization the dilation from the scaled Fock
space to the original one takes \(P(z)\) to
\(P(z/\sqrt s)\). It intertwines both unitary groups since
scalar dilation commutes with their actions and respects
the splittings. A homogeneous degree-\(d\) test therefore
contributes \(s^{-d}\) to the ratio of the diagonal
integrals. Multiply over the edges and embeddings. Subtract
the mean of the log vector of the product of the generators
to obtain the bound. Averaging over the Galois group makes
every fixed coordinate equal to this mean, giving the
geometric-mean assertion. Detailed proofs are in
Section \ref{nav:VIII-08}.
\end{proof}

In particular the factor from the generators for parallel
degrees is determined by their ideal norms. If the degrees
are nonparallel, fixed ideals alone allow arbitrarily large
and arbitrarily small positive real factors. Indeed there
is a positive unit whose logarithm has nonzero scalar product
with the degree vector, by the spanning argument. Multiplying
one generator by its integer powers gives both limits.
The calculation is in the fixed vacuum norm. The scale of
the original real integrals, the designated rational tensor
and the valuations of its projected period remain necessary
for the desired prime estimate.

\subsection{The last real projection and an integral coefficient}\label{nav:VIII-24}
At the last forward step both groups are compact at infinity.
Put \(n=n_{m-1}\), \(N'=n_m=n+t\). For a fixed real place
write \(p_1>\cdots>p_n\) for \(P_{m-1}\) in
Theorem~\ref{thm:tower170}, and let \(a_T=-e_m+h_{m-1}\).
In a reference orthonormal complex matrix frame use the
Fock inner product with density
\(e^{-\pi\operatorname{tr}(Z^*Z)}\) and probability Haar
measure on \(U(n)\) and \(U(N')\).

\begin{theorem}[A projected last theta norm in integral coordinates]
\label{thm:last182}
The source and target, after removal of the compact vacuum
characters, have the same partition
\[
       \lambda_i=p_i-(n+1)/2+i-a_T-N'/2 \quad(1\le i\le n)
\]
with zeros appended on the target. Write \(d=|\lambda|\) and
\(D_n=\dim S_\lambda(\CC^n)\). Fix real-place unit lines
\(u,v\) in these two complex representations, and write
\(Z_v(P;u)\) for the real local Rallis integral of a test
after projection to the target line. Among the reference
unit monomials of degree \(d\) there is one with
\[
   (M D_n)^{-1}\le Z_v(P;u)\le D_n^{-1},\qquad
             M=\binom{d+nN'-1}{nN'-1}.
\]
The sum of the logarithms of \(MD_n\) over the real
places is bounded by a power of \(x\) on the common
intervals. In particular the composed complex operator in
Theorem~\ref{thm:common178} can still detect every surviving
column on its real input line after fixing any one target
line at each real place and choosing the last tests accordingly.

Take now the matching highest lines and let
\(H_\lambda=\prod_i\Delta_i^{\lambda_i-\lambda_{i+1}}\)
in the leading minors, with \(\lambda_{n+1}=0\).
For a polynomial \(Q\) with algebraic integer coefficients,
column degrees \(\lambda_i\), and invariance under upper
unitriangular column operations, write \(J=\binom{1_n}0\).
At any coefficient embedding, with any scale \(s>0\) relative
to the reference real frame, the projected squared norm is
\begin{equation}
 Z_v^{(s)}(Q;u)
   = |Q(J)|^2\,
     \frac{(\pi s)^{-d}}{D_n^2}
            \prod_{i=1}^n\frac{(\lambda_i+n-i)!}{(n-i)!}.
                                                        \label{eq:last182}
\end{equation}
Thus for \(Q(J)\ne0\), the sum of the logarithms of the
ratios \(Z_v^{(s)}(Q;u)/Z_v^{(s)}(H_\lambda;u)\)
over all coefficient embeddings is exactly
\(2\sum_{\mathfrak l}\operatorname{ord}_{\mathfrak l}(Q(J))
\log N\mathfrak l\).
\end{theorem}
\begin{proof}
The compact vacuum shifts are \(a_T+N'/2\) on the source
and \(-e_m+h_m+n/2\) on the target. Subtraction from
the highest weights of Theorem~\ref{thm:tower170} gives the
stated partition. By Schur--Weyl duality the symmetric
algebra of the matrix coordinates contains exactly one
joint summand of this pair of types, in degree \(d\).
Schur orthogonality expresses \(Z_v\) as \(1/D_n\) times
the squared norm of the projection to \(u\otimes v\) in
that summand. The normalized degree-\(d\) monomials are an
orthonormal basis of size \(M\). The sum of their projected
squared norms is the trace one. Since \(D_n\le M\) and
\(\log M\le nN'\log(2+d+nN')\), the interval weight and rank
bounds give the logarithmic estimate. Apply it with the
incoming line from the preceding tests at each real place.
Projecting the test by the target matrix unit projects the
automorphic output by equivariance. This gives positive
local factors common to the full fixed row, as required
in the proof of Theorem~\ref{thm:common178}.

The column identities put \(Q\) in the source highest line.
Its component on the target highest line is
\(Q(J)H_\lambda\), since the row degrees of the diagonal
monomial are the highest weight and \(H_\lambda(J)=1\).
The Fischer squared norm of \(H_\lambda\) is
\(D_n^{-1}\prod_i(\lambda_i+n-i)!/(n-i)!\), by the Gaussian
Gram--Schmidt computation of Lemma \ref{lem:hook86}.
The reference vacuum introduces \(\pi^{-d}\), the source
matrix unit contributes the other \(1/D_n\), and
Theorem~\ref{thm:units180} gives the factor \(s^{-d}\).
The principal ideal of \(Q(J)\) now gives the final identity
by its norm. The complete local computation, with its
realization inside the complex chain, is in
Section \ref{nav:VIII-09}.
\end{proof}

For the coefficient-one highest vector the reference
calculation also gives a uniform asymptotic at the required
precision. Write \(B_\tau\) for the offset following
\eqref{eq:letters111} and \(d_\tau\) for the last degree there.
Then
\begin{align*}
  \sum_\tau \log Z_{v_\tau}^{(s_\tau)}
                 (H_{\lambda_\tau};u_\tau)
     ={}&(k-1)\sum_\tau B_\tau(\log B_\tau-1-\log\pi)\\
        &-\sum_\tau d_\tau\log s_\tau+O_{\psi,\eps}(x^C).
\end{align*}
Indeed the first \(k-1\) parts of \(\lambda_\tau\) differ
from \(B_\tau\) by a power of \(x\), and all other parts
are polynomially bounded in \(x\). Stirling's estimate in
\eqref{eq:last182} and the Weyl dimension formula give the
display, as proved in Corollary \ref{cor:asym181}. It evaluates
the final real projection for a specified polynomial lattice.
The global period requires the polynomial supplied by the
preceding lifts and their sum over finite classes to survive
in the designated arithmetic component, with the required
height bound. Those comparisons and the remaining strict
orders in \eqref{eq:tail114} still enter \eqref{eq:window}.

\subsection{Two consecutive integral evaluations}\label{nav:VIII-25}
One step earlier the source is \(U(p,1)\) and the target
is compact. Write \(n=n_{m-2}\), \(p=n-1\),
\(L=n+t_{\rm r}\) and \(a_T=-e_{m-1}+h_{m-2}\). At the
real place, let \(p_1>\cdots>p_p\) be the plus entries of
\(P_{m-2}\). Put
\[
 \alpha_i=p_i-a_T-(p+L)/2+i,\qquad
 \beta=H_0-(t_{\rm r}+1)/2,\qquad d_-=|\alpha|+\beta.
\]
Here \(\alpha\) is a partition and \(\alpha_p=\beta>0\).
Use coordinates \((Z,W)\in M_{L,p}(\CC)\times\CC^L\)
in the unit-vacuum Gaussian frame of
Theorem~\ref{thm:last182}. Set \(J=\binom{1_p}0\)
and let \(e_L\) be the last coordinate vector in \(\CC^L\).

\begin{theorem}[The two terminal real integrals]
\label{thm:two186}
Take the highest line of the lowest compact source type and
the highest line in the compact intermediate representation.
Suppose \(Q(Z,W)\) has column degrees \(\alpha_i\) in
\(Z\), is invariant under its upper unitriangular column
operations, and is homogeneous of degree \(\beta\) in \(W\).
For the indicated unit lines \(u,v\), the projected local
integral on this edge is
\[
      Z_v^{(s)}(Q;u)= |Q(J,e_L)|^2 (\pi s)^{-d_-}
                     H(\alpha)\beta!/d_\sigma .
\]
Here \(d_\sigma\) is the formal degree of the source discrete
series in the chosen real Haar measure,
\(H(\alpha)=\prod_i(\alpha_i+p-i)!/
\prod_{i<j}(\alpha_i-\alpha_j+j-i)\), and \(s>0\) is the
scale relative to the reference frame. On taking \(v\) as
the input to Theorem~\ref{thm:last182}, the ratio of the
two-edge product to the same product for the coefficient-one
highest polynomials equals
\[
                     |Q(J,e_L)\,Q'(J')|^2 ,
\]
where \(Q'\) is the last test and \(J'\) its identity frame.
For algebraic integer coefficients a nonzero product of
evaluations therefore gives the exact logarithmic norm
identity of Theorem~\ref{thm:last182} with this product
inside each valuation.
\end{theorem}
\begin{proof}
The sole negative centered entry of the source is
\(-H_0\), and its least positive entry is \(H_0\).
Adding the half-sum of noncompact positive roots and
subtracting the compact half-sum gives the lowest compact
type of this discrete series. Removing the vacuum
characters gives precisely \((\alpha;-\beta)\). The
compact target after its vacuum shift has highest weight
\((\alpha_1,\ldots,\alpha_p,0,\ldots,0,-\beta)\) by
Theorem~\ref{thm:tower170}. In a monomial its target
weights are the \(Z\)-row counts minus the \(W\)-exponents.
The last target weight \(-\beta\) forces all \(W\)-variables
to be \(W_L\). The others force the highest row degrees
in \(Z\). The surviving polynomial is consequently
\(Q(J,e_L)H_\alpha(Z) W_L^\beta\), by the highest-line
calculation in Theorem~\ref{thm:last182}. It lies on the
chosen source and target lines.

Since the target is compact, its isotypic summand of the
oscillator realizes the Hilbert tensor product of the
target type and its irreducible unitary source partner.
The source is a discrete series, so formal-degree
orthogonality evaluates the projected integral as the
squared norm of this line component divided by \(d_\sigma\).
The two polynomial factors use independent Gaussian
variables, with squared Fischer norms \(H(\alpha)\)
and \(\beta!\). The dilation and the original Gaussian
give \((\pi s)^{-d_-}\). The intermediate group has
the same highest line on both edges, so multiplication
by \eqref{eq:last182} proves the second identity.
Section \ref{nav:VIII-10}
also proves that a unit monomial can detect any lowest
compact input line and any compact target line on this edge.
\end{proof}

For the coefficient-one highest polynomials, the logarithms
on these two edges have the same leading contribution
\((k-1)\sum_\tau B_\tau(\log B_\tau-1-\log\pi)\) per edge,
after extracting the exact dilation terms and the reference
formal degree on \(U(p,1)\). That degree is \(p!/\pi^p\)
for probability compact integration and ball measure
\((1-\|z\|^2)^{-(p+1)}d^{2p}z\). The total error is bounded by
a fixed power of \(x\), by Corollary \ref{cor:twomain185}.
The arithmetic input reaching these two edges, its finite
class sums and its nonvanishing at the prescribed level
remain to be controlled. Theorem~\ref{thm:both184} bounds
the cost of the finite test places on the same intervals.
The remaining global strict orders are those in
\eqref{eq:tail114}.

\subsection{An integral polynomial on every real pair}\label{nav:VIII-26}
Fix \(j\), a real place \(\tau\), and use the forward
theta quotient of Theorem~\ref{thm:tower170}. Subtract the
target splitting exponent \(a_T\) from the ordered source
parameter. Its entries \(\ell_i\) consist of \(A=k+j(t_{\rm r}+1)\)
positive plus entries followed by alternating negative entries
of signs minus, plus, starting with minus. Put \(y=s_{j+1}\),
\(z=s_j\), \(h=(t_{\rm r}+1)/2\), and set
\[
 a_i=\ell_i+i-A-h,\quad b_v=\ell_{A+2v}-h,\quad
 c_u=\ell_{A+2u-1}+h .
\]
The three partitions used below are \(\alpha=a\), the reversed
list \(\beta=(-b_y,\ldots,-b_1)\), and the reversed list
\(\gamma=(-c_z,\ldots,-c_1)\). Write \(D=|\alpha|+|\beta|+|\gamma|\).

\begin{theorem}[The matched highest lines on the full real chain]
\label{thm:harmonic188}
At every step one can use the highest lines of the lowest
compact types of its two discrete series. In the unit-vacuum
polynomial model, a product \(\mathcal P_{j,\tau}\) of three
polynomials in minors, with integer coefficients, maps nontrivially
to these lines. Its degree \(D\) is the least Fock degree of
any compact type of the source discrete series on this pair.
For homogeneous degree-\(D\) tests \(F_{j,\tau}\) of the same
highest source compact weight, the product of the squared norms
of the projected real theta vectors, divided by the product
for \(\mathcal P_{j,\tau}\), is
\begin{equation}
          \left|\prod_{\tau}\prod_{j=0}^{m-1}
                      F_{j,\tau}(\mathcal J_{j,\tau})\right|^2 ,
                                                         \label{eq:all188}
\end{equation}
where the frames \(\mathcal J_{j,\tau}\) have entries zero and one.
This also gives the ratios of the projected real Rallis factors
on each surviving cuspidal row of Theorem~\ref{thm:global176}.
\end{theorem}
\begin{proof}
The stripped lowest compact weights are \((a,b\,;\,c)\) at
the source and \((a,0^{t_{\rm r}},c\,;\,b)\) at the target.
This follows by adding the half-sum of positive roots and
subtracting twice the compact half-sum at each slot, then
removing its compact vacuum character. In the defining
source-plus and target-plus matrix take the first square
of size \(A\). Take the terminal square of size \(y\)
in the dual source-plus and target-minus matrix and the
terminal square of size \(z\) in the dual source-minus and
target-plus matrix. Reverse both index orders in a dual
square. For a partition \(\nu\) the product of the leading
\(l\)-minors with exponents \(\nu_l-\nu_{l+1}\) is highest
for its two block actions. Use this product on the three
squares with partitions \(\alpha,\beta,\gamma\) to obtain
\(\mathcal P_{j,\tau}\). It is one on the frame with these
squares the identity and every other coordinate zero.

Here is the point that ensures survival when both groups are
noncompact. Put \(\delta_i=1\) on the positive source entries
and \(-1\) on the negative tail. It is compact dominant and
pairs nonnegatively with every positive noncompact root.
Let \(C\) be the vector of compact vacuum weights and let
\(\mu\) be the actual lowest compact highest weight. For any
compact highest weight \(\nu\) occurring in the source, the
Blattner formula places a compact Weyl translate of
\(\nu+\rho_{\rm c}\) in \(\mu+\rho_{\rm c}\) plus the cone
of positive noncompact roots. Pairing with \(\delta\) and
sorting within the compact blocks gives
\[
 \|\nu-C\|_1\ge\langle\delta,\nu-C\rangle
                 \ge\langle\delta,\mu-C\rangle=D .
\]
Each oscillator variable contributes just one signed weight on
the source, so every monomial of this weight has degree at
least that norm. Both chosen compact types attain degree \(D\).
Howe's least-degree lemma for the theta quotient now preserves
their joint harmonic copy. Concretely, the full source isotypic
space is generated from its least-degree part by the target
enveloping algebra. That part must survive somewhere. Every
target type in the quotient has degree at least \(D\) by the
same lemma, so its surviving part at degree \(D\) is joint
harmonic. There is exactly one compact tensor of joint harmonics
of the indicated source type. Its highest vector is
\(\mathcal P_{j,\tau}\).

For any test in the assertion, take the indicated target torus
weight. The absolute weights on both sides sum to the total
degree. This forces every entry used into the same three
squares, with the specified row and column degrees.
In each square the source highest identities and the highest
row weight pick out the one-dimensional line of the corresponding
minor polynomial by the Cauchy decomposition. Its coefficient
is evaluation at the identity on the square. Thus the target
highest projection gives precisely
\(F_{j,\tau}(\mathcal J_{j,\tau})\mathcal P_{j,\tau}\).
The intermediate highest lines agree on consecutive edges, so
contraction with unit vectors and multiplication gives
\eqref{eq:all188}. On an occurring cuspidal row the global norm
formula and the nonzero local quotient identify the same
relative factors. The degree calculation, compact weights,
and this application of the norm formula are proved in
Section \ref{nav:VIII-11}.
\end{proof}

For integral coefficients, a nonzero product of evaluations in
\eqref{eq:all188} has the principal-ideal norm formula of
Theorem~\ref{thm:two186} after summation over coefficient
embeddings. The Fock norm of its reference polynomial is
\((\pi s)^{-D}H(\alpha)H(\beta)H(\gamma)\) at scale \(s\).
The ratio in \eqref{eq:all188} cancels the absolute real theta
normalizations. Determining those on earlier noncompact pairs,
controlling the designated global coefficient at the original
finite level, and bounding the strict orders in
\eqref{eq:tail114} are still needed in the common-interval proof.

\subsection{An absolute estimate for the projected real integrals}\label{nav:VIII-27}
Let \(\sigma=\sigma_{j,\tau}\) and let \(u\) be a unit vector
on the source highest line in Theorem~\ref{thm:harmonic188}. Write
\(d_\sigma\) for its formal degree relative to the chosen real
Haar measure. We use \(Z_e^{(s)}\) for the actual real Rallis
integral after orthogonal projection to the target compact
highest line, at dilation \(s\). Normalize the Fock vacuum to
have norm one.

\begin{theorem}[Upper bound for the absolute real product]
\label{thm:abs190}
On every real pair, for a test \(F\) satisfying the weight and
degree hypotheses of Theorem~\ref{thm:harmonic188},
\[
 0\le d_\sigma Z_e^{(s)}(F;u)\le
 |F(\mathcal J_{j,\tau})|^2(\pi s)^{-D}
                      H(\alpha)H(\beta)H(\gamma).
\]
On a surviving cuspidal row use the reference polynomials
\(\mathcal P_{j,\tau}\) on all pairs and let \(R_\infty^\circ\)
be the product of their projected real integrals. Give the
compact factors probability in the disintegration of \(U(p,q)\)
and take the usual matrix ball with measure
\(\det(I-Z^*Z)^{-(p+q)}d^{2pq}Z\) on its quotient. Then
\(R_\infty^\circ>0\) and, on the same fields and intervals,
\begin{align*}
 \log R_\infty^\circ \le {}&
     m(k-1)\sum_\tau B_\tau(\log B_\tau-1-\log\pi)\\
 &-\sum_{\tau,j} D_{j,\tau}\log s_{j,\tau}
                       +O_{\psi,\eps}(x^C).
\end{align*}
\end{theorem}
\begin{proof}
For a noncompact source, two entries of opposite compact
sign differ by at least \(2H_0\). Indeed all positive
entries are plus, and the negative tail consists of odd
negative multiples of \(H_0\) with alternating signs.
The discrete-series integrability test of Mili\v ci\'c
requires that for every noncompact root \(\xi\),
\(|(\lambda,\xi)|>\tfrac14\sum_{\zeta}|(\zeta,\xi)|\),
as a sufficient condition. The sum is over the full root
system. In type \(A_{n-1}\) the right side is \(n-1\),
and the spacing \(H_0\) exceeds the largest rank \(n_m\).
Therefore the coefficient
\(f_u(g)=d_\sigma\overline{\langle\sigma(g)u,u\rangle}\)
is integrable. This also holds at a compact source.
Schur orthogonality gives \(f_u*f_u=f_u=f_u^*\).
Consequently its integrated oscillator operator is an
orthogonal projection and
\(d_\sigma Z_e(F;u)=\|\omega(f_u)eF\|^2\le\|eF\|^2\).
The exact formula for \(eF\) and its Gaussian norm in
Theorem~\ref{thm:harmonic188} give the local inequality.

In the specified ball measure the formal degree equals
\[
 \pi^{-pq}\prod_{l=1}^q\frac{\Gamma(p+l)}{\Gamma(l)}
       \prod_{i<l}\frac{\lambda_i-\lambda_l}{l-i},
 \qquad \lambda_1>\cdots>\lambda_{p+q}.
\]
To obtain the constant, integrate the squared coefficient
of the unit constant in the scalar Bergman discrete series.
Its Jacobian factor cancels the invariant density, leaving
ball volume \(\pi^{pq}\prod_{l=1}^q\Gamma(l)/\Gamma(p+l)\).
Harish-Chandra's formal-degree formula gives the Weyl
ratio to that reference parameter. All ranks and logarithms
of the largest weights are polynomially bounded in \(x\).
Thus \(\sum_{\tau,j}|\log d_\sigma|\) is polynomially
bounded as well. On every edge the reference consists of
three minor polynomials with Gaussian norms given above.
Exactly \(k-1\) parts are \(B_\tau+O(x^{C_1})\), while all
other parts are polynomially bounded. Stirling's inequality
in the hook products gives the claimed upper bound after
summing. Positivity on an occurring cuspidal row follows from
Theorem~\ref{thm:harmonic188}. The ball integral, all measure
constants and the root estimate are given in
Section \ref{nav:VIII-12}.
\end{proof}

For each reference test the ratio to its local upper bound is
the squared norm of the orthogonal projection of a unit vector.
The two final ratios are exactly one. The earlier ratios and
the finite class sum enter the lower arithmetic comparison for
the designated global period. The bound for the original
strict global classes in \eqref{eq:tail114} is also still
needed for the common interval.

\subsection{A compact conjugation average on a tempered column}\label{nav:III-57}
Let \(L\) be a nonarchimedean local field, \(O\) its
integer ring and \(P\) the size of its residue field.
Put \(G=\mathrm{GL}_k(L)\), \(C=\mathrm{GL}_k(O)\).
Use probability Haar measure on \(C\).

\begin{lemma}[A bound independent of the compact type]\leavevmode
\label{lem:conjugation292}
Let \(\pi\) be a tempered unitary representation of \(G\).
Let \(a\in G\) have elementary divisors
\((-d_-,0,\ldots,0,d_+)\) with \(d_\pm\ge0\), and put
\(D=d_-+d_+\). For every measurable \(|\phi|\le1\) on \(C\),
\[
  \left\|\int_C\phi(h)\pi(hah^{-1})\,dh\right\|
           \le k^D P^{-(k-1)D/2}.
\]
\end{lemma}
\begin{proof}
Let \(\lambda\) be the left regular representation and
\(e_C\) its orthogonal projection onto left \(C\)-invariants.
For \(u,v\in L^2(G)\), take absolute values inside the
integral of the regular matrix coefficient and set \(x=hy\)
in its integral over \(G\). Cauchy--Schwarz in \(h\)
bounds the answer by
\[
   \int_G U(a^{-1}y)V(y)\,dy,\qquad
        U(y)=\left(\int_C|u(hy)|^2\,dh\right)^{1/2},
\]
with \(V\) defined using \(v\). Both functions are
left \(C\)-invariant and have norms \(\|u\|,\|v\|\).
Thus \( \|e_C\lambda(a)e_C\|\) bounds the operator
in the regular representation. Let \(\omega\) be the
compactly supported measure obtained by pushing \(\phi(h)\,dh\)
forward under conjugation of \(a\). For vectors in \(\pi\)
fixed by a compact open subgroup \(J\), the matrix
coefficient of \(\pi(\omega)\) is unchanged on replacing
\(\omega\) by \(e_J*\omega*e_J\), with \(e_J\) normalized
Haar measure. This convolution has compact support and a
locally constant density. Temperedness gives weak
containment in \(\lambda\), hence the norm inequality on
that density. The two regular compact averages are
contractions. Smooth vectors are dense, which proves the
same bound for every pair of vectors in \(\pi\).

For the spherical regular bound let \(\Xi(g)\) be the
coefficient of the normalized \(C\)-fixed vector in
normalized induction of the trivial character of the
diagonal torus from the upper triangular Borel. This
unitary representation has a one-dimensional fixed space
by Iwasawa decomposition. In its compact model \(\Xi\)
is an integral of a strictly positive half-modulus, so
\(\Xi>0\). It is bi-\(C\)-invariant and
\(\Xi(a^{-1})=\Xi(a)\) by unitarity and reality. Inserting
the rank-one fixed-space projection between two matrices
gives \(\int_C\Xi(g h g')\,dh=\Xi(g)\Xi(g')\).
Consequently the positive convolution kernel
\(e_C\lambda(a)e_C\) and its transpose both send the
function \(\Xi\) on the discrete set \(C\backslash G\)
to \(\Xi(a)\Xi\). All cosets have the same Haar mass.
In Schur's test, apply the weighted Cauchy--Schwarz
inequality to each row with weights \(\Xi(y)/\Xi(x)\).
Summation by columns and the transpose identity then
give operator norm at most \(\Xi(a)\).

The spherical coefficient \(\Xi(a)\) is the value at
torus parameters all equal to one of the Satake transform
of \(1_{CaC}\), divided by \(\operatorname{vol}(CaC)\).
The coefficientwise bound in the proof of
Proposition~\ref{prop:cartan290}, using the \(D\)
one-step diagonals and their inverses, bounds this numerator
by \((kP^{(k-1)/2})^D\). Its denominator there is at
least \(P^{(k-1)D}\). Division proves the estimate.
\end{proof}

\subsection{A tempered estimate with an auxiliary Hilbert space}\label{nav:III-58}
Let \(G=\mathrm{GL}_k(L)\), \(C=\mathrm{GL}_k(O_L)\)
and \(P\) be as in Lemma~\ref{lem:conjugation292}. The
first-plane orbital integral at the terminal prime involves
an additional variable recording a covector. We shall use
the following form of the compact estimate.

\begin{lemma}[Compact invariance in a tensor product]\leavevmode
\label{lem:tensor293}
Let \(\pi\) be tempered on \(G\), let \(V\) be a Hilbert
space carrying a continuous unitary action \(h\mapsto R_h\)
of \(C\), and let \(S\) be a contraction on \(V\). If
\(u,v\in\mathcal H_\pi\otimes V\) are fixed by
\(\pi(h)\otimes R_h\) for all \(h\in C\), then
\[
 |\langle(\pi(A)\otimes S)v,u\rangle|
        \le \Xi(A)\|u\|\|v\| .
\]
Here \(\Xi\) is the positive spherical function of
Lemma~\ref{lem:conjugation292}. In particular if the only
nonzero exponents of \(A\) are \(-d,d\), allowing \(d=0\),
the right side is at most \(k^{2d}P^{-(k-1)d}\|u\|\|v\|\).
\end{lemma}
\begin{proof}
Average the coefficient under simultaneous conjugation
by \(\pi(h)\otimes R_h\). Its value is unchanged, and
the integrand operator becomes
\(\pi(hAh^{-1})\otimes D(h)\) with
\(D(h)=R_h S R_h^{-1}\). Each \(D(h)\) has norm at most one.
In the left regular representation on \(L^2(G;V)\), change
the integration variable \(x=hy\). For test functions
\(v_0,u_0\), the absolute matrix coefficient of the
average is bounded by
\[
  \int_G\int_C
   \|v_0(hA^{-1}y)\|_V\,\|u_0(hy)\|_V\,dh\,dy .
\]
Cauchy--Schwarz in \(h\) gives a regular matrix coefficient
between the two scalar left \(C\)-invariant functions
constructed in Lemma~\ref{lem:conjugation292}, with the
pointwise Hilbert norm in place of absolute value. Their
norms are the original norms. The regular bound in that
lemma is \(\Xi(A)\).

To transfer the estimate, first test on finite sums of
simple tensors whose first factors are smooth. Fix a
compact open \(J\) fixing these first factors and choose
an orthonormal basis of the finite-dimensional span \(V_0\)
of all their second factors. Compress \(D(h)\) to \(V_0\).
Each matrix entry of the average now involves a compact
measure on \(G\). Convolving on the left and right by
probability Haar measure on \(J\) preserves the test
coefficient and gives a compactly supported locally
constant density in each matrix entry. Since \(\pi\) is
tempered, its integrated action is a representation of
the reduced group \(C^*\)-algebra. On a finite matrix of
densities this representation is contractive in the
regular matrix norm, by contractivity of the induced
star representation of the matrix algebra. In the
regular model the compression to \(V_0\) and the two
\(J\)-averages are all contractions. The preceding
regular bound thus applies to the test coefficient.
Density gives the result for the averaged operator on the
whole tensor product. The explicit bound for \(\Xi(A)\)
is the last estimate in Lemma~\ref{lem:conjugation292}.
\end{proof}

\subsection{Absolute size of the last restriction on the period fields}\label{sec:absolutelower349}

Fix the parameters of \eqref{eq:letters111} on the fields of
Theorem~\ref{thm:markedstrict337}. Write
\(b=\exp(x^\eps)\), \(X'=\exp(x^{\eps/3})\) and \(d=t_{\rm r}+1\).
Use normalized Fischer forms with the highest-vector square equal
to one in both endpoint modules. At each embedding write
\(c_\tau,\kappa_\tau\) for the factors of the standard restriction
from Section~\ref{sec:lowerlattice347}. The word and translate
are implicit in these two numbers.

\begin{theorem}[An absolute arithmetic bound for the restriction scalar]\leavevmode
\label{thm:abslow349}
There is one positive integer \(D_x\) on each field \(E_x\),
valid for every word and translate at this \(x\), with
\[
            \log D_x\ll x^{C_0}(1+b/X') .
\]
For each such choice \(D_x\) is divisible by the absolute
numerator and the denominator in lowest terms of each number in
the list
\[
         \prod_\tau c_\tau,\qquad
         \prod_\tau \kappa_\tau,\qquad
         \prod_\tau c_\tau\kappa_\tau,\qquad
         \prod_\tau c_\tau^2\kappa_\tau .
\]
In particular \(\log D_x=o(b/x^A)\) for every fixed \(A>0\),
uniformly as the left endpoint of the interval tends to infinity.
\end{theorem}
\begin{proof}
We first locate the large powers in one embedding. Denote the
exponents in Section~\ref{sec:lowerlattice347} by
\(a_i=\mu_i-\lambda_{i+1}\) and \(b_i'=\lambda_i-\mu_i\).
Continue to write \(b_j,e_j\) for the parameters of the repeated
tail. Put \(N=k+md=n+1\). In the initial part of the lists of
Corollary~\ref{cor:normaldegree347} one has, for \(1\le i<k\),
\[
        a_i=v_i-u_{i+1}-1/2,\qquad b_i'=u_i-v_i-1/2 .
\]
At index \(i_j=k+(j-1)d\) the differences are
\[
 \begin{split}
 a_{i_j}&=e_j-b_j-d/2,\\
 b_{i_1}'&=u_k+b_1-1/2,\qquad
 b_{i_j}'=b_j-e_{j-1}-d/2\quad(j>1).
 \end{split}
\]
Every remaining difference vanishes, because the two entries
compared there belong to the same constant block. Thus outside
the three entries
\[
                  A=a_{k-1},\qquad B=b_1',\qquad C=b_k'
\]
all exponents are bounded by a fixed polynomial in \(x\),
uniformly in the word. Indeed \(a_i=\Delta c_{i+1}\) for
\(i\le k-2\), and \(b_i'=4\Delta-1-\Delta c_i\) for
\(2\le i<k\). Both tail parameter spacings equal \(H_0\).
The three exceptions are fixed across words at the given translate.
Since \(v_{k-1}=0\), at a high place
\(A=\lceil xb/X'\rceil+(k-2)/2\), and elsewhere \(A=0\).
At every place \(B,C\ll xb+x^{C_1}\). In particular it is
possible for \(C\) to have order \(b\) at an ordinary place.

Take the two products \(X_0,P_0\) of
Lemma~\ref{lem:threeminors349} with exactly these exponents.
The highest-vector calculation gives
\[
        \kappa_\tau=\frac{q(X_0)}{q(P_0)}
             \frac{q(X_w)}{q(X_0)}
             \frac{q(P_0)}{q(p_{l_w})}.
\]
For the middle quotient form a box as in
Theorem~\ref{thm:commonfischer348}. Its lower product is \(X_0\)
and its upper exponents are the maxima of the exponents of \(X_w\)
over all words. The three indicated large powers have already
been included, so the degree difference is polynomial in \(x\).
The logarithm of two plus the upper degree is also polynomial.
The theorem gives a common integer of polynomial logarithm
divisible by both reduced parts of the middle quotient for every
word. For the last quotient use a box with lower product \(P_0\).
The exponent of \(D_i\) in \(p_{l_w}\) is \(a_i+b_i'\).
After removing the three specified powers, the same degree bound
therefore holds. Lemma~\ref{lem:threeminors349} now shows that
both parts of every \(\kappa_\tau\) divide a single integer
of logarithm \(O(x^{C_2}(1+b/X'))\). At an ordinary place the
first quotient is exactly one and a polynomial logarithm suffices.

For \(c_\tau\), the length of the interval of permitted integers
at index \(i\), after omitting \(\mu_i\), is at most
\(d_0=\sum_i a_i\). This sum is at most \(A+x^{C_3}\).
All the factors in the product for \(c_\tau\) have absolute
value at most \(n+\lambda_1-\lambda_N\).
Their logarithms are \(O(\log(2+x^C+xb))\).
For a fixed reference word the logarithm of \(|c_\tau|\) is
therefore \(O(x^{C_4}(1+b/X'))\), by counting at most
\(n^2 d_0\) factors. For every other word, the proof of
Theorem~\ref{thm:wordlower348} bounds both parts of the ratio
of these integers by divisors of one integer of polynomial
logarithm, before combining it with the Fischer factors.
Multiplying that integer by the absolute reference value gives
one common multiple of all the \(|c_\tau|\).

Take the square of each common multiple for \(|c_\tau|\), multiply
by the integer for \(\kappa_\tau\), and multiply over the
embeddings and the translates. There are polynomially many of
these pairs by Proposition~\ref{prop:field12}. Addition of prime
valuations proves the four divisibilities, with a single fixed
exponent \(C_0\) after increasing it to include the number of
pairs. Finally both \(b\) and \(X'\) exceed every fixed power
of \(x\) for fixed \(\eps>0\). This proves the little-oh estimate.
\end{proof}

\begin{corollary}[The contribution to norms in every exterior degree]\leavevmode
\label{cor:singular349}
Remove the explicitly evaluated normal-line factor \(Z\) of
Corollary~\ref{cor:normaldegree347}. All nonzero singular values
of the remaining restriction have the same square
\(s_w=\prod_\tau c_\tau^2\kappa_\tau\). For every fixed \(A>0\),
\[
       \sup_w|\log s_w|+\sum_p\sup_w |v_p(s_w)|\log p
                         =o(b/x^A)
\]
with the suprema also over translates. The same bound applies
to the logarithm of each nonzero singular value squared in an
exterior power, divided by its positive exterior degree, and
to its finite valuations with this division. At \(g_1,g_2,T\)
these valuations vanish.
\end{corollary}
\begin{proof}
The tensor product identity in Theorem~\ref{thm:lowerlattice347}
gives the scalar on the product target. All nonzero singular
values of an exterior power are products of the indicated number
of nonzero singular values, by diagonalizing in unitary bases.
Both reduced parts of \(s_w\) divide \(D_x\).
Hence \(|\log s_w|\le\log D_x\) and
\(|v_p(s_w)|\le v_p(D_x)\). The prime sum is at most
\(\log D_x\), which proves the assertion by the theorem.
The marked primes satisfy the cutoff in
Theorem~\ref{thm:terminalformula347}.
\end{proof}

\begin{proposition}[The normal line in the integral frames at the marked primes]\leavevmode
\label{prop:normalframe349}
On each lifted definite class the last restriction, written in the
integral coefficient frames of the large and small lattices, is
the standard map of Theorem~\ref{thm:lowerlattice347}. On these
full coefficient lattices it is a split surjection at every
prime above \(g_1,g_2,T\), including the places where \(Z\) has
nonzero valuation in the original rational frames.
\end{proposition}
\begin{proof}
Write \(M\) for the small unimodular Hermitian lattice. In the
lift the lattice whose Hermitian dual is used to evaluate the
coefficients is \(M\oplus\pi\mathcal O_E e_N\), with the standard
unit normal vector \(e_N\). Indeed the terminal marking on the
normal coordinate has kernel \(\pi\mathcal O_E\). Reduction at
each prime dividing \(\pi\) tests that coordinate times
\(\bar\pi\), a unit there, and these primes have product
\((\pi)\) since \(T\) splits completely. Dualizing the orthogonal
sum gives \(M\oplus z^{-1}\mathcal O_E e_N\), where
\(z=\bar\pi\). Pairing with \(\pi e_N\) tests exactly
\(z\) times the last coordinate. This also describes the
lifting condition in
\srcloc{Development/AbelianVarietyTheory/Definitions/AutomorphicMembership/ExtUHInjective}{488}.

At a finite place take a basis matrix \(G'\) for \(M\).
The corresponding basis matrix of the lifted dual lattice is
\(G=\operatorname{diag}(G',z^{-1})\).
For a column of polynomial covectors, pullback into a frame
with basis columns \(G\) uses \(\rho_\lambda(G^{\mathsf t})\).
This follows first for linear covectors by evaluation on the
basis columns, then for products, minors and the separate
determinant lines by functoriality. Write \(L\) for the
standard restriction at one embedding and \(L_z=z^{-\delta}L\)
in the rational frames, suppressing the embedding on \(z,G,G'\).
Block equivariance and the normal degree of
Corollary~\ref{cor:normaldegree347} give
\[
          L\,\rho_\lambda(G^{\mathsf t})
                =\rho_\mu((G')^{\mathsf t})\,L_z .
\]
Explicitly, each monomial retained by extraction uses the bottom
row \(\sum b_i'\) times. The determinant power contributes
the additional exponent \(\lambda_N\). Their sum is
\(\delta=\sum_i\lambda_i-\sum_i\mu_i\). The remaining
\(n\)-block commutes with extraction and intertwines under
the smaller isotypic projection. This proves the displayed equality
with its specified sign. For the testing slot the change is
\(\rho_\mu((G')^{-1})\), preserving the same bilinear value,
since transpose matrices are adjoint in the Fischer pairing.
Tensor over all embeddings above the coefficient place.
In these frames the lattice of minor polynomials and the lattice
of integral coefficients agree at the three primes by
Theorem~\ref{thm:terminaldual325} and the hook bound.
The integral sections of Theorem~\ref{thm:lowerlattice347}
therefore give the claimed splitting on each coefficient space.
\end{proof}
The proposition describes the map at an individual class. The
signed theta column restricts the values that can occur jointly
at all the lifted classes. The absolute estimate for that
restricted table and its pairing with the incoming Bessel row
is still required in Section~\ref{nav:I-28}.


\subsection{Vector evaluations and the signed theta row in the Hilbert norm}
\label{sec:rankone354}

Use the residue sets \(G_P\) and evaluations \(E_G\) of
Theorem~\ref{thm:grid354}. In the initial Hilbert space a class
has unit mass, and the coefficient metric is the same on all
of its markings. Put \(Q=T^{fk}\). The integral source has
already been specified in Theorem~\ref{thm:pointlattice353}.

\begin{proposition}[The least squared singular value of simultaneous evaluation]\leavevmode
\label{prop:gridnorm354}
Let \(q_P\) be the largest number of members of \(G_P\) on
any one residue line. If \(q_P<T-1\) for every \(P\), then
\(E_G Q^{-1/2}\mathcal F_0\) from the separately zero-sum
Hilbert space onto the grid values is surjective. Its singular
values, counted on the target, are at most one, and the
smallest squared singular value equals
\[
             \prod_{P\mid\pi}\left(1-\frac{q_P}{T-1}\right).
\]
In particular this is at least \(1-f\sqrt T/(T-1)\)
for the same grids of size at most \(\sqrt T\) per place
as in Theorem~\ref{thm:grid354}.
\end{proposition}
\begin{proof}
The normalized transform is unitary onto the corresponding
zero-sum vector space by Theorem~\ref{thm:vectorchain353}.
Consequently the evaluation map times its adjoint is the
restriction to \(G\) of the orthogonal projection onto that
space, with both row and column indices in \(G\). At one place
the projection subtracts the mean on each scalar orbit.
If \(t\) members of \(G_P\) lie on one line, their principal
block is \(I_t-\mathbf1\mathbf1^*/(T-1)\). Its constant
line has eigenvalue \(1-t/(T-1)\) and its orthogonal
complement has eigenvalue one. Distinct lines have zero
cross terms. The full matrix is the tensor product of these
local matrices and the identity on the coefficient space.
Thus its smallest eigenvalue is the displayed product and
is positive. The last bound follows by multiplying the
elementary inequalities \(0\le q_P/(T-1)\le1\) and using
\(\prod(1-y_P)\ge1-\sum y_P\).
\end{proof}

We next use vector frequencies to measure the actual output.
Fix the base tuple and any finite Hilbert space \(Z\) of
projected tests at the small terminal classes. Use its dual
Hilbert space as the range, and write
\(\mathcal P:S^\circ\to Z^\vee\) for evaluation of the signed
theta column on those tests, in the arithmetic class frames of
Theorem~\ref{thm:covectorpaths352}. For initial class \(a\)
and final class \(b\), group the paths by the tuple \(C\) of
their composite residue matrices. Let
\(W_{ab}(C):D_a\otimes\mathbb C\to Z^\vee\) be the sum of
their coefficient maps with all the class masses and terminal
contractions of that theorem. This definition includes the
coupled choices on the last two edges. The two class frames
are held fixed. For each tuple of nonzero initial vectors
\(v\) set
\[
 \begin{split}
 H_a(v)&=\sum_{b,C}\left(\prod_{P\mid\pi}
        \psi_P\big(-f_{b,P}(C_Pv_P)\big)\right)W_{ab}(C),\\
 R_a(v)&=\prod_{P\mid\pi}(1-A_P)H_a(v).
 \end{split}
\]
Here \(A_P\) averages over multiplication of \(v_P\) by all
nonzero residue scalars, and \(f_{b,P}\) is the specified
terminal covector. In \(H_a\) one may sum over all paths before
discarding zero pullbacks. The projection to \(R_a\) makes
both conventions agree. For fixed \(b\) the character in this
formula is the matrix Fourier character attached at each
place to the rank-one map \(-v_P\otimes f_{b,P}\), since its
trace against \(C_P\) is \(-f_{b,P}(C_Pv_P)\).

\begin{theorem}[The Gram operator of the prescribed signed arithmetic column]\leavevmode
\label{thm:rankone354}
On the full ordinary initial table with separate sums zero,
\[
       \mathcal P\mathcal P^*
         = Q^{-1}\sum_a\sum_{v}R_a(v)R_a(v)^* .
\]
The vector sums are over tuples with all coordinates nonzero.
In particular, restricting the inner sums to any specified sets
of vectors gives a lower bound as positive Hermitian operators
on \(Z^\vee\). The formula uses only the centered matrix
Fourier coefficients on the indicated rank-one frequencies.
For one algebraic terminal test \(F\) there is also the
following inclusion in its actual fractional ideal of values.
For any such \(v\), any \(d_P\ne1\) in \(k_P^\times\) and
any \(y\in D_a\), the ideal contains
\[
    u^{-f} z_{\max}
      \sum_{B\subseteq S_T}(-1)^{|B|}H_a(v_B)(y)(F),
    \qquad (v_B)_P=
     \begin{cases} d_P v_P& P\in B,\\ v_P& P\notin B .\end{cases}
\]
All scalars and coefficient lattices in the integral assertion
are taken over the same enlarged coefficient field as in
Theorem~\ref{thm:pointlattice353}.
\end{theorem}
\begin{proof}
Insert Fourier inversion into each term of
Theorem~\ref{thm:covectorpaths352}. For \(g_a=\mathcal F_0\phi_a\)
the resulting vector in \(Z^\vee\) is
\(Q^{-1}\sum_{a,v} H_a(v) g_a(v)\). Each \(g_a\) has separate
sums zero. Its pairing with \(A_P H_a\) therefore vanishes,
by collecting all terms of the scalar orbit. Hence \(H_a\)
can be replaced by \(R_a\). A path for which \(f_{b,P}C_P=0\)
at one place has trivial character there, and its entire
contribution is killed by \(1-A_P\).

In the normalized vector coordinates \(Q^{-1/2}g_a\) the
operator has blocks \(Q^{-1/2}R_a(v)\). Each family
\(R_a(v)^*z\), for \(z\in Z^\vee\), is already in the
separately zero-sum space. Consequently the adjoint computed
in the full array takes its values in the actual complex
source. The unit mass at the initial class and unitarity of
the normalized transform give the displayed Gram operator
by matrix multiplication. Each summand is positive, which
proves the lower bound for arbitrary subsets.

For the integral assertion prescribe the entire vector table
at \(a\) to be \(Q u^{-f}z_{\max}y\) times the tensor product
of the differences at \(v_P,d_Pv_P\). Set all other class
tables to zero. Theorem~\ref{thm:pointlattice353} proves this
is the transform of an integral ordinary input. Its evaluation
by the inversion formula cancels \(Q\) and gives the asserted
alternating sum. For an actual complex lower estimate one
must still control the signed coefficient sums \(H_a(v)\)
on the prescribed tests and compare the resulting operator
with the analytic Bessel row of Section~\ref{sec:missing}.
\end{proof}

\section{Bounds for the first-plane integral}\label{nav:g-X-1}

\subsection{Localizing the real frame integral}\label{nav:I-57}
At a real place write \(G=U(S)(\mathbb R)\),
\(U=U_\ell(\mathbb C)\) for the upper unitriangular matrix
group, and \(X=X_*\) for the full frame in
Proposition~\ref{prop:flagunmarked264}. Write \(H\) for its
pointwise stabilizer and \(R=\{g:gX\in XU\}\).
Use \(\chi\) for the character of \(U\) on the first
superdiagonal. In the second position of the mixed oscillator,
\(\rho(p_*)\) denotes evaluation, leaving a Schwartz vector
on the core. The residual action of \(g\) after the first
polarization is denoted by \(\Omega(g)\), with its splitting
and change-of-measure factors. Thus on a product \(A(x)\rho\)
the first action is \(A(g^{-1}x)\Omega(g)\rho\).

\begin{lemma}[Real localization at any flag length]\leavevmode
\label{lem:realflag265}
There is a smooth local section \(s\) of \(G\to R\backslash G\)
through the identity. Given any smooth compactly supported
\(a\) on a sufficiently small section domain, there is a
smooth compactly supported function \(A\) on the full ambient
matrix space, supported among matrices of full rank, such that
\[
 \int_U\overline{\chi(u)}\,
       [\omega(s(y))(A\otimes\rho)](Xu,p_*)\,du
     =a(y)\,[\Omega(s(y))\rho](p_*) .
\]
The same integral is zero on cosets outside that domain.
The test \(A\) works for every residual Schwartz function
\(\rho\). One may take \(a\ge0\) of integral one on the real
quotient, with support in an arbitrarily small neighborhood
of the identity coset.
\end{lemma}
\begin{proof}
The Gram matrix of \(X\) is zero or has just its last diagonal
entry nonzero. Hence \(u^*X^*h_SXu=X^*h_SX\) for every \(u\in U\).
Full frames of this Gram matrix form a single \(G\)-orbit by
Witt extension, applied to the isometry defined by matching
the columns. The orbit \(\mathcal O\) is embedded in the open
set of full-rank matrices. Indeed, at a frame \(x\) there choose
a matrix \(Z\) with \(x^*h_SZ=1\). The differential of the Gram
map applied to \(ZD/2\) is \(D\) for any Hermitian matrix \(D\).
The regular-level-set theorem proves the assertion. The group
action identifies \(H\backslash G\) smoothly with this orbit.
Indeed the orbit map of a Lie group has constant rank, and its image
here is the entire embedded level set. In particular its
differential is onto, since a smooth constant-rank map of
smaller rank cannot cover this level set.

The action \(x\mapsto xu\) is free and proper. For the latter,
if \(x\) and \(xu\) lie in fixed compact subsets of the
full-rank set, a left inverse of \(x\) can be chosen with
uniformly bounded norm, for example \((x^*x)^{-1}x^*\) for a
positive definite auxiliary metric. Thus \(u\) is bounded
in the closed matrix group \(U\). Witt extension gives
surjectivity of the map \(R\to U\) defined by \(gX=Xu\).
Its kernel is \(H\).
The group \(R\) is closed because \(XU\) is closed in the
space of matrices, using a left inverse to \(X\).
The quotient by a closed Lie subgroup has local smooth
sections, obtained by exponentiating a complement to its Lie
algebra. On such a section domain, \((y,u)\mapsto s(y)^{-1}Xu\)
is a smooth parametrization of the corresponding open part of
\(\mathcal O\). Its inverse reads \(y\) from the class of the
frame modulo \(U\), then recovers \(u\) with the left inverse.

Choose a nonnegative compactly supported smooth function
\(\alpha\) on \(U\) of Haar integral one. On this part of
\(\mathcal O\) prescribe
\[
             A(s(y)^{-1}Xu)=a(y)\chi(u)\alpha(u)
\]
and extend by zero on the rest of the orbit. This is smooth
and has compact support in the full-rank locus. Cover the
support by finitely many submanifold charts. In each chart
multiply the prescribed tangential function by a transverse
bump equal to one on the orbit. A partition of unity and
extension by zero give a smooth compact ambient extension
whose restriction to the entire orbit is as prescribed.
Substitution in the integral proves the equality and the
support claim. A smooth positive quotient density permits
normalization of any nonzero nonnegative bump to integral one.
Taking products gives the assertion at all real places
simultaneously. The construction places no bound on derivatives
as the neighborhoods shrink.
\end{proof}

\subsection{The first vector on the common intervals}\label{nav:I-65}
Use the normalized unit-frequency section \(W_1\) corresponding
to the original word and real row. Write \(\mathbb E_1\) for
integration on \(B_\infty\backslash G_\infty\) against
\(\|W_1\|^2\) divided by its integral. With \(u'\) the
Euclidean unit vector of the first isotropic line, put
\[
 D=4s^2+3,\qquad
 {\cal P}_D=\mathbb E_1\prod_\tau
                    \max(1,\|g_\tau^{-1}u'\|)^{2D}.
\]
Thus the power in this moment is \(8s^2+6\).

\begin{theorem}[The first-vector estimate for the prescribed orbit]\leavevmode
\label{thm:energy285}
For \(O_1\) of Theorem~\ref{thm:first284}, uniformly in all
prescribed words and after the same terminal compression,
\[
  \|O_1\|_{\rm op}
    \le \exp\!\left(-(16-o(1))(r-1)f x^\varepsilon\right)
                          {\cal P}_D^{1/2}.
\]
The moment is finite at every \(x\) on the common intervals.
In particular a bound \(\log{\cal P}_D=O(fs^3x^\varepsilon)\)
there would give the same exponential decay for this orbit.
\end{theorem}
\begin{proof}
At one place let \(l_1,\ldots,l_s\) be the successive
Euclidean lengths of the unit Witt vectors under \(g^{-1}\),
orthogonalizing each against the preceding ones. Set
\(l_{s+1}=Z(g)\ge1\) and put
\(N=\max(1,\max_{1\le i\le s}l_{i+1}/l_i)\).
In the block representative of Lemma~\ref{lem:negative285}
these quantities are \(|\lambda_i|^{-1}\) and
\(\|m_0^{-1}e\|\). One has \(\|m_0\|\le2Z(g)\).
For instance, compact rotations of the positive coordinates
put \(m_0 e\) in one positive direction and the negative
line. On that plane a hyperbolic rotation with parameter
\(t\) sends \(e\) to this vector up to compact phases, with
\(\cosh t=a\) from that lemma. The remaining isometry fixes
\(e\) and is Euclidean unitary. Thus \(\|m_0\|=e^{t}\le2a\)
for \(t\ge0\). The diagonal-block formula now gives
\[
 \nu(g)\le2\max(1,l_1)N^s .
\]
Indeed \(l_i\le l_1N^{i-1}\) for all \(i\), including
\(i=s+1\), while \(1/l_i\le N^{s-i+1}\) for \(i\le s\).

The flag vectors of Theorem~\ref{thm:energy153} have length
\(\sqrt2\) before translation. At unit frequency its local
maximum \(M_\tau\) therefore satisfies
\(N_\tau\le\max(1,\sqrt2 M_\tau)\).
Take \(j=sD\) in that theorem. Its derivative argument
applies with terminal vector \(e\), since it computes the
simple coefficient using exactly the vector of the chosen
frequency. The isometry for frequency one is
Corollary~\ref{cor:rel155}. On the local lowest type,
integration of \(N_\tau^{2j}\) is consequently a positive
quadratic form bounded by
\(\{C r^2(L_*+j+1)\}^{2j}\) times the norm squared, with
an absolute enlarged \(C\). The same estimate holds for
the product of these weights with this bound raised to \(f\).
In fact the restricted forms on the finite lowest types
are positive operators, and their tensor product has norm
at most the product of the norms. This covers arbitrary
vectors of the product lowest type.

Apply Cauchy--Schwarz with the probability measure
\(\mathbb E_1\) to the products of the first lengths and
of the ratios. It yields
\[
 {\cal M}_D(W_1)\le
     2^{Df}\,{\cal P}_D^{1/2}
                   \{C r^2(L_*+sD+1)\}^{sDf}.
\]
Here \(L_*\ll\sqrt r(2\exp(x^\varepsilon)+r)\) by the
lowest-type bound in Theorem~\ref{thm:localmass156}.
The logarithmic cost outside \({\cal P}_D^{1/2}\) is
\(O(fs^3x^\varepsilon)\). Also \(l_1\le\nu(g)\), so the
moment is finite by Proposition~\ref{prop:weight163}.

Insert the last estimate into Theorem~\ref{thm:low285}
and multiply by the original finite mass from
Theorem~\ref{thm:product281}. The numerical prefactor in
Theorem~\ref{thm:first284} has \(A_\sharp\) in place of
\(A_*\). Its logarithm still has the stated upper limit,
since \(\log A_\sharp\le O(r+\log(2s))\) from its formula.
The extra cost just obtained is \(o((r-1)f x^\varepsilon)\)
because \(m^5x^\varepsilon=o(k)\) on this family.
This proves the result.
\end{proof}

The source formulates the positive first-length moment in
\srcloc{Formalism/SharpRestriction/CuspSpanNormBounds/BesselShellMomentBounds}{1634}.
The calculation here gives the other length estimates needed
for this orbit on the original real row. Theorem~\ref{thm:moment286} below proves the uniform bound
for this moment. The full projected period and the remaining
strict cohomology sum in Section~\ref{sec:missing} are open.

\subsection{Exponential decay for the first rational orbits}\label{nav:I-66}
Use the fields and intervals of Section~\ref{sec:missing} and the original
word, finite level and terminal frequency. Write
\(b=\exp(x^\varepsilon)\), \(f=[E^+:\QQ]\),
\(r=k+2s+1\) and \(D=4s^2+3\). For the unit-frequency section
\(W_1\), let \({\cal P}_D\) be the normalized first-vector
moment in Theorem~\ref{thm:energy285}. Let \(V=\prod_\tau\nu_\tau\)
be the quotient height from Section~\ref{sec:absolute163}, and
let \(\mathbb E_1\) denote the probability with density proportional
to \(\|W_1\|^2\). Set
\(c_0=g_1g_2\sqrt{-q}\), and let \(O_j\) denote the
normalized orbital Gram operator for \(\gamma_{j c_0}\).
In the sum below \(j\) is a positive squarefree integer
prime to \(2q d_0 g_1g_2T\), all of whose prime factors split
completely in \(E\). The integer one is included.

\begin{theorem}[The unit-frequency moment and the first rational orbits]\leavevmode
\label{thm:moment286}
Uniformly over the prescribed words, for every integer
\(1\le A\le4r^2\) one has \(\log\mathbb E_1 V^A\le CfsA\log b\).
For the prescribed first-vector moment and orbit one has
\[
 \log{\cal P}_D=O(fs^3\log b),\qquad
 \|O_1\|_{\rm op}\le
       \exp\bigl(-(16-o(1))(r-1)f\log b\bigr).
\]
For every fixed \(0<\alpha<16\), on the same intervals,
\[
 \sum_{j\le b^\alpha}\|O_j\|_{\rm op}
       \le \exp\bigl(-(16-\alpha-o(1))(r-1)f\log b\bigr)
\]
with the specified restriction on \(j\). Both operator estimates
hold after compression to the original terminal character and
the retained coordinates.
\end{theorem}
\begin{proof}
We prove the moment estimate at a single real place first. Put
\(l(g)=\|g^{-1}U_1\|\), where \(U_1\) is the first Euclidean
unit Witt vector. Start at the prescribed row of parameter
\(\lambda\) with lowest compact type \(U\). Follow any sequence
of \(a\) lowerings with nonzero outgoing Bessel maps as in
Proposition~\ref{prop:transfer286}, reaching parameter \(\mu\).
For each \(q'\ge0\) with \(a+q'\le4r^2\) we prove
\[
 d_\mu\int_{B\backslash G}l(g)^{2q'}
       \|\Lambda_\mu(\sigma_\mu(g)w)\|^2\,d\dot g
        \le \exp(Csq'\log b)\,\|w\|^2,
\]
for every vector \(w\) whose compact types occur in
\((F^*)^{\otimes a}\otimes U\). These vectors are compact-finite
by admissibility. At \(q'=0\) this is the isometry in
Theorem~\ref{thm:rel155}. The shifted model exists at every
retained step and is normalized by the same Haar measures.

To pass from \(q'\) to \(q'+1\), take a Euclidean orthonormal
dual basis \((\zeta_d)_{d=1}^r\) in the Witt frame.
For each \(d\), tensor equivariance gives the vector identity
\[
 \zeta_d(g^{-1}U_1)\Lambda_\mu(\sigma_\mu(g)w)=
       \sum_{i}\beta_i\Lambda_{\mu-e_i}\!
       \left(\sigma_{\mu-e_i}(g)\,p_i(\zeta_d\otimes w)\right),
\]
where the sum includes only the nonzero maps. Indeed the left
side applies \(\Lambda_\mu{\rm ev}\) to \(g(\zeta_d\otimes w)\).
Insert \(1=\sum_i\iota_i p_i\), commute \(g\) through
the inclusions and apply Proposition~\ref{prop:transfer286}.
The squared absolute values of the scalar factors on the left
sum over \(d\) to \(l(g)^2\). Bound the squared norm of
each right side by \(r\) times the sum of the squared norms
of its terms. Multiply by \(d_\mu l(g)^{2q'}\), sum over
\(d\), and integrate. The induction hypothesis applies to
each projected vector by Lemma~\ref{lem:signs286}. The
formal-degree and scalar factor is bounded by
\((2\pi)^{2s}(1+2L)^2\) as in
Proposition~\ref{prop:transfer286}. For each \(d\), the
sum of the remaining squared norms is at most
\(\exp(C_1s\log b)\|w\|^2\) by Lemma~\ref{lem:signs286}.
The whole step costs at most
\[
        r^2(2\pi)^{2s}(1+2L)^2\exp(C_1s\log b)
                         \le\exp(Cs\log b).
\]
This proves the induction with a single \(C\), independent
of the path or \(q'\). In particular it bounds the moment
of every even order up to \(8r^2\) on every initial compact vector.

The positive form for \(\max(1,l)^{2A}\) on \(U\) has operator
norm at most \(1+\exp(CsA\log b)\) in the normalized model,
since \(\max(1,l)^{2A}\le1+l^{2A}\).
The operator for the product of the local weights on the
product compact type is their tensor product, as is seen by
integrating pairs of simple tensors and using polarization.
Taking the \(f\)-fold product of the bounds yields
a bound \(\exp(CfsA\log b)\) for the product first-vector
moment of order \(2A\), including arbitrary vectors of the
product compact type. In particular \(\log{\cal P}_D=O(fs^3\log b)\).
For the height, Theorem~\ref{thm:energy285} gives pointwise
\(\nu_\tau\le2\max(1,l_\tau)N_\tau^s\), where \(N_\tau\)
is the maximum of one and the consecutive flag-length ratios.
By the local derivative estimate of Theorem~\ref{thm:energy153}
and the tensor-product argument of Theorem~\ref{thm:energy285},
\[
 \mathbb E_1\prod_\tau N_\tau^{2sA}
                 \le \{Cr^2(L_*+sA+1)\}^{2sAf}.
\]
Since \(sA\) is polynomial in \(x\), Cauchy--Schwarz applied
to the two product weights gives the claimed bound for
\(\mathbb E_1 V^A\). Theorem~\ref{thm:energy285}
now yields the first operator estimate because
\(s^3=o(r-1)\) on these intervals.

For the sum, use the positive majorant which gave
Theorem~\ref{thm:energy285}. Its finite factor at \(j=1\)
is the unsigned finite mass and its real factor bounds the
integral of the product of the two section norms.
By Theorem~\ref{thm:low285}, replacing \(c_0\) by \(j c_0\)
in the real majorant multiplies the pure-power factor by
\(j^{-(r-1)f}\), and every positive logarithmic factor
weakly decreases. The moment just proved is independent of
\(j\). As calculated in the proof of Theorem~\ref{thm:tail283},
the unsigned finite mass at all added split primes increases,
relative to its value at \(j=1\), by at most
\(2j^{(2k+2s-2)f}\). That proof also bounds finite translations
and orthogonal compressions by the same majorant for the
signed operator. Consequently
\[
 \|O_j\|_{\rm op}\le
   2\,j^{(k-2)f}\exp\bigl(-(16-o(1))(r-1)f\log b\bigr)
\]
uniformly. There are at most \(b^\alpha\) terms, and
\((k-2)/(r-1)\to1\) while \((r-1)f\to\infty\).
This gives the claimed sum.
\end{proof}

Theorem~\ref{thm:sum287} sums this family over all its split
squarefree multipliers. Other rational orbits, the arithmetic
comparison of the integral period, and the uniform strict
cohomology bound on these fields still require proof.

\subsection{Two bounds for the original real orbital integral}\label{nav:I-67}
Let \(R=r-1\), \(c_0=Q\sqrt{-q}\) and \(U_0=Q\sqrt q\).
For an integer \(j\ge1\) write \(I_j^+\) for the real
orbital integral of the product of section norms for
\(\gamma_{j c_0}\), divided by its diagonal real norm,
in the original terminal frequency and word. Use the
unsigned mass \(m_f(c_0)\) at the single base parameter
as a common normalization for all these real quantities.

\begin{lemma}[Two powers of the rational multiplier]\leavevmode
\label{lem:real287}
There are positive numbers \(A_0,A_1\), independent of
\(j\), such that
\[
 m_f(c_0) I_j^+
    \le \min\{A_0 j^{-Rf},\,A_1 j^{-2Rf}\}.
\]
On the original intervals they may be chosen uniformly
in all the prescribed words with
\[
 \log A_0\le -(16-o(1))Rf\log b,\qquad
 \log A_1\le C R f s\log b,\qquad b=\exp(x^\varepsilon).
\]
\end{lemma}
\begin{proof}
Theorem~\ref{thm:low285} bounds the positive real integral
at the physical terminal frequency by using the unit-frequency
moment of the same word. Its pure-power factor at \(j c_0\)
is \(j^{-Rf}\) times that at \(c_0\), and the positive
logarithmic factors can only decrease for \(j\ge1\).
Multiply the real majorant at \(c_0\) by \(m_f(c_0)\).
Its logarithm satisfies the stated estimate for \(A_0\) by
the proofs of Theorems~\ref{thm:energy285} and
\ref{thm:moment286}, which bound this positive majorant
before taking the signed orbital pairing.

For the other power use the exact frequency substitution
of Proposition~\ref{prop:scale284}. Its quotient Haar
factor is \(T^{-Rf}\), and the parameter at each real place
has modulus \(jU_0/T\) at unit frequency. The proof of
Theorem~\ref{thm:moment283} applies at this frequency and
to this tuple of real parameters. Indeed its fiber calculation
in Theorem~\ref{thm:fiber283} uses only a nonzero imaginary
scalar at each place, and its Cauchy--Schwarz argument uses
only left invariance of the section norm. It gives
\[
 \begin{split}
 A_1={}&m_f(c_0) C_E\,T^{Rf}(2/U_0)^{2Rf}
       \left(\prod_\tau C_{r-2,h_\tau}\right)
       \mathbb E_1\prod_\tau l(g_\tau)^{2R},\\
 &\hspace{14mm}l(g_\tau)=\|g_\tau^{-1}U_1\|.
 \end{split}
\]
Here \(C_E\le1\) by Proposition~\ref{prop:measure283}, and
\(C_{r-2,h_\tau}\le\pi^{r-1}/(r-2)!\) by
Theorem~\ref{thm:fiber283}. The length \(l(g_\tau)\) is
bounded by the minimum coset norm \(\nu_\tau(g_\tau)\)
because \(B\) fixes the Euclidean unit vector \(U_1\).
Since \(2R\le4r^2\), the height estimate of
Theorem~\ref{thm:moment286} bounds the logarithm of the
expectation by \(O(fsR\log b)\). Finally
\(\log m_f(c_0)=O(fr\log x)\), by
Theorem~\ref{thm:product281} and the conductor bound,
and \(\log T=O(s\log b)\) by Theorem~\ref{thm:prime210}.
As \(U_0\ge1\), the asserted upper bound follows by
multiplication. It remains valid after taking the largest
majorant among the finitely many words for that field and
interval point.
\end{proof}

\subsection{First-plane parameters over the whole real field}\label{nav:I-68}
Fix \(E,F=E^+,x\) and the word from Section~\ref{sec:missing}. A good
split place in this paragraph is a finite place of \(F\)
split in \(E/F\), away from \(2q d_0 Q T\), at which the
ordinary row and column levels are the full lattice
stabilizers. Let \(\mathcal A\) be the set of all nonzero
\(a\in\mathcal O_F\) such that \((a)\) is squarefree and
supported at these places. All generators are included,
in particular every unit of \(F\). Write \(O(a)\) for
the normalized ordinary orbital Gram operator at
\(\gamma_{Q\sqrt{-q}\,a}\). This parameter is skew under
conjugation of \(E/F\), as required in
Proposition~\ref{prop:plane281}.

\begin{theorem}[The sum over integral generators with squarefree split support]\leavevmode
\label{thm:field288}
On the original power intervals one has
\[
   \sum_{a\in\mathcal A}\|O(a)\|_{\rm op}
       \le \exp\bigl(-(16-o(1))(r-1)f x^\varepsilon\bigr).
\]
This holds uniformly in the prescribed words and lowest-type
unit vectors, with the original terminal frequency, including
after projection to its terminal character and retained
coordinates.
\end{theorem}
\begin{proof}
Put \(R=r-1,\ h=s-1\), \(B_x=4(k+s+1)^2\), and
let \(\omega_F((a))\) count the distinct prime ideals dividing
\((a)\). At \(v\mid(a)\), put \(P=Nv\). In the first split
component the parameter \(Q\sqrt{-q}\,a\) has valuation one.
The computation in Theorem~\ref{thm:split281} works over the
local integer ring \(\mathcal O_v\) by counting vectors,
covectors and scalars over its residue field and lifting them
modulo a single uniformizer. Proposition~\ref{prop:lattice282}
then gives precisely \(a_P\) identity terms and \(u_P\)
terms per neighbor in the notation of
Proposition~\ref{prop:split287}. Each phase has modulus at
most one after spherical averaging. The Satake bound of
Proposition~\ref{prop:satake287} applies at this split place
with residue size \(P\). The proof of
Proposition~\ref{prop:split287} therefore bounds this local
operator by \(B_xP^R\), also when the residue degree over
the rational prime exceeds one. At all other finite places
multiplication by \(a\) is by a unit. The unchanged
valuation tests give at most the original unsigned mass
there, independently of that unit. With \(c_0=Q\sqrt{-q}\)
the finite bound is consequently
\[
 m_f(c_0)\, B_x^{\omega_F((a))}\,|N_{F/\QQ}a|^R .
\]

Combine this with the signed real factorization of
Proposition~\ref{prop:places288} by the same unfolding as in
Theorem~\ref{thm:sum287}. Put \(t_\tau=|\tau a|\).
The product of \(t_\tau^{-R}\) cancels the finite norm
power exactly. For \(t_\tau<1\) use the first real estimate
and retain its logarithmic factor. For \(t_\tau\ge1\)
use both real estimates. After removing \(t_\tau^{-R}\),
their minimum is at most
\(a_0(a_1/a_0)^\theta t_\tau^{-\theta R}\)
for any \(0\le\theta\le1\). Since \(a_1/a_0\ge1\),
multiplying over the real places gives
\[
 \begin{split}
 \|O(a)\|_{\rm op}\le {}&
   m_f(c_0)C_E\,a_0^{(1-\theta)f}a_1^{\theta f}
   B_x^{\omega_F((a))} H(a)^{-\theta R}\\
   &{}\cdot\prod_\tau(1+2\log^+(1/t_\tau))^h .
 \end{split}
\]
Here \(H(a)=\prod_\tau\max(1,t_\tau)\).
Every prime ideal has norm at least two, so
\(2^{\omega_F((a))}\le |N_{F/\QQ}a|\le H(a)\).
The logarithm of the last product is at most
\(2hL_-(a)\le2h\log H(a)\), by Lemma~\ref{lem:height288}.
Choose
\[
      \theta=\frac{\log B_x/\log2+2h+3}{R}.
\]
This lies in \((0,1)\) for sufficiently large \(x\).
The remaining height factor is at most \(H(a)^{-3}\), which
Lemma~\ref{lem:height288} sums over the entire ring,
including all choices of generator, at cost \(8^f\).

Finally \(\theta s=O((s\log r+s^2)/R)=o(1)\) by the
original rank inequalities. Also \(C_E\le1\) and
\(\log m_f(c_0)=O(fr\log x)=o(Rf\log b)\).
The two bounds for \(\log a_i\) in
Proposition~\ref{prop:places288} now yield the claimed
exponent. Orthogonal compression preserves every norm
upper bound. Distinct parameters here label distinct
rational terms, since their first minors are distinct
as in Proposition~\ref{prop:plane281}.
\end{proof}

Theorem~\ref{thm:codimension294} gives the first-plane estimate at all finite supports and valuations.

\subsection{Squarefree first-plane parameters on the original intervals}\label{nav:I-69}
On the fields and levels of Section~\ref{sec:missing} put
\(\mathcal O_E^-=\{\xi\in\mathcal O_E:\bar\xi=-\xi\}\).
For \(0\ne\xi\in\mathcal O_E^-\) there is a unique integral
ideal \(\mathfrak b(\xi)\) of \(F=E^+\) such that
\(\xi\mathcal O_E=\mathfrak b(\xi)\mathcal O_E\).
Indeed \(E/F\) is unramified at finite places by
Theorem~\ref{thm:norm91}, and at a split place its two
valuations on \(\xi\) agree. These valuations specify
the exponents of the ideal in \(F\). Let \(\mathcal S\)
consist of the nonzero skew integers for which
\(\mathfrak b(\xi)\) is squarefree and supported away from
\(QT\) at places where the row and column levels are the
full lattice stabilizers. Both split and inert primes are
allowed, as are different generators and mixed prime supports.
Write \(O_\xi\) for the normalized orbital Gram operator of
the first-plane element \(\gamma_{Q\xi}\).

\begin{theorem}[The contribution of all good squarefree skew parameters]\leavevmode
\label{thm:skew289}
Uniformly on the original power intervals and prescribed words,
with \(R=r-1,\ f=[F:\QQ]\), one has
\[
           \sum_{\xi\in\mathcal S}\|O_\xi\|_{\rm op}
             \le \exp\bigl(-(16-o(1))Rf x^\varepsilon\bigr).
\]
The same estimate holds on the original terminal-character
projection and the retained coordinates.
\end{theorem}
\begin{proof}
At \(v\mid\mathfrak b(\xi)\), write \(P=Nv\). The signed
finite operator has norm at most \(B_xP^R\), with
\(B_x=4(k+s+1)^2\). At an inert place this is
Proposition~\ref{prop:sphere289}. At a split place apply
the full-level residue count of Theorem~\ref{thm:split281}
and the spherical compression of Proposition~\ref{prop:split287}.
These use a uniformizer as the first component of \(Q\xi\)
and the actual residue cardinality \(P\), independent of
its residue characteristic or degree. The character enters
with modulus at most one in the compression. At every other
place the valuation is the minimal one of
Theorem~\ref{thm:finite281}, including above \(Q\) and at the
terminal graph. The relative column is the identity there,
with unsigned mass one. The finite operator is thus bounded
in norm by
\[
               B_x^{\omega(\mathfrak b)}(N\mathfrak b)^R ,
               \qquad \mathfrak b=\mathfrak b(\xi).
\]
Here \(\omega\) counts prime ideal factors. The product
decomposition follows by unfolding as in
Theorem~\ref{thm:sum287}, since every \(\gamma_{Q\xi}\)
commutes with the column and has the same stabilizer \(J\)
by Proposition~\ref{prop:plane281}. Its first minor \(Q\xi\)
distinguishes the rational terms.

Choose one extension to \(E\) of each real embedding \(\tau\)
of \(F\), and put \(t_\tau=|\tau\xi|\). The ordinary absolute
ideal norm gives
\[
        N\mathfrak b=\prod_\tau t_\tau,\qquad
        H(\xi)=\prod_\tau\max(1,t_\tau).
\]
Indeed \(N_{E/\QQ}(\xi)=\prod_\tau t_\tau^2=(N\mathfrak b)^2\).
The real factor is \(C_E\) times the product of the signed
local factors of Proposition~\ref{prop:places288} with real
multipliers \(\tau(\xi/\sqrt{-q})\). At a single place its
two bounds therefore take the form
\[
  a'_0 t_\tau^{-R}(1+2\log^+(1/t_\tau))^{h},
       \qquad a'_1 t_\tau^{-2R},\qquad h=s-1 .
\]
One may take \(a'_0=a_0q^{R/2}(1+\log q)^h\) and increase
\(a'_1=a_1 q^R\) to \(\max(a'_0,a_1q^R)\) if necessary.
This uses \(1+2\log^+(\sqrt q/t)
\le(1+\log q)(1+2\log^+(1/t))\).
Since \(\log q=O(\log x)\) by Theorem~\ref{thm:conductor232},
these constants still satisfy
\(\log a'_0\le-(16-o(1))R\log b\) and
\(\log a'_1\le C R s\log b\) for \(b=\exp(x^\varepsilon)\).

The factor \(\prod_\tau t_\tau^{-R}\) cancels the entire
finite norm power. For \(t_\tau<1\) use the first local
estimate. For \(t_\tau\ge1\) use the geometric mean of
the two estimates with exponents \(1-\theta,\theta\).
The inequalities
\[
  2^{\omega(\mathfrak b)}\le N\mathfrak b\le H(\xi),
  \qquad\sum_\tau\log^+(1/t_\tau)\le\log H(\xi)
\]
follow from the norm identity and \(N\mathfrak b\ge1\).
The logarithmic factors together are at most \(H(\xi)^{2h}\).
Taking \(\theta=(\log B_x/\log 2+2h+3)/R\) therefore gives
\[
       \|O_\xi\|_{\rm op}
            \le C_E(a'_0)^{(1-\theta)f}(a'_1)^{\theta f}
                          H(\xi)^{-3}.
\]

We sum over all skew integers directly. For our fixed complex
embeddings their imaginary parts are \(f\) real coordinates.
Two distinct skew integers in the same product of half-open
unit intervals would have a nonzero skew difference with
absolute norm strictly below one. Thus each such product
contains at most one. In the box indexed by integers
\(j_\tau=\lfloor\log_2\max(1,t_\tau)\rfloor\) there are at
most \(6^f2^{\sum j_\tau}\) points, by the interval count in
Lemma~\ref{lem:height288}, and \(H(\xi)\ge2^{\sum j_\tau}\).
Summing gives at most \(8^f\) for the sum of \(H(\xi)^{-3}\)
over all nonzero skew integers. In particular no choice of
generator for an ideal in \(F\) occurs in this bound.
Finally \(C_E\le1\) and \(\theta s=o(1)\) by the rank
inequalities in Theorem~\ref{thm:field288}. This proves the
stated exponent. Orthogonal compression decreases the norm.
\end{proof}

Theorem~\ref{thm:completeplane293} treats repeated prime factors and arbitrary valuations at the marked primes for the full skew family.

\subsection{A sum over the actual algebraic integers}\label{nav:III-54}
For a totally real field \(F\) of degree \(f\) and
\(a\in\mathcal O_F\setminus\{0\}\), put
\[
 H(a)=\prod_{\tau:F\hookrightarrow\RR}\max(1,|\tau a|),
 \qquad
 L_-(a)=\sum_\tau\max(0,-\log|\tau a|).
\]
The height distinguishes the different archimedean sizes of
generators of the same ideal. These sizes occur in the real
orbital integrals of Section~\ref{sec:missing}.

\begin{lemma}[A uniform height sum including all unit multiples]\leavevmode
\label{lem:height288}
One has \(L_-(a)\le\log H(a)\). For any \(\delta>1\),
\[
 \sum_{a\in\mathcal O_F\setminus\{0\}} H(a)^{-\delta}
       \le \left(\frac{6}{1-2^{1-\delta}}\right)^f .
\]
In particular the sum at \(\delta=3\) is at most \(8^f\).
\end{lemma}
\begin{proof}
The difference \(\log H(a)-L_-(a)\) equals
\(\log|N_{F/\QQ}a|\), which is nonnegative because the norm
is a nonzero rational integer. For the sum, assign \(a\)
to the tuple \(t_\tau=\lfloor\log_2\max(1,|\tau a|)\rfloor\).
In this tuple's box one has \(|\tau a|<2^{t_\tau+1}\).
Cover the corresponding interval on each real line by at
most \(2^{t_\tau+2}+2\le 6\,2^{t_\tau}\) half-open unit
intervals. Two distinct algebraic integers cannot belong
to one product of such intervals. Their difference would
have all its real absolute values strictly less than one,
contradicting the norm inequality. Thus the box contains
at most \(6^f2^{\sum_\tau t_\tau}\) algebraic integers,
while \(H(a)\ge2^{\sum_\tau t_\tau}\). Sum the resulting
\(f\) geometric series over the nonnegative integers
\(t_\tau\) to obtain the assertion.
\end{proof}

\subsection{A spherical matrix coefficient at arbitrary distance}\label{nav:III-55}
At a full-level place \(v\) of the prescribed automorphic
column write \(K_H\) for the lattice stabilizer and
\(P=Nv\). For a translation by \(a\) write \(M(a)\) for
its compression to the \(K_H\)-fixed vectors. Translate
on the same prescribed real type as in
Proposition~\ref{prop:satake287}. Normalize \(K_H\) to
have volume one.

\begin{proposition}[Spherical compression for a pair of elementary divisors]\leavevmode
\label{prop:cartan290}
For a split column position with exponents
\((-d_-,0,\ldots,0,d_+)\), \(d_\pm\ge0\), one has
\[
      \|M(a)\|\le k^{d_-+d_+}
                    P^{-(k/2-1)(d_-+d_+)} .
\]
For an inert column with the same exponents and
\(d_-=d_+=d\), one has
\(\|M(a)\|\le(k+1)^d P^{-(k-2)d}\).
\end{proposition}
\begin{proof}
In the split case elementary divisors give the displayed
diagonal representative \(a\). In the inert case the
hyperspecial Cartan decomposition gives representatives
in a hyperbolic basis, with diagonal exponents
\((\nu_1,\ldots,\nu_{k/2},-\nu_{k/2},\ldots,-\nu_1)\).
Their elementary divisors as matrices over \(E_w\)
determine them up to signed permutations, so the pair
in question is the \(d\)-th power of the representative
in Proposition~\ref{prop:sphere289}.
\href{https://math.umd.edu/~tjh/satake24.pdf}
{Haines and Rostami, Cartan decomposition, Section 1}.

Let \(\phi_a=1_{K_H a K_H}\), and let \(S\) be its normalized
Satake transform in the upper triangular Borel. Dividing the
eigenvalue of \(\phi_a\) by its degree
\(\operatorname{vol}(K_H a K_H)\) gives the eigenvalue of
\(M(a)\), by averaging over both compact factors.
The transform of a nonnegative function has nonnegative
coefficients by its unipotent integral. It preserves
convolution as multiplication of Laurent polynomials.
Thus if \(a=a_1\cdots a_j\) is a product of diagonal
elements, the coefficients of \(S\phi_a\) are at most
those of \(\prod_i S\phi_{a_i}\). Indeed
\(\phi_{a_1}*\phi_{a_2}\ge\phi_{a_1a_2}\), since at
\(a_1a_2\) one may integrate over \(a_1K_H\), which has
volume one. Bi-invariance and induction
give the assertion for the product.

At a split place use \(d_+\) index-\(P\) translations and
\(d_-\) inverse translations. Their Satake polynomials
are \(P^{(k-1)/2}\sum_i X_i\) and
\(P^{(k-1)/2}\sum_i X_i^{-1}\). At the absolute values of
the parameters of any prescribed constituent each is at
most \(kP^{k/2}\), by Proposition~\ref{prop:satake287}.
Consequently the numerator of the normalized eigenvalue is
bounded by \((kP^{k/2})^{d_-+d_+}\). At an inert place use
\(d\) copies of the neighbor translation. Its polynomial in
Proposition~\ref{prop:sphere289} has nonnegative coefficients
and is at most \((k+1)P^k\) at the absolute values of the
prescribed parameters. This gives \(((k+1)P^k)^d\) for
the corresponding numerator.

For the denominator let \(\delta\) be the Borel modulus.
At the displayed diagonal \(a\), chosen with the negative
exponent first, the coefficient of the corresponding
monomial in \(S\phi_a\) is at least \(\delta(a)^{1/2}\)
because the integral contains all integral upper unipotents.
The coefficient at the opposite Weyl translate is the
same. Evaluation on the trivial representation multiplies
this last coefficient by \(\delta(a)^{1/2}\), since the
inducing character is \(\delta^{-1/2}\) and the longest
Weyl element inverts the modulus. All contributions there
are nonnegative. The degree, which is the trivial
eigenvalue, is therefore at least \(\delta(a)\).
In the split case the upper triangular entries give
\(\delta(a)=P^{(k-1)(d_-+d_+)}\). In the inert case
\(\delta(a)=P^{2(k-1)d}\) by the first-line calculation of
Proposition~\ref{prop:sphere289}. Dividing the respective
numerator bounds by these powers proves the claim on each
spherical constituent and hence on their orthogonal sum.
\end{proof}

\subsection{The factorial in the real first-plane estimate}\label{nav:III-56}
For \(0\ne\xi\in\mathcal O_E^-\) put
\(t_\tau=|\tau\xi|\) and \(H(\xi)=\prod_\tau\max(1,t_\tau)\).
Write \(d=k+1\) and \(h=s-1\). The first-plane estimate of
Theorem~\ref{thm:low285} includes the volume of a complex
ball of dimension \(d\). Its factorial is needed when
marked primes consume the full factor \(Q^{Rf}\).

\begin{lemma}[A uniform bound with the factorial retained]\leavevmode
\label{lem:factorial291}
With \(\mathfrak m\) and \(O_\xi\) as in
Proposition~\ref{prop:marked291} and Theorem~\ref{thm:skew289},
there is a constant \(C\) independent of the interval
point, word and \(\xi\) such that
\[
 \|O_\xi\|_{\rm op}\le
    (N\mathfrak m)^{R-1} H(\xi)^{-3}
          Q^{-Rf}(d!)^{-f}\exp(C R f).
\]
\end{lemma}
\begin{proof}
At the original frequency the first estimate in
Theorem~\ref{thm:low285}, applied locally as in
Proposition~\ref{prop:places288}, bounds the signed factor by
\[
 Q^{-R} L\,t_\tau^{-R}
             (1+2\log^+(1/t_\tau))^h,\quad
 L= 2^R A_\sharp \exp(C s(4s^2+3)\log b)
                  (1+2\log^+(2T/Q))^h .
\]
This follows by separating the positive logarithm of
\(2T/(Qt_\tau)\) into its constant part and its \(t_\tau\)
part as in that proposition. The explicit constants in
Proposition~\ref{prop:shell285} give
\[
 A_\sharp=(2s^2+1)(2+3(s+1))\,
                 2^{1+2h}(4\sqrt2)^d\pi^{r-2}/d! .
\]
The rank formulas in Theorem~\ref{thm:rangenew50} give
\(s^3\log b=o(k)\), and Theorem~\ref{thm:prime210} gives
\(\log(1+2\log^+(2T/Q))=O(\log x)\). Consequently
\(\log L\le-\log(d!)+C R\), and we may use
\(L=\exp(C R)/d!\) as a common larger constant.

Lemma~\ref{lem:power291} gives a second bound
\(Q^{-R} U t_\tau^{-3R}\) with \(U=\exp(C'R s\log b)\).
Here the proof of that lemma supplies \(Q^{-3R}\), so
\(Q^{-R}\) can be taken outside both bounds. Increase \(C'\)
if necessary to ensure \(U\ge L\).
The product of the signed real factors has just the
additional field constant \(C_E\le1\) by
Proposition~\ref{prop:places288}. Remove \(t_\tau^{-R}\)
from each factor, cancelling
\((N\mathfrak g\,N\mathfrak m)^R\) in
Proposition~\ref{prop:marked291} by the norm identity of
Theorem~\ref{thm:skew289}. Leave
\((N\mathfrak m)^{R-1}\) aside. The remaining finite
factor is at most \(H(\xi)^{2s+1+\log C_k/\log2}\).
All logarithmic factors cost at most \(H(\xi)^{2h}\),
since \(\prod t_\tau\ge1\) bounds the sum of the
negative logarithms by \(\log H(\xi)\).

For \(t_\tau\ge1\) interpolate the two real estimates with
exponent \(\theta\) on the second. At \(t_\tau<1\) use the
first. Take \(2R\theta=2s+1+\log C_k/\log 2+2h+3\).
The extra power \(H(\xi)^{-2R\theta}\) then leaves exactly
\(H(\xi)^{-3}\) in the majorant. The remaining per-place
constant after taking out \(Q^{-R}\) is at most
\(L^{1-\theta}U^\theta\). Its logarithm is at most
\[
  -\log(d!)+ C R+\theta\{\log(d!)+C'Rs\log b\}.
\]
Indeed \(0<\theta<1\) on the specified large intervals.
By the ranks in Theorem~\ref{thm:rangenew50},
\(\log k=O(\log x)=o(s)\), \(s^3\log b=o(k)\), and
\(\log(d!)\le d\log d\). Thus
\(\theta\log(d!)=O((s+\log k)\log k)=o(R)\) and
\(\theta R s\log b=O((s+\log k)s\log b)=o(R)\).
Multiplication over the \(f\) places proves the estimate.
\end{proof}

\subsection{A weighted bound for the real orbital integral}\label{nav:V-47}
Use \(D=2r-2\) and the measures of
Proposition~\ref{prop:measure283}. For a nonzero real Bessel
section of the original word, put
\[
 M_D(W)=
 \frac{\int_{B_\infty\backslash G_\infty}\|W(g)\|^2
                 \prod_\tau R(g_\tau)^D\,d\dot g}
      {\int_{B_\infty\backslash G_\infty}\|W(g)\|^2\,d\dot g}.
\]
All norms \(R\) use the same unit first vector as in
Theorem~\ref{thm:fiber283}.

\begin{theorem}[A common moment for the first-plane real integrals]\leavevmode
\label{thm:moment283}
For every rational skew parameter \(c\ne0\) of
Proposition~\ref{prop:plane281}, the real factor of
Theorem~\ref{thm:real282} satisfies
\[
 |\mathcal I_\infty(\gamma_c;W)|
 \le C_E\, M_D(W)
       \prod_\tau \left(\frac{2}{|\tau c|}\right)^D C_{n,h_\tau}.
\]
The right side also bounds the integral of the product of the
two real norms divided by the same diagonal normalization.
The moment is finite for the original lowest-type sections
and is independent of the rational parameter \(c\).
\end{theorem}
\begin{proof}
The change of opposite vector made in Theorem~\ref{thm:fiber283}
preserves the orbital integral and the integral of norms.
Indeed left multiplication by the corresponding central
unipotent normalizes \(J_\infty\) and preserves its Haar
measure, so it preserves the quotient measure. It acts on
\(W\) on the left by the trivial character. We may therefore
use the representatives \(\gamma'\) of that theorem
simultaneously at all real places.

Write \(R_1(g)=\prod_\tau R(g_\tau)\) and
\(R_2(g)=R_1((\gamma')^{-1}g)\), using the product element
\(\gamma'\). They are positive and invariant on the left
under \(J_\infty\). Apply Cauchy--Schwarz to the product of
norms on \(J_\infty\backslash G_\infty\) with weight
\((R_1/R_2)^D\) in the first squared integral and its inverse
in the second. In the first integral disintegrate over
\(B_\infty\backslash G_\infty\). The fiber \(J_\infty
\backslash B_\infty\) is \(U_\infty\) in coordinates \(Jn\)
with the measure of Proposition~\ref{prop:measure283}.
The norm \(\|W\|\) and \(R_1\) are constant along that fiber.
The integral of \(R_2^{-D}\) is exactly the product in the
claimed bound, times \(C_E\), by
Theorem~\ref{thm:fiber283}. In the second squared integral
left translate \(g\) by \(\gamma'\). This preserves quotient
measure because \(\gamma'\) centralizes \(J_\infty\).
The resulting first fiber integral has parameter \(-c\)
and hence has exactly the same value. The geometric mean
of these two equal expressions proves the bound.

For the sections in question the calculation is integrable.
For every \(b'\in B(F_\tau)\), \(R(g_\tau)\le
\|(b'g_\tau)^{-1}\|_{\rm op}=\|b'g_\tau\|_{\rm op}\),
since \(b'\) fixes \(u'\) and inverse and original operator
norms agree in signed coordinates. Taking the infimum
bounds \(R_1\) by the height \(V\) of
Proposition~\ref{prop:weight163}. That proposition gives
boundedness of \(W\) and integrability of \(\|W\|V^L\),
where \(L=8(f+1)(r+1)^3>D\). Thus the positive moment in
the statement is finite. All disintegrations with nonnegative
integrands also follow directly from Tonelli's theorem.
\end{proof}

\subsection{The weighted volume of the first-plane fiber}\label{nav:V-48}
Set \(h=s-1,\ d=k+1,\ n=d+2h\), and use
\(\delta=|C|\) and \(y=\gamma(C)^{-1}n(w,z)g\).
Here \(g\) is the representative from Lemma~\ref{lem:negative285}.
Let
\[
 K=2s^2,\qquad
 B_*=(2+3(s+1))\,2^{\,1+2h}(4\sqrt2)^d\,\pi^n/d!.
\]

\begin{proposition}[Weighted sections of the unipotent fiber]\leavevmode
\label{prop:shell285}
For every \(R\ge1\), with \(H=R\nu(g)\), one has
\[
 \int_{\nu(y)\le R}
       \left(\frac{{\cal A}(g)}{{\cal A}(y)}\right)^d\,dw\,dt
 \le B_*\,Z(g)^2 H^K\,
         \delta^{-(n+1)}(1+2\log^+(1/\delta))^h .
\]
\end{proposition}
\begin{proof}
The exterior product defining \({\cal A}\) has unit norm
before translation. On any coset attaining the minimum
this gives
\(\nu^{-s-1}\le{\cal A}\le\nu^{s+1}\), using the equal
norms of a unitary matrix and its inverse. The ratio
\(\theta={\cal A}(y)/{\cal A}(g)\) therefore belongs to
\([H^{-s-1},H^{s+1}]\). Divide this into its intersections
with \((2^{j-1},2^j]\), \(j\in\ZZ\). There are at most
\(2+2(s+1)\log H/\log2\le(2+3(s+1))H\) nonempty pieces.

On such a piece Lemma~\ref{lem:negative285} bounds the
squared radius in the coordinates \(x-\eta v_+/a,\eta/a\)
by \(2\sqrt2\,2^j/\delta\). The inverse change of coordinates
has real Jacobian \(a^2\). Multiplying the ball volume by
the upper bound \(2^{(1-j)d}\) on \(\theta^{-d}\) gives
\[
       a^2\,\frac{\pi^d}{d!}(4\sqrt2/\delta)^d .
\]
This is independent of \(j\). On the same set
\(n^{-1}\gamma(C)b\) has norm at most \(H\) for some \(b\in B\),
by the argument of Theorem~\ref{thm:norm284}.
For the \(h\) complex pairs the bound of
Theorem~\ref{thm:minor284} gives
\(\pi^{2h}(2H^{s+1}/\delta)^{2h}
       (1+2\log^+(1/\delta))^h\).
The real \(t\)-coordinate ranges over an interval of length
at most \(2H/\delta\). All these containing sets may be
multiplied for an upper bound even though the original
conditions couple the variables. Multiply by the number
of pieces. The total exponent of \(H\) is
\(1+2h(s+1)+1=2s^2\). Finally \(a^2\le Z(g)^2\).
\end{proof}

\subsection{The first flag coordinate of a tensor projection}\label{nav:V-49}
At one real place use the Euclidean orthonormal basis, in order,
\[
 U_1,\ldots,U_s,\quad w_1,\ldots,w_k,e,\quad V_s,\ldots,V_1 ,
 \qquad U_j=u_j/\sqrt2,\quad V_j=\sqrt2\,v_j .
\]
Here the \(w_a\) form an orthonormal basis of the positive
column and \(u_j,v_j\) are the Witt vectors of Section~\ref{sec:char94}.
The row has rank \(r=k+2s+1\). Write \(S_\mu\) for the smooth
discrete series with decreasing integral Harish--Chandra parameter
\(\mu\) and the signs of \eqref{eq:src109}. Write \(F^*\) for
the dual of its group's defining representation. In this paragraph
adjacent gaps of \(\mu\) exceed two. Let \(P_i\) be the projection
in \(F^*\otimes S_\mu\) onto the summand with parameter
\(\mu-e_i\), where \(e_i\) in a parameter denotes a coordinate
unit vector. The coherent character and splitting argument of
Lemma~\ref{lem:translate163} applies with the dual defining
representation, whose weights are the \(-e_i\). The canonical
smooth direct sum follows from the Casselman--Wallach globalization
theorem. Thus the projections intertwine the group actions
on smooth vectors. See also Bernstein and Krötz,
\href{https://arxiv.org/abs/0812.1684}{\emph{Smooth Fréchet Globalizations of Harish--Chandra Modules}}.

Put \(\omega=h(U_1,\cdot)\), the last dual basis vector, and
write \({\rm ev}\) for evaluation at \(U_1\) in the dual
factor. Both are fixed by \(B\). For the unit-frequency
character of Section~\ref{sec:char94}, let
\(\Lambda:S_\mu\longrightarrow E_u\) be a continuous Bessel
map to the compact column of parameter \(u=(u_a)_{a=1}^k\).

\begin{proposition}[A quantum minor on the prescribed Bessel model]\leavevmode
\label{prop:cofactor286}
There is a nonzero constant \(\eta\), common to all \(i,\mu,u\)
at these ranks, with \(|\eta|=(2\pi)^{2s}/2\), such that
\[
 \Lambda\,{\rm ev}\,P_i(\omega\otimes v)
  = \eta\,
   \frac{\prod_{a=1}^k(u_a+1/2-\mu_i)}
        {\prod_{j\ne i}(\mu_j-\mu_i)}\,\Lambda(v)
 \qquad (v\in S_\mu).
\]
\end{proposition}
\begin{proof}
Denote the complexified Lie algebra matrix units and their
actions by \(E_{ij}\). For two ordered lists \(I,J\) of length
\(t\), set \(T(z)_{ij}=z\delta_{ij}+E_{ij}\) and define
\[
 M_{I,J}(z)=\sum_{\pi\in\mathfrak S_t}{\rm sgn}(\pi)
           \prod_{a=1}^{t}T(z-a+1)_{i_{\pi(a)},j_a}.
\]
Factors act by composition in their printed order. We shall
also use the expansion that permutes the columns and has
factors in row order with arguments \(z-t+1,\ldots,z\).
Here is a verification. In \(t\) copies of the defining
space let \(P_{ab}^{\,0}\) interchange copies \(a,b\), and
let \(A_t\) be the signed permutation sum divided by \(t!\).
Use \(T_a(z)\) for the matrix \(T(z)\) in copy \(a\).
The matrix-unit commutator gives
\[
 R_{ab}(u-v)T_a(u)T_b(v)=T_b(v)T_a(u)R_{ab}(u-v),
 \qquad R_{ab}(w)=1-P_{ab}^{\,0}/w .
\]
Indeed its entries are
\((u-v)[T_{ij}(u),T_{kl}(v)]
=T_{kj}(u)T_{il}(v)-T_{kj}(v)T_{il}(u)\).
In lexicographic order of the pairs \(a<b\), the product of
the \(R_{ab}(b-a)\) equals \(t!A_t\). To see this, move the
factors involving \(t\) to the right without changing their
relative order, commuting only disjoint transpositions.
After the induction hypothesis the remaining calculation is
\[
 A_{t-1}\prod_{i<t}\left(1-\frac{(i\,t)}{t-i}\right)
             =A_{t-1}\left(1-\sum_{i<t}(i\,t)\right)=t A_t .
\]
In the expansion, a nonempty set of transpositions with least
index \(i\) lies in the left coset of \((i\,t)\) under
\(\mathfrak S_{t-1}\), as its inverse sends \(t\) to \(i\).
Its signed coefficient after antisymmetrizing contributes
to that coset a total of
\(-\prod_{i<j<t}(1+1/(t-j))/(t-i)=-1\).
The last equality in the display is the signed coset decomposition.
The exchange relation now reverses all \(t\) factors, moving
the antisymmetrizer from the left to the right. Taking entries
gives alternation in both lists. On the right the factors
occur in decreasing slot order. Reversing both index lists,
which leaves their combined sign unchanged, gives the asserted
row expansion.

For the full ordered list \(J_0=(1,\ldots,r)\), put
\(D_r(z)=M_{J_0,J_0}(z)\). It is central. In fact commuting
\(E_{ab}\) past the factors replaces each occurrence of row
\(b\) by row \(a\), and subtracts the replacements of column
\(a\) by column \(b\), by the matrix-unit relations. Repetitions
vanish, and for \(a=b\) the terms cancel. On a highest weight
vector of weight \(\lambda\) its value is
\(\prod_j(z-j+1+\lambda_j)\). In a nonidentity permutation
the last moved column has a strictly smaller row index and
its upper matrix unit kills the highest weight vector after
the diagonal factors to its right. The Harish--Chandra
isomorphism with shift \((r+1-2j)/2\) thus gives
\[
 D_r(Y+(r-1)/2)=\prod_j(Y+\mu_j)
\]
on the entire discrete series. The same calculation applies
to the compact \(k\)-column with parameter \(u\).

The split Casimir \(C\) on \(F^*\otimes S_\mu\) sends a
coordinate vector \(x=(x_j)\) to
\((-\sum_jE_{ij}x_j)_i\). This follows from
\(E_{ab}\zeta=-\zeta\circ E_{ab}\) on the dual.
Its eigenvalue on summand \(j\) is
\(-\mu_j-(r-1)/2\). Indeed twice \(C\) is the difference
between the tensor Casimir and the two separate Casimirs,
for the trace form. The row eigenvalue is
\(\sum\mu_i^2-\sum((r+1-2i)/2)^2\), lowering coordinate \(j\)
changes it by \(1-2\mu_j\), and the dual defining eigenvalue
is \(r\).

Let \(A(z)_{ij}=(-1)^{i+j}M_{J_0\setminus j,J_0\setminus i}(z)\),
with the inherited order. Expanding the last column gives
\(A(z)T(z-r+1)=D_r(z)I\). For an off-diagonal entry the
appended column repeats one already present. At
\(z=Y+(r-1)/2\), the matrix on the left multiplying vectors
first is scalar \(Y+\mu_j\) on summand \(j\). Consequently
\(A(z)\) there equals \(\prod_{l\ne j}(Y+\mu_l)\).
This is a polynomial identity, so substitution \(Y=-\mu_i\)
gives \(A(z)=\prod_{l\ne i}(\mu_l-\mu_i)P_i\).

It remains to compute the \((1,r)\) entry after applying
\(\Lambda\). In the complexified unipotent radical all upper
matrix units outside the middle square act through scalars
under \(\Lambda\). The only nonzero scalars belong to the
successive arrows in the ordered list
\(U_1,\ldots,U_s,e,V_s,\ldots,V_1\), where an arrow from
column \(b\) to row \(a\) means \(E_{ab}\) for two consecutive
entries \(a,b\) of this list. For the two arrows touching
\(e\) the absolute value is \(2\pi/\sqrt2\), and for every
other arrow it is \(2\pi\). Indeed the real differential is
\(4\pi i\,{\rm Re}\,d\), with
\(d(X)=\sum_{j<s}X_{U_j,U_{j+1}}+X_{U_s,e}/\sqrt2\).
Its complex-linear extension is
\(2\pi i(d(X)+\overline{d(-JX^*J)})\), where \(J\) interchanges
the two Witt strings, fixes the positive column and negates \(e\).
This also proves vanishing for every other upper entry in the
radical. The column matrix units act through the differentiated
compact column representation by equivariance.

The minor for the \((1,r)\) entry omits row \(V_1\) and
column \(U_1\). Remove the rows \(U_1,\ldots,U_s\) in this
order, using the row expansion. At each step the row index is
absent from the remaining columns, and every column to be tested
lies strictly to its right. Thus \(\Lambda\) on the leftmost
factor kills all terms except the column at the next arrow.
Deleting this row and that column leaves the parameter \(z\)
unchanged in the row expansion. Now take the columns
\(V_1,\ldots,V_s\) in that order, moving the chosen column to
the front by alternation. Its index is absent from the remaining
rows, which lie strictly to its left. In the column expansion
only the row at the adjoining arrow survives. Each deletion
changes \(z\) to \(z-1\). The remaining square is precisely the
positive \(k\)-column at \(z-s=Y+k/2\), of value
\(\prod_a(Y+1/2+u_a)\) under \(\Lambda\).
All permutation signs and arrow scalars, including the cofactor
sign, form the same constant \(\eta\) for each \(i\). Their
product has the asserted absolute value. Evaluation at \(U_1\)
selects row one, and insertion of \(\omega\) selects column \(r\).
This proves the formula.
\end{proof}

The quantum minor used here also appears in
\srcloc{Development/AbelianVarietyTheory/Proofs/JetHarmonics/ProjectorEntryValueUnderLinearFunctional}{24}.

\subsection{Separate real estimates for a field-valued multiplier}\label{nav:V-50}
Let \(c_0=Q\sqrt{-q}\), \(U_0=Q\sqrt q\),
\(R=r-1=k+2s\) and \(h=s-1\). For each real place
\(\tau\) choose any Haar density on \(B_\tau\backslash G_\tau\).
On \(J_\tau\backslash G_\tau\) use the compatible density
with ordinary \(dw\,dt\) on the first unipotent fiber as in
Proposition~\ref{prop:measure283}. For a real \(z\ne0\)
write \(I_\tau(z)\) for the signed real orbital integral
at parameter \((\tau c_0)z\) and the original frequency
\(\varpi_\tau\), divided by its diagonal norm. The quotient
of these densities in particular specifies the normalization
of this number.

\begin{proposition}[The local factors for independent real multipliers]\leavevmode
\label{prop:places288}
On the lowest compact type the signed normalized real factor
for a tuple \((z_\tau)\) is \(C_E\prod_\tau I_\tau(z_\tau)\),
for every nonzero row vector. There are positive
\(a_0\le a_1\), the same for every place and word, such that
for \(t=|z|>0\),
\[
 \begin{split}
 |I_\tau(z)|&\le a_0 t^{-R}
              (1+2\log^+(1/t))^h,\\
 |I_\tau(z)|&\le a_1 t^{-2R},\\
 \log a_0&\le -(16-o(1))R\log b,\qquad
 \log a_1\le C R s\log b .
 \end{split}
\]
Here \(C_E\le1\) is the field constant of
Proposition~\ref{prop:measure283}. Neither the signs nor the
relative sizes of the different \(z_\tau\) are restricted.
\end{proposition}
\begin{proof}
Polarize the numerator in the two input vectors of a fixed
real word, giving a sesquilinear form on its finite-dimensional
lowest compact type. The integrals for two different vectors
converge absolutely by the weighted Cauchy--Schwarz proof of
Theorem~\ref{thm:low285}, since each squared integral
uses the height moment of its own input. The denominator
similarly gives a positive definite form by
Corollary~\ref{cor:rel155}. The numerator and denominator are
invariant under simultaneous right translation by the maximal
compact group. Indeed one changes the quotient variable by
right multiplication. Schur's lemma now makes the numerator
a scalar multiple of the denominator on the irreducible
compact type, even if the scalar is not real.

The argument also applies to the product compact group and
the irreducible external tensor product of its lowest types.
One may therefore evaluate the global scalar on any nonzero
pure tensor. On this vector the integrands factor, the signed
numerator being absolutely integrable. The common scalar in
the quotient of the two product measures comes from the
unipotent fiber \(J_\infty\backslash B_\infty\). Its density
is \(C_E\prod_\tau dw_\tau\,dt_\tau\) by
Proposition~\ref{prop:measure283}. This proves the product
formula, independently of the local choices of base density.

Work now at a single place. The local estimates in the
proof of Theorem~\ref{thm:moment286}, before the forms at
different places are tensored, give
\(\mathbb E_{1,\tau}\nu_\tau^A\le\exp(CsA\log b)\)
for \(1\le A\le4r^2\). With \(D=4s^2+3\), use the local
positive integral bound of Theorem~\ref{thm:low285}
and the exact frequency change of
Proposition~\ref{prop:scale284}. The result is at most
\[
 A_\sharp\exp(CsD\log b)(2/U_0)^R t^{-R}
        (1+2\log^+(2T/(U_0t)))^h .
\]
The last parenthesis is at most
\((1+2\log^+(2T/U_0))(1+2\log^+(1/t))\).
Use the factor independent of \(t\) to define the first
constant. Its logarithm is at most
\(-(16-o(1))R\log b\) since \(\log Q=(16+o(1))\log b\),
\(\log A_\sharp=O(r+\log(2s))\),
\(\log(1+2\log^+(2T/U_0))=O(\log x)\) and
\(s^3=o(R)\).

For the second estimate use the same frequency substitution
and the local calculation of Theorem~\ref{thm:fiber283}
with power \(2R\). As in Lemma~\ref{lem:real287}, this
gives the constant
\[
     C_{r-2,h_\tau}\,T^R(2/U_0)^{2R}\,
                       \mathbb E_{1,\tau} l(g_\tau)^{2R}
\]
times \(t^{-2R}\). Use \(l\le\nu_\tau\),
\(C_{r-2,h_\tau}\le\pi^R/(r-2)!\),
and \(\log T=O(s\log b)\) to bound its logarithm by
\(C R s\log b\). Uniform constants may be taken in both
calculations. Increasing the second constant to
\(\exp(C R s\log b)\) ensures \(a_1\ge a_0\) on the
specified sufficiently distant intervals.
\end{proof}

\subsection{The first-plane integral under a change of terminal frequency}\label{nav:VI-50}
Keep the original finite level and all its characters. Here
the change of variables takes place at infinity. For unit
terminal frequency and a tuple \(c=(c_\tau)\) of nonzero
purely imaginary parameters, use \(\mathcal I(c;W_1)\)
for the quotient of the two integrals in
Theorem~\ref{thm:real282}, using
\(\gamma(C_\tau)\) of Theorem~\ref{thm:minor284} with
\(iC_\tau=c_\tau/2\). This definition also applies to the
terminal-frequency section \(W_\varpi\). Write \(\mathcal I^+\)
for the same quotient with the product of norms in the
numerator. The finite factors do not occur in either definition.

\begin{proposition}[Exact frequency change for an individual real orbit]\leavevmode
\label{prop:scale284}
Use corresponding sections \(W_\varpi,W_1\) from
Proposition~\ref{prop:scale151}, and write \(f=[F:\QQ]\).
For \(n=r-2\), the real factor of the original rational orbit is
\[
 \mathcal I_\infty(\gamma_c;W_\varpi)
       =T^{-(n+1)f}\,\mathcal I((\tau c/T)_\tau;W_1).
\]
The identical scaling holds for \(\mathcal I^+\).
\end{proposition}
\begin{proof}
The central unipotent change of opposite vector in
Theorem~\ref{thm:fiber283} commutes with \(J\) and acts on
each section by its trivial character. Apply this change in
the numerator to pass from the rational vector to the
opposite vector of Euclidean length \(1/\sqrt2\).
At a place with terminal parameter \(\beta=\varpi_\tau\),
let \(a\) multiply each unit first flag vector of a Witt
pair by \(\beta\) and its partner by \(\bar\beta^{-1}\).
It fixes the column and \(e\). Conjugation by \(a\) sends
\(\gamma(C_\tau)\) to \(\gamma(C_\tau/|\beta|^2)\).
It normalizes \(B\) and its subgroup \(J\) on the middle
space. Let \(j_B,j_J\) be its absolute Haar moduli there.

For either subgroup \(L=B,J\), substitution on its left
quotient gives
\(\int_{L\backslash G}F(ag)\,d\dot g
    =j_L\int_{L\backslash G}F(g)\,d\dot g\).
For instance take \(F(g)=\int_L F_0(lg)\,dl\), change
\(l\) to \(ala^{-1}\) inside the integral, and use left
invariance of Haar on \(G\). Local quotient sections give
the formula for all integrable functions, and then for
nonnegative functions by monotone convergence.
Since \(B=J\ltimes U\), \(j_B/j_J\) is the modulus on \(U\).
Write \(a_M\) for the middle block of \(a\). On \(U\) the
coordinates transform as
\[
         w\longmapsto\bar\beta\,a_M w,\qquad
         t\longmapsto |\beta|^2t .
\]
Here \(|\det_{\CC}a_M|=1\). Thus \(j_B/j_J=|\beta|^{2n+2}\).
The scalar \(j_B^{-1/2}\) on each section in the numerator
contributes \(j_B^{-1}\), whereas the substitution on
\(J\backslash G\) contributes \(j_J\).
The denominator is preserved by the isometry on
\(B\backslash G\). Multiply the ratios over the real places
and use \(|\varpi_\tau|^2=T\). The same argument works before
the real vectors are contracted.
\end{proof}

\subsection{The real integral for one summand of the projected Gram form}\label{nav:VII-50}
Fix one original word \(\rho\), its unit lowest-type row vector
and the corresponding \(V_\rho\)-valued Bessel section \(W\)
of Theorem~\ref{thm:ordinary163}. Use the right-invariant seed
of equation~\eqref{eq:seed127} at \(K^0\). Its column
inner product is normalized so that the diagonal orbital
term is the identity. Put \(\gamma=\gamma_{Q\sqrt{-q}}\)
and let \(J\) be its stabilizer from
Proposition~\ref{prop:plane281}.

Choose Haar measures as follows. In \(U\), the radical of the
first-line stabilizer, give every \(U_L\) of
Theorem~\ref{thm:finite281} finite local volume one.
Normalize the remaining real Haar factor so that \([U]\)
has volume one. Use this product measure with that on \(J\)
to measure \(B=J\ltimes U\), and take compatible quotient
measures everywhere. Both \([J]\) and \([B]\) can then be
given volume one: their volume ratio is the volume of
\([U]\), by integration first on the normal unipotent group.
Write \(G_\infty=G(F\otimes_\QQ\RR)\) and use the analogous
notation for its subgroups. With inner product linear in
the first entry, set
\[
 \mathcal I_\infty(\gamma;W)=
 \frac{\displaystyle\int_{J_\infty\backslash G_\infty}
       \langle W(t),W(\gamma^{-1}t)\rangle\,d\dot t}
      {\displaystyle\int_{B_\infty\backslash G_\infty}
                     \|W(t)\|^2\,d\dot t}.
\]
The numerator integrand is well defined since \(\gamma\)
commutes with \(J\) and the two \(J_\infty\)-actions on the
sections are the same and unitary.

\begin{theorem}[One rational term in the projected cuspidal Gram form]\leavevmode
\label{thm:real282}
The numerator is absolutely convergent and the denominator
is positive. The mixed orbital Gram term at \(\gamma\) for
ordinary finite columns \(\ell,\ell'\) is
\[
             \mathcal I_\infty(\gamma;W)\,
                    \langle\ell,F_q^{\mathrm{col}}\ell'\rangle .
\]
In particular its contribution to the normalized signed trace
on any nonzero terminal space of Theorem~\ref{thm:range239},
including the retained-coordinate version, is
\(\operatorname{Re}\mathcal I_\infty(\gamma;W)\) times a real
number between \(\prod_{v\mid q} L_{b_v}\) and
\(\prod_{v\mid q} U_{b_v}\).
\end{theorem}
\begin{proof}
In the orbital integral unfolded in Theorem~\ref{thm:ordinary163}
integrate first on \([J]=J(F)\backslash J(\mathbb A_F)\) and
then on \(J(\mathbb A_F)\backslash G(\mathbb A_F)\).
The absolute majorant in that theorem applies to the norms
of the sections before contraction by any column covector.
For the fixed rational orbit this majorant has finite integral.
It is constant in absolute value under the common left
\(J\)-action, and has a finite right coset of positive measure
on which both finite supports occur, for example the identity
coset. The integral of the product of the two real norms on
\(J_\infty\backslash G_\infty\) is therefore finite. The
denominator is the squared norm of the nonzero real section.
These facts justify all the following disintegrations.

The unipotent radical of \(J\) has the same Bessel character on
the two sides. Its factors cancel in the product. The integral
on \([J]\) therefore reduces to the common column argument
on \([H]\), with its Haar probability. On each finite cell
write \(t_f=n_f\), modulo the left \(J\)-action and the
right level, as in Theorem~\ref{thm:finite281}. If the
second column is represented by \(h'\), the finite
covectors in the integrand are evaluated at \(h_f\) and
\(h_f h'\). The local phases are one by
Theorem~\ref{thm:phase282}. On the real factors both vectors
are translated by the same \(\rho(h_\infty)\), since the
relative element commutes with \(H\). Integrate also over
right translations of \(h_\infty\) by the compact real
group \(H_\infty\), with probability measure. This preserves
the integral over \([H]\). Schur orthogonality on the
irreducible word contracts the real vectors to
\(\langle W(t),W(\gamma^{-1}t)\rangle/\dim\rho\);
the remaining pairing is that of the two column covectors
at \(h_f\) and \(h_f h'\).

For the identity orbital term the same calculation has
\(\|W(t)\|^2/\dim\rho\) on \(B_\infty\backslash G_\infty\)
and the identity on the finite columns. The choices of
measure give the same probability on the common column
argument in both calculations. Dividing by the diagonal
normalization cancels the dimension and leaves the displayed
ratio of real integrals. Each ratio of finite stabilizer
measures is exactly the one counted in
Theorem~\ref{thm:finite281}. Summing the translated column
pairings accordingly gives \(F_q^{\mathrm{col}}\) by
Theorem~\ref{thm:phase282}. Finally average over an
orthonormal basis in the chosen projected space.
Theorem~\ref{thm:bounds282} bounds each diagonal matrix
element between the two stated products.
\end{proof}

Thus this orbit has an explicitly positive finite multiplier
in the projected trace. Its sign in that trace is determined
by the real integral. Controlling the full trace requires
estimates for the real integrals and the sum over rational
orbits, with the arithmetic normalization of Section~\ref{sec:missing}.

\subsection{Integrating over the first isotropic line at infinity}\label{nav:VII-51}
At one real place put \(R(g)=\|g^{-1}u'\|\), in the frame
of Proposition~\ref{prop:measure283}. The Bessel group fixes
\(u'\), so this function descends to its left quotient.
Let \(n=p_0+q_0=r-2\), with \(p_0,q_0\) the signatures
on \(M\), and write \(h(w,w)=|w_+|^2-|w_-|^2\).
For \(D=2n+2\) set
\[
 C_{n,h}=\int_{\CC^n\times\RR}
   (1+|w|^2+t^2+h(w,w)^2/4)^{-D/2}\,dw\,dt .
\]
Let \(\gamma'\) be the first-plane transvection formed
with the opposite vector \(v'/\sqrt2\) and parameter
\(c\in i\RR\setminus\{0\}\). It is conjugate to the
original \(\gamma_c\) by an element of the center of \(U\)
at this place. That center has trivial Bessel character
since \(s\ge2\), and it commutes with \(J\).

\begin{theorem}[The real integral on a first-line fiber]\leavevmode
\label{thm:fiber283}
With the unscaled coordinate measure of the preceding
display, for every \(g\in G(F_\tau)\) one has
\[
 \int_{\CC^n\times\RR}
       R\bigl((\gamma')^{-1}n(w,z)g\bigr)^{-D}\,dw\,dt
                 = (|c|/2)^{-D} C_{n,h}.
\]
The same identity holds if \(c\) is replaced by \(-c\), and
\[
      2\,4^{-D/2}\frac{\pi^n}{n!}
             \le C_{n,h}\le\frac{\pi^{n+1}}{n!}.
\]
\end{theorem}
\begin{proof}
Decompose \(g=n_0 a m k_0\), where \(k_0\) preserves the
signed Euclidean metric, \(m\in U(M)\), and \(a\) multiplies
\(u',v'\) by \(\lambda,\bar\lambda^{-1}\). To see the
decomposition, the maximal compact group acts transitively
on isotropic lines: the positive and negative parts of a
unit isotropic vector have equal lengths and can be rotated
independently. After removing a right compact element, \(g\)
thus belongs to the parabolic stabilizing the first line.
Its diagonal blocks have the indicated Levi form, and its
remaining factor is the displayed radical. Right translation
of the integration variable absorbs \(n_0\) by invariance
of the unipotent measure. Now put \(R=|\lambda|^{-1}\).

The first vector under \(\gamma'\) goes to
\(u'+(c/2)v'\). Applying \(n(w,z)^{-1}\) gives
\((1+(c/2)\bar z)u'-(c/2)w+(c/2)v'\).
Write \(c/2=iC\) with real \(C\ne0\). Remove \(k_0\) by
Euclidean norm invariance and make the changes of variables
\(w\mapsto m^{-1}w\), \(t\mapsto t+1/C\).
They have Jacobian one, since a unitary matrix has complex
determinant of modulus one. The squared norm after applying
\(g^{-1}\), in the resulting variables, is
\[
  C^2\{R^{-2}+|w|^2+R^2(t^2+h(w,w)^2/4)\}.
\]
Substitute \(w=w_1/R,\ t=t_1/R^2\). The Jacobian is
\(R^{-2n-2}\), and taking the power \(-D/2\) supplies the
factor \(R^D\). These powers cancel at \(D=2n+2\).

On \(|w|\le1,\ |t|\le1\) the denominator before taking
the power is at most four. This gives the lower bound.
For the upper bound omit the nonnegative term \(h(w,w)^2/4\).
The integral of \((1+|y|^2)^{-n-1}\) in \(2n+1\) real
variables is \(\pi^{n+1}/n!\): insert the gamma integral
for the power, integrate the Gaussian in \(y\), and use
\(\Gamma(1/2)=\sqrt\pi\). This proves finiteness and the
stated bound.
\end{proof}

\subsection{A full-flag estimate at small real parameter}\label{nav:VII-52}
Write \(\nu_\tau(g)=\inf_{b\in B_\tau}\|bg\|_{\rm op}\)
and \(V(g)=\prod_\tau\nu_\tau(g_\tau)\), as in
Proposition~\ref{prop:weight163}. For a nonzero section put
\[
 {\cal M}_a(W)=
 \frac{\int_{B_\infty\backslash G_\infty}\|W(g)\|^2
                         V(g)^a\,d\dot g}
      {\int_{B_\infty\backslash G_\infty}\|W(g)\|^2\,d\dot g}.
\]
In the notation of Theorem~\ref{thm:minor284} let
\(A_*=(K_*+1)2^{n+1}\pi^n/d!\), and keep the real Haar
constant \(C_E\) of Proposition~\ref{prop:measure283}.

\begin{theorem}[A bound uniform down to zero real parameter]\leavevmode
\label{thm:norm284}
Let \(iC_\tau=c_\tau/2\), \(\delta_\tau=|C_\tau|>0\).
For the lowest-type sections of the prescribed word,
\[
 |\mathcal I(c;W)|\ \le\ \mathcal I^+(c;W)
 \ \le\ C_E A_*^f\,{\cal M}_{2K_*+1}(W)
      \prod_\tau\delta_\tau^{-(n+1)}
               (1+2\log^+(1/\delta_\tau))^{s-1}.
\]
The moment is finite, and it is common to all \(c\).
\end{theorem}
\begin{proof}
Start at one real place, suppressing \(\tau\).
The matrix norm is at least one on \(G\) and is unchanged by
inversion. The minimum in \(\nu\) is attained. Indeed a
bounded sequence in \(Bg\) has a convergent subsequence in
the ambient matrix space. The relation defining the unitary
group survives the limit and implies invertibility, and the
subgroup \(B\) is closed. In a fiber \(J\backslash B=U\)
we may therefore use a representative with \(\|g\|=\nu(g)\).
Changing this representative only translates \(U\) or
conjugates it by \(J\), preserving measure. The latter
group commutes with \(\gamma(C)\).

For \(y=\gamma(C)^{-1}n(w,z)g\) with \(\nu(y)\le R\), choose
\(b_0\in B\) with \(\|b_0y\|\le R\). Then
\(n(w,z)^{-1}\gamma(C)b_0^{-1}=g(b_0y)^{-1}\)
has norm at most \(R\nu(g)\). Apply
Theorem~\ref{thm:minor284} with exactly this bound.
For \(p=K_*+1\), integrate the resulting volume estimate
against \(pR^{-p-1}\,dR\) over \([1,\infty)\).
Because \(\nu(y)\ge1\), Tonelli gives
\[
 \int_U \nu(y)^{-p}\,dw\,dt
   \le A_*\,\nu(g)^{K_*}\delta^{-(n+1)}
                           (1+2\log^+(1/\delta))^{s-1}.
\]
Here \(p\int_1^\infty R^{K_*-p-1}\,dR=p\).
The same estimate holds with \(-C\).

Return to all places and put \(V_1=V(g)\),
\(V_2=V(\gamma^{-1}g)\). In the integral of the product
of norms on \(J_\infty\backslash G_\infty\) apply
Cauchy--Schwarz with weight \((V_1/V_2)^p\) in the
first squared integral and its reciprocal in the second.
In the first, integrate on \(J_\infty\backslash B_\infty\)
with product measure \(C_E\prod_\tau dw_\tau\,dt_\tau\).
The first section norm and \(V_1\) descend to the base.
The fiber estimates bound the remaining base weight by the
stated product times \(V(g)^{p+K_*}\).
In the other squared integral translate by \(\gamma\);
this centralizes \(J_\infty\) and switches \(C\) to \(-C\).
The two upper bounds coincide. Divide by the diagonal
squared norm to obtain the result. Finally
\(2K_*+1\le 8(f+1)(r+1)^3\). Proposition~\ref{prop:weight163}
bounds the section and proves integrability of its norm
times this power of \(V\).
\end{proof}

\subsection{A height moment independent of the column dimension}\label{nav:VII-53}
Use the positive integral \(\mathcal I^+\) from
Proposition~\ref{prop:scale284}, first at one common real
frequency, with \(\delta_\tau=|c_\tau|/2\). Put
\(D=4s^2+3\), \(p=2s^2+1\), \(A_\sharp=p B_*\),
where \(B_*\) is from Proposition~\ref{prop:shell285}.
For the product height \(V\), write \({\cal M}_D(W)\)
as in Theorem~\ref{thm:norm284}.

\begin{theorem}[An orbital bound of order \(4s^2+3\) in the height]\leavevmode
\label{thm:low285}
For each nonzero lowest-type section one has
\[
 \mathcal I^+(c;W)\le
 C_E A_\sharp^f {\cal M}_D(W)
 \prod_\tau \delta_\tau^{-(r-1)}
              (1+2\log^+(1/\delta_\tau))^{s-1}.
\]
At terminal frequency \(\varpi\), use \(W_1\) in the moment
and replace the product by
\(\prod_\tau (2/|\tau c|)^{r-1}
      (1+2\log^+(2T/|\tau c|))^{s-1}\)
for the original orbit \(\gamma_c\).
\end{theorem}
\begin{proof}
At one place integrate the estimate of
Proposition~\ref{prop:shell285} against
\(pR^{-p-1}\,dR\) for \(R\ge1\). As \(p=K+1\),
this bounds the integral over \(U\) of
\(({\cal A}(g)/{\cal A}(y))^d\nu(y)^{-p}\)
by \(A_\sharp Z(g)^2\nu(g)^K\) times the local parameter
factor. The choice of coset representative is harmless.
Both \({\cal A}\) and \(\nu\) descend to \(B\backslash G\);
a change by \(J\) preserves the fiber integral since it
commutes with \(\gamma\), and the other changes translate
the unipotent variable.

Apply Cauchy--Schwarz on \(J_\infty\backslash G_\infty\)
as in Theorem~\ref{thm:norm284}, now with weight
\[
 \prod_\tau
    \frac{{\cal A}(g_\tau)^d\,\nu_\tau(g_\tau)^p}
         {{\cal A}((\gamma^{-1}g)_\tau)^d\,
                         \nu_\tau((\gamma^{-1}g)_\tau)^p}.
\]
After fiber integration, remove \(C_E A_\sharp^f\) and
the parameter factors. The remaining weight on the base is at most
\[
             \prod_\tau Z(g_\tau)^2\nu_\tau(g_\tau)^{p+K}.
\]
Since \(Z(g)\le\nu(g)\) by Lemma~\ref{lem:negative285},
this weight costs at most \(V^{p+K+2}=V^D\).
In the second squared integral the substitution by \(\gamma\)
gives the same bound with \(-C_\tau\). Normalize by the
diagonal norm. The moment is finite by
Proposition~\ref{prop:weight163}. The assertion in original
frequency follows from Proposition~\ref{prop:scale284},
including its exact quotient-measure factor.
\end{proof}

\subsection{The real Haar measure in the first-plane integral}\label{nav:VIII-49}
Write \(n=r-2\), \(F=E^+\), \(f=[F:\QQ]\), and let
\(\Delta_E,\Delta_F\) denote the absolute discriminants.
In the signed Euclidean metric the first vector \(u=e_++e_-\)
has length \(\sqrt2\). Put
\(u'=u/\sqrt2\), \(v'=(e_+-e_-)/\sqrt2\) at infinity.
They have Euclidean length one and Hermitian pairing one.
Use the signed orthonormal coordinates of the original
space on the remaining block \(M\).

\begin{proposition}[Haar normalization on the first unipotent radical]\leavevmode
\label{prop:measure283}
Give the finite groups \(U_L\) volume one as in
Theorem~\ref{thm:real282} and give \([U]\) volume one. Write
the real unipotent in the \(u',M,v'\) frame as \(n(w,z)\)
of Theorem~\ref{thm:finite281}, with \(z=-h(w,w)/2+it\).
With ordinary complex area and real length at each place,
the resulting real Haar measure is
\[
 C_E\prod_{\tau}dw_\tau\,dt_\tau,\qquad
 C_E=2^{-f}\Delta_E^{-n/2}\Delta_F^{-1/2}.
\]
\end{proposition}
\begin{proof}
Start in the rational integral frame \(u,M,v\) of
Proposition~\ref{prop:plane281}. The quotient of \(U\) by
its center is the additive \(E\)-space \(M\); its center is
the additive trace-zero \(F\)-line. Locally the integral
subgroup maps onto \(L_M\). For \(w\in L_M\) one must solve
\(z+\bar z=-h(w,w)\) integrally. In a split algebra one
can prescribe either summand of the trace. At the other
places two is a unit, by Proposition~\ref{prop:tracelat253}.
The kernel is the trace-zero integer lattice.
Rational surjectivity follows by taking \(z=-h(w,w)/2\).
Product measure in \(w\) and the imaginary part of \(z\)
is Haar at infinity, since translations have triangular
Jacobian with diagonal one. Integration over the central
quotient and then the additive quotient of \(M\) multiplies
the two covolumes. Indeed their rational fundamental sets
lift to a fundamental set for \(U\), by first reducing \(w\)
and then its central coordinate.

The additive covolume of \(M\), with its finite lattices
assigned volume one, is \((2^{-f}\sqrt{\Delta_E})^n\).
For the standard integer lattice this is the discriminant
formula in complex area. The change to the terminal graph
lattice in the middle block acts on the two split halves
by a matrix and its inverse dual, hence changes the product
of additive volumes by one. The other finite lattices
in this block are the original ones.

Write \(E^-=F\sqrt{-d_0}\). The valuations of \(d_0\) in
\(F\) are even because it is a square up to sign in \(E\)
and \(E/F\) is unramified at every finite place. Thus there
is an integral ideal \(\mathfrak a\) of \(F\) with
\(\mathfrak a^2=d_0\mathcal O_F\). Testing valuations gives
\(\mathcal O_E\cap E^-=\sqrt{-d_0}\,\mathfrak a^{-1}\).
In imaginary-part coordinates its covolume is
\(d_0^{f/2}\sqrt{\Delta_F}/N\mathfrak a=\sqrt{\Delta_F}\).
The normalization in rational coordinates is the reciprocal
of the product of the two covolumes.

The opposite vector \(v\) of the integral frame differs from
\(v'/\sqrt2\) by a purely imaginary multiple of \(u\).
The upper unipotent changing one to the other is central
in \(U\), so this replacement leaves its coordinates unchanged.
Rescaling the first vector to \(u'\) and the opposite vector
to \(v'\) then rescales \(w\) by \(\sqrt2\) and \(z\) by two.
Their real Jacobian is \(2^{n+1}\) at each place. Dividing
the rational-coordinate normalization by \(2^{(n+1)f}\)
gives the displayed value.
\end{proof}

\subsection{The terminal minor and the size of a real unipotent fiber}\label{nav:VIII-50}
At a real place use the Euclidean unit Witt pair \(u',v'\) of
Proposition~\ref{prop:measure283}. In their perpendicular use
Euclidean unit Witt pairs \(p_i,q_i\), \(1\le i\le h=s-1\),
for the remaining isotropic flag, the negative unit terminal
vector \(e\), and the positive column of dimension \(k\).
These vectors form a Euclidean orthonormal basis, with
Hermitian pairings \(h_V(p_i,q_i)=1\). We use \(h_V\)
conjugate linear in its first argument. Set \(d=k+1\) and
\(n=2h+d=r-2\). Write
\[
 w=\sum_{i=1}^h(a_i p_i+b_i q_i)+x+\eta e,\qquad
 z=-h_V(w,w)/2+it .
\]
Here \(x\) belongs to the positive column and \(t\in\RR\).
Use ordinary real length and complex area. For a real
\(C\ne0\), let \(\gamma(C)\) fix the middle block and \(v'\),
and send \(u'\) to \(u'+iC v'\). Put \(\delta=|C|\),
\(\log^+ a=\max(0,\log a)\), and \(K_*=1+(s+1)n\).

\begin{theorem}[The full decorated flag on a real fiber]\leavevmode
\label{thm:minor284}
For \(H\ge1\), the set of \((w,t)\) for which
\(\|n(w,z)^{-1}\gamma(C)b\|_{\rm op}\le H\) for some \(b\in B\)
has volume at most
\[
  \frac{2^{n+1}\pi^n}{d!}\,H^{K_*}\delta^{-(n+1)}
              (1+2\log^+(1/\delta))^h .
\]
\end{theorem}
\begin{proof}
Put \(g=n(w,z)^{-1}\gamma(C)\).
In the block order \(u',M,v'\), the first column of \(g\)
is \((A,-iCw,iC)\) with \(A=1+iC\bar z\).
For \(p=p_i\) or \(p=e\) one has
\(gp=p+h_V(w,p)u'\).
Every \(b\in B\) adds earlier vectors to later ones in
the list \(u',p_1,\ldots,p_h,e\), with all diagonal entries
one. Consequently it fixes every initial exterior product.
For the first \(j+1\) columns, \(0\le j\le h\), take the
coefficient of that product along the same Euclidean unit
exterior vector. It is
\[
 D_j=A+iC\sum_{i=1}^j\bar b_i a_i,\qquad |D_j|\le H^{j+1}.
\]
Indeed the indicated submatrix has lower right block \(I_j\),
first row \((A,\bar b_1,\ldots,\bar b_j)\) and lower first
column \(-iC(a_1,\ldots,a_j)^{\rm t}\). The exterior norm
of \(gb\) in this degree is at most \(H^{j+1}\).
Appending \(e\) changes the coefficient to
\(D_e=D_h-iC|\eta|^2\), with \(|D_e|\le H^{s+1}\).
Since
\[
 h_V(w,w)=2\operatorname{Re}\sum_i\bar b_i a_i
                      +\|x\|^2-|\eta|^2,
 \qquad
 \operatorname{Im}D_e=-\frac{C}{2}(\|x\|^2+|\eta|^2),
\]
the full list bounds the squared length in \(d\) complex
coordinates by \(P=2H^{s+1}/\delta\). The volume of this ball
is \(\pi^d P^d/d!\), by polar integration.

For each \(i\), \(D_i-D_{i-1}=iC\bar b_i a_i\) bounds
\(|a_i b_i|\) by \(P\). The first column gives
\(|a_i|,|b_i|\le H/\delta=L_0\). One pair with these three
bounds has volume at most
\(\pi^2P^2[1+2\log^+(L_0^2/P)]\).
If \(L_0^2\le P\) the two disks suffice. Otherwise integrate
\(\pi\min(L_0^2,P^2/\rho^2)\) against \(2\pi\rho\,d\rho\)
from zero to \(L_0\), splitting at \(P/L_0\).
As \(s\ge2\), \(L_0^2/P\le1/\delta\).
Finally, \(\operatorname{Re} A=1+Ct\) puts \(t\) in an
interval of length \(2H/\delta\). Multiplication of these
volume bounds gives the claim. In particular, indefinite
cross terms of the middle space do not cost an extra power
of \(1/\delta\) in each positive column coordinate.
\end{proof}

\subsection{A negative test vector for the terminal minor}\label{nav:VIII-51}
Use the real coordinates of Theorem~\ref{thm:minor284}.
Let \(P\) be the exterior product of the \(s\) Euclidean
unit vectors in the isotropic flag, in order. Put
\[
 \alpha(g)=\|g^{-1}P\|,\qquad
 {\cal A}(g)=\|g^{-1}(P\wedge e)\|,\qquad
 Z(g)={\cal A}(g)/\alpha(g).
\]
Each norm is calculated in its exterior power. All three
functions are invariant under \(B\) on the left. In particular
\(Z(g)\) is the distance from \(g^{-1}e\) to the span of
the transformed flag.

\begin{lemma}[An ellipsoid from a negative unit vector]\leavevmode
\label{lem:negative285}
In every \(B\)-coset there is a representative \(g\) of
norm \(\nu(g)\), a positive-column vector \(v_+\), and
\(a\ge1\) with \(\|v_+\|^2=a^2-1\), \(Z(g)^2=2a^2-1\),
such that for every \(w,t\), \(y=\gamma(C)^{-1}n(w,z)g\)
satisfies
\[
 \left\|x-\eta v_+/a\right\|^2+|\eta|^2/a^2
       \le \frac{2\sqrt2}{|C|}\,
                            \frac{{\cal A}(y)}{{\cal A}(g)} .
\]
Here \(x,\eta\) are the positive and terminal coordinates of \(w\).
\end{lemma}
\begin{proof}
The maximal compact group is transitive on the isotropic
flags of the given dimensions. Indeed the projections of
an isotropic subspace into the positive and negative signed
spaces are injective and give an isometry between their
images. Orthonormal bases adapted to the projected flags
extend to orthonormal bases of the two signed spaces.
Removing a right compact factor thus puts a group element
in the flag parabolic. Its diagonal blocks on the two
isotropic strings are scalars \(\lambda_i\) and
\(\bar\lambda_i^{-1}\) on their successive quotients.
The middle block is \(m_0\in U(k,1)\). Removing the
upper unipotent factor on the left gives a representative
consisting of these diagonal blocks followed by the right
compact factor. This is allowed because that unipotent
radical lies in \(B\). On multiplying on the left by \(B\),
the projections onto the diagonal blocks still have the
same norms, since the possible extra middle factor is
unitary for the Euclidean norm. The block diagonal
representative therefore attains the minimum \(\nu(g)\).

In this representative put \(l_i=|\lambda_i|^{-1}\).
Then \(\alpha(g)=\prod_i l_i\), \(Z(g)=\|m_0^{-1}e\|\).
Multiply \(m_0e\) by a phase to write it as \(v_*=v_++ae\)
with \(a>0\). Its Hermitian norm is \(-1\).
The last diagonal entry of \(m_0^{-1}\) is the conjugate
of that of \(m_0\), by the unitary relation for the signed
form. Thus \(Z(g)^2=2a^2-1\), and \(\|m_0^{-1}v_*\|=1\).

Let \(P^-\) be the ordered product of the \(s\) unit opposite
Witt vectors. Pair \(P^-\wedge v_*\) with
\(n^{-1}\gamma(C)(P\wedge e)\) by the determinant of
Hermitian pairings. Divide this determinant by
\(h_V(v_*,e)=-a\) and call the result \(D_*\).
Using the first column and the other flag columns from
Theorem~\ref{thm:minor284} gives
\[
 D_*=D_h-iC|\eta|^2+iC\bar\eta\,h_V(v_+,x)/a .
\]
Indeed the row for \(v_*\) has first entry
\(-iC h_V(v_*,w)\), middle entries zero and final entry \(-a\).
Expansion using the identity block on the \(p_i\) gives the formula.
Consequently \(-2\operatorname{Im}D_*/C\) equals the
left side of the asserted inequality, by completing the square.

To bound \(D_*\), apply \(g^{-1}\) to both exterior vectors
in the Hermitian pairing. The first then has Euclidean
norm \(1/\alpha(g)\), using \(\|m_0^{-1}v_*\|=1\) and the
opposite scalings. The second has norm \({\cal A}(y)\).
A signed-unitary Hermitian pairing on an exterior power is
bounded by the product of Euclidean norms, since its sign
matrix is Euclidean unitary. It follows that
\(|D_*|\le {\cal A}(y)/(a\alpha(g))\).
Use \(Z(g)/a\le\sqrt2\) to conclude.
\end{proof}

\section{Matrix coefficients and Hermite moments}\label{nav:g-X-2}

\subsection{Integral vectors near a short isotropic line}\label{nav:III-38}
Use the initial unitary group \(H=U(Y)\) of Theorem~\ref{thm:short246}
and its self-dual lattice \(\Lambda=\Lambda_{Y,g}\) at \(g\).
Let \(r=\dim_EY\), \(f=[E^+:\QQ]\), \(d=2fr\), and fix a cyclotomic
modulus \(M\) for \(E\). Write \(D_E=|\operatorname{disc}E|\),
\({\rm rd}(E)=D_E^{1/(2f)}\), and
\(\|v\|^2=\sum_\tau|v|_{\tau,g}^2\).
Regard the lattice in this Euclidean real space by its diagonal
embedding. Let \(q_*\ge1\) be an upper bound for the dimension of
a rational totally isotropic subspace of \(Y\). Let \(\lambda_1\)
be the least nonzero Euclidean length in \(\Lambda\).

\begin{lemma}[Short spanning vectors for a rational flag and its perpendicular]
\label{lem:lattices247}
If \(b_g<1\), then
\(\lambda_1\le2\sqrt{f\,{\rm rd}(E)}\,b_g^{1/f}\).
If \(\lambda_1<1/2\), there are a nonzero totally isotropic
\(E\)-subspace \(W\subset Y\) and \(0<a<a'\le1/2\) with
\[
          a/a'\le(2\lambda_1)^{1/q_*}.
\]
Vectors in \(\Lambda\) of length at most \(a\) span \(W\) over
\(E\), and every integral vector outside \(W\) has length at least
\(a'\). The real space \(W_\infty\) has a spanning set of integral
vectors of lengths at most \(Ma\). If \(dMa\le a'\), the real
space \((W^{\perp_h})_\infty\) also has a spanning set of integral
vectors of lengths at most \(Md^d/a'\).
\end{lemma}
\begin{proof}
Choose a line attaining \(b_g<1\), which exists by the finiteness
argument before Theorem~\ref{thm:packing246}. For a generator \(w\)
let \(\mathfrak c=\mathfrak c_g(w)\). The scalars taking \(w\) into
\(\Lambda\) form \(\mathfrak c^{-1}\). Its covolume in the product of
complex embedding planes is \(2^{-f}D_E^{1/2}/N\mathfrak c\).
Consider the product of disks in these planes with radii
\(2\sqrt{{\rm rd}(E)}\,b_g^{1/f}/|w|_{\tau,g}\).
Its volume is \((4\pi)^fD_E^{1/2}/N\mathfrak c\), since
\(\prod_\tau|w|_{\tau,g}=b_g\sqrt{N\mathfrak c}\).
Minkowski's convex body theorem supplies a nonzero scalar of
\(\mathfrak c^{-1}\) in this product. Multiplication of \(w\)
by that scalar proves the first bound.

List the distinct lengths below \(1/2\) at which the \(E\)-span
of the integral vectors up to that length increases. The first is
\(\lambda_1\). Every vector in this list's spans is generated by
vectors whose lines have height less than one, since their content
is integral and their norm at every real block is less than \(1/2\).
All these spans are totally isotropic by Theorem~\ref{thm:short246}.
There are consequently at most \(q_*\) lengths in the list. Adjoin
\(1/2\) at the end. The product of each length divided by its
successor is \(2\lambda_1<1\). One ratio \(a/a'\) is at most
\((2\lambda_1)^{1/q_*}\). Take \(W\) to be the span through \(a\).

There are \([E:\QQ]\) linearly independent integers of \(E\)
bounded at every embedding by \(M\). Indeed the relative traces to
\(E\) of the powers of \(\zeta_M\) span \(E\) over \(\QQ\), since the
powers span the cyclotomic field and trace is onto. Each trace is a
sum of at most \(M\) roots of unity. Multiplying an \(E\)-basis
of the short vectors by these integers gives a real basis of
\(W_\infty\) of lengths at most \(Ma\). Rounding coefficients in
this basis approximates any point of \(W_\infty\) within \(dMa/2\)
by an integral vector of \(W\).

Let \(L\) be the orthogonal Euclidean projection of \(\Lambda\) onto
\(W_\infty^{\perp_{\rm eucl}}\). This is a lattice because \(W_\infty\)
is spanned by lattice vectors. A basis of the saturated subgroup
\(\Lambda\cap W_\infty\) extends to a basis of \(\Lambda\).
Lift a nonzero vector of \(L\) to \(\Lambda\), and subtract the
integral approximation to its component in \(W_\infty\).
The resulting vector is outside \(W\) and has length at least \(a'\).
If \(dMa\le a'\), the triangle inequality shows that the projected
vector has length at least \(a'/2\).

For a Euclidean lattice \(L\) of rank \(j\), let
\(\lambda_i(L)\) be its successive minima for the unit ball.
One has \(\lambda_1(L)\lambda_j(L^*)\le j^j\). Indeed
\(\lambda_i(L)\lambda_{j+1-i}(L^*)\ge1\), since the indicated independent
vectors on the two sides cannot all pair to zero, and their pairings
are integers. Minkowski's second theorem bounds the product of
all these inequalities' left sides by \(4^j/\operatorname{vol}(B_j)^2\).
The unit ball contains the cube of side \(2/\sqrt j\), giving \(j^j\).
Thus \(L^*\) has independent spanning vectors of lengths at most
\(2d^d/a'\).

In the adapted frames let \(J\) be the orthogonal signature operator,
so that half the trace of \(h\) on rational vectors is the Euclidean
pairing with \(J\) in one slot. Self-duality over \(\mathcal O_E\)
and the definition of the codifferent \(\mathfrak D_E^{-1}\) give
\(\Lambda^*=2J(\mathfrak D_E^{-1}\Lambda)\). Indeed the trace dual
condition against all multiples in \(\mathcal O_E\) says precisely
that the Hermitian pairings lie in the codifferent. The Euclidean
dual of \(L\) is \(\Lambda^*\cap W_\infty^{\perp_{\rm eucl}}\).
Consequently applying \(J/2\) to its spanning vectors gives
vectors of \(\mathfrak D_E^{-1}\Lambda\) perpendicular to \(W\)
for \(h\). Here orthogonality to the whole realification of an
\(E\)-subspace tests both real and imaginary parts in every block.

Finally \(M\mathfrak D_E^{-1}\subset\mathcal O_E\).
The different of \(\QQ(\zeta_M)\) is generated by
\(\Phi_M'(\zeta_M)\), since \(\ZZ[\zeta_M]\) is its full integer
ring. Differentiation of \(T^M-1\) shows that this generator divides
\(M\). Transitivity of the different and contraction of extended
ideals give \(M\mathcal O_E\subset\mathfrak D_E\).
Multiplying the preceding vectors by \(M\) gives the asserted
integral spanning set of \((W^{\perp_h})_\infty\).
\end{proof}

\subsection{The second moment transverse to the stabilizer}\label{nav:III-39}
Use the real groups, parameters and Haar measures of
Theorems~\ref{thm:tensor249} and~\ref{thm:capelli250}. All inner products
in this argument are linear in their first variable. Use Euclidean length
in signed orthonormal coordinates and \(W(g)=\|g\|_{\rm op}\). Endow
\(H\backslash G\) with the quotient of the chosen Haar measures.

\begin{theorem}[Detection of the discrete summand in a bounded region]
\label{thm:transverse250}
Suppose the isometric embedding \(j:\sigma\to\pi|_H\) has
multiplicity one. For all unit vectors \(v_0,w\) in the respective
lowest compact types there is an absolute constant \(C\) such that
\[
 d_\pi\int_G |g^{-1}z|^2 |\langle\pi(g)v_0,jw\rangle|^2\,dg
      \le T_0^{C n} .
\]
Consequently, with an absolute \(A\ge1\), there exists \(g_0\) with
\(W(g_0)\le T_0^{A n}\) and \(\|P\pi(g_0)v_0\|\ge T_0^{-A n^3}\).
\end{theorem}
\begin{proof}
Write \(V_i\) for the smooth tensor summand with parameter
\(\lambda-e_i\), and \(\pi_i\) for its representation. Its invariant
signed form restricts to \(\epsilon_i\langle\ ,\ \rangle_i\) with
\(\epsilon_i\in\{1,-1\}\), in the normalization of
Theorem~\ref{thm:tensor249}. The functional \(\ell_i(x)=\langle x_z,jw\rangle\)
is defined on \(V_i\). We first prove that it extends to the Hilbert
completion with squared norm \(|\gamma_i|\). This point is needed because
\(w\) is smooth under \(H\), whereas our squared integral is over \(G\).

Choose a smooth vector \(a\) of \(G\) with \(c=\langle jw,a\rangle\ne0\),
using density. Put \(y_i=P_i I_z a\) and
\(p_w(h)=d_\sigma\overline{\langle\sigma(h)w,w\rangle}\). This
coefficient is integrable by the decay and volume bounds of
Theorem~\ref{thm:tensor249}. Formal degree orthogonality gives a
self-adjoint convolution idempotent. The argument of
Theorem~\ref{thm:compact248} identifies its operator on \(\pi|_H\)
with the orthogonal projection onto \(\mathbb C jw\), by multiplicity
one. Its operator \(Q_i\) on the Hilbert space of \(\pi_i|_H\) is
also an orthogonal projection. No multiplicity assumption for that
restriction is needed. If \(\delta_z\) is the sign of the coordinate
\(z\), then for smooth \(x\in V_i\),
\[
 \begin{aligned}
 \epsilon_i\langle x,Q_i y_i\rangle_i
  &=\int_H p_w(h)[(\pi\otimes (\mathbb C^n)^*)(h)x,I_z a]\,dh\\
  &=\delta_z\langle\pi|_H(p_w)x_z,a\rangle
    =\delta_z c\,\ell_i(x).
 \end{aligned}
\]
In the first line orthogonality of the tensor summands permits replacing
\(y_i\) by \(I_z a\) in the signed pairing. In the second line the
stabilizer fixes the last coordinate and its covector. Both integrals
are absolutely convergent by unitarity in the respective Hilbert
spaces. Taking \(x=y_i\) and using Theorem~\ref{thm:capelli250} gives
\(\epsilon_i\|Q_i y_i\|_i^2=\delta_z\gamma_i |c|^2\).
The preceding identity now represents \(\ell_i\) by a Hilbert vector
of squared norm \(|\gamma_i|\), also when this quantity vanishes.

For the coordinate covector \(e^l\), take the last component of the
translated tensor \(v_0\otimes e^l\). The dual defining action gives
\[
 (g^{-1})_{lz}\langle\pi(g)v_0,jw\rangle
       =\sum_i \ell_i\bigl(\pi_i(g)P_i(v_0\otimes e^l)\bigr).
\]
Cauchy--Schwarz for this sum of \(n\) terms and Schur orthogonality in
each summand bound the squared integral times \(d_\pi\) by
\(n\sum_i (d_\pi/d_{\pi_i}) |\gamma_i|\,\|P_i(v_0\otimes e^l)\|_i^2\).
The degree ratio is at most \(2^{n-1}\) by the Vandermonde argument
in Proposition~\ref{prop:moment249}. The projector estimate on the
lowest compact type in Theorem~\ref{thm:tensor249} gives a squared
norm at most \(T_0^{C n}\), and Theorem~\ref{thm:capelli250} bounds
\(|\gamma_i|\) by \((2B+1)^{n-1}\). Summing over \(l\) proves the
transverse moment estimate.

The unweighted squared integral is \(1/d_\pi\). Take \(R^2=2T_0^{C n}\)
for a constant valid in that estimate. At least half the integral is
supported on \(|g^{-1}z|\le R\). This region is the inverse image
of a set \(\mathcal B_R\) in \(H\backslash G\). On a fiber \(Hg\)
formal degree orthogonality for \(\sigma\) integrates the unweighted
square to \(\|P\pi(g)v_0\|^2/d_\sigma\). It remains to bound the
quotient volume. Each coset in \(\mathcal B_R\) has a representative
with \(W(g)\le2R\). Indeed a compact rotation and a boost in one
signed plane carry \(z\) to \(g^{-1}z\), as in the proof following
Theorem~\ref{thm:ratio248}. The norm of that isometry and of its inverse
is at most \(2|g^{-1}z|\). In a definite space a rotation suffices.
For every such representative, \(U_{1/4}g\subset\{W\le3R\}\), where
the small ball is taken in \(H\). Quotient integration and the volume
bounds in Theorem~\ref{thm:tensor249} give
\[
 \frac{d_\pi}{d_\sigma}\operatorname{vol}(\mathcal B_R)
 \le\frac{d_\pi\operatorname{vol}_G\{W\le3R\}}
           {d_\sigma\operatorname{vol}_H(U_{1/4})}
 \le T_0^{C n^3} .
\]
Here \((3R)^{4n^2}\le T_0^{C n^3}\). Integrating the fiber identity
over \(\mathcal B_R\) now shows that at some coset
\(\|P\pi(g)v_0\|^2\ge (4T_0^{C n^3})^{-1}\), since a smaller
value throughout would contradict the half-mass lower bound. This
norm is invariant under changing representatives on the left by \(H\).
Choose the representative of norm at most \(2R\) and enlarge the absolute
constant in the two asserted estimates.
\end{proof}

\subsection{The norm of the initial Poincar\'e intertwiner}\label{nav:IV-37}
Fix a pure unit \(v_0\) in the lowest compact type of the physical
core discrete series \(\pi\). Write
\(\ell_a(F)=\mathcal B(R(a)F,u)\) on \(H(v_0)\), where
\(a\in G_\infty\) and \(u\) is the prescribed unit
lowest-type row in Theorem~\ref{thm:first250}. Let \(d_\pi\)
and \(d_\sigma\) be the product real formal degrees of the
core and row respectively. Use product Haar measures, counting
measure on rational points, and keep the original
\(K_f=K(\mathfrak m)\) and \(J=\iota^{-1}K_f\).

\begin{proposition}[A translated first restriction with an absolute lower bound]\leavevmode
\label{prop:poincarenorm261}
There is \(a\in G_\infty\) satisfying
\(W(a_\tau)\le T_0^{C n^3}\) for all \(\tau\) such that
\[
 \|\ell_a\|_*\ \ge\
 \frac{\operatorname{vol}(J)}{d_\sigma}
 \left(\frac{d_\pi}{\operatorname{vol}(K_f)}\right)^{1/2}
       \|u\|_\infty\,\exp[-P_{\psi,\eps}(x)] .
\]
Here a polynomial depending only on the fixed family and \(\eps\)
works uniformly on the original intervals, for every such row.
\end{proposition}
\begin{proof}
Choose \(g\) with \(|u(g)|>\|u\|_\infty/2\), and put \(c=\iota(g)\).
The Poincar\'e series in Theorem~\ref{thm:first250} is \(Y_b(b)\)
in the notation of Proposition~\ref{prop:corecusp251}. We
record the size of the compact averaging used to form \(b\).
The transverse estimate in Theorem~\ref{thm:transverse250}
starts at any prescribed local lowest unit in the conjugate
core, so its conjugate in the physical core can be chosen to
be the local factor of \(v_0\). The direction selection in
Theorem~\ref{thm:direction249} translates it by an element of
norm at most \(T_0^{C n^3}\). Compact translates have norm one.
Finally the compact projection in Theorem~\ref{thm:compact248}
integrates against a matrix unit in the lowest row type \(S\).
For unit \(e,w\in S\), its density is
\((\dim S)\overline{\langle\sigma(k)e,w\rangle}\) in a
choice of coefficient convention and has \(L^1\) norm at most
\(\sqrt{\dim S}\) by compact orthogonality. The mixed matrix
unit in the proof of Theorem~\ref{thm:direction249} permits an
arbitrary prescribed lowest row vector. Complex conjugation
preserves these bounds. Thus in the physical core one has
\[
 b=\int\pi(a)v_0\,d\nu(a),\qquad
 \|\nu\|_{\rm TV}\le T_0^{Cfn^2},\qquad
 \sup_{a\in\operatorname{supp}\nu}W(a_\tau)\le T_0^{C n^3}
\]
with compact support. For the variation bound we used
\(\log\dim S\ll fn^2\log T_0\), as follows from the compact
Weyl dimension formula in the proof of Theorem~\ref{thm:first250}.

We bound the norm of \(Y_b(v_0)\). In this calculation use
the coefficient convention of Proposition~\ref{prop:corecusp251}
and put \(f_0(t)=\langle\pi(t)v_0,b\rangle\). This Poincar\'e
series has just one contributing finite double coset. Its squared
norm is the level volume times the integral of the squared
periodization of \(f_0\) on \(\Gamma_c\backslash G_\infty\).
Indeed at finite representative \(c_f\) the contributing
rational points are exactly those defining \(\Gamma_c\);
translating the real variable by \(c_\infty^{-1}\) gives the
stated quotient. Haar measure on the finite right coset
contributes \(\operatorname{vol}(K_f)\), with rational
stabilizers counted in the quotient by \(\Gamma_c\).
Opening the square and unfolding one of the two sums gives
\[
 \|Y_b(v_0)\|_2^2=\operatorname{vol}(K_f)
   \sum_{\gamma\in\Gamma_c}\int_{G_\infty}
                f_0(t)\overline{f_0(\gamma t)}\,dt .
\]
For absolute convergence let \(L=8n^2+2n+9\) be the decay
exponent of Theorem~\ref{thm:tensor249}. Since \(v_0\) is a
pure lowest tensor and \(\nu\) has the displayed support and
variation, that theorem and submultiplicativity give
\(|f_0(t)|\le T_0^{C f n^5}w(t)^{-L}\).
Now \(w(t)w(\gamma t)\ge w(\gamma)\). For \(p=4n^2+2\) the
product of absolute integrands is bounded by
\(T_0^{C f n^5}w(\gamma)^{-p}w(t)^{-(L-p)}\).
The last weight is integrable and the sum of the preceding
weights is finite by Lemma~\ref{lem:returns261}. This
justifies unfolding.

Formal-degree orthogonality gives the absolute value of each
integral as \(d_\pi^{-1}|\langle\pi(\gamma)b,b\rangle|\).
Apply the same decay estimate on both compact translates in
the two instances of \(b\). It gives
\(|\langle\pi(\gamma)b,b\rangle|\le T_0^{C f n^5}w(\gamma)^{-L}\),
again using \(W(a_\tau)\le T_0^{C n^3}\) at each end. The sum
bound of Lemma~\ref{lem:returns261} therefore yields
\[
 \|Y_b(v_0)\|_2\le
 \left(\frac{\operatorname{vol}(K_f)}{d_\pi}\right)^{1/2}
                     \exp[P_1(x)]
\]
for a polynomial \(P_1\) uniform at the chosen points.

By Theorem~\ref{thm:corecopy251}, \(Y_b\) extends to a
bounded intertwiner and \(Y_b(v_0)\) belongs to \(H(v_0)\).
The compact integral of its translates by \(\nu\) is exactly
\(Y_b(b)\). Its pairing with \(u\) is the integral of
\(\ell_a(Y_b(v_0))\) against \(\nu\). To see this directly,
the compact translates of the smooth bounded cusp form are
uniformly bounded on the automorphic quotient and \(u\) is
integrable on \([H]\). Fubini gives the assertion. Theorem
~\ref{thm:first250} bounds this last pairing below in modulus by
\(\operatorname{vol}(J)d_\sigma^{-1}\|u\|_\infty
\exp[-Cfn^5\log b]\). It follows that \(Y_b(v_0)\ne0\).
The continuous function of \(a\) has a maximum in modulus
on the compact support, at least the pairing modulus divided
by \(\|\nu\|_{\rm TV}\). Divide by \(\|Y_b(v_0)\|_2\)
and combine the displayed estimates to obtain the assertion.
\end{proof}

\subsection{The height at a large value of a cusp form}\label{nav:VI-35}
Use the lattice, ranks and modulus \(M\) of Lemma~\ref{lem:lattices247}.
Let \(m_0\ge1\) be a rational integer. On \(H=U(Y)\) let
\(K_H(m_0)\) be the stabilizer of the original finite lattice acting
trivially modulo \(m_0\). For a tuple of real Lie algebra elements
\(A=(A_\tau)_\tau\) put \(\|A\|_1=\sum_\tau\|A_\tau\|_{\rm HS}\),
where the Hilbert--Schmidt norms use the adapted frames at the identity.
A cusp form has integral zero along the unipotent radical of each proper
rational parabolic subgroup, at every right translate.

\begin{theorem}[A uniform height bound from cuspidality]
\label{thm:cusp247}
Let \(u\) be a nonzero bounded smooth cuspidal automorphic function on
\(H\), right invariant under \(K_H(m_0)\), and put \(U=\|u\|_\infty\).
Suppose \(|u(s\exp A)-u(s)|\le L\|A\|_1\) for all adelic \(s\) and
all real \(A\). Write \(\kappa=m_0M^2(2d)^{d+5}\), where
\(d=2f\dim_EY\). At any point with \(|u(g)|>U/2\),
\[
  b_g\ \ge
   \left(\frac{[4\kappa(1+L/U)]^{-q_*}}
                     {4\sqrt{f\,{\rm rd}(E)}}\right)^f .
\]
The bound uses the original field, lattice and level of the form.
\end{theorem}
\begin{proof}
Set \(D_0=[4\kappa(1+L/U)]^{-1}\). The case \(b_g=1\) is immediate.
For \(b_g<1\), Lemma~\ref{lem:lattices247} shows it suffices to prove
\(2\lambda_1\ge D_0^{q_*}\). Suppose otherwise, and put
\(\theta=(2\lambda_1)^{1/q_*}<D_0\). The lemma gives an isotropic
subspace \(W\) with \(a/a'\le\theta\), so \(dMa\le a'\). Let \(N\)
be the unipotent radical of the proper parabolic stabilizing \(W\) in
\(H\). We will show that every point of \(N(\mathbb A_F)g\), up to
left rational translation and the right finite level, can be reached
from \(g\) by two real exponentials of total norm at most
\(\kappa\theta\).

Choose a rational Witt decomposition \(Y=W\oplus D\oplus W'\), with
\(W'\) isotropic and dual to \(W\) under \(h\). Denote by
\(\mathfrak z\) the central Lie subspace of maps supported in the
block \(W'\to W\) and skew adjoint there. The elements of \(N\) act
as the identity on \(W\), \(W^{\perp_h}/W\) and \(Y/W^{\perp_h}\).
In this order of the three blocks their matrices are upper unitriangular.
In the Lie algebra the \(W'\to D\) block is determined by the negative
adjoint of the \(D\to W\) block. The quotient by \(\mathfrak z\) is
therefore the additive space of the latter blocks, and any product of
three strictly upper block matrices is zero.

For integral \(w\in W\) and \(v\in W^{\perp_h}\) define
\[
          X_{w,v}(s)=w h(v,s)-v h(w,s).
\]
This is a rational, integral, skew-adjoint endomorphism in the Lie
algebra of \(N\). Its square is \(-h(v,v)w h(w,\cdot)\).
Use for \(w\) the spanning vectors of length at most \(Ma\) and for
\(v\) those of length at most \(Md^d/a'\) from the lemma. Their
operators span the quotient over \(\mathbb R\). Indeed on \(D\) the
first block is \(w h(v,\cdot)\) and these rank-one maps span all
complex matrices in each real block. The spanning sets cover the
realifications, so real bilinearity also isolates the individual
places. Taking both vectors from the short spanning set in \(W\)
gives operators spanning \(\mathfrak z_\infty\). In bases of \(W\)
and \(W'\) paired by \(h\), these blocks are \(w v^*-v w^*\), which
span the skew-Hermitian matrices. Choosing coordinate vectors and
their imaginary multiples gives the diagonal and both real directions
of every off-diagonal pair. Select a basis among the images for each
of the two spaces. Each selection has at most \(d^2\) members.

For each chosen \(X\), use the element \(t\in\mathcal O_E\) with
\(t+\bar t=1\) from Theorem~\ref{thm:short246} to form
\[
           \gamma_X=1+m_0X+t m_0^2 X^2 .
\]
Since \(X^*=-X\) and \(X^3=0\), multiplication by the adjoint gives
the identity. Both \(\gamma_X\) and its adjoint preserve \(\Lambda\)
and act trivially modulo \(m_0\Lambda\). Hence \(\gamma_X\in N(F)\)
and its finite component belongs to \(g_f K_H(m_0)g_f^{-1}\).
Its logarithm is \(m_0X+(t-1/2)m_0^2X^2\). The second term belongs
to \(\mathfrak z\), as it is skew adjoint and supported in the last
block. For a centrally supported generator \(X^2=0\).

Take any \(v\in N(\mathbb A_F)\). Left multiplication by a rational
point first puts its finite component inside \(g_f K_H(m_0)g_f^{-1}\).
Indeed \(N(F)\) is dense in \(N(\mathbb A_{F,f})\). Exponential and
logarithm here are inverse polynomials of degree at most two on the
strictly upper blocks. They identify the rational Lie algebra and its
finite adeles with the corresponding unipotent groups, preserving
integral coordinates outside finitely many primes. Density follows
from finite adelic approximation in the additive \(F\)-vector space.
Further left multiplication by integer powers of the \(\gamma_X\)
preserves this finite condition. Modulo \(\exp\mathfrak z_\infty\)
the real group is additive. Round the coordinates in the chosen
basis of the elements \(m_0 X\) in that quotient. The resulting real
component has the form \(\exp(A_0)\exp(Z_0)\), where \(A_0\) is a sum
of the corresponding \(m_0 X\)'s with coefficients in \([-1/2,1/2]\)
and \(Z_0\in\mathfrak z_\infty\). Now round in the central basis,
using the integral powers of the central generators. This also puts
every coefficient of \(Z_0\) in \([-1/2,1/2]\) in that basis.
This argument includes the case of a zero-dimensional quotient.

Conjugate \(A_0,Z_0\) by \(g_\infty^{-1}\) and call the results
\(A,Z\). The norm of each conjugated \(X_{w,v}\) is bounded by
\(2\|w\|\|v\|\) for the Euclidean lengths at \(g\), since the
Hilbert--Schmidt norm of a rank-one map is the product of the two
lengths and Cauchy--Schwarz applies in the sum over places. Thus
\[
   \|A\|_1+\|Z\|_1
      \le d^2m_0M^2(d^d a/a'+a^2)\le\kappa\theta .
\]
Here \(a^2\le a/a'\) as \(a,a'\le1/2\). Left automorphy and right
finite invariance now give \(u(vg)=u(g\exp A\exp Z)\). The Lipschitz
bound at two successive points yields \(|u(vg)-u(g)|\le L\kappa\theta\).
The quotient \(N(F)\backslash N(\mathbb A_F)\) is compact, as follows
from its central filtration by additive adelic quotients. Integrating
against Haar probability and using the vanishing constant term gives
\(|u(g)|\le L\kappa\theta<U/2\), a contradiction. Combining the
resulting lower bound for \(\lambda_1\) with the first inequality of
Lemma~\ref{lem:lattices247} proves the assertion.
\end{proof}

\subsection{A derivative bound for the lowest compact type}\label{nav:VI-36}
Give \(U(p,q)\), \(r=p+q\), the left invariant metric determined at
the identity by the Hilbert--Schmidt inner product on its real Lie algebra.
In an automorphic row, fixing all coordinates except this real
factor gives a function on the whole Lie group with the same bound
in supremum norm. Thus an estimate for this metric can be used at
every finite level.

\begin{proposition}[A uniform derivative estimate for a bounded row]
\label{prop:gradient248}
Suppose a bounded smooth automorphic row \(u\) at every real place
belongs to a discrete series with Harish-Chandra parameters of
absolute value at most \(B\ge1\), and to its lowest compact type.
Put \(U=\|u\|_\infty\) and \(r=\dim_E Y\). Then at every adelic point
and for every real Lie algebra tuple \(A\),
\[
  |u(g\exp A)-u(g)|
       \le 40\sqrt r\,(B+2r) U\sum_\tau\|A_\tau\|_{\rm HS} .
\]
The same result holds for any single compact type at each place
whose highest-weight coordinates have absolute value at most \(B+r\).
\end{proposition}
\begin{proof}
Write \(\Delta\) for the Laplace--Beltrami operator, with negative
spectrum. In a left invariant orthonormal frame it is the sum of
squares, because the group is unimodular. In the Cartan decomposition
let \(X_i\) be an orthonormal basis of the compact Lie algebra and
\(Y_j\) one of its perpendicular. The trace form
\(\operatorname{Re}\operatorname{tr}(AB)\) has signs minus and plus
on these two spaces. Its Casimir is therefore
\(\Omega=-\sum_i X_i^2+\sum_jY_j^2\), and the compact Casimir is
\(\Omega_K=-\sum_iX_i^2\). Thus \(-\Delta=2\Omega_K-\Omega\).
Fix all other arguments of the row. Its function on this group
satisfies \(-\Delta u=\lambda u\), where \(\lambda\ge0\) since the
sum of negatives of squares acts positively in the unitary series.

In the diagonal coordinates for this trace form, the Casimir on
infinitesimal character \(\nu\) is
\(\sum_i\nu_i^2-\sum_i\rho_i^2\), with \(\rho\) half the sum of
positive roots. This follows from the Harish-Chandra isomorphism for
the quadratic Casimir. On the compact type of highest weight \(\mu\)
one uses \(\nu=\mu+\rho_K\), separately on \(U(p)\) and \(U(q)\).
For the lowest compact type, the discrete-series lowest-weight
formula is \(\mu=\nu_G+\rho_{\rm nc}-\rho_K\), in the chamber of
its Harish-Chandra parameter \(\nu_G\). Each coordinate of
\(\rho_{\rm nc}-\rho_K\) has absolute value at most \((r-1)/2\),
as each of the other coordinates contributes a single signed half.
Thus \(|\mu_i|\le B+r\). The same bound is part of the hypothesis
for the other types. Since every coordinate of a half-sum is bounded
by \(r/2\), the two Casimir formulas give
\(0\le\lambda\le 10r(B+2r)^2\).

We give the metric estimate including its rank dependence. For
left invariant vectors the bracket has norm at most twice the
product of their norms, by matrix multiplication. The Koszul
formula therefore gives \(\|\nabla_XY\|\le3\|X\|\|Y\|\). Substituting
in the definition of curvature gives
\(\|R(X,Y)Z\|\le24\|X\|\|Y\|\|Z\|\). Taking a trace in dimension
\(r^2\) yields \(\operatorname{Ric}\ge-K\) with \(K=24r^2\).
This invariant metric is complete.

Let \(P_t\) be its heat semigroup with generator \(\Delta\).
The local reverse Poincar\'e inequality under \(\operatorname{Ric}\ge-K\)
reads, for a bounded real function \(a\),
\[
 |\nabla P_t a|^2\le
   \frac{K}{1-e^{-2Kt}}\big(P_t(a^2)-(P_ta)^2\big).
\]
For clarity, the semigroup gradient inequality obtained from Bochner's
formula is \(P_s|\nabla v|^2\ge e^{-2Ks}|\nabla P_s v|^2\).
Differentiate \(P_s((P_{t-s}a)^2)\), apply the gradient inequality
to \(v=P_{t-s}a\), and integrate \(0<s<t\). This gives the displayed
formula for compactly supported smooth data. On a complete manifold
with Ricci curvature bounded below the heat semigroup is conservative.
Smooth compact cutoffs of a bounded smooth function converge under
\(P_t\) locally with first derivatives for \(t>0\), by bounded
convergence for the heat kernel and interior parabolic estimates.
This gives the same formula here. This semigroup interpolation is also given in
\href{https://doi.org/10.1112/jlms.12452}{Kopfer
and Sturm, Section 2.1}.

For a bounded real eigenfunction of \(-\Delta\) of eigenvalue
\(\lambda\) one has \(P_ta=e^{-t\lambda}a\). For example, apply
It\^o's formula to \(e^{s\lambda}a(B_s)\) for the nonexplosive
diffusion with generator \(\Delta\). The stopped local martingale is
bounded on every finite time interval, so passage to the limit
preserves its expectation. Take \(t=(1+\lambda+K)^{-1}\).
Convexity of \(e^{-2s}\) on \([0,1]\) gives
\(1-e^{-2Kt}\ge (1-e^{-2})Kt\). Hence
\(\|\nabla a\|_\infty\le 3\sqrt{1+\lambda+K}\,\|a\|_\infty\).
For a complex row take the real part after multiplication by a unit
scalar chosen to make its derivative in a given direction real.
The bound for \(\lambda\) gives \(3\sqrt{1+\lambda+K}\le
40\sqrt r(B+2r)\). Integrate along the exponential path at one
place and vary the places successively. This proves the assertion.
\end{proof}

\subsection{A Sobolev bound for both terminal-period transfers}\label{nav:VI-42}
Fix an edge from \(i\) to \(i+1\) in the construction of
Theorem~\ref{thm:fullperiod266}, with \(i\ge2\). Fix the input
cusp forms \(B\) and \(A\) on the horizontal source and the
next core, respectively. Keep a single allowed original
finite test and, on a Jacobi edge, its distinguished line test.
Write \(P(\phi)\) for the outgoing complex period as a
functional of the real horizontal test in the original
polarization. In real position coordinates at all infinite
places together let \(d\) be the dimension, \(h_\nu\) the
normalized product Hermite functions, and \(\mathcal N\) the
number operator with eigenvalue \(|\nu|=\sum\nu_j\).

\begin{proposition}[An explicit order for the period functional]\leavevmode
\label{prop:order268}
With \(s=2d+2\), for every real Schwartz test,
\[
            |P(\phi)|\le M\|(1+\mathcal N)^s\phi\|_2 .
\]
One may take \(M\) to be the finite lattice sums and integrals
of the absolute input forms given below. The order \(s\)
is independent of their weights, the finite levels and all
previous choices of Hermite degrees. For a prescribed finite
list of marked terms one may use the maximum of their constants.
\end{proposition}
\begin{proof}
Use the defining integral before unfolding, as in
Propositions~\ref{prop:besorbit265} and~\ref{prop:jacorbit266}.
Its variables are \(g\) on the horizontal source, \(h\) on
the core, and \(n\) on the target flag radical. Apart from
unit characters its integrand is the product of \(B(g)A(h)\),
the horizontal theta at \((g,nh)\), and, on a Jacobi edge,
the conjugate line theta. Write \(Q(n,h)\) for the absolute
value of the last factor, and put \(Q=1\) on a Bessel edge.
The flag quotient has a compact fundamental set. For the two
reductive quotients take finitely many Siegel sets and
representatives of the finite components at the given compact
open levels. The finite arguments of the horizontal oscillator
then lie in a finite union of compact sets. On such sets its
action on a fixed finite test takes only finitely many values,
because the test has an open stabilizer. Subdivide the
integration cover accordingly into pieces \(\mathcal D_\lambda\).
On a piece the translated finite test is \(\phi_{f,\lambda}\).
Put
\[
 H_\lambda=\sum_{y\ {\rm rational\ position}}
         |\phi_{f,\lambda}(y)|(1+|y_\infty|)^{-(d+1)} .
\]
It is finite. In fixed rational coordinates the finite
support lies inside a common inverse denominator times the
integral lattice. The number of points of that real lattice
in a ball of radius \(R\ge1\) is \(O_\lambda(R^d)\), by
packing disjoint balls. Summing over dyadic shells proves
convergence.

Let \(S(g,nh)\) be the real symplectic matrix acting on the
horizontal oscillator and set
\(W=1+\|S(g,nh)\|_{\rm op}+\|S(g,nh)^{-1}\|_{\rm op}\).
There is a constant \(c_d\) in normalized real oscillator
coordinates for which the following choice follows from the estimate below.
\[
 M=c_d\sum_\lambda H_\lambda
       \int_{\mathcal D_\lambda}|B(g)A(h)|\,Q(n,h)
                          W^{2d+2}\,dg\,dh\,dn .
\]
Fixed coordinate changes can equivalently be included in the
matrix norms and in \(H_\lambda\). The integrals are finite. On the compact unipotent set the matrix norms have polynomial
growth on the two Siegel sets. The fixed line theta has
polynomial growth there by its lattice sum, as in
Proposition~\ref{prop:jacorbit266}. The two smooth
compact-finite cusp forms have rapid decrease in the
respective unbounded chamber variables. These bounds also
justify using an overlapping integration cover to majorize
the absolute period.

It remains to prove the real estimate leading to this formula.
Put \(a=\lfloor d/2\rfloor+1\). By Fourier inversion and
Cauchy--Schwarz, the sup norm of a Schwartz function is
bounded by a dimension-dependent constant times the sum of
its \(L^2\) derivative norms through order \(a\). Indeed
\((1+|\xi|^2)^{-a}\) is integrable on \(\mathbb R^d\).
Apply this estimate to \(x^\gamma f\) for
\(|\gamma|\le d+1\). Expanding
\((1+\sum |x_k|)^{d+1}\) bounds
\(\sup_x(1+|x|)^{d+1}|f(x)|\) by a fixed sum of \(L^2\)
norms of words of length at most \(d+1+a\le2d+2\) in
coordinates and derivatives.

Take \(f=\omega(S)\phi\). Moving a word through \(\omega(S)\)
replaces each of its generators by a linear combination
of the \(2d\) generators, with coefficients bounded by a
fixed multiple of \(W\). This follows by differentiating the
covariance of the Schr\"odinger representation. Each resulting
word is a fixed combination of words in creation and annihilation
operators. A word of length \(u\) sends \(h_\nu\) to a single
shifted basis element times a scalar of modulus at most
\((1+|\nu|+u)^{u/2}\), or to zero. Distinct surviving indices
remain distinct under that fixed shift. Parseval thus bounds
its norm on \(\phi\) by a dimension-dependent constant times
\(\|(1+\mathcal N)^s\phi\|_2\) for \(u\le s\).
Combining the finitely many estimates defines \(c_d\) and
gives the real supremum bound with the additional factor
\(W^s\). The horizontal rational theta sum on each piece
is at most \(H_\lambda\) times this supremum. Integration
proves the proposition.
\end{proof}

\subsection{A compactly supported operator fixing the incoming row}\label{nav:VI-43}
Let \(G\) be a product of the real unitary groups of a horizontal
row, \(K\) a maximal compact subgroup and \(g=\dim_{\mathbb R}G\).
On its Lie algebra take the invariant real trace form on each
matrix factor. It is negative on \(\mathfrak k\), including
the center, and positive on the noncompact Cartan complement
\(\mathfrak p\). Reversing the sign on \(\mathfrak k\) gives a
positive inner product. Use its left-invariant, right-\(K\)-invariant
metric and its Haar volume. For a different fixed Haar measure
rescale convolution densities accordingly. Let \(\Delta\)
be the nonnegative Laplacian. For an orthonormal basis adapted
to the two summands the associated left-invariant vector fields
give
\[
       \Delta=-\sum_j X_j^2=-\Omega_G+2C_K,\qquad
       C_K=-\sum_{X_j\in\mathfrak k}X_j^2.
\]
Indeed unimodularity makes each first-order divergence zero.
Here \(\Omega_G\) is the Casimir for the signed trace form.

\begin{proposition}[Compact convolution fixing finitely many types]\leavevmode
\label{prop:kernel270}
Let a unitary representation have a fixed infinitesimal
character on its smooth vectors, or be a Hilbert sum of
copies with that character. For a finite nonempty list
\(\mathcal T\) of its \(K\)-types write
\(\lambda_\tau=-\Omega_G(\pi)+2C_K(\tau)\).
Put \(M=\sum_{\tau\in\mathcal T}\dim\tau\),
\(H=1+\max_{\tau\in\mathcal T}\sqrt{\lambda_\tau}\).
There is a smooth \(b\) supported in the fixed compact set
\(\overline{B(1,1)}K\) such that \(\pi(b)\) is the identity
on the listed types. For each integer \(l\ge0\),
\[
              \|b\|_{C^l}\le C_{G,l}M H^{2(l+2g+2)} .
\]
\end{proposition}
\begin{proof}
On a smooth vector \(v\) of type \(\tau\),
\(-\sum d\pi(X_j)^2v=\lambda_\tau v\). Pairing with \(v\)
in the unitary inner product proves \(\lambda_\tau\ge0\).
Choose a fixed even nonnegative smooth function of integral
one supported in \((-1,1)\). Its cosine transform \(m(z)\)
is real, even and rapidly decreasing, and \(m(0)=1\).
Choose \(0<c\le1\) such that \(m(z)\ge1/2\) for \(|z|\le c\).
Let \(R=c/H\) and \(k_R=m(R\sqrt\Delta)\delta_1\).
The metric is complete. By the finite propagation theorem
for the wave equation the kernel here is supported in
\(\overline{B(1,R)}\). Indeed the multiplier is the integral
of the cosine wave operators for times of absolute value
at most \(R\). The kernel is real by this formula.
\href{https://www.ihes.fr/~gromov/wp-content/uploads/2018/08/536.pdf}
{Cheeger, Gromov and Taylor, Section 1}.

We detail its derivative estimate. Local elliptic regularity
on nested coordinate balls about the identity and Sobolev
point evaluation give
\(
 |u(1)|\le C\|(1+\Delta)^{g+1}u\|_2
\)
for smooth compactly supported \(u\). The same inequality
on the operator domain follows by local regularity there,
bounding every lower power in \(L^2\) by the highest shifted
power. Hence \((1+\Delta)^{-(g+1)}\delta_1\in L^2\), by
duality. On the fixed compact ball of radius two, the same
interior estimates followed by Sobolev embedding bound the
\(C^l\) norm of a function supported in the unit ball by
\(C_l\|(1+\Delta)^{l+g+1}u\|_2\). By the spectral theorem,
with \(j=l+2g+2\),
\[
 \|(1+\Delta)^{l+g+1}k_R\|_2
 \le \|(1+\Delta)^{-(g+1)}\delta_1\|_2
       \sup_{z\ge0}(1+z^2)^j|m(Rz)|
 \le C_l R^{-2j}.
\]
This also proves smoothness by taking all \(l\).

For a smooth \(v\) of type \(\tau\) and any smooth dual
Hilbert vector \(w\), \(F(y)=\langle\pi(y)v,w\rangle\),
linear in \(v\), satisfies \(\Delta F=\lambda_\tau F\).
The wave solution with initial data \((F,0)\) is locally
\(\cos(t\sqrt{\lambda_\tau})F\). For its value at the
identity through time \(R\), one may first cut off \(F\)
to compact support with the cutoff equal to one near the
closed radius-\(R\) ball. Uniqueness and finite propagation
give the stated solution at that point. Integrating in time
and using the real self-adjoint kernel gives
\(\int_G k_R(y)F(y)\,dy=m(R\sqrt{\lambda_\tau})F(1)\).
It follows that \(\pi(k_R)\) multiplies the listed type by
this positive scalar. Approximation extends the conclusion
to its full Hilbert isotypic space, since smooth vectors
projected onto a compact type are dense there.

With \(dk\) of mass one on \(K\), take the measure
\[
  d\mu(k)=\sum_{\tau\in\mathcal T}
       \frac{(\dim\tau)\overline{\chi_\tau(k)}}
            {m(R\sqrt{\lambda_\tau})}\,dk .
\]
Schur orthogonality gives the inverse scalar on each listed
type. It also gives \(\int_K|\chi_\tau|\,dk\le1\) by
Cauchy--Schwarz, so \(\|\mu\|\le2M\). The right convolution
\(b=k_R*\mu\) has the required support and fixes the listed
types. Right translations by \(K\) have uniformly bounded
derivatives of every fixed order in a left-invariant frame.
The estimate for \(k_R\) and \(R=c/H\) give the asserted
bound.
\end{proof}

\subsection{An exponential moment for the initial Hermite expansion}\label{nav:VI-44}
On the first even core oscillator at place \(\tau\), let
\(D=n(n+d)\) be the real position dimension, \(L_\tau\) the
energy bound for \(\phi_\tau\), and \(W(a_\tau)\ge1\) the
matrix norm of the real translate in
Lemma~\ref{lem:weightedhermite260}. Put
\[
 S_\tau=(1+16D^3)W(a_\tau)^2,\qquad
 \eta_\tau=\frac{1}{4\mathrm e S_\tau(L_\tau+3)},\qquad
 \eta=\min_\tau\eta_\tau .
\]
Here \(\mathrm e\) is the base of natural logarithms. The
indices below run over all product Hermites for the
orthogonal sum of the bottom row and its positive line at
every place.

\begin{lemma}[Exponential coefficient bound after one real translation]\leavevmode
\label{lem:moment271}
Let \(c_\alpha\) be the coefficients of the single Schwartz
test \(\omega(a)\phi^0\) of Lemma~\ref{lem:project271}.
With total degree \(|\alpha|=\sum_\tau|\alpha_\tau|\),
\[
       \sum_\alpha |c_\alpha|\exp(\eta|\alpha|)
       \ \le\ S,\qquad
       S=\prod_\tau
          2(1-\exp(-2\eta_\tau))^{-D/2}\|\phi_\tau\|_2 .
\]
In this bound each \(W(a_\tau)\) may be replaced by an upper
bound fixed in advance, using the resulting smaller
\(\eta_\tau\) throughout.
\end{lemma}
\begin{proof}
Fix a place, and write \(y=\omega(a_\tau)\phi^0_\tau\).
For every integer \(j\ge1\),
\[
       \|(1+\mathcal N)^j y\|_2
         \le [S_\tau(L_\tau+2j+1)]^j\|\phi_\tau\|_2 .
\]
Indeed, in normalized creation and annihilation operators,
the conjugate of each generator by \(\omega(a_\tau)\) has
sum of absolute linear coefficients at most \(4DW(a_\tau)\).
This is the covariance computation of
Lemma~\ref{lem:weightedhermite260}. Expand
\((1+\sum_l a_l^*a_l)^j\) after substitution. The sum of
the absolute word coefficients is at most \(S_\tau^j\).
The words have length at most \(2j\). Each step changes the
energy by at most one, with operator norm at most the square
root of one plus the current energy bound. Their norms on
\(\phi^0_\tau\) are therefore bounded by
\((L_\tau+2j+1)^j\|\phi_\tau\|_2\).

For \(j\ge1\), \(L_\tau+2j+1\le (L_\tau+3)j\) and
\(j!\ge(j/\mathrm e)^j\). Consequently the series with
terms \((2\eta_\tau)^j(1+\mathcal N)^j y/j!\) converges in
\(L^2\) with sum of norms at most \(2\|\phi_\tau\|_2\),
including the constant term. Each Hermite coordinate of
the sum is that of \(y\) times
\(\exp(2\eta_\tau(1+|\alpha_\tau|))\), by the scalar
exponential series. Cauchy--Schwarz now gives
\[
 \sum_{\alpha_\tau} e^{\eta_\tau|\alpha_\tau|}
                    |\langle h_{\alpha_\tau},y\rangle|
 \le (1-e^{-2\eta_\tau})^{-D/2}\,2\|\phi_\tau\|_2 .
\]
The squared numerical factor before the 2 is the product
of \(D\) geometric series. The archimedean test is a pure
tensor, so its absolute coefficients multiply. Multiplying
these local inequalities and using \(\eta\le\eta_\tau\)
proves the lemma.
\end{proof}

\subsection{Uniform bounds for a discrete-series matrix coefficient}\label{nav:VII-39}
Fix \(n\ge2\), \(B\ge1\), and put \(T_0=10^4 n^4(B+1)\),
\(D=4n^2+2\), \(s=2n+5\), \(L=s+2D\). For \(U(p,q)\) of rank
\(m=p+q\le n\) write \(W(h)=\|h\|_{\rm op}\) in signed orthonormal
coordinates and \(U_\delta=\{h:\|h-1\|_{\rm HS}\le\delta\}\).
Here \(W(h^{-1})=W(h)\ge1\). For a discrete series \(\rho\) let
\(d_\rho\) be its formal degree with respect to the indicated Haar
measure. Constants denoted by \(C\) below are absolute and may
increase between inequalities.

\begin{theorem}[Power decay and Haar bounds with uniform constants]
\label{thm:tensor249}
Suppose the Harish-Chandra coordinates of \(\rho\) have absolute
value at most \(B\) and consecutive gaps greater than \(24n^2+2\).
For unit vectors \(a,b\) in the lowest compact type one has
\[
 |\langle\rho(h)a,b\rangle|\le T_0^{C n^3} W(h)^{-L} .
\]
For \(S\ge1\) the same parameters satisfy
\[
 \begin{split}
 d_\rho\operatorname{vol}\{W\le S\}&\le T_0^{C n^2} S^{4m^2},\\
 d_\rho\int W(h)^{-D}\,dh&\le T_0^{C n^2},\qquad
 d_\rho\operatorname{vol}(U_{1/4})\ge T_0^{-C n^3} .
 \end{split}
\]
\end{theorem}
\begin{proof}
We track the constants in translation by the dual defining representation
\(V=(\mathbb C^m)^*\). Write \(e^l\) for its coordinate covectors,
\(E_{ab}\) for complex matrix units, and \(\lambda\) for the current
Harish-Chandra parameter. For one translation the coherent character
identity used in Lemma~\ref{lem:translate163} gives the discrete series
with parameters \(\lambda-e_i\), each once. All shifts here are in the
same regular chamber. Their infinitesimal characters are distinct.
Consequently the center splits the finite-length tensor module into
these irreducible constituents. On smooth globalizations this is also a
direct sum, by the Casselman--Wallach globalization theorem. We use
this observation for translated smooth vectors as well as compact-finite
vectors. After at most \(L<24n^2\) shifts a difference of two parameter
entries has changed by at most \(L\), so each translation used below
satisfies the same chamber condition.

On this tensor use the invariant Hermitian form \([\ ,\ ]\) given
by the Hilbert form of \(\rho\) times the signed form on \(V\).
The difference of the tensor Casimir and the two factor Casimirs,
divided by two, is the split Casimir \(A\). For the trace form its
matrix on vectors with \(m\) components is
\(A_{ab}=-d\rho(E_{ab})\), since the dual action of \(E_{ba}\) is
\(-E_{ab}\). The Casimir eigenvalue on parameter \(\nu\) is
\(\sum\nu_i^2-\sum\rho_{G,i}^2\), with \(\rho_G\) the half-sum
of positive roots. On the dual defining representation its value is
\(m\) (equivalently \(\sum E_{ab}E_{ba}\) acts by \(m\)). Thus on
summand \(i\) the eigenvalue of \(A\) is
\(\xi_i=-\lambda_i-(m-1)/2\), and its projector is
\[
       P_i=\prod_{j\ne i}\frac{A-\xi_j}{\xi_i-\xi_j} .
\]
The denominators have absolute value at least one. The split Casimir
is self-adjoint for \([\ ,\ ]\), since each Casimir is a signed sum
of squares of infinitesimal generators of an isometric action.
The real distinct eigenvalues give orthogonal summands. On each,
the nonzero restricted form is a real multiple of the unitary form.
Uniqueness follows by Schur's lemma on the irreducible Harish-Chandra
module and then by density on smooth vectors. It is nonzero by
orthogonality and nondegeneracy of the full form. Normalize the Hilbert
norm on this summand so that its square is \(|[x,x]|\). In particular
\[
 \|P_i(a\otimes e^l)\|_i^2
     = |[P_i(a\otimes e^l),a\otimes e^l]| .
\]
The absolute value on the right is at most \(\|a\|\) times the
ordinary product-space Hilbert norm of \(P_i(a\otimes e^l)\) in the
parent representation. This estimate avoids any choice of norm for a
translation intertwiner.

Here are bounds for the differential operators in this expression.
In any intermediate unitary constituent the coordinates of its parameter
are at most \(B+L\) in absolute value. The compact highest-weight
coordinates meeting our vectors are at most \(B+n+L\) in absolute
value. Initially this follows from the lowest-type formula in
Proposition~\ref{prop:gradient248}. Tensoring with a defining
dual changes the compact bound by at most one, since all weights of a
unitary compact type lie in the convex hull of the Weyl orbit of its
highest weight. The projections preserve compact types. Applying one
Lie algebra element increases the bound by at most one as well, by
the compact action on the matrix units in the adjoint representation.
On the sum of the resulting types after at most \(n\) derivatives,
the Casimir identity \(-\Delta=2\Omega_K-\Omega\) in
Proposition~\ref{prop:gradient248} gives
\[
 \sum_X\|d\rho(X)a\|^2
       \le 20n(B+4n+2L)^2\|a\|^2
\]
for a real orthonormal Lie algebra basis \(X\). The Casimir calculation
applies on each compact isotypic space and these spaces are orthogonal.
A complex matrix splits as \(X+iY\) in the real form, with both norms
at most its Hilbert--Schmidt norm, since real conjugation
\(Z\mapsto-JZ^*J\) is isometric. Successive applications show that
any word of length \(j\le n\) in matrices of Hilbert--Schmidt norm at
most one acts on the starting vector with norm at most \(T_0^{2j}\|a\|\).

More generally put \(J_m(a)\) equal to the maximum of the Hilbert
norms of the matrix-unit words of length at most \(m-1\) applied to
\(a\), including \(a\) itself. Expanding the projector gives
\[
      \|P_i(a\otimes e^l)\|_i^2
                    \le T_0^{C n}\|a\|J_m(a).
\]
Indeed each multiplication by \(A-\xi_j\) contributes at most
\(m+1\) choices to a matrix entry, one derivative or a scalar of size
at most \(B+L+n\). Taking the product-space norm adds at most
\(\sqrt m\). This proof of the last bound works for every smooth
\(a\). For our compact types the preceding word estimate bounds
its right side by \(T_0^{C n}\|a\|^2\).

To obtain decay, multiply the coefficient by an entry
\((h^{-1})_{la}\). Up to a sign this gives the coefficient for the
signed form between the tensor vectors with covectors \(e^l,e^a\).
Decomposing them gives a sum of at most \(m\) signed unitary
coefficients, with each pair of vector norms growing in product by
at most \(T_0^{C n}\). Repeat \(L\) times using the same indices.
The number and the norms of all terms together are bounded by
\((m T_0^{C n})^L\). At least one entry of \(h^{-1}\) has absolute
value at least \(W(h)/m\). Since \(L=O(n^2)\), this proves the first
assertion.

For a lowest unit vector \(a\) and \(h\in U_\delta\), take the
matrix logarithm at the identity. For \(\delta\le1/4\) it has
Hilbert--Schmidt norm at most \(2\delta\), by its power series, and
lies in the real Lie algebra by uniqueness of the logarithm near one
and \(Jh^*J=h^{-1}\). With \(\delta=(8T_0^2)^{-1}\), integrating
the first derivative along this exponential gives
\(|\langle\rho(h)a,a\rangle|\ge1/2\). Formal degree orthogonality
therefore gives \(d_\rho\operatorname{vol}(U_\delta)\le4\).
For any \(0<\delta\le1/4\) the ball \(W\le S\) is covered by at most
\((1+2\sqrt m S^2/\delta)^{2m^2}\) left translates of \(U_\delta\).
Take a maximal \(\delta/S\)-separated subset in the ambient
Hilbert--Schmidt norm and compare balls of half that radius in the
real \(2m^2\)-dimensional matrix space. A center \(g\) has
\(W(g^{-1})\le S\), so distances at most \(\delta/S\) give the claimed
translate. The small choice of \(\delta\) proves the volume bound.
On shells \(2^j\le W<2^{j+1}\) their weighted contributions to the
integral of \(W^{-D}\) decrease by a factor at least \(2^{D-4m^2}\)
in the resulting geometric majorant. This gives the second estimate.
Finally use the same cover with \(\delta=1/4\) and the first assertion
for \(a=b\) to obtain
\(1/d_\rho=\int |\langle\rho(h)a,a\rangle|^2 dh
 \le T_0^{C n^3}\operatorname{vol}(U_{1/4})\), by summing the shells
with exponent \(2L>4m^2+1\). This proves the lower bound.
\end{proof}

\subsection{Absolute Hermite coefficients of a real translate}\label{nav:VII-41}
Consider the first even core edge \(C_0,C_1\), of ranks \(n,n+d\)
with \(d=n+4\). At a real place its symplectic space is the
underlying real space of a complex matrix space of dimension
\(D=n(n+d)\). The Schr\"odinger position space has real dimension
\(D\). Use signed orthonormal Hermitian bases on both groups.
Choose the corresponding positive oscillator structure respecting
\(C_0=Y\perp Ez\), and the orthonormal product Hermite basis
\((h_\alpha)_{\alpha\in\mathbb N^D}\) in these coordinates.
Write \(\mathcal N h_\alpha=|\alpha|h_\alpha\) for the number
operator and put
\[
 s=\lfloor D/2\rfloor+1,\qquad
 A(D,L)=\big[(1+16D^3)(L+2s+1)\big]^s .
\]
Here the energy of a polynomial Fock vector is its largest
nonzero Hermite degree.

\begin{lemma}[A summable bound on the first real oscillator]\leavevmode
\label{lem:weightedhermite260}
Let \(\phi\) have energy at most \(L\) and let \(G=U(C_0)_\tau\).
For \(W(g)=\|g\|_{\rm op}\) in the defining signed coordinates,
\[
 \sum_\alpha|\langle h_\alpha,\omega(g)\phi\rangle|
 \le \sqrt{\zeta(2)}\,\|(1+\mathcal N)^s\omega(g)\phi\|_2
 \le \sqrt{\zeta(2)}\,A(D,L) W(g)^{2s}\|\phi\|_2 .
\]
For the bottom core discrete series \(\pi\), every compact-finite
unit vector \(u\) has finite weighted coefficient integral
\[
 M_\tau(u)=d_\pi\int_G
       |\langle u,\pi(g)u\rangle| W(g)^{2s}\,dg .
\]
Here \(d_\pi\) is the formal degree in the chosen measure.
If \(u\) belongs to the lowest compact type, then
\(M_\tau(u)\le T_0^{C n^3}\), with \(T_0\) of
Theorem~\ref{thm:tensor249} and an absolute \(C\).
\end{lemma}
\begin{proof}
For degree \(j\) there are at most \((1+j)^{D-1}\) indices,
since the first \(D-1\) entries determine the last. As
\(2s\ge D+1\), the sum of \((1+|\alpha|)^{-2s}\) is at most
\(\zeta(2)\). Cauchy--Schwarz with these weights proves the
first inequality.

Use the usual creation and annihilation operators \(a_i^*,a_i\)
in the Hermite basis, so \(\mathcal N=\sum_i a_i^*a_i\).
Conjugation by \(\omega(g)\) substitutes the real linear action
of \(g\) or its inverse on the Heisenberg generators. This
follows by differentiating the defining covariance of the Weil
representation on each one-parameter Heisenberg subgroup.
In the Hilbert--Schmidt metric on the matrix space, composition
on the \(C_0\) factor has norm at most \(W(g)\), since
\(\|g^{-1}\|_{\rm op}=W(g)\).
The common positive scalar in the symplectic form scales the
two orthonormal Heisenberg coordinates together. In the
normalized complex combinations defining \(a_i,a_i^*\) each
coefficient of the substitution is therefore bounded by
\(2W(g)\). Circle factors in a group splitting disappear in
conjugation.

Write \(b_i=\omega(g)^{-1}a_i\omega(g)\). Each has sum of
absolute linear coefficients at most \(4DW(g)\). Expanding
\((1+\sum_i b_i^* b_i)^s\) gives words of length at most \(2s\)
in \(a_i,a_i^*\), with sum of absolute coefficients at most
\((1+16D^3)^s W(g)^{2s}\), since \(W(g)\ge1\). On the space of
energies at most \(j\) each individual operator has norm at
most \(\sqrt{j+1}\), by its action on the Hermite basis.
The norm of each word on \(\phi\) is at most
\((L+2s+1)^s\|\phi\|_2\). Unitarity and the triangle inequality
prove the second bound.

The core gaps exceed \(24n^2+2\) in the proof of
Theorem~\ref{thm:first250}. Since \(2s\le2n^2+4n+2\),
the integer \(2s+4n^2+1\) is below those gaps for \(n\ge2\).
Use this exponent for the coefficient in
Lemma~\ref{lem:translate163} and integrate the remaining
power \(W^{-(4n^2+1)}\) by its Haar bound. This proves
finiteness for every compact-finite \(u\). For a unit vector
in the lowest type, Theorem~\ref{thm:tensor249} gives decay
with exponent \(8n^2+2n+9\) and prefactor \(T_0^{C n^3}\).
After subtracting \(2s\) the exponent is still at least
\(4n^2+2\). The formal-degree Haar bound in the same theorem
therefore gives the stated uniform estimate.
\end{proof}

\subsection{A matrix-norm moment for two smooth vectors}\label{nav:IX-35}
Keep the notation \(T_0,W,J_m,d_\rho\) of Theorem~\ref{thm:tensor249}.
The moment below permits one vector to be the projection of a smooth
vector from a larger unitary representation. Such a projection can
have infinitely many compact types.

\begin{proposition}[A second moment for the full real matrix norm]
\label{prop:moment249}
Under the parameter hypotheses of Theorem~\ref{thm:tensor249}, for
all smooth vectors \(a,b\) in \(\rho\),
\[
 d_\rho\int_{U(p,q)}\|h\|_{\rm HS}^2
                 |\langle\rho(h)a,b\rangle|^2\,dh
       \le T_0^{C n} J_m(a)\|a\|J_m(b)\|b\| .
\]
The constant is absolute, including for compact signatures.
\end{proposition}
\begin{proof}
For each pair of indices \(l,k\), use the smooth decomposition and the
signed invariant form of Theorem~\ref{thm:tensor249}. Apart from the sign
of the \(k\)-th coordinate,
\[
  (h^{-1})_{lk}\langle\rho(h)a,b\rangle
   =\sum_i \epsilon_i
      \langle\rho_i(h) P_i(a\otimes e^l), P_i(b\otimes e^k)\rangle_i ,
\]
where the sign on a summand is \(\epsilon_i\). The right side is an
identity of smooth coefficients in the corresponding globalization.
Write \(d_i\) for the formal degree of summand \(i\). The formal
degree formula on this fixed group is a common positive Haar factor
times the absolute Vandermonde product of the Harish-Chandra entries.
On lowering \(\lambda_i\) by one, each of its \(m-1\) Vandermonde
factors changes by a ratio at least \(1/2\), because each original
difference has absolute value at least two. Thus \(d_\rho/d_i\le2^{m-1}\).
The Cauchy--Schwarz inequality for the sum over \(i\), followed by
formal degree orthogonality in each \(\rho_i\), bounds the integral
of the absolute square by
\[
        m\sum_i d_i^{-1}\|P_i(a\otimes e^l)\|_i^2
                            \|P_i(b\otimes e^k)\|_i^2 .
\]
Use the smooth-vector projector estimate of Theorem~\ref{thm:tensor249}
and sum over \(l,k\). The factor \(m^4 2^{m-1}\) is absorbed by
\(T_0^{C n}\). Since \(h^{-1}=Jh^*J\), the sum of the squared matrix
entries of \(h^{-1}\) is exactly \(\|h\|_{\rm HS}^2\).
\end{proof}

\subsection{A finite Hermite cutoff along the horizontal chain}\label{nav:IX-42}
Keep one of the fields in Theorem~\ref{thm:fullperiod266}
and the same compact-finite input forms on an edge \(i\ge2\).
Use the real section through the identity of
Lemma~\ref{lem:realflag265}. In local coordinates \(y=0\)
there write \(J(y)\) for the scalar fiber integral before
insertion of the first-position cutoff, as in
Theorems~\ref{thm:bestransfer265} and~\ref{thm:jactransfer266}.
Then \(J(0)\) is the nonzero preceding period on an odd-index
edge and the preceding period times \(\|q\|_2^2\) on an
even-index edge. Write \(a\) for the real dimension of the
section. Choose a closed coordinate ball of radius
\(0<\rho_0\le1\) inside it. Put
\[
 L=\sup_{|y|\le\rho_0}\|\nabla J(y)\|,\qquad
 \rho=\min\{\rho_0,\ |J(0)|/[2(1+L)]\}.
\]
Let \(c_f,v_f\) be the given finite base scalar and orbit
measure on this edge, and \(T\) the number of original marked
terms in its finite partition. Put
\(\eta=c_fv_f|J(0)|/(2T)\). These three finite quantities
are those in the corresponding transfer theorem.

\begin{theorem}[Degree and size of an individual real test]\leavevmode
\label{thm:cutoff268}
Let \(d,s,M\) be as in Proposition~\ref{prop:order268}, with
the maximum taken over the marked terms of this edge. There
are finite positive chart and test constants \(C_0,C_{s+1}\)
described below. Set
\[
 K=\max\!\left\{1,\left\lceil
 \frac{2MC_{s+1}}{\eta\,\rho^{\,a/2+2(s+1)}}\right\rceil\right\}.
\]
One original marked term and one normalized product Hermite
test \(h_\nu\) of total degree \(|\nu|\le K\) satisfy
\[
       |P_{i+1}(h_\nu)|\ge
        \frac{\eta\,\rho^{a/2}}{2 C_0\binom{K+d}{d}^{1/2}} .
\]
This estimate can be applied successively from the nonzero
second restriction through the prescribed raw terminal Bessel
period, retaining the same terminal integral coefficient.
\end{theorem}
\begin{proof}
The function \(J\) is continuously differentiable on the
closed coordinate ball. A derivative acts on the smoothly
translated row by a first Lie algebra derivative, or on the
translated residual oscillator test. On a compact section
these operations use bounded combinations of differentiated
compact-finite cusp forms and bounded Schwartz seminorms.
The differentiated cusp forms still have rapid decrease.
The integrable chamber majorants in the proofs of
Propositions~\ref{prop:besorbit265} and~\ref{prop:jacorbit266}
therefore apply after differentiation, giving the assertion
by dominated convergence. In particular \(L\) is finite.
Throughout the ball of radius \(\rho\) the difference
\(|J(y)-J(0)|\) is at most \(|J(0)|/2\).

We record the real seminorms in the localization. Choose
once a nonnegative smooth product bump of integral one
supported inside the open unit ball in the section coordinates.
Dilate it by \(\rho\) and divide by the smooth positive
quotient-density factor, giving it quotient integral one.
Write \(\Phi_\rho\) for its ambient horizontal Schwartz test.
In the orbit chart \((y,u)\) used in
Lemma~\ref{lem:realflag265}, multiply the section bump by
the fixed upper-group bump and the prescribed character in
\(u\). Use upper support inside a fixed compact coordinate
patch. Extend to a tubular chart by a fixed transverse bump
of value one at zero. The core, line and opposite-position factors in the remaining
coordinates are fixed Schwartz functions. One may shrink \(\rho_0\) at
the outset to fit the patch and use a finite family of such
charts. All coordinate maps and their required derivatives
are bounded on the supports in the frame coordinates.

For every fixed integer \(j\ge0\) this construction gives
\[
       \|(1+\mathcal N)^j\Phi_\rho\|_2
                    \le C_j\,\rho^{-a/2-2j}.
\]
Indeed the section bump and its derivatives of order \(v\)
cost \(O(\rho^{-a-v})\), by the chain rule and the fixed
density factor. The volume of the section support is
\(O(\rho^a)\). The operator \((1+\mathcal N)^j\) is a sum
of polynomial-coefficient differential words of length at
most \(2j\). The other frame variables stay on fixed compact
supports and the weighted residual derivatives are square
integrable. Taking \(L^2\) norms gives the indicated power.
Equivalently \(C_j\) can be taken as the sum of the resulting
coordinate derivative and Jacobian bounds times the weighted
norms of the fixed factors. A fixed metaplectic change to the
original polarization replaces each word by a combination
of words of the same length, by the covariance used in
Proposition~\ref{prop:order268}. Include its matrix
coefficients in \(C_j\). These constants are independent of
the chosen radius.

The two orbital transfer formulas now give modulus at least
\(c_fv_f|J(0)|/2\) before their finite partition, by
integrating against this nonnegative bump. Some individual
marked term has modulus at least \(\eta\), since the partition
coefficients have modulus one. Fix it. For the spectral
projection \(E_K\) onto all Hermite degrees at most \(K\),
Parseval gives
\[
 |P_{i+1}((1-E_K)\Phi_\rho)|
 \le \frac{M}{K+1}\|(1+\mathcal N)^{s+1}\Phi_\rho\|_2
 \le\eta/2 .
\]
Hence the period on \(E_K\Phi_\rho\) has modulus at least
\(\eta/2\). There are \(\binom{K+d}{d}\) multi-indices in
this range, by the count of weak compositions with a slack
coordinate. Cauchy--Schwarz bounds the sum of the absolute
Hermite coefficients by the square root of this number times
\(\|\Phi_\rho\|_2\). This proves the stated individual bound.

The selected output is compact-finite and cuspidal by the
first-occurrence statement of Theorem~\ref{thm:realrow257},
so it is an admissible input to the next edge. For the final
raw character, apply the rational involution of
Lemma~\ref{lem:topchar266} to the smooth final test
\emph{before} its Hermite projection. It preserves the absolute
period, takes a marked coset to a marked coset, and at infinity
is a fixed metaplectic action. Include this action in \(C_j\)
and take \(M\) for the raw period. The same estimates now
select a single test in the original terminal coordinates.
This proves the last assertion by finite induction, starting
with Theorem~\ref{thm:lowesttwo261}.
\end{proof}

The bounds use the input forms through \(M,L,J(0)\) and the
chosen charts and residual tests through \(C_j\). Their growth
along the varying chain, and the passage from this complex
period to the integral projection of
Corollary~\ref{cor:original200}, remain to be controlled.

\subsection{The real distribution of a terminal Bessel coefficient}\label{nav:VII-05}
\label{sec:terminal151}

Keep the unitary space of Section~\ref{sec:char94}, with
\(r=k+2s+1\), and set \(B=N\rtimes U(W_0)\). At each real
place, sections for the frequency \(\beta\) transform by
\(\chi_\beta\) on \(N\) and by the column representation on
\(U(W_0)\). We use the invariant quotient measure on
\(B_\infty\backslash G_\infty\). The next formulas apply to
nonnegative integrals, including infinite values.

\begin{proposition}[Change of the terminal frequency and the height]
\label{prop:scale151}
At a single real place let \(a_\beta\) multiply the \(u_i\)
by \(\beta\in\CC^\times\), the \(v_i\) by
\(\bar\beta^{-1}\), and fix the middle space pointwise. Put
\(J_\beta=|\beta|^{2s(r-s)}\). The operation
\[
          W(g)\longmapsto J_\beta^{-1/2}W(a_\beta g)
\]
is a right-equivariant isometry from square-integrable
sections of frequency one to those of frequency \(\beta\).
For \(\beta=\varpi_\tau\) simultaneously at all real places,
write \(U_\varpi\) for its product. For every nonnegative
Borel function \(b\) and every real seed section \(W\),
\begin{equation}
 \int b(H(g))\|(U_\varpi W)(g)\|^2\,d\dot g
          =\int b(\sqrt T H(g))\|W(g)\|^2\,d\dot g .
     \label{eq:scale151}
\end{equation}
Here both sides use the same height of Section~\ref{sec:success135}.
\end{proposition}
\begin{proof}
Conjugation by \(a_\beta\) fixes the first \(s-1\) simple
coordinates of \(N\) and multiplies its last coordinate by
\(\beta\). It centralizes \(U(W_0)\). Thus left translation
of the argument gives the indicated character. It commutes
with every right translation. To compute the measure, use the
decomposition \(L\oplus W\oplus L^\vee\). In the Lie algebra
of \(B\), the compact block and the strictly upper triangular
block on \(L\) are unchanged. The block from \(W\) to \(L\)
has complex dimension \(s(k+1)\) and is multiplied by
\(\beta\). Its opposite partner is determined by skew
adjointness. The block from \(L^\vee\) to \(L\) consists of
skew-Hermitian \(s\)-by-\(s\) matrices of real dimension
\(s^2\), and is multiplied by \(|\beta|^2\). Therefore the
absolute determinant on the Lie algebra of \(B\) is
\(|\beta|^{2s(k+1)+2s^2}=J_\beta\).

Both groups are unimodular. For a compactly supported scalar
function \(v\) on \(G\), put \(v_B(g)=\int_B v(bg)\,db\).
In \(v_B(a_\beta g)\) change \(b=a_\beta b'a_\beta^{-1}\).
Quotient integration and left invariance of Haar measure on
\(G\) give
\(\int v_B(a_\beta g)\,d\dot g=J_\beta\int v_B(g)\,d\dot g\).
Local sections and a partition of unity give the same formula
for every nonnegative function on the quotient. Finally,
\((a_\beta g)^{-1}u_i=\beta^{-1}g^{-1}u_i\) for every
\(i\), so \(h_\tau(a_\beta g)=|\beta|^{-1}h_\tau(g)\).
All \(|\varpi_\tau|=\sqrt T\). Combining these formulas and
the scalar normalization proves \eqref{eq:scale151}. Replacing
\(\beta\) by its inverse gives the inverse isometry.
\end{proof}

Two other quotient functions will locate a Bessel coefficient
with large frequency. At each place use Euclidean distance
to a complex subspace and put
\[
 A_\tau(g)=\operatorname{dist}(g_\tau^{-1}u_s,g_\tau^{-1}L_{s-1}),
 \qquad Z_\tau(g)=\operatorname{dist}(g_\tau^{-1}e,g_\tau^{-1}L).
\]
Here \(L_0=0\). Left multiplication by \(B_\infty\) preserves
the corresponding affine spaces, so the functions descend.
One has \(0<A_\tau\le h_\tau\) and \(Z_\tau\ge1\), the latter
because \(e+L\) has Hermitian norm \(-1\). Define
\[
        H_{\rm end}(g)=\Bigl(\prod_\tau A_\tau(g)Z_\tau(g)\Bigr)^{1/f}.
\]

\begin{lemma}[Terminal derivatives on a region of the real quotient]
\label{lem:terminal151}
Let \(W\) be a smooth section of frequency \(\varpi\).
At each place choose a real Frobenius-orthonormal basis of
\(\mathfrak u(k+s,s+1)\). For \(j\ge1\), let \(\mathcal I_j\)
consist of words \(I\) with \(j\) right derivatives at each
place in these basis directions, in a fixed order of the
places and in every order within each place. On any measurable
region \(\Omega\) contained in \(H_{\rm end}\le U\),
\begin{equation}
 \|\mathbf1_\Omega W\|_2
 \le \left(\frac{U}{2\pi\sqrt T}\right)^{fj}
           \left(\sum_{I\in\mathcal I_j}
                      \|\mathbf1_\Omega D_I W\|_2^2\right)^{1/2}.
     \label{eq:terminal151}
\end{equation}
The inequality also holds in the seed norm, with all its finite
class coordinates and the same fixed frequency.
\end{lemma}
\begin{proof}
Fix \(g\) and a real place. Choose \(u'\in u_s+L_{s-1}\)
and \(e'\in e+L\) that realize the two displayed distances
after applying \(g_\tau^{-1}\). Orthogonal projection gives
the minimizing representatives. For a scalar \(c\) of modulus
one define
\[
       Y(w)=c\,u' h(e',w)-\bar c\,e'h(u',w).
\]
This is skew adjoint. It kills \(L\), sends \(L^\perp\) to
\(L\), and sends the whole space into \(L^\perp\).
It belongs to the Lie algebra of \(N\) since it acts as zero
on every graded piece of the completed flag. Its first
\(s-1\) simple coordinates are zero, and the \(u_s\)
coefficient of \(Y(e)\) is \(-c\). Choose the phase so that
\(-c\varpi_\tau=\sqrt T\). Then
\(\chi_{\varpi,\tau}(\exp(tY))=\exp(4\pi i\sqrt T\,t)\)
for real \(t\), by the trace from \(\CC\) to \(\RR\).

Put \(X=g_\tau^{-1}Yg_\tau\). The Frobenius norm of a
rank-one map \(v h(w,\cdot)\) is \(\|v\|\,\|w\|\), since the
signature matrix is Euclidean unitary. It follows that
\(\|X\|_{\rm F}\le 2A_\tau(g)Z_\tau(g)\).
Freeze this \(X\) while differentiating at the given \(g\).
The identity \(g_\tau\exp(tX)=\exp(tY)g_\tau\) gives the
factor \((4\pi i\sqrt T)^j\) on taking \(j\) derivatives.
Apply this at all places, whose actions commute. Express the
resulting product of derivatives in the fixed bases. The
Euclidean norm of its coefficient vector is
\(\prod_\tau\|X_\tau\|_{\rm F}^{\,j}\), by expansion of the
tensor powers. Cauchy--Schwarz gives the pointwise version
of \eqref{eq:terminal151}; integration gives the result.
The coefficients are treated as fixed at each evaluation
point, so this argument introduces no derivatives of the
minimizing representatives. Summing the squared inequalities
with the finite class inner product gives the last assertion.
\end{proof}

\begin{theorem}[Projection on the automorphic quotient]
\label{thm:project151}
Use a single terminal frequency \(\varpi\) and any finite set
of the selected words and columns. Let \(\mathcal H_\sigma\)
be the subspace of automorphic \(L^2\) in the isotypic space
of a prescribed real discrete series \(\sigma\) and in its
specified compact type. It is the image of the two commuting
orthogonal isotypic projections. Let \(L_*\ge1\) bound the
compact Lie norms and the square roots of the absolute
Casimir scalars in Lemma~\ref{lem:cas96} on this type.
For \(R>0\) put \(B(R)=1\) when \(R\le R_0\) and
\(B(R)=(1+16(R^2/Q_0)^f)^{rs}\) when \(R>R_0\).
For every seed supported in \(H<R,\ H_{\rm end}\le U\),
and every integer \(j\ge1\),
\begin{equation}
 \|\operatorname{proj}_{\mathcal H_\sigma}P\Phi\|
       \le B(R)
      \left(\frac{C r^2(L_*+fj+1)U}{\sqrt T}\right)^{fj}
                      \|\Phi\|_{\rm seed}
       \label{eq:project151}
\end{equation}
with an absolute \(C\).
\end{theorem}
\begin{proof}
Take first a smooth vector \(F\) in a finite sum of the given
copies. The functional \(\Phi\mapsto\langle P\Phi,F\rangle\)
on compact test seeds has a locally square-integrable representing
section \(W_F\) in the seed inner product. Indeed restriction to
any bounded height is bounded by Theorem~\ref{thm:sectors149}.
The Riesz representatives agree on overlaps by testing on
sections there. On \(H<R_0\) the summation map is isometric.
Together these bounds give
\(\|\mathbf1_{H<R} W_F\|_{\rm seed}\le B(R)\|F\|\).
Use complex inner products linear in the first argument,
so that \(W_F\) has the same covariance as the test seeds.

For a compact test seed and real \(a\), ordinary summation
commutes with \(R(a)\). Changing variables in the pairing
therefore gives \(R_B(a)W_F=W_{R_G(a)F}\) as distributions.
Differentiation gives \(D_IW_F=W_{D_I F}\). In particular
all these weak derivatives are locally square-integrable,
and on \(H<R\) their norms are at most \(B(R)\|D_I F\|\).
One may use smooth representatives in
Lemma~\ref{lem:terminal151}: in a local bundle chart the
right Lie directions span the tangent space, so all coordinate
derivatives of each order are linear combinations with smooth
coefficients of such words of at most that order. The local
Sobolev embedding theorem in each finite order gives a smooth
representative. The fixed finite class and word coordinates
carry unitary fibre metrics and satisfy the same argument.

There are \(r^{2fj}\) words in \(\mathcal I_j\), since the
real dimension at each place is \(r^2\). By
Lemma~\ref{lem:cas96}, each \(\|D_I F\|\) is at most
\(\{C r(L_*+fj+1)\}^{fj}\|F\|\).
In \eqref{eq:terminal151} use
\(\Omega=\{H<R,H_{\rm end}\le U\}\) and the preceding
derivative bounds. This gives
\(\|\mathbf1_\Omega W_F\|_{\rm seed}\) bounded by the factor
in \eqref{eq:project151} times \(\|F\|\). Pairing against
\(\Phi\) gives the assertion for such test vectors \(F\).
Finite sums from copies of \(\sigma\) in this compact type
are dense in \(\mathcal H_\sigma\), by the definition of
the isotypic subspace and the bounded compact-type projector.
Within each copy compact-finite vectors of the discrete series
are smooth. The constant is uniform on these copies. Taking
the supremum over their unit sphere gives \eqref{eq:project151}.
\end{proof}

\subsection{Consecutive frame lengths in the Bessel energy estimate}\label{nav:VII-06}
\label{sec:ratios153}

At a real place \(\tau\) use the flag of
Section~\ref{sec:terminal151} at \(g_\tau\). In this section set
\[
 v_i=g_\tau^{-1}u_i\ (1\le i\le s),\qquad
 v_{s+1}=g_\tau^{-1}(\varpi_\tau e),\qquad
 V_i=\langle v_1,\ldots,v_i\rangle,\qquad V_0=0 .
\]
Define \(a_{i,\tau}=\operatorname{dist}(v_i,V_{i-1})>0\)
for \(1\le i\le s+1\). These are functions of the real seed
coset. The first \(s\) are its Gram--Schmidt lengths. The last
is \(\sqrt T Z_\tau\ge\sqrt T\). With the indices starting at
zero in the archive, these are exactly the functions in
\srcloc{Formalism/SharpRestriction/CuspSpanNormBounds/RelevantZeroGramTerms}{299}.
Put
\[
 M_\tau(g)=\max_{1\le i\le s}\frac{a_{i+1,\tau}}{a_{i,\tau}},
       \qquad M(g)=\left(\prod_\tau M_\tau(g)\right)^{1/f}.
\]

\begin{lemma}[A derivative for each consecutive pair]
\label{lem:roots153}
For each \(g_\tau\) and each \(i\le s\) there is a real
unitary Lie direction \(X_i\) with
\(\|X_i\|_{\rm F}\le2a_{i,\tau}/a_{i+1,\tau}\) such that
every smooth Bessel section of frequency \(\varpi\) satisfies
\[
         W(g_\tau\exp(tX_i))=\exp(4\pi i t)\,W(g_\tau)
                 \qquad(t\in\RR).
\]
Here the other arguments of the section are held fixed.
\end{lemma}
\begin{proof}
Write \(J\) for the signature matrix, so \(h(y,w)=y^*Jw\),
\(J^*=J\), and \(J^2=1\). Let \(P_V\) denote Euclidean
orthogonal projection. At this point and this index set
\[
    u=(1-P_{V_{i-1}})v_i,\qquad
    t_0=(1-P_{V_i})v_{i+1},\qquad
    w=Jt_0/\|t_0\|^2 .
\]
Thus \(u\in v_i+V_{i-1}\), \(\|u\|=a_{i,\tau}\),
\(h(w,V_i)=0\), \(h(w,v_{i+1})=1\), and
\(\|w\|=a_{i+1,\tau}^{-1}\). Define
\[
                X_i y=u h(w,y)-w h(u,y).
\]
It is skew adjoint and has Frobenius norm at most
\(2\|u\|\|w\|\). Since \(u\in V_s\) and \(V_s\) is isotropic,
this map kills \(V_i\), sends \(V_j\) to \(V_i\) for
\(i<j\le s\), and sends \(V_s^{\perp_h}\) to \(V_s\).
By skew adjointness it also lowers the dual flag. Thus
\(g_\tau X_i g_\tau^{-1}\) lies in the Bessel unipotent
Lie algebra. If \(i<s\), the coefficient of \(v_i\) in
\(X_i v_{i+1}\) modulo \(V_{i-1}\) is one. The other simple
coordinates, including the terminal one, are zero: in the
later steps the image lies in the earlier \(V_i\).
If \(i=s\), all the ordinary simple coordinates vanish and
\(X_s(\varpi_\tau g_\tau^{-1}e)=u\). The terminal coordinate
relative to \(e\) is then \(\varpi_\tau^{-1}\), which is
multiplied by \(\varpi_\tau\) in the actual character.
In both cases the character on the exponential of the
conjugate is \(\exp(4\pi i t)\). Left covariance proves the
identity. The factor \(a_{i+1,\tau}^{-1}\) in the norm comes
from the vector representing the dual functional on the next
quotient line; this calculation allows \(Z_\tau\) to be
arbitrarily large.
\end{proof}

\begin{theorem}[Joint energy of the consecutive ratios]
\label{thm:energy153}
Let \(W\) be a smooth real section, possibly with finite class
coordinates. For every \(j\ge1\), pointwise and hence on every
measurable region \(\Omega\),
\begin{equation}
     (2\pi)^{fj}\|\mathbf1_\Omega M^{fj} W\|_2
       \ \le\
        \left(\sum_{I\in\mathcal I_j}
                   \|\mathbf1_\Omega D_I W\|_2^2\right)^{1/2}.
     \label{eq:energy153}
\end{equation}
Here \(\mathcal I_j\) uses the fixed Frobenius bases in
Lemma~\ref{lem:terminal151}. In particular at a single place,
in an isometrically realized square-integrable real Bessel
model, a compact-type vector \(W\) with the derivative parameter
\(L_*\) of Lemma~\ref{lem:cas96} satisfies
\[
      \int M_\tau^{2j}\|W\|^2\,d\dot g
              \le \{C r^2(L_*+j+1)\}^{2j}\,\|W\|_2^2 .
\]
\end{theorem}
\begin{proof}
At a given point choose an index attaining the maximum at
each real place. Freeze the resulting \(X_i\) and take \(j\)
derivatives there, at all the places. Lemma~\ref{lem:roots153}
gives a factor \((4\pi i)^{fj}\). The coefficient vector of
the resulting words in the fixed bases has Euclidean norm
at most \(\prod_\tau(2/M_\tau)^j=2^{fj}/M^{fj}\).
Pointwise Cauchy--Schwarz gives precisely the inequality,
with the norms replaced by the pointwise fibre norms.
The right hand side uses the same list of words at every
point, including when the maximizing index changes.
Squaring and integrating proves \eqref{eq:energy153}.
For the one-place assertion there are \(r^{2j}\) words.
The hypothesis on the realization transfers the Hilbert
derivative estimate of Lemma~\ref{lem:cas96} to each of them.
This proves the moment bound. The consecutive ratios are the
ones entering the energy estimate in
\srcloc{Formalism/SharpRestriction/CuspSpanNormBounds/BesselShellMomentBounds}{1642}.
\end{proof}

\begin{corollary}[Uniformity in the terminal direction]
\label{cor:ang153}
In Theorem~\ref{thm:project151}, suppose the seed is supported
in \(H<R,\ M\ge V\), where \(V>0\). Then
\[
   \|\operatorname{proj}_{\mathcal H_\sigma}P\Phi\|
      \le B(R)\left(\frac{Cr^2(L_*+fj+1)}{V}\right)^{fj}
                                         \|\Phi\|_{\rm seed}.
\]
In fact at every point, if \(h_i=(\prod_\tau a_{i,\tau})^{1/f}\)
and \(Z=(\prod_\tau Z_\tau)^{1/f}\), one has
\[
            M^{s-i+1}h_i\ \ge\ \sqrt T\, Z\qquad(1\le i\le s).
\]
In particular \(M\ge\sqrt T\, Z/h_s\ge\sqrt T/H\).
\end{corollary}
\begin{proof}
For a unit test vector \(F\) in the required discrete series
and compact type use its local representative \(W_F\) from
Theorem~\ref{thm:project151}. On \(H<R\), the \(L^2\) norm of
\(D_I W_F\) is bounded by \(B(R)\|D_I F\|\).
Equation~\eqref{eq:energy153} on the support region and the
word count from that theorem give the bound on its adjoint.
Duality gives the stated estimate. For the geometric assertion,
multiply the consecutive ratios from \(i\) through \(s\) at
one place. Their product is \(a_{s+1,\tau}/a_{i,\tau}\);
each is at most \(M_\tau\). Take the geometric mean over the
places and use \(Z_\tau\ge1\) and \(h_s\le H\).
\end{proof}

\subsection{The real Bessel norm and the formal degree}\label{nav:VII-07}
\label{sec:relative155}

Fix one real place. Write \(G=U(V)\), \(C=U(W_0)\), and
\(B=N\rtimes C\) as in Section~\ref{sec:terminal151}.
The column space \(W_0\) is definite, so \(C\) is compact.
Let \(\rho\) be the prescribed irreducible representation of
\(C\) on \(E_\rho\) and \(\sigma\) the discrete series of the row.
Give \(C\) probability Haar measure. Choose Haar measures on
\(N\) and \(G\) and use their quotient measures throughout.
Write \(d_\sigma\) for the formal degree with respect to this
choice on \(G\). Inner products in this section are linear in
the first argument.

Put \(\mathcal G=G\times C\) and
\(\mathcal H=\{(nc,c):n\in N,\ c\in C\}\).
On \(\mathcal H\) the character \(\xi(nc,c)=\chi_\beta(n)\)
is unitary. This is the real unitary Gan--Gross--Prasad Bessel
pair. The larger space adds \(s\) hyperbolic planes and the
line of \(e\) to the smaller space. Its character is obtained
from the successive flag coordinates and the last pairing
with \(e\). The change to a nonzero terminal frequency
\(\beta\) is the conjugation in Proposition~\ref{prop:scale151}.
Changing the sign of both Hermitian forms, when needed for the
sign convention on the line, preserves the unitary groups.

\begin{lemma}[Averaging a Schwartz function over the period subgroup]
\label{lem:rel155}
Let \(\mathcal C(\mathcal G)\) be the Harish-Chandra Schwartz
algebra. For \(f\) in this algebra set
\[
    Qf(g)=\int_{\mathcal H} f(hg)\overline{\xi(h)}\,dh,\qquad
    f^*(g)=\overline{f(g^{-1})}.
\]
Then \(Qf\in L^2(\mathcal H\backslash\mathcal G,\xi)\).
With \((f_1*f_2)(h)=\int_{\mathcal G}f_1(y)f_2(y^{-1}h)\,dy\)
one has
\begin{equation}
 \langle Qf,Qf'\rangle
  =\int_{\mathcal H}(f*f'^{*})(h)\overline{\xi(h)}\,dh .
 \label{eq:rel155}
\end{equation}
\end{lemma}
\begin{proof}
For smooth compactly supported functions, expand the convolution
as \(\int_{\mathcal G} f(hg)\overline{f'(g)}\,dg\).
Integrating in \(h\) first gives \(Qf(g)\).
Disintegrating \(dg\) along left \(\mathcal H\)-orbits and using
\(Qf(h_0g)=\xi(h_0)Qf(g)\) gives the left side of
\eqref{eq:rel155}. The integrations are absolutely convergent
in this case by compact support.

For the passage to Schwartz functions we use the restriction
estimate of Beuzart-Plessis for this Bessel pair. Lemma 6.5.1(ii)
of \cite{XX-XV-IV-bp155} gives
\(\int_{\mathcal H}\Xi_{\mathcal G}(h)(1+\ell(h))^{-a}\,dh
<\infty\) for sufficiently large \(a\), where \(\ell\) is
the logarithmic norm and \(\Xi_{\mathcal G}\) is the elementary
spherical function. Thus integration on \(\mathcal H\) against
\(\overline\xi\) is continuous in a Schwartz seminorm.
The same holds after any fixed right translation. The
Harish-Chandra Schwartz algebra has continuous convolution
and involution and contains \(C_c^\infty\) densely
\cite[Sections 1.5 and 2.6]{XX-XV-IV-bp155}. Approximate \(f\) by
\(f_j\in C_c^\infty\) there. Applying \eqref{eq:rel155} to
the difference \(f_j-f_i\) shows that \(Qf_j\) is Cauchy in
the quotient Hilbert space. The translated restriction estimate
gives pointwise convergence to the displayed \(Qf\).
An almost everywhere convergent subsequence of the Hilbert
space limit identifies the limits. Polarization and continuity
now give \eqref{eq:rel155}.
\end{proof}

\begin{theorem}[The normalized Bessel model of a discrete series]
\label{thm:rel155}
Suppose that the continuous Bessel functionals on the smooth
representation \(\pi=\sigma\boxtimes\rho^\vee\) form a line.
There is an intertwining map on smooth row vectors
\(\Lambda:\sigma^\infty\to E_\rho\), unique up to a complex
scalar of absolute value one under the normalization below,
such that
\begin{equation}
  d_\sigma\int_{B\backslash G}\!
        \|\Lambda(\sigma(g)v)\|^2\,d\dot g=\|v\|^2 .
  \label{eq:isom155}
\end{equation}
Its norm and phase convention up to this scalar are determined
by the Haar measures and the matrix coefficient integral
over \(\mathcal H\).
\end{theorem}
\begin{proof}
For compact-finite smooth vectors \(a,b\) of \(\pi\) set
\(f_{a,b}(g)=\langle\pi(g)a,b\rangle\) and
\[
           E(a,b)=Qf_{a,b}(1).
\]
The group \(\mathcal G\) has compact center, and \(\pi\)
is a discrete series. Its compact-finite matrix coefficients
belong to \(\mathcal C(\mathcal G)\). Equivalently, in
Harish-Chandra inversion the smoothing rank-one operators on
this discrete representation are a Schwartz spectral component
supported at its isolated atom. Inverting that component gives
the coefficients, up to formal degree and inversion of the
group argument \cite[Theorem 2.6.1]{XX-XV-IV-bp155}.

The continuous relative character on smooth vectors in
\cite[Sections 7.1--7.2]{XX-XV-IV-bp155} therefore agrees with the
ordinary integral defining \(E\) on these vectors and on their
translates. In the convention with \(\pi(h)\) in the second
argument, substitute \(h^{-1}\) for \(h\). Theorem 7.2.1
there gives nonvanishing of this form because the continuous
\((\mathcal H,\xi)\)-functional is nonzero. The form is
Hermitian by the same substitution. In each argument its
covariance is that of the functional or its conjugate. Choosing
a generator \(\ell'\) of the line, we consequently have
\(E(a,b)=c\,\ell'(a)\overline{\ell'(b)}\) on smooth vectors
with \(c\) real and nonzero. Indeed fixing the second argument
first gives a multiple of \(\ell'\), and Hermitian symmetry
determines its dependence on that argument.

For compact-finite vectors \(v,u\), Schur orthogonality on
\(\mathcal G\) gives
\[
    (f_{v,u}*f_{v,u}^{*})(h)
       =d_\pi^{-1}\|v\|^2 f_{u,u}(h).
\]
To see the order of the entries, write the first factor in the
integral as \(\langle\pi(g)v,\pi(h)^{-1}u\rangle\) and the
conjugate factor as \(\overline{\langle\pi(g)v,u\rangle}\).
Lemma~\ref{lem:rel155} now gives
\(\|Qf_{v,u}\|^2=d_\pi^{-1}\|v\|^2 E(u,u)\).
Choose \(u\) with \(\ell'(u)\ne0\), possible by density of
the compact-finite vectors in the smooth topology. Positivity
forces \(c>0\). Put \(\ell=\sqrt c\,\ell'\). At every point
\(g\) one has
\[
       Qf_{v,u}(g)=E(\pi(g)v,u)
            =\ell(\pi(g)v)\overline{\ell(u)} .
\]
Canceling \(|\ell(u)|^2\) gives the square norm
\(d_\pi^{-1}\|v\|^2\) for the scalar function
\(g\mapsto\ell(\pi(g)v)\). Approximate an arbitrary smooth
\(v\) by compact-finite vectors in its smooth topology.
The identity gives convergence in \(L^2\), and continuity
of \(\ell\) identifies the limit pointwise along a subsequence.

Define \(\Lambda\) by
\(\ell(v\otimes z)=z(\Lambda(v))\) for \(z\in E_\rho^\vee\).
It transforms by \(\chi_\beta\) on \(N\) and by \(\rho\) on
\(C\). The quotient of the product pair identifies with
\(N\backslash G\) through \((g,c)\mapsto Nc^{-1}g\).
With probability Haar measure on \(C\), quotient integration
first along \(N\) and then along \(C\) gives exactly the
measures on \(N\backslash G\) and \(B\backslash G\) fixed
above. For \(v\otimes z\) the scalar square norm is thus
\[
   \int_{B\backslash G}\int_C
      |z(\rho(c)\Lambda(\sigma(g)v))|^2\,dc\,d\dot g
     =\frac{\|z\|^2}{\dim\rho}
        \int_{B\backslash G}\|\Lambda(\sigma(g)v)\|^2\,d\dot g .
\]
The equality is compact Schur orthogonality. The product
formal degree is \(d_\pi=d_\sigma\dim\rho\), by the same
orthogonality on the two factors. The scalar identity gives
\eqref{eq:isom155}. The line of functionals and this norm fix
the modulus of the scaling.
\end{proof}

\begin{corollary}[Isometry at every real place of the period field]
\label{cor:rel155}
For every word on the row of
Theorem \ref{thm:src109},
and for \(\beta=1\) or \(\varpi_\tau\), the map
\[
        v\longmapsto \sqrt{d_{\sigma_\tau}}\,
            \Lambda_\tau(\sigma_\tau(\,\cdot\,)v)
\]
extends to an isometric right-equivariant embedding in the local
seed Hilbert space. Its tensor product over the real places
retains this property. In Theorem~\ref{thm:energy153} the
local derivative moments therefore apply to these embeddings
on the whole real quotient.
\end{corollary}
\begin{proof}
The cited packet calculation gives exactly the one-dimensional
space required in Theorem~\ref{thm:rel155}, for each column
and each nonzero terminal character. The formula commutes with
right translations by the representation law and extends by
density. Products of the local norm identities give the assertion
on pure tensors, then on finite sums by polarization and on their
completion. For a smooth input, differentiating the equivariance
identity in the Hilbert norm transfers every derivative to the
same isometric realization. This establishes the local norm
condition at
\srcloc{Formalism/SharpRestriction/CuspSpanNormBounds/BesselShellMomentBounds}{1625}.
\end{proof}

\section{Transverse pairs and higher-rank blocks}\label{nav:g-X-3}

\subsection{The large-height contribution for every single reversed block}\label{nav:I-70}
For \(2\le j\le s\) let \(\mathcal B_j\) consist of the
relevant supported pairs of Proposition~\ref{prop:block295},
counting each rational orbit once. Let
\(H(\xi)=\prod_\tau\max(1,|\tau\xi|)\), with the scalar
in its reversed cross block written \(c=Q\xi\).

\begin{theorem}[A height cutoff uniform in the block length]\leavevmode
\label{thm:blocktail295}
For \(H_0\ge1\) and \(1\le\ell\le j\), on the original
fields and intervals,
\[
   \sum_{\substack{w\in\mathcal B_j\\H(\xi)\ge H_0}}
        \|O_w\|_{\rm op}
   \le \frac{\exp(CRj\ell f s\log b)}
               {Q^{(3\ell R-2r)jf}\,H_0^{3j\ell R-4j-3}} .
\]
In particular for every fixed \(D>0\) there is \(M_D\)
such that the sum over all \(2\le j\le s\) with
\(\log H(\xi)\ge M_D f s\log b\) is at most
\(\exp(-D R f s\log b)\). These bounds hold after the
original orthogonal projections. The remaining labels
in these one-block families form a finite set at each
point of the interval.
\end{theorem}
\begin{proof}
Take the power \(A_1=3\ell R\) in Lemma~\ref{lem:blockphysical295}
at a place with \(|\tau\xi|<1\), and \(A_2=6\ell R\) at
each other place. The product of all inverse powers of
\(|\tau c|\) is at most
\[
       (Q^f A_\xi)^{-3j\ell R}\,H(\xi)^{-3j\ell R}.
\]
Here \(Q\ge1\) and \(A_\xi=\prod_\tau|\tau\xi|\).
Such independent local choices give a bound for the
actual product orbital integral. Indeed the diagonal
sesquilinear forms and the polarized signed forms
commute with the real right maximal compact group on
its given irreducible lowest type. They are scalar
on that type, by Schur's lemma. Evaluate them on a
decomposable vector, use product measures, and bound
each signed factor by the local positive integral, as
in Proposition~\ref{prop:places288}. Alternatively the
weighted Cauchy--Schwarz proof uses the product of the
chosen local height weights. The field factor is at
most one by Lemma~\ref{lem:blockphysical295}.
Multiply by Proposition~\ref{prop:blockfinite295}.
The remaining power of \(A_\xi\) is \(-(3\ell R-2r)j\).
It is negative for the ranks in question, and \(A_\xi\ge1\)
by integrality. This leaves the individual estimate
\[
   \exp(CRj\ell f s\log b)\,
           Q^{-(3\ell R-2r)jf}H(\xi)^{-3j\ell R}.
\]

For fixed \(\xi\), the algebraic integer
\(\eta=T\xi^j\lambda\) has prescribed
\(\eta\bar\eta=T^2(\xi\bar\xi)^j\).
As in Theorem~\ref{thm:codimension294}, its ideal divides
the ideal of this product, the number of divisors is at
most the absolute norm, and generators for a given ideal
differ by a root of unity. Thus at most
\(8f^2 T^{4f} H(\xi)^{4j}\) choices of \(\lambda\)
occur. Each gives at most one label for the specified
cross block. The factor \(8f^2 T^{4f}\) is absorbed in
the exponential by Theorem~\ref{thm:prime210}.
The sum of \(H^{-3}\) over all nonzero integers of \(E\)
is at most \(27^f\) by the disk count of
Theorem~\ref{thm:codimension294}. On \(H\ge H_0\) take
the remaining negative powers at \(H_0\), yielding the
first inequality.

For large \(x\), \(3jR-4j-3\ge2jR\) for every
indicated \(j\). Set \(\ell=1\), take \(M_D>C+D+1\),
and sum the resulting geometric series in \(j\) to obtain
the second assertion. Orthogonal compression decreases
the norms. Finally bounded \(H\) bounds every complex
coordinate of \(\xi\), so only finitely many integers
and then finitely many relative scalars occur.
\end{proof}

Theorem~\ref{thm:codimension294} treats all heights at \(j=1\). The bounded-norm calculation below gives a second estimate when \(j\ge2\).

\subsection{The full contribution up to a bound on the absolute norm}\label{nav:I-71}
For each \(2\le j\le s\) let \(\mathcal B_j\) be the one-block
family of Theorem~\ref{thm:blocktail295}, with \(c=Q\xi\).
The absolute norm and the height of \(\xi\) can be very
different on the prescribed varying fields. The following
sum includes every unit in \(\mathcal O_E\) allowed by the
support conditions, at every block length.

\begin{theorem}[Single reversed blocks of bounded absolute norm]\leavevmode
\label{thm:blocknorm296}
Put \(d=k+1\). Uniformly for \(0\le\vartheta\le1\), one has
\[
  \sum_{\substack{w\in\mathcal B_j\\
           |N_{E/\mathbb Q}\xi|^{1/2}\le Q^{\vartheta f}}}
      \|O_w\|_{\rm op}
   \ \le\ (d!)^{-jf}\exp(Crjf)\,Q^{-(1-\vartheta)rjf}.
\]
There is no restriction on the complex sizes of the conjugates
of \(\xi\), on the relative scalar, or on its valuations beyond
the stated norm and support conditions. The assertion holds on
the original fields, words, levels and intervals and after the
prescribed orthogonal projections.
\end{theorem}
\begin{proof}
For a single label multiply the bounds of
Theorem~\ref{thm:blockfactorial296} at all real places.
Their product bounds the real operator relative to the
diagonal, since the positive norms dominate the polarized
forms on the given irreducible lowest type, as in
Theorem~\ref{thm:blocktail295}. The rational measure factor is
at most one by Lemma~\ref{lem:blockphysical295}.
The product of \(|\tau c|\) is \(Q^f A_\xi\), where
\(A_\xi=|N_{E/\mathbb Q}\xi|^{1/2}\).
Proposition~\ref{prop:blockmarks296} therefore gives
\[
  \|O_w\|_{\rm op}
   \le (d!)^{-jf}\exp(Crjf)
                Q^{-(B_j-j^2)f} A_\xi^{\,2rj-B_j}.
\]
Here \(B_j-rj=j^2(4s-j)-2js-j^2\). For \(j\ge2\) and
\(s\ge j\) this is at least \(j^2\), since
\(s(4-2/j)\ge j+2\). It is also \(O(j^2s)=o(rj)\)
uniformly. Thus \(2rj-B_j>0\) for large \(x\).
Substitute \(Q^{\vartheta f}\) for \(A_\xi\). The remaining
negative exponent of \(Q^f\) is at least
\((1-\vartheta)rj\), since its excess is
\((1+\vartheta)(B_j-rj)-j^2\ge0\).

To sum over all heights take \(H_*=\exp(Mfs\log b)\) with
a fixed sufficiently large \(M\). Below this height the disk
estimate \(\sum_{\xi\ne0}H(\xi)^{-3}\le27^f\) from
Theorem~\ref{thm:codimension294} gives at most
\(27^f H_*^3\) integers. For each there are at most
\(8f^2T^{4f}H_*^{4j}\) admissible scalars, by
Theorem~\ref{thm:blocktail295}. The logarithm of the resulting
count is \(O_M(jfs\log b)=o(rjf)\).

For the heights \(H(\xi)\ge H_*\) use
Theorem~\ref{thm:blocktail295} with \(\ell=1\).
Since \(3j(r-1)-4j-3\ge 2j(r-1)\), choosing \(M\)
greater than its absolute cost constant gives a tail at most
\(\exp(-jrfs\log b)\), even without the norm restriction.
This is at most \((d!)^{-jf}Q^{-rjf}\) for sufficiently
large \(x\). Indeed \(\log(d!)\le d\log d\),
\(\log Q=O(\log b)\) by the marked-prime choice and
\(\log d=O(\log x)=o(\log b)\), while \(s\) tends to
infinity. Combining the two sums proves the estimate.
Orthogonal projections are contractions.
\end{proof}

Theorem~\ref{thm:allblocks297} removes the norm restriction by estimating the finite column operator at every prime.

\subsection{The real orbital integral for every transverse pair}\label{nav:I-72}
Consider the remaining transverse positions on the fields and
intervals of Section~\ref{sec:missing}, so \(\mathcal G>0\). At one real
place let \({\cal P}_\tau\) be the integral of the product
of the two section norms divided by their diagonal norms,
as in Theorem~\ref{thm:regularreal299}. On the common group
use probability on its compact factor and ordinary additive
measure on the entries of its internal unitriangular blocks
in the normalized monomial frames. At the original frequency
these are coordinates in the decorated frames with terminal
vector \(z=\varpi e\). Let \(J_i\) be the block endpoints for
the extended pair, \(i(p)\) the index of the block containing
\(p\), and \(u_p=d-J_{i(p)}\). Use \(\delta_p\) from
Proposition~\ref{prop:extend300} and \(R\) from
Proposition~\ref{prop:slice300}.

\begin{theorem}[A real bound including the last nonsingleton block]\leavevmode
\label{thm:realtrans300}
Put \(\kappa=k-d\), \(E_0=\kappa+2ds+s^2\), with \(d=s+1\).
On sufficiently distant common intervals, for every pair of
prescribed words and \(\eta_i\ge0\) with
\(\sum_i J_i\eta_i\le 200d(d+\log_2(2(k+1)))\), one has
\[
 {\cal P}_\tau\ \le\
 \frac{\exp(Ckd)\,R^{-2\kappa}\,\delta_d^{-E_0}}
      {(k!)^s\,T^{s(k+d)}}
 \prod_{p=1}^d\delta_p^{-2u_p}
 \prod_i\max(1,\delta_{J_i})^{-\eta_i}.
\]
This holds also for \(D_s=0\). In that case
\(R=\sqrt2\) and \(\delta_d=|\tau D_d|/(2T)^d\).
\end{theorem}
\begin{proof}
We start at unit frequency as in Proposition~\ref{prop:extend300}.
With exterior products in the extended space set
\(\alpha_0(g)=\|\bigwedge^d(g^{-1}P_0)\|\). This is positive
and left \(B\)-invariant, since the action on the frame is
upper unitriangular. Let \(g\) minimize the operator norm in
its original \(B\)-coset and put \(N=\nu(g)\). Extending an
original isometry by the identity on \(t\) preserves this norm,
which is at least one. In the fiber above \(g\) write
\(y=w^{-1}b g\), \(b\in J\backslash B\), using the representatives
in Proposition~\ref{prop:slice300}. The columns of \(y^{-1}P_0\)
are
\[
        g^{-1}(Z_0^\dagger,-U,I_d)^{\rm col}\,A^\dagger C .
\]
Indeed \(b_0^{-1}b=\operatorname{diag}(A,h,A^{-*})n(U,Z_0)\)
and \(b_0^{-1}wP_0=P_-C\).

On \(\nu(y)\le H\), choose a minimizing left translate of \(y\).
The columns before \(g^{-1}\) in this formula then have operator
norm at most \(M=NH\) after right multiplication by an upper
unitriangular matrix \(b_1\). This uses the equality of the
operator norms of an isometry and its inverse in the signed
Euclidean coordinates. In particular \(\|A^\dagger C b_1\|\le M\).
Apply the minor identity of Lemma~\ref{lem:pivots298} to this
matrix, permuting rows within each block and removing pivot
phases. In column \(p\) the \(u_p\) entries strictly between
blocks have moduli bounded by \(M^p/\delta_p\), and
\(\delta_p\le M^p\). The operation relating these entries to
the coordinates of \(A\) permutes and conjugates them and
preserves volume. Thus the \(A\)-volume is at most
\(\pi^{\sum u_p}M^{2\sum p u_p}\prod_p\delta_p^{-2u_p}\).
All singular values of \(A^\dagger C b_1\) are at most \(M\)
and their product is \(\delta_d\). Its inverse consequently has
norm at most \(M^{d-1}/\delta_d\). The column formula gives
\(\|U\|,\|Z_0\|\le L=M^d/\delta_d\).

For the terminal determinant, factor the same \(g\) in the
extended unitary group as
\[
        g=n(U_g,Z_g)\operatorname{diag}(A_g,h_g,A_g^{-*})k_g
\]
where \(k_g\) is Euclidean unitary. To obtain this factorization
move \(g^{-1}\operatorname{span}(P_0)\) onto the standard
maximal isotropic space by a compact isometry. Its projections
to the two signed spaces identify isometric subspaces, and their
orthonormal bases extend to bases of the signed spaces. The
remaining parabolic element has the displayed block form by the
unitary equations. In particular \(\alpha_0(g)=|\det A_g|^{-1}\).
Write \(n(U,Z_0)n(U_g,Z_g)=n(U+U_g,Z')\). Its block relation gives
\(-Z'-Z'^\dagger=(U+U_g)^\dagger(U+U_g)\). In the formula for
\(y^{-1}P_0\), remove the left compact factor \(k_g^{-1}\).
The leading \(d\)-square is now \(A_g^{-1}Z'^\dagger A^\dagger C\).
Its determinant modulus is bounded by the full exterior norm
\(\alpha_0(y)\). Also, if \(-Z'-Z'^\dagger=K\ge0\), then
\(\det K\le 2^d|\det Z'|\). For \(K>0\) conjugate \( -Z'\)
by \(K^{-1/2}\). Its real part is \(I/2\) and all its singular
values are at least \(1/2\) by Cauchy--Schwarz. For singular
\(K\) the inequality is automatic. We have proved
\[
 \det((U+U_g)^\dagger(U+U_g))
       \le \frac{2^d}{\delta_d}\frac{\alpha_0(y)}{\alpha_0(g)} .
\]

Integrate \((\alpha_0(g)/\alpha_0(y))^\kappa\) on this shell.
At fixed \(A\) the last column of \(U\) has the sphere law
of Proposition~\ref{prop:slice300}. The first \(s\) columns
lie in balls of radius \(L\), and the free skew entries have
volume at most \(\exp(Cs^2)L^{s^2}\) since the skew part has
norm at most \(L\). Enlarge these ranges independently and use
Lemma~\ref{lem:spheregram300}, translated by \(U_g\).
The quotient density supplies \(2^s\). From
\(k!/(k-i-1)!\le k^{i+1}\) and \(d\log k=o(k)\) on these
intervals, all its fixed numerical factors and \(2^{d\kappa}\)
are at most \(\exp(Ckd)/(k!)^s\) after the factorials in the
lemma are included. The shell integral is therefore bounded by
\[
 \frac{e^{Ckd}}{(k!)^s}\,
 R^{-2\kappa}\delta_d^{-E_0}\prod_p\delta_p^{-2u_p}
       M^{K_0},\qquad
 K_0=d(2ds+s^2)+2\sum_p p u_p .
\]
For endpoints with \(\delta_{J_i}\ge1\), multiply by
\((M^{J_i}/\delta_{J_i})^{\eta_i}\ge1\) on the shell.
Write \(K_1\) for the enlarged exponent of \(M\).

Integrate the shell inequality against \(pH^{-p-1}\,dH\)
for \(H\ge1\) with \(p=K_1+1\). By Tonelli this inserts
\(\nu(y)^{-p}\) and leaves \(pN^{K_1}\).
Apply Cauchy--Schwarz to the two original section norms on
\(J\backslash G\) with squared-integral weights
\[
       \left(\alpha_0(g)/\alpha_0(w^{-1}g)\right)^\kappa
       \left(\nu(g)/\nu(w^{-1}g)\right)^p
       \quad\hbox{and its reciprocal}.
\]
In the first squared integral the section norm descends to
\(B\backslash G\) and the preceding calculation integrates
its fiber. The remaining normalized moment is of order
\(2K_1+1=O(d^3+d\log k)\). The local estimate in
Theorem~\ref{thm:moment286}, rounding the exponent upwards,
costs \(\exp(O(sd^3\log b))\le\exp(Ckd)\) by
Theorem~\ref{thm:rangenew50}. For the reciprocal integral
translate by \(w\) and interchange the two frames. Their cross
matrix becomes \(C^*\), with the same pivot moduli in each
reversed block and the same \(R\). Common unipotent coordinates
in the second frame are transported by \(D\mapsto C^{-1}D^{-*}C\).
Inside each block the pivot scalar cancels, and reversal, phase
changes and unitriangular inversion preserve ordinary Haar
measure. The compact common factor has unit mass in both
frames. Thus the second disintegration has exactly the required
normalization and uses the diagonal moment for the second word.
This proves the estimate at unit frequency.

Finally conjugate from the physical frequency as in
Proposition~\ref{prop:scale284}. On \(B\) the change scales
the \((k+1)s\) rectangular complex entries by modulus
\(\sqrt T\), the \(s^2\) real center coordinates by \(T\),
and the \(s\)-flag upper entries with modulus one.
Its Haar modulus is \(T^{s(k+d)}\).
On \(J\) it carries the original monomial-frame coordinates
to the unit ones with modulus one. Indeed all \(d\) vectors of the
first extended decorated frame have acquired the same scalar
\(\sqrt2\,\tau\varpi\) relative to the unit frame, by
Proposition~\ref{prop:extend300}. Thus the unitriangular
coefficients in those frames agree entry by entry, and the
compact probabilities are preserved. The change of normalized
section integrals contributes the common modulus divided by
the \(B\)-modulus, by the quotient integration formula in
Proposition~\ref{prop:scale284}. This gives the displayed
power of \(T\).
\end{proof}

This estimates a real factor of the original orbital integral
even when the isotropic cross determinant vanishes. For an
integral projected period one must still combine the bound
with the finite operators and the adelic measure of the common
\(U(k-1)\). In particular the determinant entering the estimate
is \(|\tau\Delta|/(2T)^d\), and integrality of \(\Delta/Q^d\)
alone gives no uniform lower bound of one for this number.

\subsection{The archimedean integral for a whole isotropic block}\label{nav:III-60}
At one real place use Euclidean unit Witt vectors and
write \(l_j(g)\) for the Euclidean norm of the exterior
product of the first \(j\) vectors after \(g^{-1}\).
In these coordinates let \(C_0\) be the invertible
cross block for an exchange as in Proposition~\ref{prop:block295}.
Use \(W\in\mathbb C^{n\times j}\) and a skew-Hermitian
\(j\)-square matrix \(V\) as coordinates on the unipotent
radical. In its upper right block put
\(Z=-W^\dagger H W/2+V\), where \(H\) is the signed
Hermitian matrix on \(M_j\) and \(\dagger\) is ordinary
conjugate transpose. The three upper blocks in the
unipotent are \(-W^\dagger H,W,Z\). Give \(W\) ordinary
complex area entry by entry and \(V\) imaginary diagonal
length and complex area above the diagonal. These are Haar
coordinates by the upper triangular multiplication law.

\begin{theorem}[The matrix volume on a first-block fiber]\leavevmode
\label{thm:matrixfiber295}
For \(A=2\alpha\) with \(\alpha>n+3j/2-1\), the integral
of \(l_j(w^{-1}u_0g)^{-A}\) over this radical equals
\[
        |\det C_0|^{-A}\,l_j(g)^{A-2(n+j)}\,K_{n,j,A}.
\]
Its constant satisfies
\[
 K_{n,j,2\alpha}\le
 \frac{\pi^{nj+j^2/2}}{2^{j(j-1)/2}}
 \prod_{i=1}^j
   \frac{\Gamma(\alpha-n-j/2-i+1)}
        {\Gamma(\alpha-i+1)} .
\]
\end{theorem}
\begin{proof}
The maximal compact group acts transitively on isotropic
\(j\)-planes because their positive and negative projections
are isometric copies with matched orthonormal bases.
Thus one may write
\(g=u_1\operatorname{diag}(D,m,D^{-*})k_0\), where
\(m\) is unitary for \(H\) and \(k_0\) is Euclidean
unitary. Right translation in \(U\) absorbs \(u_1\).
The columns defining the transformed exterior vector
then have blocks
\[
       (D^{-1}Z^\dagger,\ -m^{-1}W,\ D^\dagger)^{\rm col}
                    C_0 .
\]
Here the notation means stacking the three blocks.
The \(D\)-dependent change \(W=m W_1 D^\dagger\),
\(V=D V_1D^\dagger\) gives \(Z=D Z_1D^\dagger\).
The Jacobian is \(|\det D|^{2(n+j)}\). Indeed the
complex rectangular block gives the power \(2n\), as
\(|\det m|=1\). Congruence on skew-Hermitian matrices
gives the power \(2j\), by diagonalizing the positive
part of \(D\) and multiplying the real diagonal and
complex off-diagonal scalings. The common right matrix
on the stacked columns is \(D^\dagger C_0\). Since
\(l_j(g)=|\det D|^{-1}\), this proves the first identity
with
\[
 K_{n,j,2\alpha}
  =\int\det(I+W^\dagger W+ZZ^\dagger)^{-\alpha}\,dW\,dV .
\]

For a positive Hermitian \(j\)-square matrix \(B\) use
the analogous independent entry measure. Set
\(\Gamma_j(z)=\int_{B>0}e^{-\operatorname{tr}B}
(\det B)^{z-j}\,dB\) for \(z>j-1\).
Integrating its last column by the Schur complement gives
\(\Gamma_j(z)=\pi^{j-1}\Gamma(z-j+1)\Gamma_{j-1}(z)\).
In fact if \(B'\) is the preceding principal block and
\(b\) the off-diagonal column, the new positive scalar is
\(B_{jj}-b^\dagger (B')^{-1}b\), and the Gaussian integral
over \(b\) gives \(\pi^{j-1}\det B'\). Congruence by a
positive \(S^{-1/2}\) consequently gives
\[
 (\det S)^{-\alpha}
   =\Gamma_j(\alpha)^{-1}
    \int_{B>0}e^{-\operatorname{tr}(BS)}
                        (\det B)^{\alpha-j}\,dB .
\]
Insert \(S=I+W^\dagger W+ZZ^\dagger\). Tonelli's
theorem applies to the positive integrals. At fixed
\(B,W\), the Gaussian in \(Z\) integrated over the
translate of the skew-Hermitian subspace is at most its
integral on that subspace through zero. To see this,
orthogonally project the translation onto the subspace
for the real inner product associated to the positive
quadratic form \(\operatorname{tr}(BZZ^\dagger)\).
Translation inside the subspace preserves measure and
the orthogonal component supplies a factor at most one.

Conjugate \(B\) unitarily to diagonal entries \(b_i>0\).
The skew-Hermitian Gaussian gives
\(\pi^{j^2/2}(\det B)^{-1/2}
\prod_{i<i'}(b_i+b_{i'})^{-1}\). By the arithmetic and
geometric mean inequality this is at most
\(\pi^{j^2/2}2^{-j(j-1)/2}(\det B)^{-j/2}\).
The \(W\)-Gaussian gives \(\pi^{nj}(\det B)^{-n}\)
by integration over its rows. The remaining integral is
\(\Gamma_j(\alpha-n-j/2)\), finite in the stated range.
Taking the quotient of the two gamma integrals proves
the bound.
\end{proof}

\subsection{Minors controlling the entries between blocks}\label{nav:III-62}
\begin{lemma}[Quotients of minors for the internal unipotent]\leavevmode
\label{lem:pivots298}
Use \(C,z\) from Proposition~\ref{prop:composition298}.
After a signed permutation of rows, \(z^*C b_1\), for any
upper unitriangular \(b_1\), becomes \(L D b_1\), where
\(D_{pp}=c_{i(p)}\) is diagonal and \(L\) is block lower
unitriangular with identity diagonal blocks. Its strict
off-block entries are, up to sign, the conjugates of the
independent entries in \(z\), each used once. If an index
\(\ell\) belongs to a later block than \(p\), then
\[
  L_{\ell p}G_p=
   \det (LDb_1)_{(1,\ldots,p-1,\ell),(1,\ldots,p)} .
\]
Also \(\det (LDb_1)_{[p],[p]}=G_p\).
\end{lemma}
\begin{proof}
Write \(C=S D\) with \(S\) the block diagonal product of
the signed reversals. Premultiplication by \(S^{-1}\)
gives \(L=S^{-1}z^*S\), of the stated shape. In the first
\(p\) columns, upper triangularity of \(b_1\) permits only
rows \(1,\ldots,p\) of that matrix to enter. Its leading
square has determinant one. The corresponding part of \(D\)
has determinant \(G_p\). On selecting the first \(p-1\)
rows and row \(\ell\) of \(L\), its final column is zero
except for \(L_{\ell p}\), and the preceding leading
determinant is one. This proves the formula. Selecting the
first \(p\) rows instead gives the last assertion.
\end{proof}

\subsection{A spherical integral for the two terminal vectors}\label{nav:III-63}
\begin{lemma}[The average of a negative power of the terminal pairing]\leavevmode
\label{lem:angle299}
Let \(a_g=|h(e,m_g e)|\ge1\), \(B_\tau=|\tau\beta|>1\)
and use the probability measure on \(h_0\) in
Proposition~\ref{prop:middle299}. For an integer
\(1\le p\le k-1\),
\[
 \int a_1^{-p}\,dh_0
    \le (a_gB_\tau)^{-p}
        \frac{\Gamma(k)\Gamma(k-p)}{\Gamma(k-p/2)^2}
    \le e^{Ck}(a_gB_\tau)^{-p}.
\]
\end{lemma}
\begin{proof}
After multiplying the two terminal vectors by phases, their
negative coordinates are \(a_g,B_\tau\). Their positive
parts have lengths \(\sqrt{a_g^2-1}\) and \(\sqrt{B_\tau^2-1}\).
The second positive direction is uniform on the complex
unit sphere. Consequently \(a_1\) has the distribution of
\(a_gB_\tau|1-\rho z|\), with \(z\) the first coordinate
of a uniform unit vector in \(\mathbb C^k\) and
\[
               \rho=\sqrt{1-a_g^{-2}}\sqrt{1-B_\tau^{-2}}<1.
\]
This formula also applies to \(a_g=1\).

Write \((a)_m=a(a+1)\cdots(a+m-1)\), with \((a)_0=1\).
Independent rotations give
\(\mathbb E[z^m\bar z^n]=0\) for \(m\ne n\) and
\(\mathbb E|z|^{2m}=m!/(k)_m\).
For the latter formula take independent complex Gaussians
with density \(\pi^{-1}e^{-|u|^2}\). Their direction is uniform,
independent of their radius. The numerator moment of the first
Gaussian is \(m!\) and the \(2m\)-th radial moment is
\(\Gamma(k+m)/\Gamma(k)\), by polar integration.
With \(\alpha=p/2\), the absolutely convergent binomial series
for \((1-\rho z)^{-\alpha}\) and its conjugate now give
\[
 \mathbb E|1-\rho z|^{-p}
   =\sum_{m\ge0}\frac{(\alpha)_m^2}{m!(k)_m}\rho^{2m}
   \le \frac{B(\alpha,k-2\alpha)}{B(\alpha,k-\alpha)} .
\]
For the inequality substitute the beta integral with density
\(x^{\alpha-1}(1-x)^{k-\alpha-1}/B(\alpha,k-\alpha)\)
for the ratio \((\alpha)_m/(k)_m\), by taking its \(m\)-th
moment. Sum the other binomial series with \(\rho=1\)
inside this positive integral. It gives the additional
factor \((1-x)^{-\alpha}\), which is integrable since
\(k-2\alpha>0\). The resulting beta ratio is the first
gamma quotient in the assertion. At \(\rho=1\) the
positive series increases termwise with \(p\). Its largest
value in this range is attained at \(p=k-1\).
Stirling's formula bounds
\(\Gamma(k)/\Gamma((k+1)/2)^2\) by \(e^{Ck}\).
\end{proof}

\subsection{A Gram integral with one column on a sphere}\label{nav:III-64}
\begin{lemma}[A spherical column and translated ball cutoffs]\leavevmode
\label{lem:spheregram300}
Let \(k>d=s+1\ge2\), \(L,R>0\), and put \(\kappa=k-d\).
In a \(k\)-by-\(d\) complex matrix \(A\), integrate each of the
first \(s\) columns over a ball of radius \(L\) with arbitrary
center, using ordinary volume. Integrate its last column over
a sphere of radius \(R\) with arbitrary center, using uniform
probability. Then
\[
 \int \det(A^\dagger A)^{-\kappa}\,dA
 \le R^{-2\kappa}\pi^{ks}L^{2ds}
       \prod_{i=1}^s\frac{1}{i!(k-i-1)!(d-i)} .
\]
\end{lemma}
\begin{proof}
First fix the last column \(v\). The event \(v=0\) has sphere
measure zero. Starting with \(v\), add the other columns in their
given order. Their successive perpendicular lengths and \(\|v\|\)
have squared product \(\det(A^\dagger A)\) by Gram--Schmidt.
The columns added so far are independent almost everywhere by
Fubini, since every proper complex subspace has zero volume.

Integrate the added columns in reverse order. For the \(i\)-th
added column its parallel component lies in complex dimension
\(i\), and its perpendicular component in dimension \(k-i\).
The original ball is contained in the product of the two
projected balls of radius \(L\). The parallel volume is
\(\pi^iL^{2i}/i!\). A translate of the perpendicular ball
integrates \(\|z\|^{-2\kappa}\) to at most its value on the
centered ball. Indeed each superlevel set of the integrand
is a centered ball and its intersection with a ball of fixed
radius is largest when their centers coincide. Integrating the
superlevel volumes proves the claim. Polar integration gives
the other factor
\(\pi^{k-i}L^{2(d-i)}/((k-i-1)!(d-i))\).
These bounds do not depend on the preceding columns.

It remains to integrate \(\|v\|^{-2\kappa}\). In real dimension
\(n=2k\), the average of \(\|u+R\omega\|^{2-n}\) over unit
directions \(\omega\) is at most \(R^{2-n}\), for every \(u\).
Here is a direct proof. Away from \(\|u\|=R\) the average is
radial and harmonic by differentiation under the integral,
since the radial Laplacian of \(r^{2-n}\) vanishes on \(r>0\).
Solutions of \(f''+(n-1)f'/r=0\) have the form
\(a+b r^{2-n}\). In the interior boundedness at zero and the
value at zero give the constant \(R^{2-n}\).
In the exterior the asymptotic \(r^{2-n}(1+o(1))\) at
infinity gives \(r^{2-n}\). On the intervening sphere Fatou's
lemma applied from inside gives the asserted upper bound.
Raise this nonnegative integrand to the power
\(\kappa/(k-1)\in(0,1)\). Jensen's inequality bounds the new
average by \(R^{-2\kappa}\). Multiplication of the factors
proves the result.
\end{proof}

\subsection{Integration of the inverse Gram determinant}\label{nav:V-52}
\begin{lemma}[A rectangular determinant integral with a column cutoff]\leavevmode
\label{lem:blockgram296}
Let \(d\ge j\ge1\) be integers and \(L>0\). Give the
\(d\)-by-\(j\) complex matrix \(Y\) ordinary entry measure, and
restrict each column to the Euclidean ball of radius \(L\).
Then
\[
  \int \det(Y^\dagger Y)^{-(d-j)}\,dY
  \le
  \pi^{dj}L^{2j^2}
  \prod_{i=1}^j
      \frac{1}{(i-1)!(d-i)!(j-i+1)}.
\]
\end{lemma}
\begin{proof}
Outside a set of measure zero all initial columns are independent.
Indeed, having fixed \(i-1<d\) independent columns, their complex
span has zero \(2d\)-dimensional measure. Decompose the \(i\)-th
column into components parallel and perpendicular to this span.
Write \(t_i\) for the length of the perpendicular component.
Gram--Schmidt orthogonalization gives
\(\det(Y^\dagger Y)=\prod_i t_i^2\).

Integrate the last column first and continue in reverse order.
At each step enlarge the ball to the product of the two balls
of radius \(L\) for its orthogonal components. The parallel
part has complex dimension \(i-1\) and volume
\(\pi^{i-1}L^{2(i-1)}/(i-1)!\). For the other part put
\(v=d-i+1\). Polar integration gives
\[
 \int_{\|z\|\le L}\|z\|^{-2(d-j)}dz
    =\frac{\pi^v L^{2(j-i+1)}}{(v-1)!(j-i+1)}.
\]
In fact the sphere at radius one has area \(2\pi^v/(v-1)!\),
as follows by integrating \(e^{-\|z\|^2}\) in Cartesian and
polar coordinates. The radial power is integrable since
\(j-i+1>0\). The product of these two factors does not
depend on the preceding columns. Tonelli's theorem and
multiplication over \(i\) give the asserted bound.
\end{proof}

\subsection{The mixed real orbital integral with two terminal lines}\label{nav:V-54}
Retain the real quotient in Proposition~\ref{prop:middle299}.
For lowest-type sections \(W_1,W_2\) of prescribed words
use the positive real quantity
\[
 {\cal P}_\tau(w)=
 \frac{\displaystyle\int_{J\backslash G}
          \|W_1(g)\|\,\|W_2(w^{-1}g)\|\,d\dot g}
      {\displaystyle\prod_{\ell=1}^2
        \left(\int_{B\backslash G}\|W_\ell(g)\|^2\,d\dot g\right)^{1/2}} .
\]
Take nonzero sections, ordinary coordinate measures on the
unipotents and probability measure on each compact middle
group. The sections in the numerator and denominator have
the same physical terminal frequency. Set \(d_0=k+1\),
\(E_s=d_0-s+3s^2\), and \(a_*=-D_{s+1}/Q^s=T\beta a_t\).
The congruences in Proposition~\ref{prop:shift203} imply
that \(a_*\) is integral at every finite place.

\begin{theorem}[Real decay in the separation of the negative lines]\leavevmode
\label{thm:regularreal299}
For \(\eta_i,\eta_*\ge0\) with
\(\sum_iJ_i\eta_i+(s+1)\eta_*\le 200s(s+1+\log_2(2(k+1)))\),
the local integral satisfies the bound
\[
\begin{split}
 {\cal P}_\tau(w)
 &\le \frac{e^{Crs}}{(d_0!)^s}\,
    |\tau\beta|^{-(d_0-3s)}
       |\tau G_s|^{-E_s}\prod_{p=1}^s|\tau G_p|^{-2u_p}\\
 &\hspace{1.3cm}\cdot
     \max(1,|\tau a_*|)^{-\eta_*}
     \prod_{i=1}^t\max(1,|\tau a_i|)^{-\eta_i}.
\end{split}
\]
The two words may coincide or be distinct. The constant is
uniform in their choice, in the relative middle vector and
in the prefix composition.
\end{theorem}
\begin{proof}
Work first at unit terminal frequency on
\(y=w^{-1}h_0\widehat z u_0g\), with \(g\) minimal.
Use the notation of Theorem~\ref{thm:compositionreal298},
now \(j=s\) and \(h=0\). On the shell \(\nu(y)\le H_0\)
write \(M_0=NH_0\). The first \(s\) columns still have
blocks \((Z^\dagger,-W,I_s)^{\rm col}\Omega b_1\), since
the rotations and \(m\) act only on the middle space.
Thus all bounds there for the entries between blocks, and
the norm bound \(M_0^s/\delta_{[s]}\) for \(W,Z,V\),
continue to hold.

Set \(b_1'=d_0-s\). In bounding
\(({\cal A}(g)/{\cal A}(y))^{b_1'}\) on this shell,
Lemma~\ref{lem:twoneg299} contributes the factor
\[
       (2^s)^{b_1'}\delta_{[s]}^{-b_1'}
                        a_1^{2s}(Z_g/a_1)^{b_1'}
\]
before the inverse Gram determinant integral. Its column
radius has an additional factor \(N\), since
\(\|Y\|\le 3N\|W\|\). By Lemma~\ref{lem:blockgram296}
the cost of this radius increase is \(N^{2s^2}\).
All surviving powers of \(\delta_{[p]}\) and \(M_0\)
are the same as in Theorem~\ref{thm:compositionreal298}.

Integrate the factors involving \(a_1,Z_g\) over the actual
compact fiber.
Put \(p_1=b_1'-2s=d_0-3s\), which lies between \(1\)
and \(k-1\) on the large intervals. Since
\(Z_g=\sqrt{2a_g^2-1}\le\sqrt{2}\,a_g\), Lemma~\ref{lem:angle299}
gives
\[
  \int a_1^{2s}(Z_g/a_1)^{b_1'}\,dh_0
       \le e^{Ck}a_g^{2s}|\tau\beta|^{-p_1}
       \le e^{Ck}N^{2s}|\tau\beta|^{-p_1}.
\]
This is valid before contracting the column vectors or
summing finite cells.

In addition to the endpoint tests in
Theorem~\ref{thm:compositionreal298}, the shell bounds
\(|\tau\beta|\delta_{[s]}\le M_0^{s+1}\). Indeed take
the Hermitian pairing of \(P\wedge e\) with its translated
image. Its modulus for the exchange is this product by
block diagonality. Left and right upper changes of
decorated frame preserve the exterior vector and its
pairing, as do the compact column rotations. Its Euclidean
pairing bound on a representative of norm at most \(M_0\)
is \(M_0^{s+1}\). Insert its inverse power \(\eta_*\)
with the compensating factor \(M_0^{(s+1)\eta_*}\)
whenever \(|\tau a_*|\ge1\). Treat the other endpoint
tests at the places \(|\tau a_i|\ge1\) as before.

The squared-integral weights in Cauchy--Schwarz are the
ratio of the full exterior norms to power \(b_1'\) and the
ratio of the heights to power one larger than the total
shell exponent of \(H_0\), and their reciprocals.
Integrate the shell bounds as in
Theorem~\ref{thm:compositionreal298}. The extra moments
of \(N\) above have exponent \(2s^2+2s\). The endpoint
tests cost at most the weighted sum in the hypothesis
in each shell. All required base moments therefore have
exponent \(O(s^3+s\log(k+1))\), within
Theorem~\ref{thm:moment286}. Their logarithmic cost is
\(O(s^4\log b)=o(rs)\). In the second squared integral
translate by \(w\). The inverse has the same absolute
prefix products and its middle terminal coefficient is
\(\bar\beta\), by the unitary relation for \(m\).
It normalizes \(J\) by Proposition~\ref{prop:middle299}.
The diagonal moment estimate applies to the second word
with the same constants.

The terminal-frequency dilation fixes the middle space,
so it commutes with \(m\) and the compact quotient.
Scale \(\delta_{[p]}\) to \(|\tau G_p|/(2T)^p\) and
use the quotient factor \(T^{-s(n+s)}\) as in
Theorem~\ref{thm:compositionreal298}, with \(n=k+1\).
For the additional terminal test,
\(|\tau\beta|\,|\tau G_s|/(2T)^s
=Q^s|\tau a_*|/(T(2T)^s)\).
Its extra constant has logarithm \(O((s+1)\eta_*\log(2T))\),
within \(O(rs)\). This proves the asserted estimate with
the stated normalizations.
\end{proof}

The real estimate applies to the smaller common column when
\(D_s\ne0\). Its finite orbital factors and the normalization
of the adelic common \(U(k-1)\) still enter the global period.
When \(D_s=0\) the final reversed block contains an isotropic
and a terminal vector, and the representative above does not
cover it.

\subsection{Two distinct negative terminal lines}\label{nav:VI-56}
Return to the remaining transverse pairs with common positive
column of rank \(k-1\). Here we consider those for which the
isotropic cross determinant \(D_s\) is nonzero. Then
\(\mathcal G=|D_{s+1}|^2-T^2|D_s|^2\) is totally positive,
where \(|z|^2=z\bar z\). In the normal form of
Lemma~\ref{lem:forms203} the terminal entry is a singleton
block, since both \(s\) and \(s+1\) are block endpoints.
Use the composition and products \(G_p\) of
Proposition~\ref{prop:composition298} for the isotropic
prefix of length \(s\). Its signed determinant convention
gives \(G_s=D_s\).

\begin{proposition}[A rational representative and its compact real quotient]\leavevmode
\label{prop:middle299}
Write \(\beta=-D_{s+1}/(T D_s)\). One can represent the
pair by an exchange with prefix cross block \(C\) as in
Proposition~\ref{prop:composition298} and middle block
\(m\in U(W_0\perp Ee)\), with \(h(e,e)=-1\), such that
\[
             me=v+\beta e,\qquad h(v,v)=\beta\bar\beta-1,
                     \qquad v\in W_0 .
\]
Here \(m\) fixes \(v^\perp\cap W_0\) pointwise. The common
stabilizer \(J\) is the product of the diagonal prefix
unitriangular groups and \(U(v^\perp\cap W_0)\), and the
exchange normalizes \(J\) with modulus one. At a real place
the extra quotient in \(J\backslash B\) is the unit sphere
of the positive column, with its rotation-invariant probability
measure when both compact groups have mass one.
\end{proposition}
\begin{proof}
After the two upper unitriangular changes of frame, the
terminal vectors are perpendicular to both isotropic lists,
since the last block of the cross matrix is a singleton.
The two isotropic spaces pair perfectly and span a nondegenerate
space of dimension \(2s\). The map fixing the original
isotropic vectors and the original terminal vector, and
sending the second isotropic list to the standard opposite
list multiplied by \(C\), is an isometry on their span.
Witt extension gives a rational isometry of the ambient
space fixing the first decorated frame. The second terminal
vector divided by \(\varpi\) is now a negative unit in
\(W_0\perp Ee\). Its \(e\)-coordinate is \(\beta\) by
the determinant ratio and \(\varpi\bar\varpi=T\).
The norm identity gives the asserted length of \(v\).
In particular \(|\tau\beta|>1\) at every real place.

Put \(t=e-(v+\beta e)\) and \(b_0=h(t,e)=\bar\beta-1\).
This is nonzero. On the middle space take
\[
             m(z)=z-t\,h(t,z)/b_0 .
\]
Expanding the Hermitian pairing shows this is an isometry,
because \(h(t,t)=b_0+\bar b_0\). It sends \(e\) to the
required vector and fixes the indicated perpendicular.
Using it as the middle block therefore realizes both
decorated frames.

An element of \(J\) stabilizes the two opposite isotropic
spaces and is block diagonal on their sum and its perpendicular.
The prefix calculation is the one in
Proposition~\ref{prop:composition298}. The middle part
fixes \(e\) and \(v+\beta e\), so it fixes \(v\) as well.
This gives the stated group. The middle block commutes with
that group, and the exchange acts by inverse adjoints and
reversal inside the prefix blocks. These operations preserve
Haar measure there. With \(\widehat z,u_0\) as in that
proposition, representatives on the real fiber can thus be
taken in the order \(h_0\widehat z u_0\), with \(h_0\)
in the middle \(U(k)\) modulo the left stabilizer \(U(k-1)\).
The pushforward by \(h_0\mapsto h_0^{-1}v/\|v\|\) is the
unique invariant probability on the sphere. Conjugation
by these compact matrices preserves the other Haar factors.
\end{proof}

\subsection{Higher decay at the physical terminal frequency}\label{nav:VII-55}
At a real place \(\tau\), write \(l(g)=\|g^{-1}u'\|\) for
the length of the first Euclidean unit isotropic vector
in Theorem~\ref{thm:fiber283}. Use the same quotient Haar
densities as in Proposition~\ref{prop:places288}. Write
\(\mathcal I_{\tau,\xi}\) for the signed local numerator divided by
its diagonal norm at parameter \(Q\tau\xi\) and the original
terminal character. Thus \(\mathcal I_{\tau,\xi}=I_\tau(\tau(\xi/\sqrt{-q}))\)
in the notation of that proposition.

\begin{lemma}[A real estimate beyond twice the first-plane exponent]\leavevmode
\label{lem:power291}
Let \(R=r-1=n+1\), \(n=\dim_{\CC} M\), and
\(t=|\tau\xi|>0\). On the prescribed words,
\[
       |\mathcal I_{\tau,\xi}|\le a_2\,t^{-3R},\qquad
                       \log a_2\le C R s\log b .
\]
The constant is uniform on the common intervals and across
the real places. The same bound holds for the normalized
local integral of the product of the two section norms.
\end{lemma}
\begin{proof}
First use unit terminal frequency and let \(c_\tau\ne0\)
be the imaginary local orbital parameter. In the first
unipotent coordinates of Theorem~\ref{thm:fiber283} put
\(L=l(g)\). For \(A\ge2R\) its change of variables gives
\[
 \int l((\gamma')^{-1}n(w,z)g)^{-A}\,dw\,dv
     =(|c_\tau|/2)^{-A} L^{A-2R} C_{n,\tau,A} .
\]
Here \(v\) is the real central coordinate. In fact before
dilation the squared length is
\((|c_\tau|/2)^2\{L^{-2}+|w|^2+
L^2(v^2+h_M(w,w)^2/4)\}\), after the unit-Jacobian middle
isometry and central translation. Substituting \(w=y/L\),
\(v=u/L^2\) gives the Jacobian \(L^{-2n-2}\).
The remaining integral has integrand
\((1+|y|^2+u^2+h_M(y,y)^2/4)^{-A/2}\), so
\(C_{n,\tau,A}\le C_{n,\tau,2R}\le\pi^R/n!\).
Changing the sign of \(c_\tau\) has the same result.

On \(J(F_\tau)\backslash G_\tau\), the quotient by the
common stabilizer of the two flags, apply Cauchy--Schwarz as in
Theorem~\ref{thm:moment283}, with weight
\((l(g)/l((\gamma')^{-1}g))^A\) in the first squared
integral and its reciprocal in the second.
Disintegrating the first over \(B_\tau\backslash G_\tau\)
multiplies its diagonal density by \(l(g)^A\) and by the
fiber integral just computed. The product contains the
power \(l(g)^{2A-2R}\). For the second integral translation
by \(\gamma'\) replaces \(c_\tau\) by \(-c_\tau\), because
\(\gamma'\) centralizes the stabilizer. The two bounds
therefore agree. The integrands after squaring are
nonnegative and Tonelli's theorem applies.

Take \(A=3R\). The normalized moment required is
\(\mathbb E_{1,\tau}l(g)^{4R}\le\exp(CRs\log b)\)
by the local proof of Theorem~\ref{thm:moment286}, since
\(4R\le4r^2\). This also proves convergence. Now use
Proposition~\ref{prop:scale284} at the physical frequency.
Its Jacobian contributes \(T^{-R}\), and the unit-frequency
parameter has modulus \(Qt/T\). Thus one may take the
uniform majorant for
\[
     (\pi^R/n!)\,T^{2R}(2/Q)^{3R}
                           \mathbb E_{1,\tau}l(g)^{4R}
\]
as \(a_2\). Theorem~\ref{thm:prime210} gives
\(\log T=O(s\log b)\), which proves the stated size bound.
\end{proof}

\subsection{The real orbital estimate with a factorial for each exchanged vector}\label{nav:VII-56}
Put \(d=k+1,\ h=s-j,\ n=2h+d\), and define
\[
 D_0=d-j,\qquad K_j=j^2(4s-j),\qquad
                         B_j=jD_0+K_j .
\]
Here \(K_j\) is an exponent of a local height. Keep the positive
orbital integral divided by the diagonal norm, with the
coordinate measures in Lemma~\ref{lem:blockphysical295}.

\begin{theorem}[The real factor at the original terminal frequency]\leavevmode
\label{thm:blockfactorial296}
For \(2\le j\le s\) its value at \(\tau\) is at most
\[
          (d!)^{-j}\exp(Crj)\,|\tau c|^{-B_j}.
\]
The constant is uniform in the word, real place, block length,
and point on the prescribed intervals.
\end{theorem}
\begin{proof}
Work first at unit terminal frequency, with
\(\delta=|c_\tau|/2\). Take \(g\) minimal in its Bessel coset,
put \(N=\nu(g)\), and use \(y=w^{-1}u_0g\). For \(H_0\ge1\)
consider the part of the radical on which \(\nu(y)\le H_0\).
There is \(b_0\in B\) with \(\|u_0^{-1}wb_0\|\le NH_0\).
Indeed choose a representative \(b_0^{-1}y\) of norm at most
\(H_0\), take its inverse and multiply by \(g\). For these
unitary matrices the inverse has the same operator norm.
On the first plane \(b_0\) has upper unitriangular matrix
\(b_1\). The first \(j\) columns of \(u_0^{-1}wb_0\) have
blocks \(Z^\dagger,-W,I_j\), each multiplied on the right by
\(C_0 b_1\). Put \(L=NH_0/\delta\). The last block gives
\(\|b_1\|\le L\), hence \(L\ge1\) on the set in question.
The product of the singular values of \(b_1\) is one, so
\(\|b_1^{-1}\|\le L^{j-1}\). The first two blocks now give
\[
                   \|Z\|,\ \|W\|,\ \|V\|\le L^j .
\]
For the center we used \(V=(Z-Z^\dagger)/2\). An empty set
contributes zero throughout.

Put \(U(g,y)={\cal A}(g)/{\cal A}(y)\).
Lemma~\ref{lem:blockell296} gives, almost everywhere,
\[
 U(g,y)^{D_0}\le(2^j\sqrt{2})^{D_0}\delta^{-jD_0}
                                     \det(Y^\dagger Y)^{-D_0}.
\]
Changing the coordinates \(X,\eta\) to \(Y\) gives
\(dX\,d\eta=a^{2j}dY\) by triangularity. Its columns have
length at most \(3L^j\), since \(\|v_+\|/a\le1\).
Apply Lemma~\ref{lem:blockgram296} with this radius.
The entries of \(A,B\) lie in disks of radius \(L^j\).
The skew center lies in the product of its \(j\) real
intervals and \(j(j-1)/2\) complex disks of that radius.
The degree in \(L^j\) from these factors and the integral in
\(Y\) is
\[
                4hj+j^2+2j^2=4hj+3j^2 .
\]
Thus the bound on the set \(\nu(y)\le H_0\) is
\[
 \int U(g,y)^{D_0}\,du_0
 \ \le\ (d!)^{-j}e^{Crj}\,a^{2j}\,
             \delta^{-B_j}(NH_0)^{K_j}.
\]
In estimating the constants we used \(d!/(d-i)!\le d^i\)
for \(1\le i\le j\). The resulting logarithmic cost
\(O(j^2\log d)\), and all powers of \(2,3,\pi\) from the
balls and determinant, are \(O(rj)\) on these intervals.

Take \(p=K_j+1\). Integrate the last bound against
\(pH_0^{-p-1}\,dH_0\) for \(H_0\ge1\). Since \(\nu(y)\ge1\),
Tonelli gives the same shell bound with weight
\(\nu(y)^{-p}\), with \(H_0^{K_j}\) replaced by \(p\).
In the Cauchy--Schwarz inequality on \(J\backslash G\) use
\(U(g,w^{-1}g)^{D_0}(\nu(g)/\nu(w^{-1}g))^p\) for the first
squared section norm and the reciprocal for the second.
Disintegrate over \(B\backslash G\) as in
Lemma~\ref{lem:blockphysical295}. The preceding computation
bounds the inner integral, and \(a\le\nu(g)\). The remaining
diagonal moment has exponent \(2K_j+2j+1\) on \(\nu(g)\).
The reciprocal weight gives the same estimate on changing
variables by \(w\). Indeed \(w^{-1}\) is an exchange of the
same form with the same singular values of its cross block,
and \(w\) preserves the measure on \(J\).

The moment exponent is at most \(4r^2\) for large \(x\).
The local proof of Theorem~\ref{thm:moment286} bounds its
logarithmic cost by \(O(j^2s^2\log b)=o(rj)\), since
\(j\le s\) and \(s^3\log b=o(k)\) by
Theorem~\ref{thm:rangenew50}. This proves the unit-frequency
bound. At the physical frequency replace \(\delta\) by
\(|\tau c|/(2T)\) and multiply by \(T^{-j(n+j)}\), as
in Lemma~\ref{lem:blockphysical295}. The exponent
\(B_j-j(n+j)=K_j-2js\) on \(T\) is \(O(j^2s)\), and
\(\log T=O(s\log b)\) by Theorem~\ref{thm:prime210}.
Finally \(2^{B_j}\le e^{Crj}\) on the same intervals.
All of these costs fit in the stated constant.
\end{proof}

\subsection{The sum over all single reversed blocks}\label{nav:VII-57}
\begin{theorem}[All coefficients and all heights for one reversed block]\leavevmode
\label{thm:allblocks297}
On the fields and intervals of Section~\ref{sec:missing}, for every
prescribed word and every \(2\le j\le s\),
\[
                 \sum_{w\in\mathcal B_j}\|O_w\|_{\rm op}
                       \le ((k+1)!)^{-jf}\exp(Crjf).
\]
Here \(\mathcal B_j\) includes every relevant supported
single reversed block of length \(j<s\), and for \(j=s\)
every such block in the transverse full common column case.
The estimate holds at the original level and terminal character.
Together with Theorem~\ref{thm:codimension294} it gives a total
at most \(((k+1)!)^{-f}\exp(Crf)\) for these positions
including \(j=1\). It remains valid after all prescribed
orthogonal projections.
\end{theorem}
\begin{proof}
Put \(d=k+1,\ R=r-1,\ H=H(\xi)\), and use the exponent
\(B_j=j(d-j)+j^2(4s-j)\) of Theorem~\ref{thm:blockfactorial296}.
For \(j\ge2\), its bounds in the proof of
Theorem~\ref{thm:blocknorm296} give \(kj\le B_j\le 2rj\)
for large \(x\). At a real place put \(t=|\tau\xi|\).
For every \(t>0\) we have the two majorants
\[
       L_0(Qt)^{-B_j},\qquad U_0(Qt)^{-G},\qquad
                  G=6jR,
\]
where \(L_0=(d!)^{-j}\exp(C_1rj)\) by
Theorem~\ref{thm:blockfactorial296} and
\(U_0=\exp(C_2rjs\log b)\) by Lemma~\ref{lem:blockphysical295}
with \(\ell=1\). Increase \(C_2\) to have \(U_0\ge L_0\).
They bound the positive integral with the same measures
and terminal frequency. The product bounds polarized
coefficients relative to the diagonal, with field factor at
most one, by Theorem~\ref{thm:blocktail295}.

Use \(L_0(Qt)^{-B_j}\) at the places \(t<1\). At all other
places take the geometric interpolation with weight
\[
               \theta=\frac{\Delta_j+4j+3}{G-B_j}
\]
on the second estimate. By Proposition~\ref{prop:blockfinite297}
and \(G-B_j\ge 3rj\) for sufficiently large \(r\),
\(\theta=O((s+\log(k+1))/r)<1\), and \(\theta>0\).
Since \(Q\ge1\), the product is at most
\[
    L_0^f (U_0/L_0)^{\theta f}
         (Q^f A_\xi)^{-B_j} H^{-\Delta_j-4j-3}.
\]
Multiply by Proposition~\ref{prop:blockfinite297}. Because
\(B_j\ge kj\) and \(1\le A_\xi\le H\), the factors remaining
from \(Q^f A_\xi\) cost at most \(Q^{f\Delta_j}H^{\Delta_j}\).
The logarithm of \(Q^{\Delta_j}\) is \(O(rj)\), using
\(\log Q=O(\log b)\). For the interpolation constant,
\(\log(d!)\le d\log d\) gives
\[
     \theta\log(U_0/L_0)
          =O\bigl(j(s+\log(k+1))(s\log b+\log d)\bigr)
          =o(rj).
\]
Here \(\log(k+1)=O(\log x)=o(s)\) and
\(s^3\log b=o(k)\) on the common intervals. Thus each single
label costs at most
\[
                  (d!)^{-jf}\exp(Crjf)\,H^{-4j-3}.
\]
For fixed \(\xi\), the number of allowed \(\lambda\) is at most
\(8f^2T^{4f}H^{4j}\) by Theorem~\ref{thm:blocktail295},
and each pair of parameters gives at most one label.
The sum of \(H^{-3}\) over the nonzero integers of \(E\)
is at most \(27^f\) by Theorem~\ref{thm:codimension294}.
Since \(\log T=O(s\log b)\), the extra constants fit into
\(\exp(Crjf)\). This proves the individual block-length
bound. For large \(x\), \((d!)^{-f}\exp(Crf)\) tends to
zero since \(r\asymp k\) and \(\log(d!)\gg k\log k\).
The geometric sum in \(j\ge2\), and the separate
codimension-one result, give the asserted total. Compression
by orthogonal projections decreases operator norms.
\end{proof}

For a prefix with several reversed blocks, Proposition~\ref{prop:composition298} identifies the additional unipotent coordinates. Their real and finite contributions are compared in Theorem~\ref{thm:allflags298}.

\subsection{The real integral for all blocks in a composition}\label{nav:VII-58}
For a composition of \(j\) put \(h=s-j,\ d=k+1\) and
\(E_j=d-j+4hj+3j^2\). Keep the ordinary unipotent coordinate
measures at a single real place \(\tau\). The integral below
is the positive orbital integral relative to the same
diagonal norm as in Theorem~\ref{thm:blockfactorial296}.

\begin{theorem}[Decay in the terminal and intermediate prefix products]\leavevmode
\label{thm:compositionreal298}
Let \(2\le j\le s\). For real \(\eta_i\ge0\) such that
\(\sum_i J_i\eta_i\le200j(s+1+\log_2(2(k+1)))\), the
integral at the original frequency is bounded by
\[
 \frac{\exp(Crj)}{(d!)^j}\,
 |\tau G_j|^{-E_j}
 \prod_{p=1}^j|\tau G_p|^{-2u_p}
 \prod_{i=1}^t\max(1,|\tau a_i|)^{-\eta_i}.
\]
All constants are uniform in the composition and the
field and interval parameters.
\end{theorem}
\begin{proof}
Work first with unit terminal frequency. Put
\(\delta_{[p]}=|\tau G_p|/2^p\). On the fiber above a
minimal representative \(g\), let \(N=\nu(g)\) and
\(y=w^{-1}\widehat z u_0g\). On \(\nu(y)\le H_0\) there
exists an element of \(B\) whose upper unitriangular action
\(b_1\) on the \(j\)-plane gives the following bounds.
With \(\Omega=z^*\tau C/2\), the three blocks of the
first \(j\) columns in the calculation of
Theorem~\ref{thm:blockfactorial296} are
\((Z^\dagger,-W,I_j)^{\rm col}\Omega b_1\), of operator
norm at most \(NH_0\). Here one multiplies the inverse of
a minimal representative of \(y\) on the left by \(g\),
and uses the last identity in Proposition~\ref{prop:composition298}.
Put \(M_0=NH_0\). Every \(p\)-minor of the last block is
at most \(M_0^p\), by the exterior-power norm. In the
notation of Lemma~\ref{lem:pivots298} this implies
\[
 \delta_{[p]}\le M_0^p,\qquad
 |L_{\ell p}|\le M_0^p/\delta_{[p]}
                \quad (i(\ell)>i(p)).
\]
The volume of all possible \(z\) is therefore at most
\(\pi^{\sum u_p} M_0^{2\sum p u_p}
\prod_p\delta_{[p]}^{-2u_p}\), since there are \(u_p\)
such entries in column \(p\), each with ordinary complex area.

Next use \(|\det(\Omega b_1)|=\delta_{[j]}\) and its
norm bound to get
\(\|(\Omega b_1)^{-1}\|\le M_0^{j-1}/\delta_{[j]}\).
The norms of \(W,Z,V\) are therefore at most \(M_0^j/\delta_{[j]}\).
Fix \(z\). In the proof of Lemma~\ref{lem:blockell296}
remove the right matrix \(\Omega\)
from the first \(j\) columns. This
multiplies the pairing determinant by \(\det\Omega\)
and leaves the identical Schur complement and Hermitian
part in that proof. Consequently the bound on
\(\det(Y^\dagger Y)\) uses \(\delta_{[j]}^{-1}\) in
place of \(\delta^{-j}\).
Apply Lemma~\ref{lem:blockgram296} with column radius
\(3M_0^j/\delta_{[j]}\) and use the entry bounds on
the two residual rectangular blocks and the skew center.
The combined degree in \(M_0^j/\delta_{[j]}\) is
\(4hj+3j^2\), as in Theorem~\ref{thm:blockfactorial296}.
It follows that the integral of
\(({\cal A}(g)/{\cal A}(y))^{d-j}\) on this shell,
now over both \(z\) and \(u_0\), is at most
\[
 \frac{e^{Crj}a^{2j}}{(d!)^j}
 M_0^{K}\,\delta_{[j]}^{-E_j}
               \prod_p\delta_{[p]}^{-2u_p},
 \qquad K=j(4hj+3j^2)+2\sum_p p u_p.
\]
Here \(a\le N\) is the scalar of the negative test
vector. All fixed numerical factors from the \(z\)-disks
cost \(e^{O(j^2)}\).

At this place let \(I\) consist of the endpoints with
\(|\tau a_i|\ge1\). On the nonempty shell
\(M_0^{J_i}\delta_{[J_i]}^{-1}\ge1\) for each of these
endpoints. We may multiply the last majorant by
\(\prod_{i\in I}(M_0^{J_i}/\delta_{[J_i]})^{\eta_i}\).
Put \(K'=K+\sum_{i\in I}J_i\eta_i\) and \(p_0=K'+1\).
Integrate the shell bound against \(p_0H_0^{-p_0-1}\,dH_0\)
on \([1,\infty)\). As in Theorem~\ref{thm:blockfactorial296},
apply weighted Cauchy--Schwarz on the full \(J\backslash G\),
using the ratio of the two full terminal exterior norms to
power \(d-j\) and the ratio of the two quotient heights to
power \(p_0\). The second squared integral, on changing
variables by \(w\), uses \(w^{-1}\), whose block coefficients
are unit signs times \(\bar c_i\). All the absolute prefix
products in its shell bound are therefore the same.
The remaining normalized base moment has exponent
\(2K'+2j+1=O(j^2s+j\log(k+1))\). For a nonintegral
exponent take the next integer, as \(\nu\ge1\).
This is less than \(4r^2\) for large \(x\), and
Theorem~\ref{thm:moment286} bounds its logarithmic cost
by \(O(j^2s^2\log b)=o(rj)\).

Finally replace \(\delta_{[p]}\) by
\(|\tau G_p|/(2T)^p\). The frequency dilation acts trivially
on \(z\), since it is scalar on the prefix. Its factor on
the quotient density is therefore \(T^{-j(n+j)}\) as
in Lemma~\ref{lem:blockphysical295}. In the factors for
\(i\in I\) use \(G_{J_i}=Q^{J_i}a_i\) and \(Q\ge1\).
The exponent of \(T\) left over is at most
\[
 jE_j+2\sum_p p u_p+\sum_{i\in I}J_i\eta_i-j(n+j)
                           = O(j^2s+j\log(k+1)).
\]
Its logarithmic cost is \(o(rj)\) because
\(\log T=O(s\log b)\) and \(s^3\log b=o(k)\).
The powers of two cost \(e^{O(rj)}\), proving the assertion.
\end{proof}

\subsection{The terminal determinant for a reversed block}\label{nav:VIII-55}
At a real place keep the Euclidean Witt basis used in
Theorem~\ref{thm:matrixfiber295}. Write \(h=s-j\), \(d=k+1\) and
\(\delta=|c_\tau|/2\) at unit terminal frequency. The cross block
is \(C_0\), with \(C_0/\delta\) Euclidean unitary. For the full
isotropic exterior vector \(P\) and the negative unit terminal
vector \(e\), put
\[
  \alpha(g)=\|g^{-1}P\|,\qquad
  {\cal A}(g)=\|g^{-1}(P\wedge e)\|.
\]
Both descend to the left Bessel quotient. In the \(h\) residual
Witt pairs, the definite column and the terminal line, decompose
the rectangular radical coordinate into row blocks
\[
             W=pA+qB+X+e\eta .
\]
Here \(A,B\) are \(h\) by \(j\), \(X\) is \(k\) by \(j\)
and \(\eta\) is one by \(j\). Denote the skew-Hermitian center
coordinate by \(V\) and the upper right block by \(Z\).

\begin{lemma}[A Gram determinant in the positive and terminal coordinates]\leavevmode
\label{lem:blockell296}
Choose the minimal representative \(g\) of its Bessel coset
as in Lemma~\ref{lem:negative285}. There exist \(a\ge1\) and
a positive-column vector \(v_+\) of squared length \(a^2-1\).
For \(y=w^{-1}u_0g\) form the \(d\)-by-\(j\) matrix
\[
               Y=\begin{pmatrix} X-v_+\eta/a\\ \eta/a\end{pmatrix}.
\]
Then
\[
 \det(Y^\dagger Y)\le
 2^j\sqrt{2}\,\delta^{-j}\frac{{\cal A}(y)}{{\cal A}(g)},
 \qquad a\le\nu(g).
\]
These conclusions hold for every residual scalar \(\lambda\)
in Proposition~\ref{prop:block295}.
\end{lemma}
\begin{proof}
Use the full-flag parabolic and remove its upper unipotent
on the left, retaining a possible right compact factor.
As shown in Lemma~\ref{lem:negative285}, the middle block is
\(m_0\in U(k,1)\). A phase multiple of \(m_0e\) is
\(v_*=v_++ae\). We have
\({\cal A}(g)/\alpha(g)=\sqrt{2a^2-1}\) and
\(\|g^{-1}(P^-\wedge v_*)\|=1/\alpha(g)\), where \(P^-\)
is the exterior product of all the opposite unit flag vectors.
In particular \(a\le\|m_0\|\le\nu(g)\).

Pair \(P^-\wedge v_*\) with \(u_0^{-1}w(P\wedge e)\)
by the determinant of Hermitian pairings. In the first \(j\)
columns remove the right factor \(C_0\). Their blocks are
now \(Z^\dagger,-W,I_j\). Each residual column has a factor
\(\lambda\) of modulus one. After removing them, the pairing
matrix is
\[
 \begin{pmatrix}
 Z^\dagger & B^\dagger &-\eta^\dagger\\
 -A & I_h &0\\
 -(v_+^\dagger X-a\eta)&0&-a
 \end{pmatrix}.
\]
This follows by applying \(u_0^{-1}\) to the residual flag
vectors. Its top row block on a middle vector \(z\) is
\(W^\dagger H z\) for the middle Hermitian form \(H\).
The Schur complement of \(\operatorname{diag}(I_h,-a)\) is
\[
 D=Z^\dagger+B^\dagger A-\eta^\dagger\eta
                                  +\eta^\dagger v_+^\dagger X/a.
\]
Applying \(g^{-1}\) to both exterior vectors preserves their
Hermitian pairing. Its modulus is at most the product of
their Euclidean norms, since the sign matrix of the Hermitian
form on the exterior power is Euclidean unitary. The factors
removed from its determinant have modulus \(a\delta^j\). Thus
\[
          |\det D|\le {\cal A}(y)/(a\delta^j\alpha(g)).
\]
Now \(Z+Z^\dagger=-W^\dagger H W\). Substituting the four
row blocks of \(W\) and using \(\|v_+\|^2=a^2-1\) gives
\(-D-D^\dagger=Y^\dagger Y\). For \(S=Y^\dagger Y/2>0\),
congruence by \(S^{-1/2}\) changes \(-D\) into the identity
plus a skew-Hermitian matrix. Its determinant has modulus at
least one, by diagonalizing the skew-Hermitian part. It follows
that \(2^{-j}\det(Y^\dagger Y)\le|\det D|\), also when \(Y\)
has deficient rank because the left side then vanishes. Finally
\(\sqrt{2a^2-1}/a\le\sqrt2\), which proves the assertion.
\end{proof}

\subsection{The Gram bound for distinct negative vectors}\label{nav:VIII-56}
At unit terminal frequency write a minimal representative
\(g\) as in Lemma~\ref{lem:negative285}, with middle block
\(m_g\in U(k,1)\). Put \(v_*=m_ge\), \(Z_g=\|m_g^{-1}e\|\),
\(N=\nu(g)\). For \(y=w^{-1}h_0\widehat z u_0g\), set
\[
  e_1=h_0^{-1}me,\qquad a_1=|h(v_*,e_1)|,\qquad
             \Omega=z^*\tau C/2.
\]
Write \(W\) for the \((k+1)\)-by-\(s\) rectangular block
of \(u_0\) and \(Z\) for its upper right block. As before
\({\cal A}(g)=\|g^{-1}(P\wedge e)\|\) for the full isotropic
exterior vector \(P\).

\begin{lemma}[Change of rectangular coordinates for the two terminal lines]\leavevmode
\label{lem:twoneg299}
There is a complex linear change from \(W\) to \(Y\), at
fixed \(g,h_0\), with
\[
 dW=a_1^{2s}dY,\qquad \|Y\|\le3N\|W\|,\qquad
 \det(Y^\dagger Y)\le
    \frac{2^s Z_g}{a_1|\det\Omega|}
                     \frac{{\cal A}(y)}{{\cal A}(g)} .
\]
\end{lemma}
\begin{proof}
Let \(H=\operatorname{diag}(I_k,-1)\). The first \(s\)
columns of \(u_0^{-1}\widehat z^{-1}h_0^{-1}w\) are
\((Z^\dagger,-W,I_s)^{\rm col}\Omega\). The column
for \(e\) is \((W^\dagger H e_1,e_1,0)^{\rm col}\).
Pair their exterior product against \(P^-\wedge v_*\),
where \(P^-\) is the product of the opposite isotropic
vectors. After removing \(\Omega\), the terminal diagonal
entry is \(h(v_*,e_1)\), whose modulus \(a_1\) is at
least one. In fact \(m_g^{-1}e_1\) is a negative unit
and the modulus of its \(e\)-coordinate is \(a_1\).
The Schur complement is
\[
 D=Z^\dagger+
       W^\dagger H e_1\,h(v_*,W)/h(v_*,e_1).
\]
The transformed test exterior vector has norm
\(1/\alpha(g)\), with \(\alpha(g)=\|g^{-1}P\|\), and
\({\cal A}(g)/\alpha(g)=Z_g\), as in
Lemma~\ref{lem:blockell296}. Bounding the pairing by the
Euclidean norms therefore gives
\[
             |\det D|\le\frac{{\cal A}(y)}
                          {a_1|\det\Omega|\alpha(g)} .
\]

Write \(m_g^{-1}W=(X',\eta')^{\rm col}\) and take a phase
multiple of \(m_g^{-1}e_1\) with positive final coordinate,
say \(v_1+a_1e\). The phase cancels in the quotient defining
\(D\). Then \(\|v_1\|^2=a_1^2-1\). Since
\(Z+Z^\dagger=-W^\dagger H W\), substitution and completion
of the square give
\[
       -D-D^\dagger=Y^\dagger Y,\qquad
       Y=\begin{pmatrix}X'-v_1\eta'/a_1\\ \eta'/a_1\end{pmatrix}.
\]
Indeed the cross terms are
\(-X'^\dagger v_1\eta'/a_1\) and its adjoint, and the
coefficient on \(\eta'^\dagger\eta'\) is one.
Thus \(2^{-s}\det(Y^\dagger Y)\le|\det D|\), as in
Lemma~\ref{lem:blockell296}. The two matrix changes have
absolute complex determinants one and \(a_1^{-1}\) on
each column, respectively. Finally \(\|m_g^{-1}\|\le N\)
and \(\|v_1\|/a_1\le1\). This gives the Jacobian, the
norm bound and the claimed determinant inequality.
\end{proof}

\subsection{The orbit measure with a fixed positive vector}\label{nav:VIII-57}
Use the unit-frequency coordinates of Proposition~\ref{prop:extend300}
at one real place. Their Hermitian form on the three blocks
has identity off-diagonal blocks pairing the two \(d\)-spaces
and identity on the positive \(k\)-space. Let \(\widetilde B\)
stabilize the decorated flag of \(P_0\) in the extended group,
and \(B\) be its subgroup fixing \(t\). The latter is precisely
the original Bessel subgroup. Write \(\gamma_0=1/\sqrt2\).
Then \(t=(\gamma_0{\bf e}_d,0,\gamma_0{\bf e}_d)\).
The common stabilizer of the original pair is denoted by \(J\).

\begin{proposition}[Coordinates and Haar measure on the original fiber]\leavevmode
\label{prop:slice300}
Use unit mass on compact groups and ordinary additive measures
on the entries of the common unitriangular groups in the
monomial frames. There are coordinates on \(J\backslash B\)
consisting of a block upper unitriangular matrix \(A\) with
identity diagonal blocks, \(s\) arbitrary vectors in \(\CC^k\),
a skew-Hermitian \(s\)-square, and one point on a sphere of
radius
\[
                R=\sqrt{2\,\tau\mathcal G}/|\tau\Delta| .
\]
Their quotient measure is \(2^s\) times ordinary additive
measures and sphere probability, for the original Euclidean
Witt measures on \(B\). When \(D_s=0\) the radius is \(\sqrt2\).
\end{proposition}
\begin{proof}
Every element of \(\widetilde B\) has a unique expression
\[
 \operatorname{diag}(A,h,A^{-*})\,n(U,Z_0),\qquad
 n(U,Z_0)=
 \begin{pmatrix}I&-U^\dagger&Z_0\\0&I&U\\0&0&I\end{pmatrix},
 \qquad Z_0+Z_0^\dagger=-U^\dagger U .
\]
Here initially \(A\) is any upper unitriangular \(d\)-square,
and \(h\in U(k)\). This follows from flag preservation and
the three block equations for the Hermitian form. At \(t\)
the condition that the displayed matrix fix \(t\) gives
\(U_d=0,\ Z_0{\bf e}_d=A^{-1}{\bf e}_d-{\bf e}_d\).
For any \(A,h\) these equations can be solved, for example
by taking \(U=0\) and completing the last column to a
skew-Hermitian matrix. Thus \(B\) realizes every upper
unitriangular change of its extended frame. Left and right
changes within \(B\) allow the two flags to have cross
matrix \(C\) from Proposition~\ref{prop:extend300}. Denote
the resulting unit-frequency relative isometry by \(w\).

There is \(b_0\in\widetilde B\) fixing the first frame
pointwise and taking \(P_- C\) to \(wP_0\). Indeed these
second frames have the same cross matrix with \(P_0\),
and both are totally isotropic. The resulting isometry
of the two nondegenerate spans extends over their positive
perpendiculars. Write \(t'=b_0^{-1}t=(a,m,\gamma_0{\bf e}_d)\).
Pairing \(t'\) with \(P_0\) gives its last block as displayed.
Pairing it with \(P_-C\) gives \(C^*a=\gamma_0{\bf e}_d\),
since \(w\) fixes \(t\). In particular \(a\) lies on the
first coordinate line of the last reversed block. Its positive
middle part satisfies \(\|m\|^2=\tau\mathcal G/|\tau\Delta|^2\)
by Proposition~\ref{prop:extend300}.

Left multiplication by \(b_0^{-1}\) sends \(B\) onto the
elements of \(\widetilde B\) taking \(t\) to \(t'\).
In the displayed coordinates their defining equations are
\[
           U_d=h^{-1}m/\gamma_0,\qquad
           Z_0{\bf e}_d=A^{-1}a/\gamma_0-{\bf e}_d .
\]
The last diagonal entry has the required real part because
\(2\gamma_0\operatorname{Re}a_d+\|m\|^2=1\).
Consequently the first \(s\) columns of \(U\) are free,
as is the leading \(s\)-square of
\((Z_0-Z_0^\dagger)/2\). These data determine every remaining
entry of \(Z_0\) by its Hermitian-part equation.

In these two opposite frames the common stabilizer consists
of block diagonal matrices
\(\operatorname{diag}(D,h_0,D^{-*})\). Indeed preserving
both isotropic spaces eliminates both off-diagonal radicals.
The conditions that \(D\) and \(C^{-1}D^{-*}C\) be upper
unitriangular allow just the entries within individual
diagonal blocks of \(D\), as in
Proposition~\ref{prop:composition298}. They fix \(a\) and
the last dual coordinate vector. Fixing \(t'\) thus requires
exactly \(h_0m=m\). In particular \(h_0\) runs over \(U(k-1)\).
Taking the left quotient removes the diagonal blocks from
\(A\) and replaces \(h\) by \(h^{-1}m/\gamma_0\) on its sphere.

For the density, use product measure on the free entries just
described, on the upper entries of \(A\), and on \(h\).
Left multiplication by a fixed element of \(\widetilde B\)
transports this product between its two fibers whenever the
target vector changes. It left translates \(A,h\). For fixed
\(A,h\) it translates \(U\) by a fixed matrix, by conjugating
the left unipotent through \(\operatorname{diag}(A,h,A^{-*})\).
For fixed \(A,h,U\) it translates the free skew entries as well,
by the multiplication law for \(n(U,Z_0)\). All three changes
have Jacobian one. In particular the product at \(t'=t\) is
Haar on \(B\). To compare it with the original Witt normalization
it suffices to compare at the identity. The \(ks\) positive
rectangular entries, the \(s\)-square skew center, the compact
factor, and the upper entries among the first \(s\) indices
agree to first order. If \(\eta_i\) is the negative middle
coordinate in original rectangular column \(i\), the image of
\(p_d=(t+e)/\sqrt2\) has first-order coefficient
\(\bar\eta_i/\sqrt2\) on \(p_i\). Thus
\(d^{\,2}A_{id}=d^{\,2}\eta_i/2\). This proves the factor
\(2^s\). The block decomposition of \(A\) has product Haar
measure by triangular multiplication. Dividing by the measures
on the common groups gives the stated quotient. All these
calculations are on the fiber of the original group fixing \(t\).
\end{proof}

\subsection{The actual Bessel integral at the displaced real flag}\label{sec:displacedreal344}
Take the integral ratios of Theorem~\ref{thm:dominantflag344} and
the period in Proposition~\ref{prop:dominantexchange344}. Denote
the ordinary integral with probability Haar measures on $[N]$ and
$[H]$ by $P_{\chi_*}(A,\phi)$. A scalar attached to the period
normalization multiplies this integral. Fix a real place $\tau$
and its chosen complex embedding $\sigma$. For the one-parameter
root of Lemma~\ref{lem:root94}, in the original signed coordinate
plane, write
\[
       \chi(n_\tau(t))=\exp(4\pi i\,\operatorname{Re}(\alpha_\tau t)),
          \qquad Q_\tau=|\alpha_\tau\,\sigma(c_s)|.
\]
The coefficient $\alpha_\tau$ is nonzero since this is the simple
root adjacent to the marked middle line. This definition retains
the scalar in the given additive character.

\begin{theorem}[Decay on the fixed row type at the new archimedean frequency]\leavevmode\label{thm:displacedreal344}
Use a smooth square-integrable row $A$ with the level, compact type
and Casimir bounds of Theorems~\ref{thm:rest96} and~\ref{thm:decay96}.
Let $r$ be the rank, $d=[F:\QQ]r^2$, $I$ the index of its finite
invariance subgroup in the standard lattice level, $B_0$ the height
bound in the restriction estimate and $L$ the parameter for
the derivatives in Lemma~\ref{lem:cas96}. For every integer $k\ge0$,
\[
 \begin{split}
 |P_{\chi_*}(A,\phi)|
 &\le \sqrt I\,\exp\{C d\log(2d B_0)\}
       \frac{\{Cr(L+k+d+1)\}^{k+d}}{(4\pi Q_\tau)^k}
                                  \|A\|_2\,\|\phi\|_2 .
 \end{split}
\]
The constant $C$ is absolute. On the common intervals take the
finite level and real parameter bounds of Theorem~\ref{thm:raw97}.
The rational diagonal in Theorem~\ref{thm:dominantflag344} can
also be chosen with $Q_\tau\ge T$. With this choice, uniformly
for sufficiently large $x$ and all the indicated rows and columns,
\[
               |P_{\chi_*}(A,\phi)|
                    \le \exp(-cQ_\tau/r)\|A\|_2\,\|\phi\|_2
\]
with $c>0$ absolute. The logarithmic upper bound for $2gW$ in
that theorem increases by at most
$O(1+\log T+\log^+(1/|\alpha_\tau|))$.
\end{theorem}
\begin{proof}
In the three-dimensional coordinate block $u_s,e,v_s$ the
diagonal scales $u_s$ by $c_s$, fixes $e$ and scales $v_s$ by
$\bar c_s^{-1}$. Conjugating the matrix $n_\tau(t)$ gives
$n_\tau(\sigma(c_s)t)$, including the quadratic term in its
upper corner. Put $\xi=\overline{\alpha_\tau\sigma(c_s)}/Q_\tau$.
Along $b\mapsto n_\tau(b\xi)$, $b\in\RR$, the conjugated character
is $\exp(4\pi i b Q_\tau)$. This root commutes with the column
group. Right translation in $[N]$ and differentiation $k$ times
therefore give
\[
        P_{\chi_*}(X_\xi^k A,\phi)
                    =(4\pi i Q_\tau)^k P_{\chi_*}(A,\phi).
\]
Differentiation is valid on the compact period domains, where the
smooth row and all the derivatives in use are bounded. Cauchy--Schwarz
on the two probability spaces bounds the absolute value on the
left by $\|\phi\|_2\sup_{n,h}|(X_\xi^k A)(nh)|$.
Theorem~\ref{thm:rest96} bounds that supremum by the prefactor
$\sqrt I\exp\{C d\log(2d B_0)\}$ times the largest $L^2$ norm
of derivatives through order $d$ of $X_\xi^k A$. The generator
$X_\xi$ has Frobenius norm two, as in Theorem~\ref{thm:decay96},
since $|\xi|=1$. Thus the words have length at most $k+d$ and
Lemma~\ref{lem:cas96} bounds each norm by
$\{Cr(L+k+d+1)\}^{k+d}\|A\|_2$. This proves the first inequality
for the given global automorphic integral.

In the construction of Theorem~\ref{thm:dominantflag344}, keep the
same congruence class of the rational integer $a$ modulo $J$.
Add $2gT/|\alpha_\tau|$ to the lower bound there for $a$, and take
the first representative above the new bound. The logarithmic
increase asserted in the statement follows from $a\le J$ plus
the new lower bound. The same estimate of the conjugates of
$\theta_s$ still gives $|\sigma(c_s)|\ge W/2=a/(2g)$. This proves
$Q_\tau\ge T$ while retaining all the finite conditions.

Fix an absolute $\kappa>0$ small enough and set
$k=\lfloor\kappa Q_\tau/r\rfloor$. The level and weight hypotheses
give $L\le\exp(C' x^\varepsilon)$, polynomial bounds on $r,d$,
and a polynomial bound for $\log I+\log B_0$.
Lemma~\ref{lem:size97} gives
$\log T\ge(16-o(1))m x^\varepsilon$ with $m\to\infty$.
Thus $Cr(L+k+d+1)/(4\pi Q_\tau)\le1/2$ for large $x$, uniformly
for $Q_\tau\ge T$. The $k$ powers contribute $2^{-k}$. The remaining
logarithm, apart from the input norms, is at most
\[
             \tfrac12\log I+C d\log(2d B_0)+O(d\log(2Q_\tau)).
\]
This is $o(Q_\tau/r)$ uniformly in the same range. Indeed the
first two terms grow polynomially and $\log(2Q)/Q$ decreases
for $Q\ge T$ for large $T$. Absorbing the multiplier proves the
second inequality.
\end{proof}

For example, suppose a nonzero algebraic period $p$ in a number
field $J_0$ has conjugates $a_\nu P_{\chi_{*,\nu}}(A_\nu,\phi_\nu)$,
all satisfying this estimate with $Q_{\tau,\nu}\ge T$ and common
rank $r$. Its denominator ideal $\mathfrak d(p)$ then satisfies
\[
 \frac{\log N\mathfrak d(p)+
      \sum_\nu\log(|a_\nu|\,\|A_\nu\|_2\,\|\phi_\nu\|_2)}
      {[J_0:\QQ]}\ \ge\ cT/r .
\]
Indeed $|N(p)|\ge(N\mathfrak d(p))^{-1}$ because its numerator
ideal is integral. Multiplying the analytic bounds over the
embeddings proves the assertion. Applying it to the designated
theta column would require identifying that algebraic period and
its conjugates and determining the scalars $a_\nu$.

\clearpage
\section{References}\label{nav:bibliography-XX}
\begingroup\def\section*#1{}\begin{thebibliography}{99}\small
\bibitem{XX-XV-ichino328} A. Ichino, \emph{Theta lifting for discrete series representations of real unitary groups}, 2020. \href{https://arxiv.org/abs/2002.09148}{Preprint}.
\bibitem{XX-XV-IV-bp155} R. Beuzart-Plessis, \emph{A local trace formula for the Gan--Gross--Prasad conjecture for unitary groups: the archimedean case}, Ast\'erisque \textbf{418} (2020). \href{https://arxiv.org/pdf/1506.01452}{Sections 1.5, 2.6 and 7.2, Lemma 6.5.1(ii), and Theorem 7.2.1}.
\end{thebibliography}\endgroup
\end{document}
